% Generated human-review packet template.  Source records, Lean statements,
% and raw-source-to-Spec screening fields are substituted by
% scripts/review_dashboard_packet.py.
% Do not hand-edit generated packets; annotate the PDF or reconcile notes into
% the paper-local human-review log.
\documentclass[10pt]{article}
\usepackage[margin=0.43in,footskip=0.2in]{geometry}
\usepackage{fontspec}
% Use a Unicode-complete main face as well as a Unicode-complete code face.
% Reviewer reasons and source labels can contain exact Greek symbols outside
% verbatim blocks; silently dropping one would corrupt the review packet.
\setmainfont{DejaVu Serif}
% The broad DejaVu Sans Unicode coverage preserves Lean's mathematical glyphs
% (including U+2216 set difference and superscript identifiers) in verbatim
% review blocks.  Its proportional metrics are preferable to silently missing
% exact semantic text from a narrower monospace font.
\setmonofont{DejaVu Sans}
\usepackage{hyperref}
\usepackage{amssymb}
\usepackage{fvextra}
\usepackage{enumitem}
\setlength{\parindent}{0pt}
\setlength{\parskip}{0.15em}
\linespread{0.92}
\hypersetup{colorlinks=true,linkcolor=blue,urlcolor=blue,breaklinks=true}
\DefineVerbatimEnvironment{ReviewVerbatim}{Verbatim}{breaklines=true,breakanywhere=true,fontsize=\scriptsize}
\newcommand{\reviewerbox}{%
  \par\noindent\fbox{\parbox{0.96\linewidth}{\rule{0pt}{0.38in}}}\par
}
\newcommand{\reviewmatch}[1]{%
  \par\noindent\textbf{Reviewer check:} $\square$\; #1\par
  \noindent\emph{If left unchecked, explain the discrepancy or uncertainty below.}\par
}

\begin{document}
\fontsize{9}{10}\selectfont

\begin{center}
{\LARGE Human review packet: MendelsonWhang1990PriorityPricing}\\[0.35em]
\small Generated from byte-pinned source excerpts, the expanded PaperInterface
specifications, and hash-bound raw-source-to-Spec screening records.
\end{center}

\noindent\textbf{Paper:} MendelsonWhang1990PriorityPricing\\
\textbf{Paper id:} \texttt{MendelsonWhang1990PriorityPricing}\\
\textbf{Source version:} \parbox[t]{0.76\linewidth}{\raggedright Operations Research 38(5), September--October 1990, pp. 870--883}\\
\textbf{Source record:} \texttt{audit/paper\_statement\_map.json}\\
\textbf{Rows:} 5\\
\textbf{Generated:} 2026-09-07\\
\textbf{Source URL:} \href{https://pubsonline.informs.org/doi/10.1287/opre.38.5.870}{official source}





\section*{How to use this packet}
For each source claim, compare the exact source excerpts with the expanded
PaperInterface specification. Mark up this PDF or fill the blank reviewer box in the TeX source;
return the annotated file for reconciliation into the paper-local human-review
log.

\section*{Open the interactive dashboard}
The interactive dashboard is an optional alternative to using this PDF; it
presents the same review surface rather than an additional review requirement.
After forking and cloning (or pulling) AppliedModelingLib, run the following from the
repository root. The cache command is needed only the first time in a new clone.
\begin{ReviewVerbatim}
cd <your-repository-checkout>
lake exe cache get
python3 scripts/review_dashboard.py --paper MendelsonWhang1990PriorityPricing --serve
\end{ReviewVerbatim}
Open \url{http://127.0.0.1:8765/} in a browser and keep that terminal running
while reviewing. The dashboard records saved annotations in this paper's local
reviewer-owned review record.

\section*{Contents}
\begin{itemize}[leftmargin=1.2em]
\item \hyperlink{governing-declaration-section-5ef5ef0364b6939c}{Governing Lean declarations (86; supporting context)}
\item \hyperlink{library-section-994978aceb8dd183}{Semantic library prerequisites (23; not paper claims)}
\begin{itemize}[leftmargin=1.2em]
\item \hyperlink{library-prerequisite-dba97dbb9d8c7c91}{Applied\allowbreak{}Modeling\allowbreak{}Lib.\allowbreak{}Probability.\allowbreak{}Poisson\allowbreak{}Process.\allowbreak{}multiclass\allowbreak{}Stationary\allowbreak{}Poisson\allowbreak{}Work\allowbreak{}Measure}
\item \hyperlink{library-prerequisite-4b8788f2f9337ce3}{Applied\allowbreak{}Modeling\allowbreak{}Lib.\allowbreak{}Probability.\allowbreak{}stabilized\allowbreak{}Finite\allowbreak{}Replay\allowbreak{}Response}
\item \hyperlink{library-prerequisite-755cdb28ceaf46d2}{Applied\allowbreak{}Modeling\allowbreak{}Lib.\allowbreak{}Queueing.\allowbreak{}Multiclass\allowbreak{}Stationary\allowbreak{}Poisson\allowbreak{}Work\allowbreak{}Class\allowbreak{}Tagged\allowbreak{}Sample}
\item \hyperlink{library-prerequisite-2373fe691c32d999}{Applied\allowbreak{}Modeling\allowbreak{}Lib.\allowbreak{}Queueing.\allowbreak{}Nonpreemptive\allowbreak{}Priority\allowbreak{}Job}
\item \hyperlink{library-prerequisite-e68cb5fce009ec53}{Applied\allowbreak{}Modeling\allowbreak{}Lib.\allowbreak{}Queueing.\allowbreak{}Nonpreemptive\allowbreak{}Priority\allowbreak{}Work\allowbreak{}State}
\item \hyperlink{library-prerequisite-837de7f0cc179814}{Applied\allowbreak{}Modeling\allowbreak{}Lib.\allowbreak{}Queueing.\allowbreak{}admit\allowbreak{}Nonpreemptive\allowbreak{}Priority\allowbreak{}Job}
\item \hyperlink{library-prerequisite-96200a424cf39dee}{Applied\allowbreak{}Modeling\allowbreak{}Lib.\allowbreak{}Queueing.\allowbreak{}advance\allowbreak{}Nonpreemptive\allowbreak{}Priority\allowbreak{}Work\allowbreak{}State}
\item \hyperlink{library-prerequisite-d49e277073e052d4}{Applied\allowbreak{}Modeling\allowbreak{}Lib.\allowbreak{}Queueing.\allowbreak{}total\allowbreak{}Nonpreemptive\allowbreak{}Priority\allowbreak{}Work\allowbreak{}Jobs}
\item \hyperlink{library-prerequisite-74e6067041d9c8dc}{Applied\allowbreak{}Modeling\allowbreak{}Lib.\allowbreak{}Queueing.\allowbreak{}run\allowbreak{}Nonpreemptive\allowbreak{}Priority\allowbreak{}Arrival\allowbreak{}Trace}
\item \hyperlink{library-prerequisite-6e4e64d50452f911}{Applied\allowbreak{}Modeling\allowbreak{}Lib.\allowbreak{}Queueing.\allowbreak{}canonical\allowbreak{}Stationary\allowbreak{}Priority\allowbreak{}Class\allowbreak{}Tagged\allowbreak{}Finite\allowbreak{}Window\allowbreak{}State}
\item \hyperlink{library-prerequisite-bb4b8c8010d5959b}{Applied\allowbreak{}Modeling\allowbreak{}Lib.\allowbreak{}Queueing.\allowbreak{}canonical\allowbreak{}Stationary\allowbreak{}Priority\allowbreak{}Class\allowbreak{}Tagged\allowbreak{}Future\allowbreak{}Arrival\allowbreak{}Jobs}
\item \hyperlink{library-prerequisite-cc744f6520cc7734}{Applied\allowbreak{}Modeling\allowbreak{}Lib.\allowbreak{}Queueing.\allowbreak{}clear\allowbreak{}Nonpreemptive\allowbreak{}Priority\allowbreak{}Completion\allowbreak{}Ledger}
\item \hyperlink{library-prerequisite-795b7af1c661bc8d}{Applied\allowbreak{}Modeling\allowbreak{}Lib.\allowbreak{}Queueing.\allowbreak{}finite\allowbreak{}Nonpreemptive\allowbreak{}Priority\allowbreak{}Queue\allowbreak{}Wait}
\item \hyperlink{library-prerequisite-886cbd948a5fd8b1}{Applied\allowbreak{}Modeling\allowbreak{}Lib.\allowbreak{}Queueing.\allowbreak{}nonpreemptive\allowbreak{}Priority\allowbreak{}Recorded\allowbreak{}Completion\allowbreak{}Time}
\item \hyperlink{library-prerequisite-3702583c174721c0}{Applied\allowbreak{}Modeling\allowbreak{}Lib.\allowbreak{}Queueing.\allowbreak{}stationary\allowbreak{}Priority\allowbreak{}Class\allowbreak{}Tagged\allowbreak{}Job}
\item \hyperlink{library-prerequisite-098447dd1f7443fd}{Applied\allowbreak{}Modeling\allowbreak{}Lib.\allowbreak{}Queueing.\allowbreak{}stationary\allowbreak{}Priority\allowbreak{}Class\allowbreak{}Tagged\allowbreak{}Finite\allowbreak{}Replay\allowbreak{}Post\allowbreak{}Arrival\allowbreak{}State}
\item \hyperlink{library-prerequisite-945b0f6ad93a6ca5}{Applied\allowbreak{}Modeling\allowbreak{}Lib.\allowbreak{}Queueing.\allowbreak{}stationary\allowbreak{}Priority\allowbreak{}Class\allowbreak{}Tagged\allowbreak{}Finite\allowbreak{}Replay\allowbreak{}Post\allowbreak{}Arrival\allowbreak{}Response\allowbreak{}Time}
\item \hyperlink{library-prerequisite-c3ca5cbd067c9957}{Applied\allowbreak{}Modeling\allowbreak{}Lib.\allowbreak{}Queueing.\allowbreak{}stationary\allowbreak{}Priority\allowbreak{}Class\allowbreak{}Tagged\allowbreak{}Palm\allowbreak{}Measure}
\item \hyperlink{library-prerequisite-6b2b5577562811ed}{Applied\allowbreak{}Modeling\allowbreak{}Lib.\allowbreak{}Queueing.\allowbreak{}stationary\allowbreak{}Priority\allowbreak{}Class\allowbreak{}Tagged\allowbreak{}Response\allowbreak{}Time}
\item \hyperlink{library-prerequisite-8f1ec496ea0db7a0}{Applied\allowbreak{}Modeling\allowbreak{}Lib.\allowbreak{}Queueing.\allowbreak{}stationary\allowbreak{}Priority\allowbreak{}Class\allowbreak{}Tagged\allowbreak{}Work\allowbreak{}Requirement}
\item \hyperlink{library-prerequisite-e35fa775c8012225}{Applied\allowbreak{}Modeling\allowbreak{}Lib.\allowbreak{}Queueing.\allowbreak{}stationary\allowbreak{}Priority\allowbreak{}Class\allowbreak{}Tagged\allowbreak{}Queue\allowbreak{}Wait}
\item \hyperlink{library-prerequisite-1fb8084d5335bf95}{Applied\allowbreak{}Modeling\allowbreak{}Lib.\allowbreak{}Queueing.\allowbreak{}stationary\allowbreak{}Priority\allowbreak{}Work\allowbreak{}Requirement}
\item \hyperlink{library-prerequisite-ea60ce4afba27cd7}{Applied\allowbreak{}Modeling\allowbreak{}Lib.\allowbreak{}Queueing.\allowbreak{}stationary\allowbreak{}Priority\allowbreak{}Work\allowbreak{}Requirement\_\allowbreak{}has\allowbreak{}Law}
\end{itemize}
\item \hyperlink{paper-prerequisite-section-5ef5ef0364b6939c}{Paper-specific semantic prerequisites (19; not paper claims)}
\begin{itemize}[leftmargin=1.2em]
\item \hyperlink{paper-prerequisite-60ffddd9622d7d65}{Mendelson\allowbreak{}Whang1990\allowbreak{}Priority\allowbreak{}Pricing.\allowbreak{}priority\allowbreak{}Time\allowbreak{}Dependent\allowbreak{}Price}
\item \hyperlink{paper-prerequisite-ae05cc1b41c33470}{Mendelson\allowbreak{}Whang1990\allowbreak{}Priority\allowbreak{}Pricing.\allowbreak{}priority\allowbreak{}Time\allowbreak{}Linear\allowbreak{}Coefficient}
\item \hyperlink{paper-prerequisite-1840b9f7b498b425}{Mendelson\allowbreak{}Whang1990\allowbreak{}Priority\allowbreak{}Pricing.\allowbreak{}priority\allowbreak{}Time\allowbreak{}Quadratic\allowbreak{}Coefficient}
\item \hyperlink{paper-prerequisite-d3f248a0395e1c76}{Mendelson\allowbreak{}Whang1990\allowbreak{}Priority\allowbreak{}Pricing.\allowbreak{}heterogeneous\allowbreak{}Priority\allowbreak{}Time\allowbreak{}Dependent\allowbreak{}Price}
\item \hyperlink{paper-prerequisite-49957d909aa24bc5}{Mendelson\allowbreak{}Whang1990\allowbreak{}Priority\allowbreak{}Pricing.\allowbreak{}priority\allowbreak{}Queueing\allowbreak{}Time}
\item \hyperlink{paper-prerequisite-69133735a09302ef}{Mendelson\allowbreak{}Whang1990\allowbreak{}Priority\allowbreak{}Pricing.\allowbreak{}homogeneous\allowbreak{}Priority\allowbreak{}Queueing\allowbreak{}Time}
\item \hyperlink{paper-prerequisite-3db4e2685735dcc0}{Mendelson\allowbreak{}Whang1990\allowbreak{}Priority\allowbreak{}Pricing.\allowbreak{}homogeneous\allowbreak{}Priority\allowbreak{}Price}
\item \hyperlink{paper-prerequisite-9c41d2efaf9290fc}{Mendelson\allowbreak{}Whang1990\allowbreak{}Priority\allowbreak{}Pricing.\allowbreak{}homogeneous\allowbreak{}Priority\allowbreak{}Sojourn\allowbreak{}Time}
\item \hyperlink{paper-prerequisite-53f58e4a79f3cd3c}{Mendelson\allowbreak{}Whang1990\allowbreak{}Priority\allowbreak{}Pricing.\allowbreak{}homogeneous\allowbreak{}Priority\allowbreak{}Expected\allowbreak{}Cost}
\item \hyperlink{paper-prerequisite-0efbdff1257f8766}{Mendelson\allowbreak{}Whang1990\allowbreak{}Priority\allowbreak{}Pricing.\allowbreak{}homogeneous\allowbreak{}Priority\allowbreak{}Population}
\item \hyperlink{paper-prerequisite-6c93de2469b53cb2}{Mendelson\allowbreak{}Whang1990\allowbreak{}Priority\allowbreak{}Pricing.\allowbreak{}individual\allowbreak{}Demand\allowbreak{}Condition}
\item \hyperlink{paper-prerequisite-7dbaca4145f3a4a5}{Mendelson\allowbreak{}Whang1990\allowbreak{}Priority\allowbreak{}Pricing.\allowbreak{}net\allowbreak{}Value}
\item \hyperlink{paper-prerequisite-54c5682d350b5d3d}{Mendelson\allowbreak{}Whang1990\allowbreak{}Priority\allowbreak{}Pricing.\allowbreak{}priority\allowbreak{}Cheating\allowbreak{}Penalty}
\item \hyperlink{paper-prerequisite-476808203e3008c1}{Mendelson\allowbreak{}Whang1990\allowbreak{}Priority\allowbreak{}Pricing.\allowbreak{}priority\allowbreak{}Time\allowbreak{}Dependent\allowbreak{}Expected\allowbreak{}Cost}
\item \hyperlink{paper-prerequisite-84dca75fd93aa845}{Mendelson\allowbreak{}Whang1990\allowbreak{}Priority\allowbreak{}Pricing.\allowbreak{}stable\allowbreak{}Priority\allowbreak{}Flow\allowbreak{}Domain}
\item \hyperlink{paper-prerequisite-3ae6b5293fffafb7}{Mendelson\allowbreak{}Whang1990\allowbreak{}Priority\allowbreak{}Pricing.\allowbreak{}priority\allowbreak{}Time\allowbreak{}Dependent\allowbreak{}Pricing\allowbreak{}Optimal\allowbreak{}And\allowbreak{}Incentive\allowbreak{}Compatible\allowbreak{}At}
\item \hyperlink{paper-prerequisite-33c9dc55ddc7d5b6}{Mendelson\allowbreak{}Whang1990\allowbreak{}Priority\allowbreak{}Pricing.\allowbreak{}source\allowbreak{}Delay\allowbreak{}Cost\allowbreak{}Conditions}
\item \hyperlink{paper-prerequisite-9c435671d506207c}{Mendelson\allowbreak{}Whang1990\allowbreak{}Priority\allowbreak{}Pricing.\allowbreak{}source\allowbreak{}Value\allowbreak{}Function\allowbreak{}Conditions}
\item \hyperlink{paper-prerequisite-816477dc3af70c36}{Mendelson\allowbreak{}Whang1990\allowbreak{}Priority\allowbreak{}Pricing.\allowbreak{}stationary\allowbreak{}Class\allowbreak{}Tagged\allowbreak{}Expected\allowbreak{}Priority\allowbreak{}Time\allowbreak{}Dependent\allowbreak{}Price}
\end{itemize}
\item \hyperlink{source-section-f10b5f68e4acf3dc}{Source claims (5)}
\begin{itemize}[leftmargin=1.2em]
\item \hyperlink{source-claim-0d52797e81514cb8}{theorem\allowbreak{}One\allowbreak{}Externality\allowbreak{}Pricing\allowbreak{}Spec}
\item \hyperlink{source-claim-86fb5cbbc8012c1d}{appendix\allowbreak{}Priority\allowbreak{}Queueing\allowbreak{}Time\allowbreak{}Spec}
\item \hyperlink{source-claim-35dff0bacecfebeb}{theorem\allowbreak{}Two\allowbreak{}Priority\allowbreak{}Incentive\allowbreak{}Compatibility\allowbreak{}Spec}
\item \hyperlink{source-claim-a044a878df73f034}{theorem\allowbreak{}Three\allowbreak{}Priority\allowbreak{}Time\allowbreak{}Dependent\allowbreak{}Pricing\allowbreak{}Spec}
\item \hyperlink{source-claim-324a284201e724f2}{theorem\allowbreak{}Four\allowbreak{}Cheating\allowbreak{}Penalty\allowbreak{}Monotonicity\allowbreak{}Spec}
\end{itemize}
\end{itemize}

\clearpage
\hypertarget{governing-declaration-section-5ef5ef0364b6939c}{}
\section*{Governing Lean declarations}
These exact graph-bound declaration bodies are retained by the displayed specifications or reviewed prerequisites. They are supporting context and do not add source-result rows or semantic judgments.
\hypertarget{governing-declaration-fb5ec66edd06d906}{}
\subsection*{\small\ttfamily\raggedright Applied\allowbreak{}Modeling\allowbreak{}Lib.\allowbreak{}Probability.\allowbreak{}Palm.\allowbreak{}Shift\allowbreak{}Invariant\allowbreak{}Probability\allowbreak{}Law}
\textbf{Declaration location:} reusable library
\paragraph{Lean declaration body}
\begin{ReviewVerbatim}
field Pbase:
{Ω : Type u_1} →
  [inst : MeasurableSpace Ω] →
    AppliedModelingLib.Probability.Palm.ShiftInvariantProbabilityLaw Ω → MeasureTheory.Measure Ω

field isProbability:
∀ {Ω : Type u_1} [inst : MeasurableSpace Ω] (self : AppliedModelingLib.Probability.Palm.ShiftInvariantProbabilityLaw Ω),
  MeasureTheory.IsProbabilityMeasure self.Pbase

field shift:
{Ω : Type u_1} →
  [inst : MeasurableSpace Ω] → AppliedModelingLib.Probability.Palm.ShiftInvariantProbabilityLaw Ω → ℝ → Ω → Ω

field shift_zero:
∀ {Ω : Type u_1} [inst : MeasurableSpace Ω] (self : AppliedModelingLib.Probability.Palm.ShiftInvariantProbabilityLaw Ω),
  self.shift 0 = id

field shift_add:
∀ {Ω : Type u_1} [inst : MeasurableSpace Ω] (self : AppliedModelingLib.Probability.Palm.ShiftInvariantProbabilityLaw Ω)
  (s t : ℝ), self.shift (s + t) = self.shift s ∘ self.shift t

field shift_preserving:
∀ {Ω : Type u_1} [inst : MeasurableSpace Ω] (self : AppliedModelingLib.Probability.Palm.ShiftInvariantProbabilityLaw Ω)
  (t : ℝ), MeasureTheory.MeasurePreserving (self.shift t) self.Pbase self.Pbase
\end{ReviewVerbatim}


\hypertarget{governing-declaration-0d8c93dcf696a606}{}
\subsection*{\small\ttfamily\raggedright Applied\allowbreak{}Modeling\allowbreak{}Lib.\allowbreak{}Probability.\allowbreak{}Poisson\allowbreak{}Process.\allowbreak{}interarrival}
\textbf{Declaration location:} reusable library
\paragraph{Lean declaration body}
\begin{ReviewVerbatim}
type:
ℕ → (ℕ → ℝ) → ℝ

value:
fun (n : ℕ) (a : ℕ → ℝ) => a n
\end{ReviewVerbatim}


\hypertarget{governing-declaration-544406615ce2d8dc}{}
\subsection*{\small\ttfamily\raggedright Applied\allowbreak{}Modeling\allowbreak{}Lib.\allowbreak{}Probability.\allowbreak{}Poisson\allowbreak{}Process.\allowbreak{}arrival\allowbreak{}Time}
\textbf{Declaration location:} reusable library
\paragraph{Lean declaration body}
\begin{ReviewVerbatim}
type:
ℕ → (ℕ → ℝ) → ℝ

value:
fun (n : ℕ) (a : ℕ → ℝ) => ∑ i ∈ Finset.range (n + 1), AppliedModelingLib.Probability.PoissonProcess.interarrival i a
\end{ReviewVerbatim}

\paragraph{Direct retained dependencies}
\begin{itemize}\raggedright
\item \hyperlink{governing-declaration-0d8c93dcf696a606}{Applied\allowbreak{}Modeling\allowbreak{}Lib.\allowbreak{}Probability.\allowbreak{}Poisson\allowbreak{}Process.\allowbreak{}interarrival}
\end{itemize}
\hypertarget{governing-declaration-9fb1a8c287ed38cc}{}
\subsection*{\small\ttfamily\raggedright Applied\allowbreak{}Modeling\allowbreak{}Lib.\allowbreak{}Probability.\allowbreak{}Poisson\allowbreak{}Process.\allowbreak{}canonical\allowbreak{}Renewal\allowbreak{}Count}
\textbf{Declaration location:} reusable library
\paragraph{Lean declaration body}
\begin{ReviewVerbatim}
type:
ℝ → (ℕ → ℝ) → ℕ

value:
fun (t : ℝ) (a : ℕ → ℝ) =>
  if h : ∃ (n : ℕ), t < AppliedModelingLib.Probability.PoissonProcess.arrivalTime n a then Nat.find h else 0
\end{ReviewVerbatim}

\paragraph{Direct retained dependencies}
\begin{itemize}\raggedright
\item \hyperlink{governing-declaration-544406615ce2d8dc}{Applied\allowbreak{}Modeling\allowbreak{}Lib.\allowbreak{}Probability.\allowbreak{}Poisson\allowbreak{}Process.\allowbreak{}arrival\allowbreak{}Time}
\end{itemize}
\hypertarget{governing-declaration-570cb9faec150ef0}{}
\subsection*{\small\ttfamily\raggedright Applied\allowbreak{}Modeling\allowbreak{}Lib.\allowbreak{}Probability.\allowbreak{}Poisson\allowbreak{}Process.\allowbreak{}canonical\allowbreak{}Renewal\allowbreak{}Count\allowbreak{}LE}
\textbf{Declaration location:} reusable library
\paragraph{Lean declaration body}
\begin{ReviewVerbatim}
type:
ℝ → (ℕ → ℝ) → ℕ

value:
fun (t : ℝ) (a : ℕ → ℝ) =>
  if h : ∃ (n : ℕ), t ≤ AppliedModelingLib.Probability.PoissonProcess.arrivalTime n a then Nat.find h else 0
\end{ReviewVerbatim}

\paragraph{Direct retained dependencies}
\begin{itemize}\raggedright
\item \hyperlink{governing-declaration-544406615ce2d8dc}{Applied\allowbreak{}Modeling\allowbreak{}Lib.\allowbreak{}Probability.\allowbreak{}Poisson\allowbreak{}Process.\allowbreak{}arrival\allowbreak{}Time}
\end{itemize}
\hypertarget{governing-declaration-3729a6ac90fa3324}{}
\subsection*{\small\ttfamily\raggedright Applied\allowbreak{}Modeling\allowbreak{}Lib.\allowbreak{}Probability.\allowbreak{}Poisson\allowbreak{}Process.\allowbreak{}two\allowbreak{}Sided\allowbreak{}Gap}
\textbf{Declaration location:} reusable library
\paragraph{Lean declaration body}
\begin{ReviewVerbatim}
type:
ℤ → (ℤ → ℝ) → ℝ

value:
fun (i : ℤ) (a : ℤ → ℝ) => a i
\end{ReviewVerbatim}


\hypertarget{governing-declaration-9208635caa21f800}{}
\subsection*{\small\ttfamily\raggedright Applied\allowbreak{}Modeling\allowbreak{}Lib.\allowbreak{}Probability.\allowbreak{}Poisson\allowbreak{}Process.\allowbreak{}candidate\allowbreak{}Palm\allowbreak{}Arrival}
\textbf{Declaration location:} reusable library
\paragraph{Lean declaration body}
\begin{ReviewVerbatim}
type:
(ℤ → ℝ) → ℤ → ℝ

value:
fun (ω : ℤ → ℝ) (a : ℤ) =>
  match a with
  | Int.ofNat n => ∑ i ∈ Finset.range n, AppliedModelingLib.Probability.PoissonProcess.twoSidedGap (Int.ofNat i) ω
  | Int.negSucc n =>
    -∑ i ∈ Finset.range (n + 1), AppliedModelingLib.Probability.PoissonProcess.twoSidedGap (Int.negSucc i) ω
\end{ReviewVerbatim}

\paragraph{Direct retained dependencies}
\begin{itemize}\raggedright
\item \hyperlink{governing-declaration-3729a6ac90fa3324}{Applied\allowbreak{}Modeling\allowbreak{}Lib.\allowbreak{}Probability.\allowbreak{}Poisson\allowbreak{}Process.\allowbreak{}two\allowbreak{}Sided\allowbreak{}Gap}
\end{itemize}
\hypertarget{governing-declaration-f492363c01e8969e}{}
\subsection*{\small\ttfamily\raggedright Applied\allowbreak{}Modeling\allowbreak{}Lib.\allowbreak{}Probability.\allowbreak{}Poisson\allowbreak{}Process.\allowbreak{}suspension\allowbreak{}Carrier}
\textbf{Declaration location:} reusable library
\paragraph{Lean declaration body}
\begin{ReviewVerbatim}
type:
(ℤ → ℝ) × ℝ → Prop

value:
fun (a : (ℤ → ℝ) × ℝ) =>
  {p : (ℤ → ℝ) × ℝ | 0 ≤ p.2 ∧ p.2 < AppliedModelingLib.Probability.PoissonProcess.twoSidedGap 0 p.1} a
\end{ReviewVerbatim}

\paragraph{Direct retained dependencies}
\begin{itemize}\raggedright
\item \hyperlink{governing-declaration-3729a6ac90fa3324}{Applied\allowbreak{}Modeling\allowbreak{}Lib.\allowbreak{}Probability.\allowbreak{}Poisson\allowbreak{}Process.\allowbreak{}two\allowbreak{}Sided\allowbreak{}Gap}
\end{itemize}
\hypertarget{governing-declaration-adb3b56227ee711d}{}
\subsection*{\small\ttfamily\raggedright Applied\allowbreak{}Modeling\allowbreak{}Lib.\allowbreak{}Probability.\allowbreak{}Poisson\allowbreak{}Process.\allowbreak{}suspension\allowbreak{}Future\allowbreak{}Path}
\textbf{Declaration location:} reusable library
\paragraph{Lean declaration body}
\begin{ReviewVerbatim}
type:
(ℤ → ℝ) → ℕ → ℝ

value:
fun (a : ℤ → ℝ) (a_1 : ℕ) => AppliedModelingLib.Probability.PoissonProcess.twoSidedGap (Int.ofNat a_1) a
\end{ReviewVerbatim}

\paragraph{Direct retained dependencies}
\begin{itemize}\raggedright
\item \hyperlink{governing-declaration-3729a6ac90fa3324}{Applied\allowbreak{}Modeling\allowbreak{}Lib.\allowbreak{}Probability.\allowbreak{}Poisson\allowbreak{}Process.\allowbreak{}two\allowbreak{}Sided\allowbreak{}Gap}
\end{itemize}
\hypertarget{governing-declaration-3b441995d5597d49}{}
\subsection*{\small\ttfamily\raggedright Applied\allowbreak{}Modeling\allowbreak{}Lib.\allowbreak{}Probability.\allowbreak{}Poisson\allowbreak{}Process.\allowbreak{}suspension\allowbreak{}Gap\allowbreak{}Shift}
\textbf{Declaration location:} reusable library
\paragraph{Lean declaration body}
\begin{ReviewVerbatim}
type:
ℤ → (ℤ → ℝ) → ℤ → ℝ

value:
fun (k : ℤ) (a : ℤ → ℝ) (a_1 : ℤ) => AppliedModelingLib.Probability.PoissonProcess.twoSidedGap (a_1 + k) a
\end{ReviewVerbatim}

\paragraph{Direct retained dependencies}
\begin{itemize}\raggedright
\item \hyperlink{governing-declaration-3729a6ac90fa3324}{Applied\allowbreak{}Modeling\allowbreak{}Lib.\allowbreak{}Probability.\allowbreak{}Poisson\allowbreak{}Process.\allowbreak{}two\allowbreak{}Sided\allowbreak{}Gap}
\end{itemize}
\hypertarget{governing-declaration-5e8dbc9b6cbf2c33}{}
\subsection*{\small\ttfamily\raggedright Applied\allowbreak{}Modeling\allowbreak{}Lib.\allowbreak{}Probability.\allowbreak{}Poisson\allowbreak{}Process.\allowbreak{}suspension\allowbreak{}Past\allowbreak{}Path}
\textbf{Declaration location:} reusable library
\paragraph{Lean declaration body}
\begin{ReviewVerbatim}
type:
(ℤ → ℝ) → ℕ → ℝ

value:
fun (a : ℤ → ℝ) (a_1 : ℕ) => AppliedModelingLib.Probability.PoissonProcess.twoSidedGap (Int.negSucc a_1) a
\end{ReviewVerbatim}

\paragraph{Direct retained dependencies}
\begin{itemize}\raggedright
\item \hyperlink{governing-declaration-3729a6ac90fa3324}{Applied\allowbreak{}Modeling\allowbreak{}Lib.\allowbreak{}Probability.\allowbreak{}Poisson\allowbreak{}Process.\allowbreak{}two\allowbreak{}Sided\allowbreak{}Gap}
\end{itemize}
\hypertarget{governing-declaration-79e7fd2195dfecc3}{}
\subsection*{\small\ttfamily\raggedright Applied\allowbreak{}Modeling\allowbreak{}Lib.\allowbreak{}Probability.\allowbreak{}Poisson\allowbreak{}Process.\allowbreak{}suspension\allowbreak{}Crossing\allowbreak{}Index\allowbreak{}Past\allowbreak{}Closed}
\textbf{Declaration location:} reusable library
\paragraph{Lean declaration body}
\begin{ReviewVerbatim}
type:
ℝ → (ℤ → ℝ) × ℝ → ℤ

value:
fun (t : ℝ) (a : (ℤ → ℝ) × ℝ) =>
  if 0 ≤ a.2 + t then
    Int.ofNat
      (AppliedModelingLib.Probability.PoissonProcess.canonicalRenewalCount (a.2 + t)
        (AppliedModelingLib.Probability.PoissonProcess.suspensionFuturePath a.1))
  else
    Int.negSucc
      (AppliedModelingLib.Probability.PoissonProcess.canonicalRenewalCountLE (-(a.2 + t))
        (AppliedModelingLib.Probability.PoissonProcess.suspensionPastPath a.1))
\end{ReviewVerbatim}

\paragraph{Direct retained dependencies}
\begin{itemize}\raggedright
\item \hyperlink{governing-declaration-9fb1a8c287ed38cc}{Applied\allowbreak{}Modeling\allowbreak{}Lib.\allowbreak{}Probability.\allowbreak{}Poisson\allowbreak{}Process.\allowbreak{}canonical\allowbreak{}Renewal\allowbreak{}Count}
\item \hyperlink{governing-declaration-570cb9faec150ef0}{Applied\allowbreak{}Modeling\allowbreak{}Lib.\allowbreak{}Probability.\allowbreak{}Poisson\allowbreak{}Process.\allowbreak{}canonical\allowbreak{}Renewal\allowbreak{}Count\allowbreak{}LE}
\item \hyperlink{governing-declaration-adb3b56227ee711d}{Applied\allowbreak{}Modeling\allowbreak{}Lib.\allowbreak{}Probability.\allowbreak{}Poisson\allowbreak{}Process.\allowbreak{}suspension\allowbreak{}Future\allowbreak{}Path}
\item \hyperlink{governing-declaration-5e8dbc9b6cbf2c33}{Applied\allowbreak{}Modeling\allowbreak{}Lib.\allowbreak{}Probability.\allowbreak{}Poisson\allowbreak{}Process.\allowbreak{}suspension\allowbreak{}Past\allowbreak{}Path}
\end{itemize}
\hypertarget{governing-declaration-4ca384fa6b6dd395}{}
\subsection*{\small\ttfamily\raggedright Applied\allowbreak{}Modeling\allowbreak{}Lib.\allowbreak{}Probability.\allowbreak{}Poisson\allowbreak{}Process.\allowbreak{}palm\allowbreak{}Tagged\allowbreak{}Arrival\allowbreak{}Indices}
\textbf{Declaration location:} reusable library
\paragraph{Lean declaration body}
\begin{ReviewVerbatim}
type:
ℝ → ℝ → (ℤ → ℝ) → Finset ℤ

value:
fun (a b : ℝ) (g : ℤ → ℝ) =>
  {i ∈
      Finset.Icc (AppliedModelingLib.Probability.PoissonProcess.suspensionCrossingIndexPastClosed a (g, 0))
        (AppliedModelingLib.Probability.PoissonProcess.suspensionCrossingIndexPastClosed b (g, 0)) |
    a ≤ AppliedModelingLib.Probability.PoissonProcess.candidatePalmArrival g i ∧
      AppliedModelingLib.Probability.PoissonProcess.candidatePalmArrival g i < b}
\end{ReviewVerbatim}

\paragraph{Direct retained dependencies}
\begin{itemize}\raggedright
\item \hyperlink{governing-declaration-9208635caa21f800}{Applied\allowbreak{}Modeling\allowbreak{}Lib.\allowbreak{}Probability.\allowbreak{}Poisson\allowbreak{}Process.\allowbreak{}candidate\allowbreak{}Palm\allowbreak{}Arrival}
\item \hyperlink{governing-declaration-79e7fd2195dfecc3}{Applied\allowbreak{}Modeling\allowbreak{}Lib.\allowbreak{}Probability.\allowbreak{}Poisson\allowbreak{}Process.\allowbreak{}suspension\allowbreak{}Crossing\allowbreak{}Index\allowbreak{}Past\allowbreak{}Closed}
\end{itemize}
\hypertarget{governing-declaration-55048043b93956f0}{}
\subsection*{\small\ttfamily\raggedright Applied\allowbreak{}Modeling\allowbreak{}Lib.\allowbreak{}Probability.\allowbreak{}Poisson\allowbreak{}Process.\allowbreak{}palm\allowbreak{}Tagged\allowbreak{}Arrival\allowbreak{}Indices\allowbreak{}Right\allowbreak{}Closed}
\textbf{Declaration location:} reusable library
\paragraph{Lean declaration body}
\begin{ReviewVerbatim}
type:
ℝ → ℝ → (ℤ → ℝ) → Finset ℤ

value:
fun (a b : ℝ) (g : ℤ → ℝ) =>
  {i ∈
      Finset.Icc (AppliedModelingLib.Probability.PoissonProcess.suspensionCrossingIndexPastClosed a (g, 0))
        (AppliedModelingLib.Probability.PoissonProcess.suspensionCrossingIndexPastClosed b (g, 0)) |
    a < AppliedModelingLib.Probability.PoissonProcess.candidatePalmArrival g i ∧
      AppliedModelingLib.Probability.PoissonProcess.candidatePalmArrival g i ≤ b}
\end{ReviewVerbatim}

\paragraph{Direct retained dependencies}
\begin{itemize}\raggedright
\item \hyperlink{governing-declaration-9208635caa21f800}{Applied\allowbreak{}Modeling\allowbreak{}Lib.\allowbreak{}Probability.\allowbreak{}Poisson\allowbreak{}Process.\allowbreak{}candidate\allowbreak{}Palm\allowbreak{}Arrival}
\item \hyperlink{governing-declaration-79e7fd2195dfecc3}{Applied\allowbreak{}Modeling\allowbreak{}Lib.\allowbreak{}Probability.\allowbreak{}Poisson\allowbreak{}Process.\allowbreak{}suspension\allowbreak{}Crossing\allowbreak{}Index\allowbreak{}Past\allowbreak{}Closed}
\end{itemize}
\hypertarget{governing-declaration-2f283554547b82b0}{}
\subsection*{\small\ttfamily\raggedright Applied\allowbreak{}Modeling\allowbreak{}Lib.\allowbreak{}Probability.\allowbreak{}Poisson\allowbreak{}Process.\allowbreak{}suspension\allowbreak{}Flow}
\textbf{Declaration location:} reusable library
\paragraph{Lean declaration body}
\begin{ReviewVerbatim}
type:
ℝ → (ℤ → ℝ) × ℝ → (ℤ → ℝ) × ℝ

value:
fun (t : ℝ) (a : (ℤ → ℝ) × ℝ) =>
  have k : ℤ := AppliedModelingLib.Probability.PoissonProcess.suspensionCrossingIndexPastClosed t a;
  (AppliedModelingLib.Probability.PoissonProcess.suspensionGapShift k a.1,
    a.2 + t - AppliedModelingLib.Probability.PoissonProcess.candidatePalmArrival a.1 k)
\end{ReviewVerbatim}

\paragraph{Direct retained dependencies}
\begin{itemize}\raggedright
\item \hyperlink{governing-declaration-9208635caa21f800}{Applied\allowbreak{}Modeling\allowbreak{}Lib.\allowbreak{}Probability.\allowbreak{}Poisson\allowbreak{}Process.\allowbreak{}candidate\allowbreak{}Palm\allowbreak{}Arrival}
\item \hyperlink{governing-declaration-79e7fd2195dfecc3}{Applied\allowbreak{}Modeling\allowbreak{}Lib.\allowbreak{}Probability.\allowbreak{}Poisson\allowbreak{}Process.\allowbreak{}suspension\allowbreak{}Crossing\allowbreak{}Index\allowbreak{}Past\allowbreak{}Closed}
\item \hyperlink{governing-declaration-3b441995d5597d49}{Applied\allowbreak{}Modeling\allowbreak{}Lib.\allowbreak{}Probability.\allowbreak{}Poisson\allowbreak{}Process.\allowbreak{}suspension\allowbreak{}Gap\allowbreak{}Shift}
\end{itemize}
\hypertarget{governing-declaration-446bfa9c60d40523}{}
\subsection*{\small\ttfamily\raggedright Applied\allowbreak{}Modeling\allowbreak{}Lib.\allowbreak{}Probability.\allowbreak{}Poisson\allowbreak{}Process.\allowbreak{}suspension\allowbreak{}Good\allowbreak{}Gap\allowbreak{}Path}
\textbf{Declaration location:} reusable library
\paragraph{Lean declaration body}
\begin{ReviewVerbatim}
type:
(ℤ → ℝ) → Prop

value:
fun (ω : ℤ → ℝ) =>
  (∀ (i : ℤ), 0 < AppliedModelingLib.Probability.PoissonProcess.twoSidedGap i ω) ∧
    ∀ (k : ℤ),
      Filter.Tendsto
          (fun (n : ℕ) =>
            AppliedModelingLib.Probability.PoissonProcess.arrivalTime n
              (AppliedModelingLib.Probability.PoissonProcess.suspensionFuturePath
                (AppliedModelingLib.Probability.PoissonProcess.suspensionGapShift k ω)))
          Filter.atTop Filter.atTop ∧
        Filter.Tendsto
          (fun (n : ℕ) =>
            AppliedModelingLib.Probability.PoissonProcess.arrivalTime n
              (AppliedModelingLib.Probability.PoissonProcess.suspensionPastPath
                (AppliedModelingLib.Probability.PoissonProcess.suspensionGapShift k ω)))
          Filter.atTop Filter.atTop
\end{ReviewVerbatim}

\paragraph{Direct retained dependencies}
\begin{itemize}\raggedright
\item \hyperlink{governing-declaration-544406615ce2d8dc}{Applied\allowbreak{}Modeling\allowbreak{}Lib.\allowbreak{}Probability.\allowbreak{}Poisson\allowbreak{}Process.\allowbreak{}arrival\allowbreak{}Time}
\item \hyperlink{governing-declaration-adb3b56227ee711d}{Applied\allowbreak{}Modeling\allowbreak{}Lib.\allowbreak{}Probability.\allowbreak{}Poisson\allowbreak{}Process.\allowbreak{}suspension\allowbreak{}Future\allowbreak{}Path}
\item \hyperlink{governing-declaration-3b441995d5597d49}{Applied\allowbreak{}Modeling\allowbreak{}Lib.\allowbreak{}Probability.\allowbreak{}Poisson\allowbreak{}Process.\allowbreak{}suspension\allowbreak{}Gap\allowbreak{}Shift}
\item \hyperlink{governing-declaration-5e8dbc9b6cbf2c33}{Applied\allowbreak{}Modeling\allowbreak{}Lib.\allowbreak{}Probability.\allowbreak{}Poisson\allowbreak{}Process.\allowbreak{}suspension\allowbreak{}Past\allowbreak{}Path}
\item \hyperlink{governing-declaration-3729a6ac90fa3324}{Applied\allowbreak{}Modeling\allowbreak{}Lib.\allowbreak{}Probability.\allowbreak{}Poisson\allowbreak{}Process.\allowbreak{}two\allowbreak{}Sided\allowbreak{}Gap}
\end{itemize}
\hypertarget{governing-declaration-ab1cfe37ad8c32b0}{}
\subsection*{\small\ttfamily\raggedright Applied\allowbreak{}Modeling\allowbreak{}Lib.\allowbreak{}Probability.\allowbreak{}Poisson\allowbreak{}Process.\allowbreak{}Good\allowbreak{}Suspension\allowbreak{}State}
\textbf{Declaration location:} reusable library
\paragraph{Lean declaration body}
\begin{ReviewVerbatim}
type:
Type

value:
{ p : (ℤ → ℝ) × ℝ //
  p ∈ AppliedModelingLib.Probability.PoissonProcess.suspensionCarrier ∧
    AppliedModelingLib.Probability.PoissonProcess.suspensionGoodGapPath p.1 }
\end{ReviewVerbatim}

\paragraph{Direct retained dependencies}
\begin{itemize}\raggedright
\item \hyperlink{governing-declaration-f492363c01e8969e}{Applied\allowbreak{}Modeling\allowbreak{}Lib.\allowbreak{}Probability.\allowbreak{}Poisson\allowbreak{}Process.\allowbreak{}suspension\allowbreak{}Carrier}
\item \hyperlink{governing-declaration-446bfa9c60d40523}{Applied\allowbreak{}Modeling\allowbreak{}Lib.\allowbreak{}Probability.\allowbreak{}Poisson\allowbreak{}Process.\allowbreak{}suspension\allowbreak{}Good\allowbreak{}Gap\allowbreak{}Path}
\end{itemize}
\hypertarget{governing-declaration-1991889bac68724c}{}
\subsection*{\small\ttfamily\raggedright Applied\allowbreak{}Modeling\allowbreak{}Lib.\allowbreak{}Probability.\allowbreak{}Poisson\allowbreak{}Process.\allowbreak{}good\allowbreak{}Suspension\allowbreak{}Flow}
\textbf{Declaration location:} reusable library
\paragraph{Lean declaration body}
\begin{ReviewVerbatim}
type:
ℝ →
  AppliedModelingLib.Probability.PoissonProcess.GoodSuspensionState →
    AppliedModelingLib.Probability.PoissonProcess.GoodSuspensionState

value:
fun (t : ℝ) (a : AppliedModelingLib.Probability.PoissonProcess.GoodSuspensionState) =>
  ⟨AppliedModelingLib.Probability.PoissonProcess.suspensionFlow t ↑a,
    AppliedModelingLib.Probability.PoissonProcess.goodSuspensionFlow._proof_1 t a⟩
\end{ReviewVerbatim}

\paragraph{Direct retained dependencies}
\begin{itemize}\raggedright
\item \hyperlink{governing-declaration-ab1cfe37ad8c32b0}{Applied\allowbreak{}Modeling\allowbreak{}Lib.\allowbreak{}Probability.\allowbreak{}Poisson\allowbreak{}Process.\allowbreak{}Good\allowbreak{}Suspension\allowbreak{}State}
\item \hyperlink{governing-declaration-f492363c01e8969e}{Applied\allowbreak{}Modeling\allowbreak{}Lib.\allowbreak{}Probability.\allowbreak{}Poisson\allowbreak{}Process.\allowbreak{}suspension\allowbreak{}Carrier}
\item \hyperlink{governing-declaration-2f283554547b82b0}{Applied\allowbreak{}Modeling\allowbreak{}Lib.\allowbreak{}Probability.\allowbreak{}Poisson\allowbreak{}Process.\allowbreak{}suspension\allowbreak{}Flow}
\item \hyperlink{governing-declaration-446bfa9c60d40523}{Applied\allowbreak{}Modeling\allowbreak{}Lib.\allowbreak{}Probability.\allowbreak{}Poisson\allowbreak{}Process.\allowbreak{}suspension\allowbreak{}Good\allowbreak{}Gap\allowbreak{}Path}
\end{itemize}
\hypertarget{governing-declaration-a8385d799d0a4d16}{}
\subsection*{\small\ttfamily\raggedright Applied\allowbreak{}Modeling\allowbreak{}Lib.\allowbreak{}Probability.\allowbreak{}Poisson\allowbreak{}Process.\allowbreak{}inst\allowbreak{}Measurable\allowbreak{}Space\allowbreak{}Good\allowbreak{}Suspension\allowbreak{}State}
\textbf{Declaration location:} reusable library
\paragraph{Lean declaration body}
\begin{ReviewVerbatim}
type:
MeasurableSpace AppliedModelingLib.Probability.PoissonProcess.GoodSuspensionState

value:
{ MeasurableSet' := AppliedModelingLib.Probability.PoissonProcess.instMeasurableSpaceGoodSuspensionState._aux_1,
  measurableSet_empty := AppliedModelingLib.Probability.PoissonProcess.instMeasurableSpaceGoodSuspensionState._proof_3,
  measurableSet_compl := AppliedModelingLib.Probability.PoissonProcess.instMeasurableSpaceGoodSuspensionState._proof_4,
  measurableSet_iUnion :=
    AppliedModelingLib.Probability.PoissonProcess.instMeasurableSpaceGoodSuspensionState._proof_5 }
\end{ReviewVerbatim}

\paragraph{Direct retained dependencies}
\begin{itemize}\raggedright
\item \hyperlink{governing-declaration-ab1cfe37ad8c32b0}{Applied\allowbreak{}Modeling\allowbreak{}Lib.\allowbreak{}Probability.\allowbreak{}Poisson\allowbreak{}Process.\allowbreak{}Good\allowbreak{}Suspension\allowbreak{}State}
\item \hyperlink{governing-declaration-f492363c01e8969e}{Applied\allowbreak{}Modeling\allowbreak{}Lib.\allowbreak{}Probability.\allowbreak{}Poisson\allowbreak{}Process.\allowbreak{}suspension\allowbreak{}Carrier}
\item \hyperlink{governing-declaration-446bfa9c60d40523}{Applied\allowbreak{}Modeling\allowbreak{}Lib.\allowbreak{}Probability.\allowbreak{}Poisson\allowbreak{}Process.\allowbreak{}suspension\allowbreak{}Good\allowbreak{}Gap\allowbreak{}Path}
\end{itemize}
\hypertarget{governing-declaration-c513c48a384a886a}{}
\subsection*{\small\ttfamily\raggedright Applied\allowbreak{}Modeling\allowbreak{}Lib.\allowbreak{}Probability.\allowbreak{}Poisson\allowbreak{}Process.\allowbreak{}suspension\allowbreak{}Base\allowbreak{}Arrival}
\textbf{Declaration location:} reusable library
\paragraph{Lean declaration body}
\begin{ReviewVerbatim}
type:
AppliedModelingLib.Probability.PoissonProcess.GoodSuspensionState → ℤ → ℝ

value:
fun (p : AppliedModelingLib.Probability.PoissonProcess.GoodSuspensionState) (i : ℤ) =>
  AppliedModelingLib.Probability.PoissonProcess.candidatePalmArrival (↑p).1 i - (↑p).2
\end{ReviewVerbatim}

\paragraph{Direct retained dependencies}
\begin{itemize}\raggedright
\item \hyperlink{governing-declaration-ab1cfe37ad8c32b0}{Applied\allowbreak{}Modeling\allowbreak{}Lib.\allowbreak{}Probability.\allowbreak{}Poisson\allowbreak{}Process.\allowbreak{}Good\allowbreak{}Suspension\allowbreak{}State}
\item \hyperlink{governing-declaration-9208635caa21f800}{Applied\allowbreak{}Modeling\allowbreak{}Lib.\allowbreak{}Probability.\allowbreak{}Poisson\allowbreak{}Process.\allowbreak{}candidate\allowbreak{}Palm\allowbreak{}Arrival}
\item \hyperlink{governing-declaration-f492363c01e8969e}{Applied\allowbreak{}Modeling\allowbreak{}Lib.\allowbreak{}Probability.\allowbreak{}Poisson\allowbreak{}Process.\allowbreak{}suspension\allowbreak{}Carrier}
\item \hyperlink{governing-declaration-446bfa9c60d40523}{Applied\allowbreak{}Modeling\allowbreak{}Lib.\allowbreak{}Probability.\allowbreak{}Poisson\allowbreak{}Process.\allowbreak{}suspension\allowbreak{}Good\allowbreak{}Gap\allowbreak{}Path}
\end{itemize}
\hypertarget{governing-declaration-cf114817e77adc2f}{}
\subsection*{\small\ttfamily\raggedright Applied\allowbreak{}Modeling\allowbreak{}Lib.\allowbreak{}Probability.\allowbreak{}Poisson\allowbreak{}Process.\allowbreak{}suspension\allowbreak{}Base\allowbreak{}Arrival\allowbreak{}Indices}
\textbf{Declaration location:} reusable library
\paragraph{Lean declaration body}
\begin{ReviewVerbatim}
type:
ℝ → ℝ → AppliedModelingLib.Probability.PoissonProcess.GoodSuspensionState → Finset ℤ

value:
fun (a b : ℝ) (p : AppliedModelingLib.Probability.PoissonProcess.GoodSuspensionState) =>
  {i ∈
      Finset.Icc (AppliedModelingLib.Probability.PoissonProcess.suspensionCrossingIndexPastClosed a ↑p)
        (AppliedModelingLib.Probability.PoissonProcess.suspensionCrossingIndexPastClosed b ↑p) |
    a ≤ AppliedModelingLib.Probability.PoissonProcess.suspensionBaseArrival p i ∧
      AppliedModelingLib.Probability.PoissonProcess.suspensionBaseArrival p i < b}
\end{ReviewVerbatim}

\paragraph{Direct retained dependencies}
\begin{itemize}\raggedright
\item \hyperlink{governing-declaration-ab1cfe37ad8c32b0}{Applied\allowbreak{}Modeling\allowbreak{}Lib.\allowbreak{}Probability.\allowbreak{}Poisson\allowbreak{}Process.\allowbreak{}Good\allowbreak{}Suspension\allowbreak{}State}
\item \hyperlink{governing-declaration-c513c48a384a886a}{Applied\allowbreak{}Modeling\allowbreak{}Lib.\allowbreak{}Probability.\allowbreak{}Poisson\allowbreak{}Process.\allowbreak{}suspension\allowbreak{}Base\allowbreak{}Arrival}
\item \hyperlink{governing-declaration-f492363c01e8969e}{Applied\allowbreak{}Modeling\allowbreak{}Lib.\allowbreak{}Probability.\allowbreak{}Poisson\allowbreak{}Process.\allowbreak{}suspension\allowbreak{}Carrier}
\item \hyperlink{governing-declaration-79e7fd2195dfecc3}{Applied\allowbreak{}Modeling\allowbreak{}Lib.\allowbreak{}Probability.\allowbreak{}Poisson\allowbreak{}Process.\allowbreak{}suspension\allowbreak{}Crossing\allowbreak{}Index\allowbreak{}Past\allowbreak{}Closed}
\item \hyperlink{governing-declaration-446bfa9c60d40523}{Applied\allowbreak{}Modeling\allowbreak{}Lib.\allowbreak{}Probability.\allowbreak{}Poisson\allowbreak{}Process.\allowbreak{}suspension\allowbreak{}Good\allowbreak{}Gap\allowbreak{}Path}
\end{itemize}
\hypertarget{governing-declaration-09b9ceda015085bb}{}
\subsection*{\small\ttfamily\raggedright Applied\allowbreak{}Modeling\allowbreak{}Lib.\allowbreak{}Probability.\allowbreak{}Poisson\allowbreak{}Process.\allowbreak{}suspension\allowbreak{}Base\allowbreak{}Arrival\allowbreak{}Indices\allowbreak{}Right\allowbreak{}Closed}
\textbf{Declaration location:} reusable library
\paragraph{Lean declaration body}
\begin{ReviewVerbatim}
type:
ℝ → ℝ → AppliedModelingLib.Probability.PoissonProcess.GoodSuspensionState → Finset ℤ

value:
fun (a b : ℝ) (p : AppliedModelingLib.Probability.PoissonProcess.GoodSuspensionState) =>
  {i ∈
      Finset.Icc (AppliedModelingLib.Probability.PoissonProcess.suspensionCrossingIndexPastClosed a ↑p)
        (AppliedModelingLib.Probability.PoissonProcess.suspensionCrossingIndexPastClosed b ↑p) |
    a < AppliedModelingLib.Probability.PoissonProcess.suspensionBaseArrival p i ∧
      AppliedModelingLib.Probability.PoissonProcess.suspensionBaseArrival p i ≤ b}
\end{ReviewVerbatim}

\paragraph{Direct retained dependencies}
\begin{itemize}\raggedright
\item \hyperlink{governing-declaration-ab1cfe37ad8c32b0}{Applied\allowbreak{}Modeling\allowbreak{}Lib.\allowbreak{}Probability.\allowbreak{}Poisson\allowbreak{}Process.\allowbreak{}Good\allowbreak{}Suspension\allowbreak{}State}
\item \hyperlink{governing-declaration-c513c48a384a886a}{Applied\allowbreak{}Modeling\allowbreak{}Lib.\allowbreak{}Probability.\allowbreak{}Poisson\allowbreak{}Process.\allowbreak{}suspension\allowbreak{}Base\allowbreak{}Arrival}
\item \hyperlink{governing-declaration-f492363c01e8969e}{Applied\allowbreak{}Modeling\allowbreak{}Lib.\allowbreak{}Probability.\allowbreak{}Poisson\allowbreak{}Process.\allowbreak{}suspension\allowbreak{}Carrier}
\item \hyperlink{governing-declaration-79e7fd2195dfecc3}{Applied\allowbreak{}Modeling\allowbreak{}Lib.\allowbreak{}Probability.\allowbreak{}Poisson\allowbreak{}Process.\allowbreak{}suspension\allowbreak{}Crossing\allowbreak{}Index\allowbreak{}Past\allowbreak{}Closed}
\item \hyperlink{governing-declaration-446bfa9c60d40523}{Applied\allowbreak{}Modeling\allowbreak{}Lib.\allowbreak{}Probability.\allowbreak{}Poisson\allowbreak{}Process.\allowbreak{}suspension\allowbreak{}Good\allowbreak{}Gap\allowbreak{}Path}
\end{itemize}
\hypertarget{governing-declaration-a5c5dbd6c5850d8d}{}
\subsection*{\small\ttfamily\raggedright Applied\allowbreak{}Modeling\allowbreak{}Lib.\allowbreak{}Probability.\allowbreak{}Poisson\allowbreak{}Process.\allowbreak{}two\allowbreak{}Sided\allowbreak{}Interarrival\allowbreak{}Measure}
\textbf{Declaration location:} reusable library
\paragraph{Lean declaration body}
\begin{ReviewVerbatim}
type:
ℝ → MeasureTheory.Measure (ℤ → ℝ)

value:
fun (rate : ℝ) => MeasureTheory.Measure.infinitePi fun (x : ℤ) => ProbabilityTheory.expMeasure rate
\end{ReviewVerbatim}


\hypertarget{governing-declaration-53aa8db88cfaeec0}{}
\subsection*{\small\ttfamily\raggedright Applied\allowbreak{}Modeling\allowbreak{}Lib.\allowbreak{}Probability.\allowbreak{}Poisson\allowbreak{}Process.\allowbreak{}suspension\allowbreak{}Measure}
\textbf{Declaration location:} reusable library
\paragraph{Lean declaration body}
\begin{ReviewVerbatim}
type:
ℝ → MeasureTheory.Measure ((ℤ → ℝ) × ℝ)

value:
fun (rate : ℝ) =>
  ENNReal.ofReal rate •
    ((AppliedModelingLib.Probability.PoissonProcess.twoSidedInterarrivalMeasure rate).prod
          MeasureTheory.volume).restrict
      AppliedModelingLib.Probability.PoissonProcess.suspensionCarrier
\end{ReviewVerbatim}

\paragraph{Direct retained dependencies}
\begin{itemize}\raggedright
\item \hyperlink{governing-declaration-f492363c01e8969e}{Applied\allowbreak{}Modeling\allowbreak{}Lib.\allowbreak{}Probability.\allowbreak{}Poisson\allowbreak{}Process.\allowbreak{}suspension\allowbreak{}Carrier}
\item \hyperlink{governing-declaration-a5c5dbd6c5850d8d}{Applied\allowbreak{}Modeling\allowbreak{}Lib.\allowbreak{}Probability.\allowbreak{}Poisson\allowbreak{}Process.\allowbreak{}two\allowbreak{}Sided\allowbreak{}Interarrival\allowbreak{}Measure}
\end{itemize}
\hypertarget{governing-declaration-ad65cc28d9707a99}{}
\subsection*{\small\ttfamily\raggedright Applied\allowbreak{}Modeling\allowbreak{}Lib.\allowbreak{}Probability.\allowbreak{}Poisson\allowbreak{}Process.\allowbreak{}good\allowbreak{}Suspension\allowbreak{}Measure}
\textbf{Declaration location:} reusable library
\paragraph{Lean declaration body}
\begin{ReviewVerbatim}
type:
ℝ → MeasureTheory.Measure AppliedModelingLib.Probability.PoissonProcess.GoodSuspensionState

value:
fun (rate : ℝ) =>
  MeasureTheory.Measure.comap Subtype.val (AppliedModelingLib.Probability.PoissonProcess.suspensionMeasure rate)
\end{ReviewVerbatim}

\paragraph{Direct retained dependencies}
\begin{itemize}\raggedright
\item \hyperlink{governing-declaration-ab1cfe37ad8c32b0}{Applied\allowbreak{}Modeling\allowbreak{}Lib.\allowbreak{}Probability.\allowbreak{}Poisson\allowbreak{}Process.\allowbreak{}Good\allowbreak{}Suspension\allowbreak{}State}
\item \hyperlink{governing-declaration-a8385d799d0a4d16}{Applied\allowbreak{}Modeling\allowbreak{}Lib.\allowbreak{}Probability.\allowbreak{}Poisson\allowbreak{}Process.\allowbreak{}inst\allowbreak{}Measurable\allowbreak{}Space\allowbreak{}Good\allowbreak{}Suspension\allowbreak{}State}
\item \hyperlink{governing-declaration-f492363c01e8969e}{Applied\allowbreak{}Modeling\allowbreak{}Lib.\allowbreak{}Probability.\allowbreak{}Poisson\allowbreak{}Process.\allowbreak{}suspension\allowbreak{}Carrier}
\item \hyperlink{governing-declaration-446bfa9c60d40523}{Applied\allowbreak{}Modeling\allowbreak{}Lib.\allowbreak{}Probability.\allowbreak{}Poisson\allowbreak{}Process.\allowbreak{}suspension\allowbreak{}Good\allowbreak{}Gap\allowbreak{}Path}
\item \hyperlink{governing-declaration-53aa8db88cfaeec0}{Applied\allowbreak{}Modeling\allowbreak{}Lib.\allowbreak{}Probability.\allowbreak{}Poisson\allowbreak{}Process.\allowbreak{}suspension\allowbreak{}Measure}
\end{itemize}
\hypertarget{governing-declaration-a69bedb5b0d420c6}{}
\subsection*{\small\ttfamily\raggedright Applied\allowbreak{}Modeling\allowbreak{}Lib.\allowbreak{}Probability.\allowbreak{}Queueing.\allowbreak{}Tagged\allowbreak{}Arrival\allowbreak{}At\allowbreak{}Zero}
\textbf{Declaration location:} reusable library
\paragraph{Lean declaration body}
\begin{ReviewVerbatim}
field Ptag:
{Ω : Type u_1} →
  [inst : MeasurableSpace Ω] → AppliedModelingLib.Probability.Queueing.TaggedArrivalAtZero Ω → MeasureTheory.Measure Ω

field isProbability:
∀ {Ω : Type u_1} [inst : MeasurableSpace Ω] (self : AppliedModelingLib.Probability.Queueing.TaggedArrivalAtZero Ω),
  MeasureTheory.IsProbabilityMeasure self.Ptag

field arrivals:
{Ω : Type u_1} → [inst : MeasurableSpace Ω] → AppliedModelingLib.Probability.Queueing.TaggedArrivalAtZero Ω → Ω → ℤ → ℝ

field tag_at_zero:
∀ {Ω : Type u_1} [inst : MeasurableSpace Ω] (self : AppliedModelingLib.Probability.Queueing.TaggedArrivalAtZero Ω),
  ∀ᵐ (ω : Ω) ∂self.Ptag, self.arrivals ω 0 = 0

field arrivals_strict:
∀ {Ω : Type u_1} [inst : MeasurableSpace Ω] (self : AppliedModelingLib.Probability.Queueing.TaggedArrivalAtZero Ω),
  ∀ᵐ (ω : Ω) ∂self.Ptag, StrictMono (self.arrivals ω)
\end{ReviewVerbatim}


\hypertarget{governing-declaration-5a10835f452c324b}{}
\subsection*{\small\ttfamily\raggedright Applied\allowbreak{}Modeling\allowbreak{}Lib.\allowbreak{}Probability.\allowbreak{}Palm.\allowbreak{}target\allowbreak{}Passive\allowbreak{}Tagged\allowbreak{}Arrival\allowbreak{}At\allowbreak{}Zero}
\textbf{Declaration location:} reusable library
\paragraph{Lean declaration body}
\begin{ReviewVerbatim}
type:
{Ωtag : Type u_1} →
  {Γ : Type u_2} →
    [inst : MeasurableSpace Ωtag] →
      [inst_1 : MeasurableSpace Γ] →
        AppliedModelingLib.Probability.Queueing.TaggedArrivalAtZero Ωtag →
          AppliedModelingLib.Probability.Palm.ShiftInvariantProbabilityLaw Γ →
            AppliedModelingLib.Probability.Queueing.TaggedArrivalAtZero (Ωtag × Γ)

value:
fun {Ωtag : Type u_1} {Γ : Type u_2} [MeasurableSpace Ωtag] [MeasurableSpace Γ]
    (tagged : AppliedModelingLib.Probability.Queueing.TaggedArrivalAtZero Ωtag)
    (passive : AppliedModelingLib.Probability.Palm.ShiftInvariantProbabilityLaw Γ) =>
  { Ptag := tagged.Ptag.prod passive.Pbase,
    isProbability := AppliedModelingLib.Probability.Palm.targetPassiveTaggedArrivalAtZero._proof_1 tagged passive,
    arrivals := fun (z : Ωtag × Γ) => tagged.arrivals z.1,
    tag_at_zero := AppliedModelingLib.Probability.Palm.targetPassiveTaggedArrivalAtZero._proof_3 tagged passive,
    arrivals_strict := AppliedModelingLib.Probability.Palm.targetPassiveTaggedArrivalAtZero._proof_4 tagged passive }
\end{ReviewVerbatim}

\paragraph{Direct retained dependencies}
\begin{itemize}\raggedright
\item \hyperlink{governing-declaration-fb5ec66edd06d906}{Applied\allowbreak{}Modeling\allowbreak{}Lib.\allowbreak{}Probability.\allowbreak{}Palm.\allowbreak{}Shift\allowbreak{}Invariant\allowbreak{}Probability\allowbreak{}Law}
\item \hyperlink{governing-declaration-a69bedb5b0d420c6}{Applied\allowbreak{}Modeling\allowbreak{}Lib.\allowbreak{}Probability.\allowbreak{}Queueing.\allowbreak{}Tagged\allowbreak{}Arrival\allowbreak{}At\allowbreak{}Zero}
\end{itemize}
\hypertarget{governing-declaration-b7f1849fb6a62341}{}
\subsection*{\small\ttfamily\raggedright Applied\allowbreak{}Modeling\allowbreak{}Lib.\allowbreak{}Probability.\allowbreak{}Poisson\allowbreak{}Process.\allowbreak{}candidate\allowbreak{}Tagged\allowbreak{}Arrival\allowbreak{}At\allowbreak{}Zero}
\textbf{Declaration location:} reusable library
\paragraph{Lean declaration body}
\begin{ReviewVerbatim}
type:
(rate : ℝ) → 0 < rate → AppliedModelingLib.Probability.Queueing.TaggedArrivalAtZero (ℤ → ℝ)

value:
fun (rate : ℝ) (hrate : 0 < rate) =>
  { Ptag := AppliedModelingLib.Probability.PoissonProcess.twoSidedInterarrivalMeasure rate,
    isProbability :=
      AppliedModelingLib.Probability.PoissonProcess.isProbabilityMeasure_twoSidedInterarrivalMeasure hrate,
    arrivals := AppliedModelingLib.Probability.PoissonProcess.candidatePalmArrival,
    tag_at_zero := AppliedModelingLib.Probability.PoissonProcess.ae_candidatePalmArrival_tag_at_zero,
    arrivals_strict := AppliedModelingLib.Probability.PoissonProcess.ae_candidatePalmArrival_strictMono hrate }
\end{ReviewVerbatim}

\paragraph{Direct retained dependencies}
\begin{itemize}\raggedright
\item \hyperlink{governing-declaration-9208635caa21f800}{Applied\allowbreak{}Modeling\allowbreak{}Lib.\allowbreak{}Probability.\allowbreak{}Poisson\allowbreak{}Process.\allowbreak{}candidate\allowbreak{}Palm\allowbreak{}Arrival}
\item \hyperlink{governing-declaration-a5c5dbd6c5850d8d}{Applied\allowbreak{}Modeling\allowbreak{}Lib.\allowbreak{}Probability.\allowbreak{}Poisson\allowbreak{}Process.\allowbreak{}two\allowbreak{}Sided\allowbreak{}Interarrival\allowbreak{}Measure}
\item \hyperlink{governing-declaration-a69bedb5b0d420c6}{Applied\allowbreak{}Modeling\allowbreak{}Lib.\allowbreak{}Probability.\allowbreak{}Queueing.\allowbreak{}Tagged\allowbreak{}Arrival\allowbreak{}At\allowbreak{}Zero}
\end{itemize}
\hypertarget{governing-declaration-837dedc24e6342be}{}
\subsection*{\small\ttfamily\raggedright Applied\allowbreak{}Modeling\allowbreak{}Lib.\allowbreak{}Probability.\allowbreak{}Queueing.\allowbreak{}int\allowbreak{}Path\allowbreak{}Shift}
\textbf{Declaration location:} reusable library
\paragraph{Lean declaration body}
\begin{ReviewVerbatim}
type:
{α : Type u_1} → ℤ → (ℤ → α) → ℤ → α

value:
fun {α : Type u_1} (k : ℤ) (x : ℤ → α) (a : ℤ) => x (k + a)
\end{ReviewVerbatim}


\hypertarget{governing-declaration-34375b94749a398e}{}
\subsection*{\small\ttfamily\raggedright Applied\allowbreak{}Modeling\allowbreak{}Lib.\allowbreak{}Probability.\allowbreak{}Queueing.\allowbreak{}stationary\allowbreak{}Poisson\allowbreak{}Work\allowbreak{}Requirement}
\textbf{Declaration location:} reusable library
\paragraph{Lean declaration body}
\begin{ReviewVerbatim}
type:
AppliedModelingLib.Probability.PoissonProcess.GoodSuspensionState × (ℤ → ℝ) → ℤ → ℝ

value:
fun (z : AppliedModelingLib.Probability.PoissonProcess.GoodSuspensionState × (ℤ → ℝ)) (i : ℤ) => z.2 i
\end{ReviewVerbatim}

\paragraph{Direct retained dependencies}
\begin{itemize}\raggedright
\item \hyperlink{governing-declaration-ab1cfe37ad8c32b0}{Applied\allowbreak{}Modeling\allowbreak{}Lib.\allowbreak{}Probability.\allowbreak{}Poisson\allowbreak{}Process.\allowbreak{}Good\allowbreak{}Suspension\allowbreak{}State}
\end{itemize}
\hypertarget{governing-declaration-13879bfc74e6d5d0}{}
\subsection*{\small\ttfamily\raggedright Applied\allowbreak{}Modeling\allowbreak{}Lib.\allowbreak{}Probability.\allowbreak{}Poisson\allowbreak{}Process.\allowbreak{}multiclass\allowbreak{}Stationary\allowbreak{}Poisson\allowbreak{}Work\allowbreak{}Requirement}
\textbf{Declaration location:} reusable library
\paragraph{Lean declaration body}
\begin{ReviewVerbatim}
type:
{Class : Type u_1} →
  Class → (Class → AppliedModelingLib.Probability.PoissonProcess.GoodSuspensionState × (ℤ → ℝ)) → ℤ → ℝ

value:
fun {Class : Type u_1} (i : Class)
    (a : Class → AppliedModelingLib.Probability.PoissonProcess.GoodSuspensionState × (ℤ → ℝ)) (a_1 : ℤ) =>
  AppliedModelingLib.Probability.Queueing.stationaryPoissonWorkRequirement (a i) a_1
\end{ReviewVerbatim}

\paragraph{Direct retained dependencies}
\begin{itemize}\raggedright
\item \hyperlink{governing-declaration-ab1cfe37ad8c32b0}{Applied\allowbreak{}Modeling\allowbreak{}Lib.\allowbreak{}Probability.\allowbreak{}Poisson\allowbreak{}Process.\allowbreak{}Good\allowbreak{}Suspension\allowbreak{}State}
\item \hyperlink{governing-declaration-34375b94749a398e}{Applied\allowbreak{}Modeling\allowbreak{}Lib.\allowbreak{}Probability.\allowbreak{}Queueing.\allowbreak{}stationary\allowbreak{}Poisson\allowbreak{}Work\allowbreak{}Requirement}
\end{itemize}
\hypertarget{governing-declaration-4c9cdfa96671089f}{}
\subsection*{\small\ttfamily\raggedright Applied\allowbreak{}Modeling\allowbreak{}Lib.\allowbreak{}Probability.\allowbreak{}Queueing.\allowbreak{}timed\allowbreak{}Embedded\allowbreak{}Arrival}
\textbf{Declaration location:} reusable library
\paragraph{Lean declaration body}
\begin{ReviewVerbatim}
type:
{α : Type u_1} → AppliedModelingLib.Probability.PoissonProcess.GoodSuspensionState × (ℤ → α) → ℤ → ℝ

value:
fun {α : Type u_1} (z : AppliedModelingLib.Probability.PoissonProcess.GoodSuspensionState × (ℤ → α)) (i : ℤ) =>
  AppliedModelingLib.Probability.PoissonProcess.suspensionBaseArrival z.1 i
\end{ReviewVerbatim}

\paragraph{Direct retained dependencies}
\begin{itemize}\raggedright
\item \hyperlink{governing-declaration-ab1cfe37ad8c32b0}{Applied\allowbreak{}Modeling\allowbreak{}Lib.\allowbreak{}Probability.\allowbreak{}Poisson\allowbreak{}Process.\allowbreak{}Good\allowbreak{}Suspension\allowbreak{}State}
\item \hyperlink{governing-declaration-c513c48a384a886a}{Applied\allowbreak{}Modeling\allowbreak{}Lib.\allowbreak{}Probability.\allowbreak{}Poisson\allowbreak{}Process.\allowbreak{}suspension\allowbreak{}Base\allowbreak{}Arrival}
\end{itemize}
\hypertarget{governing-declaration-499050173dfb71ee}{}
\subsection*{\small\ttfamily\raggedright Applied\allowbreak{}Modeling\allowbreak{}Lib.\allowbreak{}Probability.\allowbreak{}Queueing.\allowbreak{}stationary\allowbreak{}Poisson\allowbreak{}Work\allowbreak{}Arrival}
\textbf{Declaration location:} reusable library
\paragraph{Lean declaration body}
\begin{ReviewVerbatim}
type:
AppliedModelingLib.Probability.PoissonProcess.GoodSuspensionState × (ℤ → ℝ) → ℤ → ℝ

value:
fun (z : AppliedModelingLib.Probability.PoissonProcess.GoodSuspensionState × (ℤ → ℝ)) (i : ℤ) =>
  AppliedModelingLib.Probability.Queueing.timedEmbeddedArrival z i
\end{ReviewVerbatim}

\paragraph{Direct retained dependencies}
\begin{itemize}\raggedright
\item \hyperlink{governing-declaration-ab1cfe37ad8c32b0}{Applied\allowbreak{}Modeling\allowbreak{}Lib.\allowbreak{}Probability.\allowbreak{}Poisson\allowbreak{}Process.\allowbreak{}Good\allowbreak{}Suspension\allowbreak{}State}
\item \hyperlink{governing-declaration-4c9cdfa96671089f}{Applied\allowbreak{}Modeling\allowbreak{}Lib.\allowbreak{}Probability.\allowbreak{}Queueing.\allowbreak{}timed\allowbreak{}Embedded\allowbreak{}Arrival}
\end{itemize}
\hypertarget{governing-declaration-04aec0e04fc92779}{}
\subsection*{\small\ttfamily\raggedright Applied\allowbreak{}Modeling\allowbreak{}Lib.\allowbreak{}Probability.\allowbreak{}Queueing.\allowbreak{}timed\allowbreak{}Embedded\allowbreak{}Suspension\allowbreak{}Flow}
\textbf{Declaration location:} reusable library
\paragraph{Lean declaration body}
\begin{ReviewVerbatim}
type:
{α : Type u_1} →
  ℝ →
    AppliedModelingLib.Probability.PoissonProcess.GoodSuspensionState × (ℤ → α) →
      AppliedModelingLib.Probability.PoissonProcess.GoodSuspensionState × (ℤ → α)

value:
fun {α : Type u_1} (t : ℝ) (a : AppliedModelingLib.Probability.PoissonProcess.GoodSuspensionState × (ℤ → α)) =>
  (AppliedModelingLib.Probability.PoissonProcess.goodSuspensionFlow t a.1,
    AppliedModelingLib.Probability.Queueing.intPathShift
      (AppliedModelingLib.Probability.PoissonProcess.suspensionCrossingIndexPastClosed t ↑a.1) a.2)
\end{ReviewVerbatim}

\paragraph{Direct retained dependencies}
\begin{itemize}\raggedright
\item \hyperlink{governing-declaration-ab1cfe37ad8c32b0}{Applied\allowbreak{}Modeling\allowbreak{}Lib.\allowbreak{}Probability.\allowbreak{}Poisson\allowbreak{}Process.\allowbreak{}Good\allowbreak{}Suspension\allowbreak{}State}
\item \hyperlink{governing-declaration-1991889bac68724c}{Applied\allowbreak{}Modeling\allowbreak{}Lib.\allowbreak{}Probability.\allowbreak{}Poisson\allowbreak{}Process.\allowbreak{}good\allowbreak{}Suspension\allowbreak{}Flow}
\item \hyperlink{governing-declaration-f492363c01e8969e}{Applied\allowbreak{}Modeling\allowbreak{}Lib.\allowbreak{}Probability.\allowbreak{}Poisson\allowbreak{}Process.\allowbreak{}suspension\allowbreak{}Carrier}
\item \hyperlink{governing-declaration-79e7fd2195dfecc3}{Applied\allowbreak{}Modeling\allowbreak{}Lib.\allowbreak{}Probability.\allowbreak{}Poisson\allowbreak{}Process.\allowbreak{}suspension\allowbreak{}Crossing\allowbreak{}Index\allowbreak{}Past\allowbreak{}Closed}
\item \hyperlink{governing-declaration-446bfa9c60d40523}{Applied\allowbreak{}Modeling\allowbreak{}Lib.\allowbreak{}Probability.\allowbreak{}Poisson\allowbreak{}Process.\allowbreak{}suspension\allowbreak{}Good\allowbreak{}Gap\allowbreak{}Path}
\item \hyperlink{governing-declaration-837dedc24e6342be}{Applied\allowbreak{}Modeling\allowbreak{}Lib.\allowbreak{}Probability.\allowbreak{}Queueing.\allowbreak{}int\allowbreak{}Path\allowbreak{}Shift}
\end{itemize}
\hypertarget{governing-declaration-174520ae711e6558}{}
\subsection*{\small\ttfamily\raggedright Applied\allowbreak{}Modeling\allowbreak{}Lib.\allowbreak{}Probability.\allowbreak{}Poisson\allowbreak{}Process.\allowbreak{}multiclass\allowbreak{}Stationary\allowbreak{}Poisson\allowbreak{}Work\allowbreak{}Flow}
\textbf{Declaration location:} reusable library
\paragraph{Lean declaration body}
\begin{ReviewVerbatim}
type:
{Class : Type u_1} →
  ℝ →
    (Class → AppliedModelingLib.Probability.PoissonProcess.GoodSuspensionState × (ℤ → ℝ)) →
      Class → AppliedModelingLib.Probability.PoissonProcess.GoodSuspensionState × (ℤ → ℝ)

value:
fun {Class : Type u_1} (t : ℝ) (a : Class → AppliedModelingLib.Probability.PoissonProcess.GoodSuspensionState × (ℤ → ℝ))
    (a_1 : Class) =>
  AppliedModelingLib.Probability.Queueing.timedEmbeddedSuspensionFlow t (a a_1)
\end{ReviewVerbatim}

\paragraph{Direct retained dependencies}
\begin{itemize}\raggedright
\item \hyperlink{governing-declaration-ab1cfe37ad8c32b0}{Applied\allowbreak{}Modeling\allowbreak{}Lib.\allowbreak{}Probability.\allowbreak{}Poisson\allowbreak{}Process.\allowbreak{}Good\allowbreak{}Suspension\allowbreak{}State}
\item \hyperlink{governing-declaration-04aec0e04fc92779}{Applied\allowbreak{}Modeling\allowbreak{}Lib.\allowbreak{}Probability.\allowbreak{}Queueing.\allowbreak{}timed\allowbreak{}Embedded\allowbreak{}Suspension\allowbreak{}Flow}
\end{itemize}
\hypertarget{governing-declaration-fd22338cf629f5f6}{}
\subsection*{\small\ttfamily\raggedright Applied\allowbreak{}Modeling\allowbreak{}Lib.\allowbreak{}Probability.\allowbreak{}Queueing.\allowbreak{}timed\allowbreak{}Embedded\allowbreak{}Suspension\allowbreak{}Product\allowbreak{}Measure}
\textbf{Declaration location:} reusable library
\paragraph{Lean declaration body}
\begin{ReviewVerbatim}
type:
{α : Type u_1} →
  [inst : MeasurableSpace α] →
    ℝ →
      MeasureTheory.Measure (ℤ → α) →
        MeasureTheory.Measure (AppliedModelingLib.Probability.PoissonProcess.GoodSuspensionState × (ℤ → α))

value:
fun {α : Type u_1} [MeasurableSpace α] (rate : ℝ) (P : MeasureTheory.Measure (ℤ → α)) =>
  (AppliedModelingLib.Probability.PoissonProcess.goodSuspensionMeasure rate).prod P
\end{ReviewVerbatim}

\paragraph{Direct retained dependencies}
\begin{itemize}\raggedright
\item \hyperlink{governing-declaration-ab1cfe37ad8c32b0}{Applied\allowbreak{}Modeling\allowbreak{}Lib.\allowbreak{}Probability.\allowbreak{}Poisson\allowbreak{}Process.\allowbreak{}Good\allowbreak{}Suspension\allowbreak{}State}
\item \hyperlink{governing-declaration-ad65cc28d9707a99}{Applied\allowbreak{}Modeling\allowbreak{}Lib.\allowbreak{}Probability.\allowbreak{}Poisson\allowbreak{}Process.\allowbreak{}good\allowbreak{}Suspension\allowbreak{}Measure}
\item \hyperlink{governing-declaration-a8385d799d0a4d16}{Applied\allowbreak{}Modeling\allowbreak{}Lib.\allowbreak{}Probability.\allowbreak{}Poisson\allowbreak{}Process.\allowbreak{}inst\allowbreak{}Measurable\allowbreak{}Space\allowbreak{}Good\allowbreak{}Suspension\allowbreak{}State}
\end{itemize}
\hypertarget{governing-declaration-10a259bddd1a9d6f}{}
\subsection*{\small\ttfamily\raggedright Applied\allowbreak{}Modeling\allowbreak{}Lib.\allowbreak{}Probability.\allowbreak{}Queueing.\allowbreak{}stationary\allowbreak{}Poisson\allowbreak{}Work\allowbreak{}Measure}
\textbf{Declaration location:} reusable library
\paragraph{Lean declaration body}
\begin{ReviewVerbatim}
type:
ℝ → MeasureTheory.Measure (AppliedModelingLib.Probability.PoissonProcess.GoodSuspensionState × (ℤ → ℝ))

value:
fun (arrivalRate : ℝ) =>
  AppliedModelingLib.Probability.Queueing.timedEmbeddedSuspensionProductMeasure arrivalRate
    (AppliedModelingLib.Probability.PoissonProcess.twoSidedInterarrivalMeasure 1)
\end{ReviewVerbatim}

\paragraph{Direct retained dependencies}
\begin{itemize}\raggedright
\item \hyperlink{governing-declaration-ab1cfe37ad8c32b0}{Applied\allowbreak{}Modeling\allowbreak{}Lib.\allowbreak{}Probability.\allowbreak{}Poisson\allowbreak{}Process.\allowbreak{}Good\allowbreak{}Suspension\allowbreak{}State}
\item \hyperlink{governing-declaration-a8385d799d0a4d16}{Applied\allowbreak{}Modeling\allowbreak{}Lib.\allowbreak{}Probability.\allowbreak{}Poisson\allowbreak{}Process.\allowbreak{}inst\allowbreak{}Measurable\allowbreak{}Space\allowbreak{}Good\allowbreak{}Suspension\allowbreak{}State}
\item \hyperlink{governing-declaration-a5c5dbd6c5850d8d}{Applied\allowbreak{}Modeling\allowbreak{}Lib.\allowbreak{}Probability.\allowbreak{}Poisson\allowbreak{}Process.\allowbreak{}two\allowbreak{}Sided\allowbreak{}Interarrival\allowbreak{}Measure}
\item \hyperlink{governing-declaration-fd22338cf629f5f6}{Applied\allowbreak{}Modeling\allowbreak{}Lib.\allowbreak{}Probability.\allowbreak{}Queueing.\allowbreak{}timed\allowbreak{}Embedded\allowbreak{}Suspension\allowbreak{}Product\allowbreak{}Measure}
\end{itemize}
\hypertarget{governing-declaration-0a1733e4d91e5dfa}{}
\subsection*{\small\ttfamily\raggedright Applied\allowbreak{}Modeling\allowbreak{}Lib.\allowbreak{}Probability.\allowbreak{}Poisson\allowbreak{}Process.\allowbreak{}multiclass\allowbreak{}Stationary\allowbreak{}Poisson\allowbreak{}Work\allowbreak{}Shift\allowbreak{}Invariant\allowbreak{}Law}
\textbf{Declaration location:} reusable library
\paragraph{Lean declaration body}
\begin{ReviewVerbatim}
type:
{Class : Type u_1} →
  [Fintype Class] →
    (arrivalRate : Class → ℝ) →
      (∀ (i : Class), 0 < arrivalRate i) →
        AppliedModelingLib.Probability.Palm.ShiftInvariantProbabilityLaw
          (Class → AppliedModelingLib.Probability.PoissonProcess.GoodSuspensionState × (ℤ → ℝ))

value:
fun {Class : Type u_1} [Fintype Class] (arrivalRate : Class → ℝ) (harrivalRate : ∀ (i : Class), 0 < arrivalRate i) =>
  { Pbase := AppliedModelingLib.Probability.PoissonProcess.multiclassStationaryPoissonWorkMeasure arrivalRate,
    isProbability :=
      AppliedModelingLib.Probability.PoissonProcess.isProbabilityMeasure_multiclassStationaryPoissonWorkMeasure
        arrivalRate harrivalRate,
    shift := AppliedModelingLib.Probability.PoissonProcess.multiclassStationaryPoissonWorkFlow,
    shift_zero :=
      AppliedModelingLib.Probability.PoissonProcess.multiclassStationaryPoissonWorkShiftInvariantLaw._proof_1,
    shift_add :=
      AppliedModelingLib.Probability.PoissonProcess.multiclassStationaryPoissonWorkShiftInvariantLaw._proof_2,
    shift_preserving := fun (t : ℝ) =>
      AppliedModelingLib.Probability.PoissonProcess.multiclassStationaryPoissonWorkFlow_measurePreserving arrivalRate
        harrivalRate t }
\end{ReviewVerbatim}

\paragraph{Direct retained dependencies}
\begin{itemize}\raggedright
\item \hyperlink{governing-declaration-fb5ec66edd06d906}{Applied\allowbreak{}Modeling\allowbreak{}Lib.\allowbreak{}Probability.\allowbreak{}Palm.\allowbreak{}Shift\allowbreak{}Invariant\allowbreak{}Probability\allowbreak{}Law}
\item \hyperlink{governing-declaration-ab1cfe37ad8c32b0}{Applied\allowbreak{}Modeling\allowbreak{}Lib.\allowbreak{}Probability.\allowbreak{}Poisson\allowbreak{}Process.\allowbreak{}Good\allowbreak{}Suspension\allowbreak{}State}
\item \hyperlink{governing-declaration-a8385d799d0a4d16}{Applied\allowbreak{}Modeling\allowbreak{}Lib.\allowbreak{}Probability.\allowbreak{}Poisson\allowbreak{}Process.\allowbreak{}inst\allowbreak{}Measurable\allowbreak{}Space\allowbreak{}Good\allowbreak{}Suspension\allowbreak{}State}
\item \hyperlink{governing-declaration-174520ae711e6558}{Applied\allowbreak{}Modeling\allowbreak{}Lib.\allowbreak{}Probability.\allowbreak{}Poisson\allowbreak{}Process.\allowbreak{}multiclass\allowbreak{}Stationary\allowbreak{}Poisson\allowbreak{}Work\allowbreak{}Flow}
\item \hyperlink{library-prerequisite-dba97dbb9d8c7c91}{Applied\allowbreak{}Modeling\allowbreak{}Lib.\allowbreak{}Probability.\allowbreak{}Poisson\allowbreak{}Process.\allowbreak{}multiclass\allowbreak{}Stationary\allowbreak{}Poisson\allowbreak{}Work\allowbreak{}Measure}
\end{itemize}
\hypertarget{governing-declaration-9bd1e8f00693c01b}{}
\subsection*{\small\ttfamily\raggedright Applied\allowbreak{}Modeling\allowbreak{}Lib.\allowbreak{}Probability.\allowbreak{}Queueing.\allowbreak{}timed\allowbreak{}Embedded\allowbreak{}Tagged\allowbreak{}Arrival\allowbreak{}At\allowbreak{}Zero}
\textbf{Declaration location:} reusable library
\paragraph{Lean declaration body}
\begin{ReviewVerbatim}
type:
{α : Type u_1} →
  [inst : MeasurableSpace α] →
    (rate : ℝ) →
      0 < rate →
        (PtagPath : MeasureTheory.Measure (ℤ → α)) →
          [MeasureTheory.IsProbabilityMeasure PtagPath] →
            AppliedModelingLib.Probability.Queueing.TaggedArrivalAtZero ((ℤ → ℝ) × (ℤ → α))

value:
fun {α : Type u_1} [MeasurableSpace α] (rate : ℝ) (hrate : 0 < rate) (PtagPath : MeasureTheory.Measure (ℤ → α))
    [MeasureTheory.IsProbabilityMeasure PtagPath] =>
  { Ptag := (AppliedModelingLib.Probability.PoissonProcess.candidateTaggedArrivalAtZero rate hrate).Ptag.prod PtagPath,
    isProbability :=
      AppliedModelingLib.Probability.Queueing.timedEmbeddedTaggedArrivalAtZero._proof_1 rate hrate PtagPath,
    arrivals := fun (z : (ℤ → ℝ) × (ℤ → α)) => AppliedModelingLib.Probability.PoissonProcess.candidatePalmArrival z.1,
    tag_at_zero := AppliedModelingLib.Probability.Queueing.timedEmbeddedTaggedArrivalAtZero._proof_3 rate PtagPath,
    arrivals_strict :=
      AppliedModelingLib.Probability.Queueing.timedEmbeddedTaggedArrivalAtZero._proof_4 rate hrate PtagPath }
\end{ReviewVerbatim}

\paragraph{Direct retained dependencies}
\begin{itemize}\raggedright
\item \hyperlink{governing-declaration-9208635caa21f800}{Applied\allowbreak{}Modeling\allowbreak{}Lib.\allowbreak{}Probability.\allowbreak{}Poisson\allowbreak{}Process.\allowbreak{}candidate\allowbreak{}Palm\allowbreak{}Arrival}
\item \hyperlink{governing-declaration-b7f1849fb6a62341}{Applied\allowbreak{}Modeling\allowbreak{}Lib.\allowbreak{}Probability.\allowbreak{}Poisson\allowbreak{}Process.\allowbreak{}candidate\allowbreak{}Tagged\allowbreak{}Arrival\allowbreak{}At\allowbreak{}Zero}
\item \hyperlink{governing-declaration-a5c5dbd6c5850d8d}{Applied\allowbreak{}Modeling\allowbreak{}Lib.\allowbreak{}Probability.\allowbreak{}Poisson\allowbreak{}Process.\allowbreak{}two\allowbreak{}Sided\allowbreak{}Interarrival\allowbreak{}Measure}
\item \hyperlink{governing-declaration-a69bedb5b0d420c6}{Applied\allowbreak{}Modeling\allowbreak{}Lib.\allowbreak{}Probability.\allowbreak{}Queueing.\allowbreak{}Tagged\allowbreak{}Arrival\allowbreak{}At\allowbreak{}Zero}
\end{itemize}
\hypertarget{governing-declaration-e9116e504d73a45c}{}
\subsection*{\small\ttfamily\raggedright Applied\allowbreak{}Modeling\allowbreak{}Lib.\allowbreak{}Probability.\allowbreak{}Queueing.\allowbreak{}stationary\allowbreak{}Poisson\allowbreak{}Work\allowbreak{}Tagged\allowbreak{}Arrival\allowbreak{}At\allowbreak{}Zero}
\textbf{Declaration location:} reusable library
\paragraph{Lean declaration body}
\begin{ReviewVerbatim}
type:
{arrivalRate : ℝ} → 0 < arrivalRate → AppliedModelingLib.Probability.Queueing.TaggedArrivalAtZero ((ℤ → ℝ) × (ℤ → ℝ))

value:
fun {arrivalRate : ℝ} (harrivalRate : 0 < arrivalRate) =>
  AppliedModelingLib.Probability.Queueing.timedEmbeddedTaggedArrivalAtZero arrivalRate harrivalRate
    (AppliedModelingLib.Probability.PoissonProcess.twoSidedInterarrivalMeasure 1)
\end{ReviewVerbatim}

\paragraph{Direct retained dependencies}
\begin{itemize}\raggedright
\item \hyperlink{governing-declaration-a5c5dbd6c5850d8d}{Applied\allowbreak{}Modeling\allowbreak{}Lib.\allowbreak{}Probability.\allowbreak{}Poisson\allowbreak{}Process.\allowbreak{}two\allowbreak{}Sided\allowbreak{}Interarrival\allowbreak{}Measure}
\item \hyperlink{governing-declaration-a69bedb5b0d420c6}{Applied\allowbreak{}Modeling\allowbreak{}Lib.\allowbreak{}Probability.\allowbreak{}Queueing.\allowbreak{}Tagged\allowbreak{}Arrival\allowbreak{}At\allowbreak{}Zero}
\item \hyperlink{governing-declaration-9bd1e8f00693c01b}{Applied\allowbreak{}Modeling\allowbreak{}Lib.\allowbreak{}Probability.\allowbreak{}Queueing.\allowbreak{}timed\allowbreak{}Embedded\allowbreak{}Tagged\allowbreak{}Arrival\allowbreak{}At\allowbreak{}Zero}
\end{itemize}
\hypertarget{governing-declaration-5d3724f4dc832d77}{}
\subsection*{\small\ttfamily\raggedright Applied\allowbreak{}Modeling\allowbreak{}Lib.\allowbreak{}Queueing.\allowbreak{}Nonpreemptive\allowbreak{}Priority\allowbreak{}Arrival\allowbreak{}Index}
\textbf{Declaration location:} reusable library
\paragraph{Lean declaration body}
\begin{ReviewVerbatim}
type:
ℕ → Type

value:
fun (n : ℕ) => (_ : Fin n) × ℤ
\end{ReviewVerbatim}


\hypertarget{governing-declaration-62a268164b6ae28b}{}
\subsection*{\small\ttfamily\raggedright Applied\allowbreak{}Modeling\allowbreak{}Lib.\allowbreak{}Queueing.\allowbreak{}Stationary\allowbreak{}Poisson\allowbreak{}Work\allowbreak{}Path}
\textbf{Declaration location:} reusable library
\paragraph{Lean declaration body}
\begin{ReviewVerbatim}
type:
Type

value:
AppliedModelingLib.Probability.PoissonProcess.GoodSuspensionState × (ℤ → ℝ)
\end{ReviewVerbatim}

\paragraph{Direct retained dependencies}
\begin{itemize}\raggedright
\item \hyperlink{governing-declaration-ab1cfe37ad8c32b0}{Applied\allowbreak{}Modeling\allowbreak{}Lib.\allowbreak{}Probability.\allowbreak{}Poisson\allowbreak{}Process.\allowbreak{}Good\allowbreak{}Suspension\allowbreak{}State}
\end{itemize}
\hypertarget{governing-declaration-0e2f4b9eb63fa294}{}
\subsection*{\small\ttfamily\raggedright Applied\allowbreak{}Modeling\allowbreak{}Lib.\allowbreak{}Queueing.\allowbreak{}empty\allowbreak{}Nonpreemptive\allowbreak{}Priority\allowbreak{}Work\allowbreak{}State}
\textbf{Declaration location:} reusable library
\paragraph{Lean declaration body}
\begin{ReviewVerbatim}
type:
{n : ℕ} → {JobId : Type u_1} → ℝ → AppliedModelingLib.Queueing.NonpreemptivePriorityWorkState n JobId

value:
fun {n : ℕ} {JobId : Type u_1} (time : ℝ) =>
  { currentTime := time, active := none, waiting := fun (x : Fin n) => [], completed := [] }
\end{ReviewVerbatim}

\paragraph{Direct retained dependencies}
\begin{itemize}\raggedright
\item \hyperlink{library-prerequisite-2373fe691c32d999}{Applied\allowbreak{}Modeling\allowbreak{}Lib.\allowbreak{}Queueing.\allowbreak{}Nonpreemptive\allowbreak{}Priority\allowbreak{}Job}
\item \hyperlink{library-prerequisite-e68cb5fce009ec53}{Applied\allowbreak{}Modeling\allowbreak{}Lib.\allowbreak{}Queueing.\allowbreak{}Nonpreemptive\allowbreak{}Priority\allowbreak{}Work\allowbreak{}State}
\end{itemize}
\hypertarget{governing-declaration-d385a019cc9e3344}{}
\subsection*{\small\ttfamily\raggedright Applied\allowbreak{}Modeling\allowbreak{}Lib.\allowbreak{}Queueing.\allowbreak{}enqueue\allowbreak{}Nonpreemptive\allowbreak{}Priority\allowbreak{}Job}
\textbf{Declaration location:} reusable library
\paragraph{Lean declaration body}
\begin{ReviewVerbatim}
type:
{n : ℕ} →
  {JobId : Type u_1} →
    AppliedModelingLib.Queueing.NonpreemptivePriorityWorkState n JobId →
      AppliedModelingLib.Queueing.NonpreemptivePriorityJob n JobId →
        AppliedModelingLib.Queueing.NonpreemptivePriorityWorkState n JobId

value:
fun {n : ℕ} {JobId : Type u_1} (state : AppliedModelingLib.Queueing.NonpreemptivePriorityWorkState n JobId)
    (job : AppliedModelingLib.Queueing.NonpreemptivePriorityJob n JobId) =>
  { currentTime := state.currentTime, active := state.active,
    waiting := Function.update state.waiting job.priority (state.waiting job.priority ++ [job]),
    completed := state.completed }
\end{ReviewVerbatim}

\paragraph{Direct retained dependencies}
\begin{itemize}\raggedright
\item \hyperlink{library-prerequisite-2373fe691c32d999}{Applied\allowbreak{}Modeling\allowbreak{}Lib.\allowbreak{}Queueing.\allowbreak{}Nonpreemptive\allowbreak{}Priority\allowbreak{}Job}
\item \hyperlink{library-prerequisite-e68cb5fce009ec53}{Applied\allowbreak{}Modeling\allowbreak{}Lib.\allowbreak{}Queueing.\allowbreak{}Nonpreemptive\allowbreak{}Priority\allowbreak{}Work\allowbreak{}State}
\end{itemize}
\hypertarget{governing-declaration-f664e388bb7b2fb8}{}
\subsection*{\small\ttfamily\raggedright Applied\allowbreak{}Modeling\allowbreak{}Lib.\allowbreak{}Queueing.\allowbreak{}finite\allowbreak{}Priority\allowbreak{}Inclusive\allowbreak{}Load}
\textbf{Declaration location:} reusable library
\paragraph{Lean declaration body}
\begin{ReviewVerbatim}
type:
{n : ℕ} → (Fin n → ℝ) → (Fin n → ℝ) → Fin n → ℝ

value:
fun {n : ℕ} (meanService arrivalRate : Fin n → ℝ) (i : Fin n) =>
  ∑ j : Fin n, if j ≤ i then meanService j * arrivalRate j else 0
\end{ReviewVerbatim}


\hypertarget{governing-declaration-f2b9a1af307e7469}{}
\subsection*{\small\ttfamily\raggedright Applied\allowbreak{}Modeling\allowbreak{}Lib.\allowbreak{}Queueing.\allowbreak{}finite\allowbreak{}Priority\allowbreak{}Inclusive\allowbreak{}Slack}
\textbf{Declaration location:} reusable library
\paragraph{Lean declaration body}
\begin{ReviewVerbatim}
type:
{n : ℕ} → (Fin n → ℝ) → (Fin n → ℝ) → Fin n → ℝ

value:
fun {n : ℕ} (meanService arrivalRate : Fin n → ℝ) (i : Fin n) =>
  1 - AppliedModelingLib.Queueing.finitePriorityInclusiveLoad meanService arrivalRate i
\end{ReviewVerbatim}

\paragraph{Direct retained dependencies}
\begin{itemize}\raggedright
\item \hyperlink{governing-declaration-f664e388bb7b2fb8}{Applied\allowbreak{}Modeling\allowbreak{}Lib.\allowbreak{}Queueing.\allowbreak{}finite\allowbreak{}Priority\allowbreak{}Inclusive\allowbreak{}Load}
\end{itemize}
\hypertarget{governing-declaration-80f4f37f124e636e}{}
\subsection*{\small\ttfamily\raggedright Applied\allowbreak{}Modeling\allowbreak{}Lib.\allowbreak{}Queueing.\allowbreak{}finite\allowbreak{}Priority\allowbreak{}Residual\allowbreak{}Work}
\textbf{Declaration location:} reusable library
\paragraph{Lean declaration body}
\begin{ReviewVerbatim}
type:
{n : ℕ} → (Fin n → ℝ) → (Fin n → ℝ) → ℝ

value:
fun {n : ℕ} (arrivalRate meanService : Fin n → ℝ) => ∑ j : Fin n, arrivalRate j * meanService j ^ 2
\end{ReviewVerbatim}


\hypertarget{governing-declaration-b9bd4f553bbc8c90}{}
\subsection*{\small\ttfamily\raggedright Applied\allowbreak{}Modeling\allowbreak{}Lib.\allowbreak{}Queueing.\allowbreak{}finite\allowbreak{}Priority\allowbreak{}Strict\allowbreak{}Load}
\textbf{Declaration location:} reusable library
\paragraph{Lean declaration body}
\begin{ReviewVerbatim}
type:
{n : ℕ} → (Fin n → ℝ) → (Fin n → ℝ) → Fin n → ℝ

value:
fun {n : ℕ} (meanService arrivalRate : Fin n → ℝ) (i : Fin n) =>
  ∑ j : Fin n, if j < i then meanService j * arrivalRate j else 0
\end{ReviewVerbatim}


\hypertarget{governing-declaration-3f5cf610a8190094}{}
\subsection*{\small\ttfamily\raggedright Applied\allowbreak{}Modeling\allowbreak{}Lib.\allowbreak{}Queueing.\allowbreak{}finite\allowbreak{}Priority\allowbreak{}Strict\allowbreak{}Slack}
\textbf{Declaration location:} reusable library
\paragraph{Lean declaration body}
\begin{ReviewVerbatim}
type:
{n : ℕ} → (Fin n → ℝ) → (Fin n → ℝ) → Fin n → ℝ

value:
fun {n : ℕ} (meanService arrivalRate : Fin n → ℝ) (i : Fin n) =>
  1 - AppliedModelingLib.Queueing.finitePriorityStrictLoad meanService arrivalRate i
\end{ReviewVerbatim}

\paragraph{Direct retained dependencies}
\begin{itemize}\raggedright
\item \hyperlink{governing-declaration-b9bd4f553bbc8c90}{Applied\allowbreak{}Modeling\allowbreak{}Lib.\allowbreak{}Queueing.\allowbreak{}finite\allowbreak{}Priority\allowbreak{}Strict\allowbreak{}Load}
\end{itemize}
\hypertarget{governing-declaration-ab8863aef2e4b99f}{}
\subsection*{\small\ttfamily\raggedright Applied\allowbreak{}Modeling\allowbreak{}Lib.\allowbreak{}Queueing.\allowbreak{}finite\allowbreak{}Nonpreemptive\allowbreak{}Priority\allowbreak{}Sojourn\allowbreak{}Time}
\textbf{Declaration location:} reusable library
\paragraph{Lean declaration body}
\begin{ReviewVerbatim}
type:
{n : ℕ} → (Fin n → ℝ) → (Fin n → ℝ) → Fin n → ℝ

value:
fun {n : ℕ} (arrivalRate meanService : Fin n → ℝ) (i : Fin n) =>
  AppliedModelingLib.Queueing.finiteNonpreemptivePriorityQueueWait arrivalRate meanService i + meanService i
\end{ReviewVerbatim}

\paragraph{Direct retained dependencies}
\begin{itemize}\raggedright
\item \hyperlink{library-prerequisite-795b7af1c661bc8d}{Applied\allowbreak{}Modeling\allowbreak{}Lib.\allowbreak{}Queueing.\allowbreak{}finite\allowbreak{}Nonpreemptive\allowbreak{}Priority\allowbreak{}Queue\allowbreak{}Wait}
\end{itemize}
\hypertarget{governing-declaration-adcd68c3e7c5c1bb}{}
\subsection*{\small\ttfamily\raggedright Applied\allowbreak{}Modeling\allowbreak{}Lib.\allowbreak{}Queueing.\allowbreak{}has\allowbreak{}Priority\allowbreak{}Waiting\allowbreak{}Job}
\textbf{Declaration location:} reusable library
\paragraph{Lean declaration body}
\begin{ReviewVerbatim}
type:
{n : ℕ} → {JobId : Type u_1} → AppliedModelingLib.Queueing.NonpreemptivePriorityWorkState n JobId → Prop

value:
fun {n : ℕ} {JobId : Type u_1} (state : AppliedModelingLib.Queueing.NonpreemptivePriorityWorkState n JobId) =>
  ∃ (i : Fin n), 0 < (state.waiting i).length
\end{ReviewVerbatim}

\paragraph{Direct retained dependencies}
\begin{itemize}\raggedright
\item \hyperlink{library-prerequisite-2373fe691c32d999}{Applied\allowbreak{}Modeling\allowbreak{}Lib.\allowbreak{}Queueing.\allowbreak{}Nonpreemptive\allowbreak{}Priority\allowbreak{}Job}
\item \hyperlink{library-prerequisite-e68cb5fce009ec53}{Applied\allowbreak{}Modeling\allowbreak{}Lib.\allowbreak{}Queueing.\allowbreak{}Nonpreemptive\allowbreak{}Priority\allowbreak{}Work\allowbreak{}State}
\end{itemize}
\hypertarget{governing-declaration-de954530ca750d3f}{}
\subsection*{\small\ttfamily\raggedright Applied\allowbreak{}Modeling\allowbreak{}Lib.\allowbreak{}Queueing.\allowbreak{}inst\allowbreak{}Decidable\allowbreak{}Eq\_\allowbreak{}applied\allowbreak{}Modeling\allowbreak{}Lib}
\textbf{Declaration location:} reusable library
\paragraph{Lean declaration body}
\begin{ReviewVerbatim}
type:
{Class : Type u_1} → (a b : Class) → Decidable (a = b)

value:
fun {Class : Type u_1} (a b : Class) => Classical.decEq Class a b
\end{ReviewVerbatim}


\hypertarget{governing-declaration-8f366849032baf38}{}
\subsection*{\small\ttfamily\raggedright Applied\allowbreak{}Modeling\allowbreak{}Lib.\allowbreak{}Queueing.\allowbreak{}multiclass\allowbreak{}Stationary\allowbreak{}Poisson\allowbreak{}Work\allowbreak{}Class\allowbreak{}Tagged\allowbreak{}Arrival}
\textbf{Declaration location:} reusable library
\paragraph{Lean declaration body}
\begin{ReviewVerbatim}
type:
{Class : Type u_2} →
  (i : Class) → AppliedModelingLib.Queueing.MulticlassStationaryPoissonWorkClassTaggedSample Class i → Class → ℤ → ℝ

value:
fun {Class : Type u_2} (i : Class)
    (a : AppliedModelingLib.Queueing.MulticlassStationaryPoissonWorkClassTaggedSample Class i) (a_1 : Class)
    (a_2 : ℤ) =>
  if hji : a_1 = i then AppliedModelingLib.Probability.PoissonProcess.candidatePalmArrival a.1.1 a_2
  else AppliedModelingLib.Probability.Queueing.stationaryPoissonWorkArrival (a.2 ⟨a_1, hji⟩) a_2
\end{ReviewVerbatim}

\paragraph{Direct retained dependencies}
\begin{itemize}\raggedright
\item \hyperlink{governing-declaration-9208635caa21f800}{Applied\allowbreak{}Modeling\allowbreak{}Lib.\allowbreak{}Probability.\allowbreak{}Poisson\allowbreak{}Process.\allowbreak{}candidate\allowbreak{}Palm\allowbreak{}Arrival}
\item \hyperlink{governing-declaration-499050173dfb71ee}{Applied\allowbreak{}Modeling\allowbreak{}Lib.\allowbreak{}Probability.\allowbreak{}Queueing.\allowbreak{}stationary\allowbreak{}Poisson\allowbreak{}Work\allowbreak{}Arrival}
\item \hyperlink{library-prerequisite-755cdb28ceaf46d2}{Applied\allowbreak{}Modeling\allowbreak{}Lib.\allowbreak{}Queueing.\allowbreak{}Multiclass\allowbreak{}Stationary\allowbreak{}Poisson\allowbreak{}Work\allowbreak{}Class\allowbreak{}Tagged\allowbreak{}Sample}
\item \hyperlink{governing-declaration-62a268164b6ae28b}{Applied\allowbreak{}Modeling\allowbreak{}Lib.\allowbreak{}Queueing.\allowbreak{}Stationary\allowbreak{}Poisson\allowbreak{}Work\allowbreak{}Path}
\item \hyperlink{governing-declaration-de954530ca750d3f}{Applied\allowbreak{}Modeling\allowbreak{}Lib.\allowbreak{}Queueing.\allowbreak{}inst\allowbreak{}Decidable\allowbreak{}Eq\_\allowbreak{}applied\allowbreak{}Modeling\allowbreak{}Lib}
\end{itemize}
\hypertarget{governing-declaration-2317cefa31cefbe1}{}
\subsection*{\small\ttfamily\raggedright Applied\allowbreak{}Modeling\allowbreak{}Lib.\allowbreak{}Queueing.\allowbreak{}multiclass\allowbreak{}Stationary\allowbreak{}Poisson\allowbreak{}Work\allowbreak{}Class\allowbreak{}Tagged\allowbreak{}Requirement}
\textbf{Declaration location:} reusable library
\paragraph{Lean declaration body}
\begin{ReviewVerbatim}
type:
{Class : Type u_1} →
  (i : Class) →
    ((ℤ → ℝ) × (ℤ → ℝ)) × ({ j : Class // j ≠ i } → AppliedModelingLib.Queueing.StationaryPoissonWorkPath) → ℝ

value:
fun {Class : Type u_1} (i : Class)
    (a : ((ℤ → ℝ) × (ℤ → ℝ)) × ({ j : Class // j ≠ i } → AppliedModelingLib.Queueing.StationaryPoissonWorkPath)) =>
  a.1.2 0
\end{ReviewVerbatim}

\paragraph{Direct retained dependencies}
\begin{itemize}\raggedright
\item \hyperlink{governing-declaration-62a268164b6ae28b}{Applied\allowbreak{}Modeling\allowbreak{}Lib.\allowbreak{}Queueing.\allowbreak{}Stationary\allowbreak{}Poisson\allowbreak{}Work\allowbreak{}Path}
\end{itemize}
\hypertarget{governing-declaration-33efae4a88ebb3d1}{}
\subsection*{\small\ttfamily\raggedright Applied\allowbreak{}Modeling\allowbreak{}Lib.\allowbreak{}Queueing.\allowbreak{}multiclass\allowbreak{}Stationary\allowbreak{}Poisson\allowbreak{}Work\allowbreak{}Class\allowbreak{}Tagged\allowbreak{}Requirement\allowbreak{}At}
\textbf{Declaration location:} reusable library
\paragraph{Lean declaration body}
\begin{ReviewVerbatim}
type:
{Class : Type u_2} →
  (i : Class) → AppliedModelingLib.Queueing.MulticlassStationaryPoissonWorkClassTaggedSample Class i → Class → ℤ → ℝ

value:
fun {Class : Type u_2} (i : Class)
    (a : AppliedModelingLib.Queueing.MulticlassStationaryPoissonWorkClassTaggedSample Class i) (a_1 : Class)
    (a_2 : ℤ) =>
  if hji : a_1 = i then a.1.2 a_2 else (a.2 ⟨a_1, hji⟩).2 a_2
\end{ReviewVerbatim}

\paragraph{Direct retained dependencies}
\begin{itemize}\raggedright
\item \hyperlink{governing-declaration-ab1cfe37ad8c32b0}{Applied\allowbreak{}Modeling\allowbreak{}Lib.\allowbreak{}Probability.\allowbreak{}Poisson\allowbreak{}Process.\allowbreak{}Good\allowbreak{}Suspension\allowbreak{}State}
\item \hyperlink{library-prerequisite-755cdb28ceaf46d2}{Applied\allowbreak{}Modeling\allowbreak{}Lib.\allowbreak{}Queueing.\allowbreak{}Multiclass\allowbreak{}Stationary\allowbreak{}Poisson\allowbreak{}Work\allowbreak{}Class\allowbreak{}Tagged\allowbreak{}Sample}
\item \hyperlink{governing-declaration-62a268164b6ae28b}{Applied\allowbreak{}Modeling\allowbreak{}Lib.\allowbreak{}Queueing.\allowbreak{}Stationary\allowbreak{}Poisson\allowbreak{}Work\allowbreak{}Path}
\item \hyperlink{governing-declaration-de954530ca750d3f}{Applied\allowbreak{}Modeling\allowbreak{}Lib.\allowbreak{}Queueing.\allowbreak{}inst\allowbreak{}Decidable\allowbreak{}Eq\_\allowbreak{}applied\allowbreak{}Modeling\allowbreak{}Lib}
\end{itemize}
\hypertarget{governing-declaration-9bede4eaca6a95fb}{}
\subsection*{\small\ttfamily\raggedright Applied\allowbreak{}Modeling\allowbreak{}Lib.\allowbreak{}Queueing.\allowbreak{}multiclass\allowbreak{}Stationary\allowbreak{}Poisson\allowbreak{}Work\allowbreak{}Rest\allowbreak{}Law}
\textbf{Declaration location:} reusable library
\paragraph{Lean declaration body}
\begin{ReviewVerbatim}
type:
{Class : Type u_1} →
  [Fintype Class] →
    (arrivalRate : Class → ℝ) →
      (∀ (j : Class), 0 < arrivalRate j) →
        (i : Class) →
          AppliedModelingLib.Probability.Palm.ShiftInvariantProbabilityLaw
            ({ j : Class // j ≠ i } → AppliedModelingLib.Queueing.StationaryPoissonWorkPath)

value:
fun {Class : Type u_1} [Fintype Class] (arrivalRate : Class → ℝ) (harrivalRate : ∀ (j : Class), 0 < arrivalRate j)
    (i : Class) =>
  AppliedModelingLib.Probability.PoissonProcess.multiclassStationaryPoissonWorkShiftInvariantLaw
    (fun (j : { j : Class // j ≠ i }) => arrivalRate ↑j)
    (AppliedModelingLib.Queueing.multiclassStationaryPoissonWorkRestLaw._proof_1 arrivalRate harrivalRate i)
\end{ReviewVerbatim}

\paragraph{Direct retained dependencies}
\begin{itemize}\raggedright
\item \hyperlink{governing-declaration-fb5ec66edd06d906}{Applied\allowbreak{}Modeling\allowbreak{}Lib.\allowbreak{}Probability.\allowbreak{}Palm.\allowbreak{}Shift\allowbreak{}Invariant\allowbreak{}Probability\allowbreak{}Law}
\item \hyperlink{governing-declaration-ab1cfe37ad8c32b0}{Applied\allowbreak{}Modeling\allowbreak{}Lib.\allowbreak{}Probability.\allowbreak{}Poisson\allowbreak{}Process.\allowbreak{}Good\allowbreak{}Suspension\allowbreak{}State}
\item \hyperlink{governing-declaration-a8385d799d0a4d16}{Applied\allowbreak{}Modeling\allowbreak{}Lib.\allowbreak{}Probability.\allowbreak{}Poisson\allowbreak{}Process.\allowbreak{}inst\allowbreak{}Measurable\allowbreak{}Space\allowbreak{}Good\allowbreak{}Suspension\allowbreak{}State}
\item \hyperlink{governing-declaration-0a1733e4d91e5dfa}{Applied\allowbreak{}Modeling\allowbreak{}Lib.\allowbreak{}Probability.\allowbreak{}Poisson\allowbreak{}Process.\allowbreak{}multiclass\allowbreak{}Stationary\allowbreak{}Poisson\allowbreak{}Work\allowbreak{}Shift\allowbreak{}Invariant\allowbreak{}Law}
\item \hyperlink{governing-declaration-62a268164b6ae28b}{Applied\allowbreak{}Modeling\allowbreak{}Lib.\allowbreak{}Queueing.\allowbreak{}Stationary\allowbreak{}Poisson\allowbreak{}Work\allowbreak{}Path}
\item \hyperlink{governing-declaration-de954530ca750d3f}{Applied\allowbreak{}Modeling\allowbreak{}Lib.\allowbreak{}Queueing.\allowbreak{}inst\allowbreak{}Decidable\allowbreak{}Eq\_\allowbreak{}applied\allowbreak{}Modeling\allowbreak{}Lib}
\end{itemize}
\hypertarget{governing-declaration-efe1ba9dba9a89a1}{}
\subsection*{\small\ttfamily\raggedright Applied\allowbreak{}Modeling\allowbreak{}Lib.\allowbreak{}Queueing.\allowbreak{}nonempty\allowbreak{}Priority\allowbreak{}Classes}
\textbf{Declaration location:} reusable library
\paragraph{Lean declaration body}
\begin{ReviewVerbatim}
type:
{n : ℕ} → (Fin n → ℕ) → Finset (Fin n)

value:
fun {n : ℕ} (backlog : Fin n → ℕ) => {i : Fin n | 0 < backlog i}
\end{ReviewVerbatim}


\hypertarget{governing-declaration-d817e39256a94da7}{}
\subsection*{\small\ttfamily\raggedright Applied\allowbreak{}Modeling\allowbreak{}Lib.\allowbreak{}Queueing.\allowbreak{}next\allowbreak{}Nonpreemptive\allowbreak{}Priority}
\textbf{Declaration location:} reusable library
\paragraph{Lean declaration body}
\begin{ReviewVerbatim}
type:
{n : ℕ} → (backlog : Fin n → ℕ) → (∃ (i : Fin n), 0 < backlog i) → Fin n

value:
fun {n : ℕ} (backlog : Fin n → ℕ) (hnonempty : ∃ (i : Fin n), 0 < backlog i) =>
  (AppliedModelingLib.Queueing.nonemptyPriorityClasses backlog).min'
    (AppliedModelingLib.Queueing.nextNonpreemptivePriority._proof_1 backlog hnonempty)
\end{ReviewVerbatim}

\paragraph{Direct retained dependencies}
\begin{itemize}\raggedright
\item \hyperlink{governing-declaration-efe1ba9dba9a89a1}{Applied\allowbreak{}Modeling\allowbreak{}Lib.\allowbreak{}Queueing.\allowbreak{}nonempty\allowbreak{}Priority\allowbreak{}Classes}
\end{itemize}
\hypertarget{governing-declaration-2d99be8c1afb689e}{}
\subsection*{\small\ttfamily\raggedright Applied\allowbreak{}Modeling\allowbreak{}Lib.\allowbreak{}Queueing.\allowbreak{}next\allowbreak{}Priority\allowbreak{}Waiting\allowbreak{}Class}
\textbf{Declaration location:} reusable library
\paragraph{Lean declaration body}
\begin{ReviewVerbatim}
type:
{n : ℕ} →
  {JobId : Type u_1} →
    (state : AppliedModelingLib.Queueing.NonpreemptivePriorityWorkState n JobId) →
      AppliedModelingLib.Queueing.hasPriorityWaitingJob state → Fin n

value:
fun {n : ℕ} {JobId : Type u_1} (state : AppliedModelingLib.Queueing.NonpreemptivePriorityWorkState n JobId)
    (hwaiting : AppliedModelingLib.Queueing.hasPriorityWaitingJob state) =>
  AppliedModelingLib.Queueing.nextNonpreemptivePriority (fun (i : Fin n) => (state.waiting i).length) hwaiting
\end{ReviewVerbatim}

\paragraph{Direct retained dependencies}
\begin{itemize}\raggedright
\item \hyperlink{library-prerequisite-2373fe691c32d999}{Applied\allowbreak{}Modeling\allowbreak{}Lib.\allowbreak{}Queueing.\allowbreak{}Nonpreemptive\allowbreak{}Priority\allowbreak{}Job}
\item \hyperlink{library-prerequisite-e68cb5fce009ec53}{Applied\allowbreak{}Modeling\allowbreak{}Lib.\allowbreak{}Queueing.\allowbreak{}Nonpreemptive\allowbreak{}Priority\allowbreak{}Work\allowbreak{}State}
\item \hyperlink{governing-declaration-adcd68c3e7c5c1bb}{Applied\allowbreak{}Modeling\allowbreak{}Lib.\allowbreak{}Queueing.\allowbreak{}has\allowbreak{}Priority\allowbreak{}Waiting\allowbreak{}Job}
\item \hyperlink{governing-declaration-d817e39256a94da7}{Applied\allowbreak{}Modeling\allowbreak{}Lib.\allowbreak{}Queueing.\allowbreak{}next\allowbreak{}Nonpreemptive\allowbreak{}Priority}
\end{itemize}
\hypertarget{governing-declaration-2ae4329b7cd91c5b}{}
\subsection*{\small\ttfamily\raggedright Applied\allowbreak{}Modeling\allowbreak{}Lib.\allowbreak{}Queueing.\allowbreak{}nonpreemptive\allowbreak{}Priority\allowbreak{}Arrival\allowbreak{}Window\allowbreak{}Indices}
\textbf{Declaration location:} reusable library
\paragraph{Lean declaration body}
\begin{ReviewVerbatim}
type:
{n : ℕ} → (Fin n → Finset ℤ) → Finset (AppliedModelingLib.Queueing.NonpreemptivePriorityArrivalIndex n)

value:
fun {n : ℕ} (indices : Fin n → Finset ℤ) => Finset.univ.sigma indices
\end{ReviewVerbatim}

\paragraph{Direct retained dependencies}
\begin{itemize}\raggedright
\item \hyperlink{governing-declaration-5d3724f4dc832d77}{Applied\allowbreak{}Modeling\allowbreak{}Lib.\allowbreak{}Queueing.\allowbreak{}Nonpreemptive\allowbreak{}Priority\allowbreak{}Arrival\allowbreak{}Index}
\end{itemize}
\hypertarget{governing-declaration-0dd8378f656d740f}{}
\subsection*{\small\ttfamily\raggedright Applied\allowbreak{}Modeling\allowbreak{}Lib.\allowbreak{}Queueing.\allowbreak{}start\allowbreak{}Next\allowbreak{}Nonpreemptive\allowbreak{}Priority\allowbreak{}Job}
\textbf{Declaration location:} reusable library
\paragraph{Lean declaration body}
\begin{ReviewVerbatim}
type:
{n : ℕ} →
  {JobId : Type u_1} →
    AppliedModelingLib.Queueing.NonpreemptivePriorityWorkState n JobId →
      AppliedModelingLib.Queueing.NonpreemptivePriorityWorkState n JobId

value:
fun {n : ℕ} {JobId : Type u_1} (state : AppliedModelingLib.Queueing.NonpreemptivePriorityWorkState n JobId) =>
  match state.active with
  | some val => state
  | none =>
    if hwaiting : AppliedModelingLib.Queueing.hasPriorityWaitingJob state then
      have i : Fin n := AppliedModelingLib.Queueing.nextPriorityWaitingClass state hwaiting;
      match state.waiting i with
      | [] => state
      | job :: tail =>
        { currentTime := state.currentTime, active := some (job, job.serviceWork),
          waiting := Function.update state.waiting i tail, completed := state.completed }
    else state
\end{ReviewVerbatim}

\paragraph{Direct retained dependencies}
\begin{itemize}\raggedright
\item \hyperlink{library-prerequisite-2373fe691c32d999}{Applied\allowbreak{}Modeling\allowbreak{}Lib.\allowbreak{}Queueing.\allowbreak{}Nonpreemptive\allowbreak{}Priority\allowbreak{}Job}
\item \hyperlink{library-prerequisite-e68cb5fce009ec53}{Applied\allowbreak{}Modeling\allowbreak{}Lib.\allowbreak{}Queueing.\allowbreak{}Nonpreemptive\allowbreak{}Priority\allowbreak{}Work\allowbreak{}State}
\item \hyperlink{governing-declaration-adcd68c3e7c5c1bb}{Applied\allowbreak{}Modeling\allowbreak{}Lib.\allowbreak{}Queueing.\allowbreak{}has\allowbreak{}Priority\allowbreak{}Waiting\allowbreak{}Job}
\item \hyperlink{governing-declaration-2d99be8c1afb689e}{Applied\allowbreak{}Modeling\allowbreak{}Lib.\allowbreak{}Queueing.\allowbreak{}next\allowbreak{}Priority\allowbreak{}Waiting\allowbreak{}Class}
\end{itemize}
\hypertarget{governing-declaration-e7772cc6b7f51cdf}{}
\subsection*{\small\ttfamily\raggedright Applied\allowbreak{}Modeling\allowbreak{}Lib.\allowbreak{}Queueing.\allowbreak{}complete\allowbreak{}Nonpreemptive\allowbreak{}Priority\allowbreak{}Work\allowbreak{}Job}
\textbf{Declaration location:} reusable library
\paragraph{Lean declaration body}
\begin{ReviewVerbatim}
type:
{n : ℕ} →
  {JobId : Type u_1} →
    AppliedModelingLib.Queueing.NonpreemptivePriorityWorkState n JobId →
      AppliedModelingLib.Queueing.NonpreemptivePriorityWorkState n JobId

value:
fun {n : ℕ} {JobId : Type u_1} (state : AppliedModelingLib.Queueing.NonpreemptivePriorityWorkState n JobId) =>
  match state.active with
  | none => state
  | some (job, snd) =>
    AppliedModelingLib.Queueing.startNextNonpreemptivePriorityJob
      { currentTime := state.currentTime, active := none, waiting := state.waiting,
        completed := (job, state.currentTime) :: state.completed }
\end{ReviewVerbatim}

\paragraph{Direct retained dependencies}
\begin{itemize}\raggedright
\item \hyperlink{library-prerequisite-2373fe691c32d999}{Applied\allowbreak{}Modeling\allowbreak{}Lib.\allowbreak{}Queueing.\allowbreak{}Nonpreemptive\allowbreak{}Priority\allowbreak{}Job}
\item \hyperlink{library-prerequisite-e68cb5fce009ec53}{Applied\allowbreak{}Modeling\allowbreak{}Lib.\allowbreak{}Queueing.\allowbreak{}Nonpreemptive\allowbreak{}Priority\allowbreak{}Work\allowbreak{}State}
\item \hyperlink{governing-declaration-0dd8378f656d740f}{Applied\allowbreak{}Modeling\allowbreak{}Lib.\allowbreak{}Queueing.\allowbreak{}start\allowbreak{}Next\allowbreak{}Nonpreemptive\allowbreak{}Priority\allowbreak{}Job}
\end{itemize}
\hypertarget{governing-declaration-fe14523475acac38}{}
\subsection*{\small\ttfamily\raggedright Applied\allowbreak{}Modeling\allowbreak{}Lib.\allowbreak{}Queueing.\allowbreak{}advance\allowbreak{}Then\allowbreak{}Admit\allowbreak{}Nonpreemptive\allowbreak{}Priority\allowbreak{}Job}
\textbf{Declaration location:} reusable library
\paragraph{Lean declaration body}
\begin{ReviewVerbatim}
type:
{n : ℕ} →
  {JobId : Type u_1} →
    ℕ →
      AppliedModelingLib.Queueing.NonpreemptivePriorityWorkState n JobId →
        AppliedModelingLib.Queueing.NonpreemptivePriorityJob n JobId →
          AppliedModelingLib.Queueing.NonpreemptivePriorityWorkState n JobId

value:
fun {n : ℕ} {JobId : Type u_1} (fuel : ℕ) (state : AppliedModelingLib.Queueing.NonpreemptivePriorityWorkState n JobId)
    (job : AppliedModelingLib.Queueing.NonpreemptivePriorityJob n JobId) =>
  AppliedModelingLib.Queueing.admitNonpreemptivePriorityJob
    (AppliedModelingLib.Queueing.advanceNonpreemptivePriorityWorkState fuel job.arrivalTime state) job
\end{ReviewVerbatim}

\paragraph{Direct retained dependencies}
\begin{itemize}\raggedright
\item \hyperlink{library-prerequisite-2373fe691c32d999}{Applied\allowbreak{}Modeling\allowbreak{}Lib.\allowbreak{}Queueing.\allowbreak{}Nonpreemptive\allowbreak{}Priority\allowbreak{}Job}
\item \hyperlink{library-prerequisite-e68cb5fce009ec53}{Applied\allowbreak{}Modeling\allowbreak{}Lib.\allowbreak{}Queueing.\allowbreak{}Nonpreemptive\allowbreak{}Priority\allowbreak{}Work\allowbreak{}State}
\item \hyperlink{library-prerequisite-837de7f0cc179814}{Applied\allowbreak{}Modeling\allowbreak{}Lib.\allowbreak{}Queueing.\allowbreak{}admit\allowbreak{}Nonpreemptive\allowbreak{}Priority\allowbreak{}Job}
\item \hyperlink{library-prerequisite-96200a424cf39dee}{Applied\allowbreak{}Modeling\allowbreak{}Lib.\allowbreak{}Queueing.\allowbreak{}advance\allowbreak{}Nonpreemptive\allowbreak{}Priority\allowbreak{}Work\allowbreak{}State}
\end{itemize}
\hypertarget{governing-declaration-2945df1a7a8ec970}{}
\subsection*{\small\ttfamily\raggedright Applied\allowbreak{}Modeling\allowbreak{}Lib.\allowbreak{}Queueing.\allowbreak{}stationary\allowbreak{}Priority\allowbreak{}Class\allowbreak{}Tagged\allowbreak{}Arrival}
\textbf{Declaration location:} reusable library
\paragraph{Lean declaration body}
\begin{ReviewVerbatim}
type:
{Class : Type u_1} →
  (i : Class) → AppliedModelingLib.Queueing.MulticlassStationaryPoissonWorkClassTaggedSample Class i → Class → ℤ → ℝ

value:
fun {Class : Type u_1} (i : Class)
    (a : AppliedModelingLib.Queueing.MulticlassStationaryPoissonWorkClassTaggedSample Class i) (a_1 : Class)
    (a_2 : ℤ) =>
  AppliedModelingLib.Queueing.multiclassStationaryPoissonWorkClassTaggedArrival i a a_1 a_2
\end{ReviewVerbatim}

\paragraph{Direct retained dependencies}
\begin{itemize}\raggedright
\item \hyperlink{library-prerequisite-755cdb28ceaf46d2}{Applied\allowbreak{}Modeling\allowbreak{}Lib.\allowbreak{}Queueing.\allowbreak{}Multiclass\allowbreak{}Stationary\allowbreak{}Poisson\allowbreak{}Work\allowbreak{}Class\allowbreak{}Tagged\allowbreak{}Sample}
\item \hyperlink{governing-declaration-8f366849032baf38}{Applied\allowbreak{}Modeling\allowbreak{}Lib.\allowbreak{}Queueing.\allowbreak{}multiclass\allowbreak{}Stationary\allowbreak{}Poisson\allowbreak{}Work\allowbreak{}Class\allowbreak{}Tagged\allowbreak{}Arrival}
\end{itemize}
\hypertarget{governing-declaration-fffbbc78aec2f9da}{}
\subsection*{\small\ttfamily\raggedright Applied\allowbreak{}Modeling\allowbreak{}Lib.\allowbreak{}Queueing.\allowbreak{}stationary\allowbreak{}Priority\allowbreak{}Class\allowbreak{}Tagged\allowbreak{}Arrival\allowbreak{}Index\allowbreak{}Key}
\textbf{Declaration location:} reusable library
\paragraph{Lean declaration body}
\begin{ReviewVerbatim}
type:
{n : ℕ} →
  (i : Fin n) →
    AppliedModelingLib.Queueing.MulticlassStationaryPoissonWorkClassTaggedSample (Fin n) i →
      AppliedModelingLib.Queueing.NonpreemptivePriorityArrivalIndex n → Lex (ℝ × Σₗ (j : Fin n), ℤ)

value:
fun {n : ℕ} (i : Fin n) (z : AppliedModelingLib.Queueing.MulticlassStationaryPoissonWorkClassTaggedSample (Fin n) i)
    (q : AppliedModelingLib.Queueing.NonpreemptivePriorityArrivalIndex n) =>
  toLex (AppliedModelingLib.Queueing.stationaryPriorityClassTaggedArrival i z q.fst q.snd, toLex q)
\end{ReviewVerbatim}

\paragraph{Direct retained dependencies}
\begin{itemize}\raggedright
\item \hyperlink{library-prerequisite-755cdb28ceaf46d2}{Applied\allowbreak{}Modeling\allowbreak{}Lib.\allowbreak{}Queueing.\allowbreak{}Multiclass\allowbreak{}Stationary\allowbreak{}Poisson\allowbreak{}Work\allowbreak{}Class\allowbreak{}Tagged\allowbreak{}Sample}
\item \hyperlink{governing-declaration-5d3724f4dc832d77}{Applied\allowbreak{}Modeling\allowbreak{}Lib.\allowbreak{}Queueing.\allowbreak{}Nonpreemptive\allowbreak{}Priority\allowbreak{}Arrival\allowbreak{}Index}
\item \hyperlink{governing-declaration-2945df1a7a8ec970}{Applied\allowbreak{}Modeling\allowbreak{}Lib.\allowbreak{}Queueing.\allowbreak{}stationary\allowbreak{}Priority\allowbreak{}Class\allowbreak{}Tagged\allowbreak{}Arrival}
\end{itemize}
\hypertarget{governing-declaration-6eb7b19311e0b293}{}
\subsection*{\small\ttfamily\raggedright Applied\allowbreak{}Modeling\allowbreak{}Lib.\allowbreak{}Queueing.\allowbreak{}stationary\allowbreak{}Priority\allowbreak{}Class\allowbreak{}Tagged\allowbreak{}Arrival\allowbreak{}Index\allowbreak{}LE}
\textbf{Declaration location:} reusable library
\paragraph{Lean declaration body}
\begin{ReviewVerbatim}
type:
{n : ℕ} →
  (i : Fin n) →
    AppliedModelingLib.Queueing.MulticlassStationaryPoissonWorkClassTaggedSample (Fin n) i →
      AppliedModelingLib.Queueing.NonpreemptivePriorityArrivalIndex n →
        AppliedModelingLib.Queueing.NonpreemptivePriorityArrivalIndex n → Prop

value:
fun {n : ℕ} (i : Fin n) (z : AppliedModelingLib.Queueing.MulticlassStationaryPoissonWorkClassTaggedSample (Fin n) i)
    (first second : AppliedModelingLib.Queueing.NonpreemptivePriorityArrivalIndex n) =>
  AppliedModelingLib.Queueing.stationaryPriorityClassTaggedArrivalIndexKey i z first ≤
    AppliedModelingLib.Queueing.stationaryPriorityClassTaggedArrivalIndexKey i z second
\end{ReviewVerbatim}

\paragraph{Direct retained dependencies}
\begin{itemize}\raggedright
\item \hyperlink{library-prerequisite-755cdb28ceaf46d2}{Applied\allowbreak{}Modeling\allowbreak{}Lib.\allowbreak{}Queueing.\allowbreak{}Multiclass\allowbreak{}Stationary\allowbreak{}Poisson\allowbreak{}Work\allowbreak{}Class\allowbreak{}Tagged\allowbreak{}Sample}
\item \hyperlink{governing-declaration-5d3724f4dc832d77}{Applied\allowbreak{}Modeling\allowbreak{}Lib.\allowbreak{}Queueing.\allowbreak{}Nonpreemptive\allowbreak{}Priority\allowbreak{}Arrival\allowbreak{}Index}
\item \hyperlink{governing-declaration-fffbbc78aec2f9da}{Applied\allowbreak{}Modeling\allowbreak{}Lib.\allowbreak{}Queueing.\allowbreak{}stationary\allowbreak{}Priority\allowbreak{}Class\allowbreak{}Tagged\allowbreak{}Arrival\allowbreak{}Index\allowbreak{}Key}
\end{itemize}
\hypertarget{governing-declaration-258d1e091ab1b39d}{}
\subsection*{\small\ttfamily\raggedright Applied\allowbreak{}Modeling\allowbreak{}Lib.\allowbreak{}Queueing.\allowbreak{}stationary\allowbreak{}Priority\allowbreak{}Class\allowbreak{}Tagged\allowbreak{}Future\allowbreak{}Arrival\allowbreak{}Indices}
\textbf{Declaration location:} reusable library
\paragraph{Lean declaration body}
\begin{ReviewVerbatim}
type:
{n : ℕ} →
  (i : Fin n) →
    AppliedModelingLib.Queueing.MulticlassStationaryPoissonWorkClassTaggedSample (Fin n) i → ℝ → Fin n → Finset ℤ

value:
fun {n : ℕ} (i : Fin n) (z : AppliedModelingLib.Queueing.MulticlassStationaryPoissonWorkClassTaggedSample (Fin n) i)
    (t : ℝ) (a : Fin n) =>
  if hji : a = i then AppliedModelingLib.Probability.PoissonProcess.palmTaggedArrivalIndicesRightClosed 0 t z.1.1
  else AppliedModelingLib.Probability.PoissonProcess.suspensionBaseArrivalIndicesRightClosed 0 t (z.2 ⟨a, hji⟩).1
\end{ReviewVerbatim}

\paragraph{Direct retained dependencies}
\begin{itemize}\raggedright
\item \hyperlink{governing-declaration-ab1cfe37ad8c32b0}{Applied\allowbreak{}Modeling\allowbreak{}Lib.\allowbreak{}Probability.\allowbreak{}Poisson\allowbreak{}Process.\allowbreak{}Good\allowbreak{}Suspension\allowbreak{}State}
\item \hyperlink{governing-declaration-55048043b93956f0}{Applied\allowbreak{}Modeling\allowbreak{}Lib.\allowbreak{}Probability.\allowbreak{}Poisson\allowbreak{}Process.\allowbreak{}palm\allowbreak{}Tagged\allowbreak{}Arrival\allowbreak{}Indices\allowbreak{}Right\allowbreak{}Closed}
\item \hyperlink{governing-declaration-09b9ceda015085bb}{Applied\allowbreak{}Modeling\allowbreak{}Lib.\allowbreak{}Probability.\allowbreak{}Poisson\allowbreak{}Process.\allowbreak{}suspension\allowbreak{}Base\allowbreak{}Arrival\allowbreak{}Indices\allowbreak{}Right\allowbreak{}Closed}
\item \hyperlink{library-prerequisite-755cdb28ceaf46d2}{Applied\allowbreak{}Modeling\allowbreak{}Lib.\allowbreak{}Queueing.\allowbreak{}Multiclass\allowbreak{}Stationary\allowbreak{}Poisson\allowbreak{}Work\allowbreak{}Class\allowbreak{}Tagged\allowbreak{}Sample}
\item \hyperlink{governing-declaration-62a268164b6ae28b}{Applied\allowbreak{}Modeling\allowbreak{}Lib.\allowbreak{}Queueing.\allowbreak{}Stationary\allowbreak{}Poisson\allowbreak{}Work\allowbreak{}Path}
\end{itemize}
\hypertarget{governing-declaration-e1de2ae4651d3437}{}
\subsection*{\small\ttfamily\raggedright Applied\allowbreak{}Modeling\allowbreak{}Lib.\allowbreak{}Queueing.\allowbreak{}canonical\allowbreak{}Stationary\allowbreak{}Priority\allowbreak{}Class\allowbreak{}Tagged\allowbreak{}Future\allowbreak{}Arrival\allowbreak{}Indices}
\textbf{Declaration location:} reusable library
\paragraph{Lean declaration body}
\begin{ReviewVerbatim}
type:
{n : ℕ} →
  (i : Fin n) →
    AppliedModelingLib.Queueing.MulticlassStationaryPoissonWorkClassTaggedSample (Fin n) i →
      ℝ → List (AppliedModelingLib.Queueing.NonpreemptivePriorityArrivalIndex n)

value:
fun {n : ℕ} (i : Fin n) (z : AppliedModelingLib.Queueing.MulticlassStationaryPoissonWorkClassTaggedSample (Fin n) i)
    (t : ℝ) =>
  (AppliedModelingLib.Queueing.nonpreemptivePriorityArrivalWindowIndices
        (AppliedModelingLib.Queueing.stationaryPriorityClassTaggedFutureArrivalIndices i z t)).sort
    (AppliedModelingLib.Queueing.stationaryPriorityClassTaggedArrivalIndexLE i z)
\end{ReviewVerbatim}

\paragraph{Direct retained dependencies}
\begin{itemize}\raggedright
\item \hyperlink{library-prerequisite-755cdb28ceaf46d2}{Applied\allowbreak{}Modeling\allowbreak{}Lib.\allowbreak{}Queueing.\allowbreak{}Multiclass\allowbreak{}Stationary\allowbreak{}Poisson\allowbreak{}Work\allowbreak{}Class\allowbreak{}Tagged\allowbreak{}Sample}
\item \hyperlink{governing-declaration-5d3724f4dc832d77}{Applied\allowbreak{}Modeling\allowbreak{}Lib.\allowbreak{}Queueing.\allowbreak{}Nonpreemptive\allowbreak{}Priority\allowbreak{}Arrival\allowbreak{}Index}
\item \hyperlink{governing-declaration-2ae4329b7cd91c5b}{Applied\allowbreak{}Modeling\allowbreak{}Lib.\allowbreak{}Queueing.\allowbreak{}nonpreemptive\allowbreak{}Priority\allowbreak{}Arrival\allowbreak{}Window\allowbreak{}Indices}
\item \hyperlink{governing-declaration-6eb7b19311e0b293}{Applied\allowbreak{}Modeling\allowbreak{}Lib.\allowbreak{}Queueing.\allowbreak{}stationary\allowbreak{}Priority\allowbreak{}Class\allowbreak{}Tagged\allowbreak{}Arrival\allowbreak{}Index\allowbreak{}LE}
\item \hyperlink{governing-declaration-258d1e091ab1b39d}{Applied\allowbreak{}Modeling\allowbreak{}Lib.\allowbreak{}Queueing.\allowbreak{}stationary\allowbreak{}Priority\allowbreak{}Class\allowbreak{}Tagged\allowbreak{}Future\allowbreak{}Arrival\allowbreak{}Indices}
\end{itemize}
\hypertarget{governing-declaration-949c09de5ada557f}{}
\subsection*{\small\ttfamily\raggedright Applied\allowbreak{}Modeling\allowbreak{}Lib.\allowbreak{}Queueing.\allowbreak{}stationary\allowbreak{}Priority\allowbreak{}Class\allowbreak{}Tagged\allowbreak{}Work\allowbreak{}Requirement\allowbreak{}At}
\textbf{Declaration location:} reusable library
\paragraph{Lean declaration body}
\begin{ReviewVerbatim}
type:
{Class : Type u_1} →
  (Class → ℝ) →
    (i : Class) → AppliedModelingLib.Queueing.MulticlassStationaryPoissonWorkClassTaggedSample Class i → Class → ℤ → ℝ

value:
fun {Class : Type u_1} (meanService : Class → ℝ) (i : Class)
    (a : AppliedModelingLib.Queueing.MulticlassStationaryPoissonWorkClassTaggedSample Class i) (a_1 : Class)
    (a_2 : ℤ) =>
  meanService a_1 * AppliedModelingLib.Queueing.multiclassStationaryPoissonWorkClassTaggedRequirementAt i a a_1 a_2
\end{ReviewVerbatim}

\paragraph{Direct retained dependencies}
\begin{itemize}\raggedright
\item \hyperlink{library-prerequisite-755cdb28ceaf46d2}{Applied\allowbreak{}Modeling\allowbreak{}Lib.\allowbreak{}Queueing.\allowbreak{}Multiclass\allowbreak{}Stationary\allowbreak{}Poisson\allowbreak{}Work\allowbreak{}Class\allowbreak{}Tagged\allowbreak{}Sample}
\item \hyperlink{governing-declaration-33efae4a88ebb3d1}{Applied\allowbreak{}Modeling\allowbreak{}Lib.\allowbreak{}Queueing.\allowbreak{}multiclass\allowbreak{}Stationary\allowbreak{}Poisson\allowbreak{}Work\allowbreak{}Class\allowbreak{}Tagged\allowbreak{}Requirement\allowbreak{}At}
\end{itemize}
\hypertarget{governing-declaration-05133ca1584020e9}{}
\subsection*{\small\ttfamily\raggedright Applied\allowbreak{}Modeling\allowbreak{}Lib.\allowbreak{}Queueing.\allowbreak{}stationary\allowbreak{}Priority\allowbreak{}Input\allowbreak{}Decidable\allowbreak{}Eq}
\textbf{Declaration location:} reusable library
\paragraph{Lean declaration body}
\begin{ReviewVerbatim}
type:
{Class : Type u_1} → (a b : Class) → Decidable (a = b)

value:
fun {Class : Type u_1} (a b : Class) => Classical.decEq Class a b
\end{ReviewVerbatim}


\hypertarget{governing-declaration-9f498af160c98bd8}{}
\subsection*{\small\ttfamily\raggedright Applied\allowbreak{}Modeling\allowbreak{}Lib.\allowbreak{}Queueing.\allowbreak{}stationary\allowbreak{}Priority\allowbreak{}Class\allowbreak{}Tagged\allowbreak{}Arrival\allowbreak{}Window\allowbreak{}Indices}
\textbf{Declaration location:} reusable library
\paragraph{Lean declaration body}
\begin{ReviewVerbatim}
type:
{Class : Type u_1} →
  (i : Class) →
    AppliedModelingLib.Queueing.MulticlassStationaryPoissonWorkClassTaggedSample Class i → ℝ → ℝ → Class → Finset ℤ

value:
fun {Class : Type u_1} (i : Class)
    (z : AppliedModelingLib.Queueing.MulticlassStationaryPoissonWorkClassTaggedSample Class i) (a b : ℝ) (j : Class) =>
  if hji : j = i then AppliedModelingLib.Probability.PoissonProcess.palmTaggedArrivalIndices a b z.1.1
  else AppliedModelingLib.Probability.PoissonProcess.suspensionBaseArrivalIndices a b (z.2 ⟨j, hji⟩).1
\end{ReviewVerbatim}

\paragraph{Direct retained dependencies}
\begin{itemize}\raggedright
\item \hyperlink{governing-declaration-ab1cfe37ad8c32b0}{Applied\allowbreak{}Modeling\allowbreak{}Lib.\allowbreak{}Probability.\allowbreak{}Poisson\allowbreak{}Process.\allowbreak{}Good\allowbreak{}Suspension\allowbreak{}State}
\item \hyperlink{governing-declaration-4ca384fa6b6dd395}{Applied\allowbreak{}Modeling\allowbreak{}Lib.\allowbreak{}Probability.\allowbreak{}Poisson\allowbreak{}Process.\allowbreak{}palm\allowbreak{}Tagged\allowbreak{}Arrival\allowbreak{}Indices}
\item \hyperlink{governing-declaration-cf114817e77adc2f}{Applied\allowbreak{}Modeling\allowbreak{}Lib.\allowbreak{}Probability.\allowbreak{}Poisson\allowbreak{}Process.\allowbreak{}suspension\allowbreak{}Base\allowbreak{}Arrival\allowbreak{}Indices}
\item \hyperlink{library-prerequisite-755cdb28ceaf46d2}{Applied\allowbreak{}Modeling\allowbreak{}Lib.\allowbreak{}Queueing.\allowbreak{}Multiclass\allowbreak{}Stationary\allowbreak{}Poisson\allowbreak{}Work\allowbreak{}Class\allowbreak{}Tagged\allowbreak{}Sample}
\item \hyperlink{governing-declaration-62a268164b6ae28b}{Applied\allowbreak{}Modeling\allowbreak{}Lib.\allowbreak{}Queueing.\allowbreak{}Stationary\allowbreak{}Poisson\allowbreak{}Work\allowbreak{}Path}
\item \hyperlink{governing-declaration-05133ca1584020e9}{Applied\allowbreak{}Modeling\allowbreak{}Lib.\allowbreak{}Queueing.\allowbreak{}stationary\allowbreak{}Priority\allowbreak{}Input\allowbreak{}Decidable\allowbreak{}Eq}
\end{itemize}
\hypertarget{governing-declaration-e2ec7ee20660f05d}{}
\subsection*{\small\ttfamily\raggedright Applied\allowbreak{}Modeling\allowbreak{}Lib.\allowbreak{}Queueing.\allowbreak{}canonical\allowbreak{}Stationary\allowbreak{}Priority\allowbreak{}Class\allowbreak{}Tagged\allowbreak{}Arrival\allowbreak{}Window\allowbreak{}Indices}
\textbf{Declaration location:} reusable library
\paragraph{Lean declaration body}
\begin{ReviewVerbatim}
type:
{n : ℕ} →
  (i : Fin n) →
    AppliedModelingLib.Queueing.MulticlassStationaryPoissonWorkClassTaggedSample (Fin n) i →
      ℝ → ℝ → List (AppliedModelingLib.Queueing.NonpreemptivePriorityArrivalIndex n)

value:
fun {n : ℕ} (i : Fin n) (z : AppliedModelingLib.Queueing.MulticlassStationaryPoissonWorkClassTaggedSample (Fin n) i)
    (a b : ℝ) =>
  (AppliedModelingLib.Queueing.nonpreemptivePriorityArrivalWindowIndices
        (AppliedModelingLib.Queueing.stationaryPriorityClassTaggedArrivalWindowIndices i z a b)).sort
    (AppliedModelingLib.Queueing.stationaryPriorityClassTaggedArrivalIndexLE i z)
\end{ReviewVerbatim}

\paragraph{Direct retained dependencies}
\begin{itemize}\raggedright
\item \hyperlink{library-prerequisite-755cdb28ceaf46d2}{Applied\allowbreak{}Modeling\allowbreak{}Lib.\allowbreak{}Queueing.\allowbreak{}Multiclass\allowbreak{}Stationary\allowbreak{}Poisson\allowbreak{}Work\allowbreak{}Class\allowbreak{}Tagged\allowbreak{}Sample}
\item \hyperlink{governing-declaration-5d3724f4dc832d77}{Applied\allowbreak{}Modeling\allowbreak{}Lib.\allowbreak{}Queueing.\allowbreak{}Nonpreemptive\allowbreak{}Priority\allowbreak{}Arrival\allowbreak{}Index}
\item \hyperlink{governing-declaration-2ae4329b7cd91c5b}{Applied\allowbreak{}Modeling\allowbreak{}Lib.\allowbreak{}Queueing.\allowbreak{}nonpreemptive\allowbreak{}Priority\allowbreak{}Arrival\allowbreak{}Window\allowbreak{}Indices}
\item \hyperlink{governing-declaration-6eb7b19311e0b293}{Applied\allowbreak{}Modeling\allowbreak{}Lib.\allowbreak{}Queueing.\allowbreak{}stationary\allowbreak{}Priority\allowbreak{}Class\allowbreak{}Tagged\allowbreak{}Arrival\allowbreak{}Index\allowbreak{}LE}
\item \hyperlink{governing-declaration-9f498af160c98bd8}{Applied\allowbreak{}Modeling\allowbreak{}Lib.\allowbreak{}Queueing.\allowbreak{}stationary\allowbreak{}Priority\allowbreak{}Class\allowbreak{}Tagged\allowbreak{}Arrival\allowbreak{}Window\allowbreak{}Indices}
\end{itemize}
\hypertarget{governing-declaration-c8dff705756357c1}{}
\subsection*{\small\ttfamily\raggedright Applied\allowbreak{}Modeling\allowbreak{}Lib.\allowbreak{}Queueing.\allowbreak{}canonical\allowbreak{}Stationary\allowbreak{}Priority\allowbreak{}Class\allowbreak{}Tagged\allowbreak{}Arrival\allowbreak{}Window\allowbreak{}Jobs}
\textbf{Declaration location:} reusable library
\paragraph{Lean declaration body}
\begin{ReviewVerbatim}
type:
{n : ℕ} →
  (Fin n → ℝ) →
    (i : Fin n) →
      AppliedModelingLib.Queueing.MulticlassStationaryPoissonWorkClassTaggedSample (Fin n) i →
        ℝ →
          ℝ →
            List
              (AppliedModelingLib.Queueing.NonpreemptivePriorityJob n
                (AppliedModelingLib.Queueing.NonpreemptivePriorityArrivalIndex n))

value:
fun {n : ℕ} (meanService : Fin n → ℝ) (i : Fin n)
    (z : AppliedModelingLib.Queueing.MulticlassStationaryPoissonWorkClassTaggedSample (Fin n) i) (a b : ℝ) =>
  List.map
    (fun (q : AppliedModelingLib.Queueing.NonpreemptivePriorityArrivalIndex n) =>
      { identifier := q, priority := q.fst,
        arrivalTime := AppliedModelingLib.Queueing.stationaryPriorityClassTaggedArrival i z q.fst q.snd,
        serviceWork :=
          AppliedModelingLib.Queueing.stationaryPriorityClassTaggedWorkRequirementAt meanService i z q.fst q.snd })
    (AppliedModelingLib.Queueing.canonicalStationaryPriorityClassTaggedArrivalWindowIndices i z a b)
\end{ReviewVerbatim}

\paragraph{Direct retained dependencies}
\begin{itemize}\raggedright
\item \hyperlink{library-prerequisite-755cdb28ceaf46d2}{Applied\allowbreak{}Modeling\allowbreak{}Lib.\allowbreak{}Queueing.\allowbreak{}Multiclass\allowbreak{}Stationary\allowbreak{}Poisson\allowbreak{}Work\allowbreak{}Class\allowbreak{}Tagged\allowbreak{}Sample}
\item \hyperlink{governing-declaration-5d3724f4dc832d77}{Applied\allowbreak{}Modeling\allowbreak{}Lib.\allowbreak{}Queueing.\allowbreak{}Nonpreemptive\allowbreak{}Priority\allowbreak{}Arrival\allowbreak{}Index}
\item \hyperlink{library-prerequisite-2373fe691c32d999}{Applied\allowbreak{}Modeling\allowbreak{}Lib.\allowbreak{}Queueing.\allowbreak{}Nonpreemptive\allowbreak{}Priority\allowbreak{}Job}
\item \hyperlink{governing-declaration-e2ec7ee20660f05d}{Applied\allowbreak{}Modeling\allowbreak{}Lib.\allowbreak{}Queueing.\allowbreak{}canonical\allowbreak{}Stationary\allowbreak{}Priority\allowbreak{}Class\allowbreak{}Tagged\allowbreak{}Arrival\allowbreak{}Window\allowbreak{}Indices}
\item \hyperlink{governing-declaration-2945df1a7a8ec970}{Applied\allowbreak{}Modeling\allowbreak{}Lib.\allowbreak{}Queueing.\allowbreak{}stationary\allowbreak{}Priority\allowbreak{}Class\allowbreak{}Tagged\allowbreak{}Arrival}
\item \hyperlink{governing-declaration-949c09de5ada557f}{Applied\allowbreak{}Modeling\allowbreak{}Lib.\allowbreak{}Queueing.\allowbreak{}stationary\allowbreak{}Priority\allowbreak{}Class\allowbreak{}Tagged\allowbreak{}Work\allowbreak{}Requirement\allowbreak{}At}
\end{itemize}
\hypertarget{governing-declaration-f2f50f38ca2b2a0b}{}
\subsection*{\small\ttfamily\raggedright Applied\allowbreak{}Modeling\allowbreak{}Lib.\allowbreak{}Queueing.\allowbreak{}total\allowbreak{}Priority\allowbreak{}Waiting\allowbreak{}Jobs}
\textbf{Declaration location:} reusable library
\paragraph{Lean declaration body}
\begin{ReviewVerbatim}
type:
{n : ℕ} → {JobId : Type u_1} → AppliedModelingLib.Queueing.NonpreemptivePriorityWorkState n JobId → ℕ

value:
fun {n : ℕ} {JobId : Type u_1} (state : AppliedModelingLib.Queueing.NonpreemptivePriorityWorkState n JobId) =>
  ∑ i : Fin n, (state.waiting i).length
\end{ReviewVerbatim}

\paragraph{Direct retained dependencies}
\begin{itemize}\raggedright
\item \hyperlink{library-prerequisite-2373fe691c32d999}{Applied\allowbreak{}Modeling\allowbreak{}Lib.\allowbreak{}Queueing.\allowbreak{}Nonpreemptive\allowbreak{}Priority\allowbreak{}Job}
\item \hyperlink{library-prerequisite-e68cb5fce009ec53}{Applied\allowbreak{}Modeling\allowbreak{}Lib.\allowbreak{}Queueing.\allowbreak{}Nonpreemptive\allowbreak{}Priority\allowbreak{}Work\allowbreak{}State}
\end{itemize}
\hypertarget{governing-declaration-971183e1bc53251f}{}
\subsection*{\small\ttfamily\raggedright Applied\allowbreak{}Modeling\allowbreak{}Lib.\allowbreak{}finite\allowbreak{}Delay\allowbreak{}Net\allowbreak{}Value}
\textbf{Declaration location:} reusable library
\paragraph{Lean declaration body}
\begin{ReviewVerbatim}
type:
{Class : Type u_1} → [Fintype Class] → (Class → ℝ → ℝ) → (Class → ℝ) → (Class → (Class → ℝ) → ℝ) → (Class → ℝ) → ℝ

value:
fun {Class : Type u_1} [Fintype Class] (value : Class → ℝ → ℝ) (delayCost : Class → ℝ)
    (waitingTime : Class → (Class → ℝ) → ℝ) (x : Class → ℝ) =>
  ∑ j : Class, (value j (x j) - delayCost j * x j * waitingTime j x)
\end{ReviewVerbatim}


\hypertarget{governing-declaration-855778ee64a02a7a}{}
\subsection*{\small\ttfamily\raggedright Applied\allowbreak{}Modeling\allowbreak{}Lib.\allowbreak{}finite\allowbreak{}Marginal\allowbreak{}Externality\allowbreak{}Price}
\textbf{Declaration location:} reusable library
\paragraph{Lean declaration body}
\begin{ReviewVerbatim}
type:
{Class : Type u_1} → [Fintype Class] → (Class → ℝ) → (Class → ℝ) → (Class → ℝ) → ℝ

value:
fun {Class : Type u_1} [Fintype Class] (delayCost flow partialWaiting : Class → ℝ) =>
  ∑ j : Class, delayCost j * flow j * partialWaiting j
\end{ReviewVerbatim}


\hypertarget{governing-declaration-806f5a55d6021334}{}
\subsection*{\small\ttfamily\raggedright Mendelson\allowbreak{}Whang1990\allowbreak{}Priority\allowbreak{}Pricing.\allowbreak{}externality\allowbreak{}Price}
\textbf{Declaration location:} paper-local
\paragraph{Lean declaration body}
\begin{ReviewVerbatim}
type:
{Class : Type u_1} → [Fintype Class] → (Class → ℝ) → (Class → ℝ) → (Class → ℝ) → ℝ

value:
fun {Class : Type u_1} [Fintype Class] (delayCost flow partialWaiting : Class → ℝ) =>
  AppliedModelingLib.finiteMarginalExternalityPrice delayCost flow partialWaiting
\end{ReviewVerbatim}

\paragraph{Direct retained dependencies}
\begin{itemize}\raggedright
\item \hyperlink{governing-declaration-855778ee64a02a7a}{Applied\allowbreak{}Modeling\allowbreak{}Lib.\allowbreak{}finite\allowbreak{}Marginal\allowbreak{}Externality\allowbreak{}Price}
\end{itemize}
\hypertarget{governing-declaration-3df5d95507228c45}{}
\subsection*{\small\ttfamily\raggedright Mendelson\allowbreak{}Whang1990\allowbreak{}Priority\allowbreak{}Pricing.\allowbreak{}paper\allowbreak{}Interface\allowbreak{}Classical\allowbreak{}Decidable\allowbreak{}Eq}
\textbf{Declaration location:} paper-local
\paragraph{Lean declaration body}
\begin{ReviewVerbatim}
type:
(α : Type u_1) → (a b : α) → Decidable (a = b)

value:
fun (α : Type u_1) (a b : α) => Classical.decEq α a b
\end{ReviewVerbatim}


\hypertarget{governing-declaration-96f4c7bfcd6904cc}{}
\subsection*{\small\ttfamily\raggedright Mendelson\allowbreak{}Whang1990\allowbreak{}Priority\allowbreak{}Pricing.\allowbreak{}priority\allowbreak{}Linear\allowbreak{}Numerator}
\textbf{Declaration location:} paper-local
\paragraph{Lean declaration body}
\begin{ReviewVerbatim}
type:
{n : ℕ} → (Fin n → ℝ) → (Fin n → ℝ) → (Fin n → ℝ) → Fin n → ℝ

value:
fun {n : ℕ} (arrivalRate delayCost meanService : Fin n → ℝ) (i : Fin n) =>
  delayCost i * arrivalRate i * AppliedModelingLib.Queueing.finitePriorityResidualWork arrivalRate meanService
\end{ReviewVerbatim}

\paragraph{Direct retained dependencies}
\begin{itemize}\raggedright
\item \hyperlink{governing-declaration-80f4f37f124e636e}{Applied\allowbreak{}Modeling\allowbreak{}Lib.\allowbreak{}Queueing.\allowbreak{}finite\allowbreak{}Priority\allowbreak{}Residual\allowbreak{}Work}
\end{itemize}
\hypertarget{governing-declaration-41afc2346e61f660}{}
\subsection*{\small\ttfamily\raggedright Mendelson\allowbreak{}Whang1990\allowbreak{}Priority\allowbreak{}Pricing.\allowbreak{}priority\allowbreak{}Sojourn\allowbreak{}Derivative}
\textbf{Declaration location:} paper-local
\paragraph{Lean declaration body}
\begin{ReviewVerbatim}
type:
{n : ℕ} → (Fin n → ℝ) → (Fin n → ℝ) → Fin n → Fin n → ℝ

value:
fun {n : ℕ} (arrivalRate meanService : Fin n → ℝ) (arrivalClass priorityClass : Fin n) =>
  (meanService arrivalClass ^ 2 *
        (AppliedModelingLib.Queueing.finitePriorityStrictSlack meanService arrivalRate priorityClass *
          AppliedModelingLib.Queueing.finitePriorityInclusiveSlack meanService arrivalRate priorityClass) -
      AppliedModelingLib.Queueing.finitePriorityResidualWork arrivalRate meanService *
        ((-if arrivalClass < priorityClass then meanService arrivalClass else 0) *
            AppliedModelingLib.Queueing.finitePriorityInclusiveSlack meanService arrivalRate priorityClass +
          AppliedModelingLib.Queueing.finitePriorityStrictSlack meanService arrivalRate priorityClass *
            -if arrivalClass ≤ priorityClass then meanService arrivalClass else 0)) /
    (AppliedModelingLib.Queueing.finitePriorityStrictSlack meanService arrivalRate priorityClass *
        AppliedModelingLib.Queueing.finitePriorityInclusiveSlack meanService arrivalRate priorityClass) ^
      2
\end{ReviewVerbatim}

\paragraph{Direct retained dependencies}
\begin{itemize}\raggedright
\item \hyperlink{governing-declaration-f2b9a1af307e7469}{Applied\allowbreak{}Modeling\allowbreak{}Lib.\allowbreak{}Queueing.\allowbreak{}finite\allowbreak{}Priority\allowbreak{}Inclusive\allowbreak{}Slack}
\item \hyperlink{governing-declaration-80f4f37f124e636e}{Applied\allowbreak{}Modeling\allowbreak{}Lib.\allowbreak{}Queueing.\allowbreak{}finite\allowbreak{}Priority\allowbreak{}Residual\allowbreak{}Work}
\item \hyperlink{governing-declaration-3f5cf610a8190094}{Applied\allowbreak{}Modeling\allowbreak{}Lib.\allowbreak{}Queueing.\allowbreak{}finite\allowbreak{}Priority\allowbreak{}Strict\allowbreak{}Slack}
\end{itemize}
\hypertarget{governing-declaration-703b87082556903f}{}
\subsection*{\small\ttfamily\raggedright Mendelson\allowbreak{}Whang1990\allowbreak{}Priority\allowbreak{}Pricing.\allowbreak{}priority\allowbreak{}Sojourn\allowbreak{}Time}
\textbf{Declaration location:} paper-local
\paragraph{Lean declaration body}
\begin{ReviewVerbatim}
type:
{n : ℕ} → (Fin n → ℝ) → (Fin n → ℝ) → Fin n → ℝ

value:
fun {n : ℕ} (arrivalRate meanService : Fin n → ℝ) (i : Fin n) =>
  AppliedModelingLib.Queueing.finiteNonpreemptivePrioritySojournTime arrivalRate meanService i
\end{ReviewVerbatim}

\paragraph{Direct retained dependencies}
\begin{itemize}\raggedright
\item \hyperlink{governing-declaration-ab8863aef2e4b99f}{Applied\allowbreak{}Modeling\allowbreak{}Lib.\allowbreak{}Queueing.\allowbreak{}finite\allowbreak{}Nonpreemptive\allowbreak{}Priority\allowbreak{}Sojourn\allowbreak{}Time}
\end{itemize}
\hypertarget{governing-declaration-65f779b7d3134146}{}
\subsection*{\small\ttfamily\raggedright Mendelson\allowbreak{}Whang1990\allowbreak{}Priority\allowbreak{}Pricing.\allowbreak{}priority\allowbreak{}Population}
\textbf{Declaration location:} paper-local
\paragraph{Lean declaration body}
\begin{ReviewVerbatim}
type:
{n : ℕ} → (Fin n → ℝ) → (Fin n → ℝ) → Fin n → ℝ

value:
fun {n : ℕ} (arrivalRate meanService : Fin n → ℝ) (i : Fin n) =>
  arrivalRate i * MendelsonWhang1990PriorityPricing.prioritySojournTime arrivalRate meanService i
\end{ReviewVerbatim}

\paragraph{Direct retained dependencies}
\begin{itemize}\raggedright
\item \hyperlink{governing-declaration-703b87082556903f}{Mendelson\allowbreak{}Whang1990\allowbreak{}Priority\allowbreak{}Pricing.\allowbreak{}priority\allowbreak{}Sojourn\allowbreak{}Time}
\end{itemize}
\hypertarget{governing-declaration-c9c37e19702a9c85}{}
\subsection*{\small\ttfamily\raggedright Mendelson\allowbreak{}Whang1990\allowbreak{}Priority\allowbreak{}Pricing.\allowbreak{}priority\allowbreak{}Schedule\allowbreak{}Sojourn\allowbreak{}Time}
\textbf{Declaration location:} paper-local
\paragraph{Lean declaration body}
\begin{ReviewVerbatim}
type:
{n : ℕ} → (Fin n → ℝ) → (Fin n → ℝ) → Equiv.Perm (Fin n) → Fin n → ℝ

value:
fun {n : ℕ} (arrivalRate meanService : Fin n → ℝ) (schedule : Equiv.Perm (Fin n)) (customer : Fin n) =>
  MendelsonWhang1990PriorityPricing.prioritySojournTime (arrivalRate ∘ ⇑schedule) (meanService ∘ ⇑schedule)
    ((Equiv.symm schedule) customer)
\end{ReviewVerbatim}

\paragraph{Direct retained dependencies}
\begin{itemize}\raggedright
\item \hyperlink{governing-declaration-703b87082556903f}{Mendelson\allowbreak{}Whang1990\allowbreak{}Priority\allowbreak{}Pricing.\allowbreak{}priority\allowbreak{}Sojourn\allowbreak{}Time}
\end{itemize}
\hypertarget{governing-declaration-ac7edc3a22d8e267}{}
\subsection*{\small\ttfamily\raggedright Mendelson\allowbreak{}Whang1990\allowbreak{}Priority\allowbreak{}Pricing.\allowbreak{}priority\allowbreak{}Schedule\allowbreak{}Delay\allowbreak{}Cost}
\textbf{Declaration location:} paper-local
\paragraph{Lean declaration body}
\begin{ReviewVerbatim}
type:
{n : ℕ} → (Fin n → ℝ) → (Fin n → ℝ) → (Fin n → ℝ) → Equiv.Perm (Fin n) → ℝ

value:
fun {n : ℕ} (arrivalRate delayCost meanService : Fin n → ℝ) (schedule : Equiv.Perm (Fin n)) =>
  ∑ customer : Fin n,
    delayCost customer * arrivalRate customer *
      MendelsonWhang1990PriorityPricing.priorityScheduleSojournTime arrivalRate meanService schedule customer
\end{ReviewVerbatim}

\paragraph{Direct retained dependencies}
\begin{itemize}\raggedright
\item \hyperlink{governing-declaration-c9c37e19702a9c85}{Mendelson\allowbreak{}Whang1990\allowbreak{}Priority\allowbreak{}Pricing.\allowbreak{}priority\allowbreak{}Schedule\allowbreak{}Sojourn\allowbreak{}Time}
\end{itemize}
\hypertarget{governing-declaration-2d648b9cf4c8cf80}{}
\subsection*{\small\ttfamily\raggedright Mendelson\allowbreak{}Whang1990\allowbreak{}Priority\allowbreak{}Pricing.\allowbreak{}priority\allowbreak{}Schedule\allowbreak{}Net\allowbreak{}Value}
\textbf{Declaration location:} paper-local
\paragraph{Lean declaration body}
\begin{ReviewVerbatim}
type:
{n : ℕ} → (Fin n → ℝ → ℝ) → (Fin n → ℝ) → (Fin n → ℝ) → (Fin n → ℝ) → Equiv.Perm (Fin n) → ℝ

value:
fun {n : ℕ} (value : Fin n → ℝ → ℝ) (arrivalRate delayCost meanService : Fin n → ℝ) (schedule : Equiv.Perm (Fin n)) =>
  ∑ customer : Fin n, value customer (arrivalRate customer) -
    MendelsonWhang1990PriorityPricing.priorityScheduleDelayCost arrivalRate delayCost meanService schedule
\end{ReviewVerbatim}

\paragraph{Direct retained dependencies}
\begin{itemize}\raggedright
\item \hyperlink{governing-declaration-ac7edc3a22d8e267}{Mendelson\allowbreak{}Whang1990\allowbreak{}Priority\allowbreak{}Pricing.\allowbreak{}priority\allowbreak{}Schedule\allowbreak{}Delay\allowbreak{}Cost}
\end{itemize}
\hypertarget{governing-declaration-f2dbc4d52705c26d}{}
\subsection*{\small\ttfamily\raggedright Mendelson\allowbreak{}Whang1990\allowbreak{}Priority\allowbreak{}Pricing.\allowbreak{}priority\allowbreak{}Successor\allowbreak{}Or\allowbreak{}Zero}
\textbf{Declaration location:} paper-local
\paragraph{Lean declaration body}
\begin{ReviewVerbatim}
type:
{n : ℕ} → (Fin n → ℝ) → Fin n → ℝ

value:
fun {n : ℕ} (value : Fin n → ℝ) (i : Fin n) => if h : ↑i + 1 < n then value ⟨↑i + 1, h⟩ else 0
\end{ReviewVerbatim}


\clearpage
\hypertarget{library-section-994978aceb8dd183}{}
\section*{Material library prerequisites}
\hypertarget{library-prerequisite-dba97dbb9d8c7c91}{}
\subsection*{\small\ttfamily\raggedright Applied\allowbreak{}Modeling\allowbreak{}Lib.\allowbreak{}Probability.\allowbreak{}Poisson\allowbreak{}Process.\allowbreak{}multiclass\allowbreak{}Stationary\allowbreak{}Poisson\allowbreak{}Work\allowbreak{}Measure}
\textbf{Source locator:} {\footnotesize\raggedright cited publication, lines 227-\allowbreak{}242\par}
\paragraph{Verbatim paper-source connection}
\begin{ReviewVerbatim}
giving rise to X = A (p) or V'(X) = 4-'(X/A) = p.                     Al. Arrival Processes

  Consider the case where the delay cost v is positive.               The arrival process of class-i jobs to the system, which
Here, the cost faced by an individual user consists of                reflects the aggregation of infinitesimal users' job flow,
two components: a direct payment of p per job, and                    follows a stationary Poisson process with a price-




                                                 This content downloaded from
                               108.30.55.84 on Mon, 24 Aug 2026 10:55:53 UTC
                                         All use subject to https://about.jstor.org/terms

[form-feed in source extraction]
                                                                                Priority Pricingfor the MIMI1 Queue / 873

dependent arrival rate Xi. The R arrival processes are                 by the system (per unit of time) to the organization,
stochastically independent.                                            where we assume for now that the system administra-

[Next verbatim source excerpt]

 Lq *)= The mean number of class-i jobs in the                          KLEINROCK, L. 1967. Optimum Bribing for Queue Posi-
         queue in steady state;                                              tion. Opns. Res. 15, 304-318.
Wi(.) = The expected waiting time of a class-i job                      KLEINROCK, L. 1976. Queueing System, Vols. I and II.
         in the queue in steady state;                                       John Wiley, New York.
 Wi ( )= The expected sojourn time of a class-i job                     KNUDSEN, N. C. 1972. Individual and Social Optimiza-
         in the system in steady state.                                      tion in a Multi-Server Queue With a General Cost-
\end{ReviewVerbatim}



\paragraph{Lean-expanded library semantic target}
\begin{ReviewVerbatim}
type:
{Class : Type u_1} →
  [Fintype Class] →
    (Class → ℝ) →
      MeasureTheory.Measure (Class → AppliedModelingLib.Probability.PoissonProcess.GoodSuspensionState × (ℤ → ℝ))

value:
fun {Class : Type u_1} [Fintype Class] (arrivalRate : Class → ℝ) =>
  MeasureTheory.Measure.pi fun (i : Class) =>
    AppliedModelingLib.Probability.Queueing.stationaryPoissonWorkMeasure (arrivalRate i)
\end{ReviewVerbatim}

\paragraph{Exact Lean library declaration}
\begin{ReviewVerbatim}
/-- One independently sampled stationary marked-Poisson input per class.
Each coordinate contains both its stationary arrival suspension and a complete
two-sided iid unit-exponential work-mark path. -/
noncomputable def multiclassStationaryPoissonWorkMeasure (arrivalRate : Class -> Real) :
    Measure (Class -> (GoodSuspensionState × (Int -> Real))) :=
  Measure.pi fun i => Queueing.stationaryPoissonWorkMeasure (arrivalRate i)
\end{ReviewVerbatim}

\paragraph{Direct retained dependencies}
\begin{itemize}\raggedright
\item \hyperlink{governing-declaration-ab1cfe37ad8c32b0}{Applied\allowbreak{}Modeling\allowbreak{}Lib.\allowbreak{}Probability.\allowbreak{}Poisson\allowbreak{}Process.\allowbreak{}Good\allowbreak{}Suspension\allowbreak{}State}
\item \hyperlink{governing-declaration-a8385d799d0a4d16}{Applied\allowbreak{}Modeling\allowbreak{}Lib.\allowbreak{}Probability.\allowbreak{}Poisson\allowbreak{}Process.\allowbreak{}inst\allowbreak{}Measurable\allowbreak{}Space\allowbreak{}Good\allowbreak{}Suspension\allowbreak{}State}
\item \hyperlink{governing-declaration-10a259bddd1a9d6f}{Applied\allowbreak{}Modeling\allowbreak{}Lib.\allowbreak{}Probability.\allowbreak{}Queueing.\allowbreak{}stationary\allowbreak{}Poisson\allowbreak{}Work\allowbreak{}Measure}
\end{itemize}
\paragraph{Recorded source-to-library screening}
\textbf{Verdict:} \texttt{matches}
\\
{\raggedright
\textbf{Reason:} A1's independent stationary class inputs are represented by the product stationary marked-\allowbreak{}Poisson work measure;\allowbreak{} work marks are the service-\allowbreak{}process carrier, not an added claim.\allowbreak{}
\par}
\reviewmatch{Matches source input}
\noindent\textbf{Reviewer annotation}\par
\reviewerbox
\clearpage
\hypertarget{library-prerequisite-4b8788f2f9337ce3}{}
\subsection*{\small\ttfamily\raggedright Applied\allowbreak{}Modeling\allowbreak{}Lib.\allowbreak{}Probability.\allowbreak{}stabilized\allowbreak{}Finite\allowbreak{}Replay\allowbreak{}Response}
\textbf{Source locator:} {\footnotesize\raggedright cited publication, lines 1118-\allowbreak{}1129\par}
\paragraph{Verbatim paper-source connection}
\begin{ReviewVerbatim}
Wi(.) = The expected waiting time of a class-i job                      KLEINROCK, L. 1976. Queueing System, Vols. I and II.
         in the queue in steady state;                                       John Wiley, New York.
 Wi ( )= The expected sojourn time of a class-i job                     KNUDSEN, N. C. 1972. Individual and Social Optimiza-
         in the system in steady state.                                      tion in a Multi-Server Queue With a General Cost-
                                                                             Benefit Structure. Econometrica 40, 515-528.
For the M/M/1 queue with nonpreemptive priority                         KREPS, D. 1990. A Course in Microeconomic Theory.
discipline (Heyman and Sobel 1982, p. 435)                                 Princeton University Press, Princeton, N.J.
                                                                        MARCHAND, M. 1974. Priority Pricing. Mgmt. Sci. 20,
   _ i( L1()
 W1CX)  = W1 Wqc1+ -+
                   Ci AR   (A-)
                      cj. (-                                                  1131-1140.

[Next verbatim source excerpt]

A4. Service Process                                                      Vi'(Xi) = pi + vi Wi (5)
The processor is modeled as an M/M/1 queueing                          where (5) is analogous to (1) with total expected class-
system with nonpreemptive priority among job                           i user cost consisting of two components: a direct
classes. The processing requirement Hi(-) of a class-i                 payment of pi per job, and the expected class-i delay
job, ri, follows an exponential distribution with finite cost, given by vi * Wi(X).
mean ci. Processing requirements are independent           The following theorem generalizes the net value
across jobs and identically distributed within job       maximizing pricing scheme, (3), to the multiclass case.
classes.

[Next verbatim source excerpt]

giving rise to X = A (p) or V'(X) = 4-'(X/A) = p.                     Al. Arrival Processes

  Consider the case where the delay cost v is positive.               The arrival process of class-i jobs to the system, which
Here, the cost faced by an individual user consists of                reflects the aggregation of infinitesimal users' job flow,
two components: a direct payment of p per job, and                    follows a stationary Poisson process with a price-




                                                 This content downloaded from
                               108.30.55.84 on Mon, 24 Aug 2026 10:55:53 UTC
                                         All use subject to https://about.jstor.org/terms

[form-feed in source extraction]
                                                                                Priority Pricingfor the MIMI1 Queue / 873

dependent arrival rate Xi. The R arrival processes are                 by the system (per unit of time) to the organization,
stochastically independent.                                            where we assume for now that the system administra-
\end{ReviewVerbatim}



\paragraph{Lean-expanded library semantic target}
\begin{ReviewVerbatim}
type:
{Ω : Type u_1} → (ℕ → Ω → ℝ) → Ω → ℝ

value:
fun {Ω : Type u_1} (finiteResponse : ℕ → Ω → ℝ) (a : Ω) =>
  Filter.limsup (fun (n : ℕ) => finiteResponse n a) Filter.atTop
\end{ReviewVerbatim}

\paragraph{Exact Lean library declaration}
\begin{ReviewVerbatim}
/-- The canonical response induced by a sequence of total finite replay
responses.  Under eventual equality it is the eventual literal value. -/
noncomputable def stabilizedFiniteReplayResponse
    (finiteResponse : Nat → Ω → ℝ) : Ω → ℝ :=
  fun ω => limsup (fun n => finiteResponse n ω) atTop
\end{ReviewVerbatim}

\paragraph{Recorded source-to-library screening}
\textbf{Verdict:} \texttt{matches}
\\
{\raggedright
\textbf{Reason:} The expanded value takes the limsup of total finite replay response observations.\allowbreak{} On the proved eventual-\allowbreak{}stabilization event this equals the literal stationary tagged response, so it is a total realization of the steady-\allowbreak{}state response quantity used in Appendix (A-\allowbreak{}1), not a replacement queue formula.\allowbreak{}
\par}
\reviewmatch{Matches source input}
\noindent\textbf{Reviewer annotation}\par
\reviewerbox
\clearpage
\hypertarget{library-prerequisite-755cdb28ceaf46d2}{}
\subsection*{\small\ttfamily\raggedright Applied\allowbreak{}Modeling\allowbreak{}Lib.\allowbreak{}Queueing.\allowbreak{}Multiclass\allowbreak{}Stationary\allowbreak{}Poisson\allowbreak{}Work\allowbreak{}Class\allowbreak{}Tagged\allowbreak{}Sample}
\textbf{Source locator:} {\footnotesize\raggedright cited publication, lines 227-\allowbreak{}242\par}
\paragraph{Verbatim paper-source connection}
\begin{ReviewVerbatim}
giving rise to X = A (p) or V'(X) = 4-'(X/A) = p.                     Al. Arrival Processes

  Consider the case where the delay cost v is positive.               The arrival process of class-i jobs to the system, which
Here, the cost faced by an individual user consists of                reflects the aggregation of infinitesimal users' job flow,
two components: a direct payment of p per job, and                    follows a stationary Poisson process with a price-




                                                 This content downloaded from
                               108.30.55.84 on Mon, 24 Aug 2026 10:55:53 UTC
                                         All use subject to https://about.jstor.org/terms

[form-feed in source extraction]
                                                                                Priority Pricingfor the MIMI1 Queue / 873

dependent arrival rate Xi. The R arrival processes are                 by the system (per unit of time) to the organization,
stochastically independent.                                            where we assume for now that the system administra-

[Next verbatim source excerpt]

 Lq *)= The mean number of class-i jobs in the                          KLEINROCK, L. 1967. Optimum Bribing for Queue Posi-
         queue in steady state;                                              tion. Opns. Res. 15, 304-318.
Wi(.) = The expected waiting time of a class-i job                      KLEINROCK, L. 1976. Queueing System, Vols. I and II.
         in the queue in steady state;                                       John Wiley, New York.
 Wi ( )= The expected sojourn time of a class-i job                     KNUDSEN, N. C. 1972. Individual and Social Optimiza-
         in the system in steady state.                                      tion in a Multi-Server Queue With a General Cost-
\end{ReviewVerbatim}



\paragraph{Lean-expanded library semantic target}
\begin{ReviewVerbatim}
type:
(Class : Type u_2) → Class → Type (max 0 u_2)

value:
fun (Class : Type u_2) (i : Class) =>
  ((ℤ → ℝ) × (ℤ → ℝ)) × ({ j : Class // j ≠ i } → AppliedModelingLib.Queueing.StationaryPoissonWorkPath)
\end{ReviewVerbatim}

\paragraph{Exact Lean library declaration}
\begin{ReviewVerbatim}
/-- The complete input carrier seen from one selected class arrival.  The
selected class uses its Palm-recentered arrival and mark paths, while every
other class retains a shifted stationary marked-Poisson path. -/
abbrev MulticlassStationaryPoissonWorkClassTaggedSample
    (Class : Type*) (i : Class) :=
  ((Int → Real) × (Int → Real)) ×
    ({j : Class // j ≠ i} → StationaryPoissonWorkPath)
\end{ReviewVerbatim}

\paragraph{Direct retained dependencies}
\begin{itemize}\raggedright
\item \hyperlink{governing-declaration-62a268164b6ae28b}{Applied\allowbreak{}Modeling\allowbreak{}Lib.\allowbreak{}Queueing.\allowbreak{}Stationary\allowbreak{}Poisson\allowbreak{}Work\allowbreak{}Path}
\end{itemize}
\paragraph{Recorded source-to-library screening}
\textbf{Verdict:} \texttt{matches}
\\
{\raggedright
\textbf{Reason:} This is the A1 selected-\allowbreak{}arrival/\allowbreak{}Palm carrier:\allowbreak{} the selected class is recentered at its arrival and other classes remain stationary.\allowbreak{}
\par}
\reviewmatch{Matches source input}
\noindent\textbf{Reviewer annotation}\par
\reviewerbox
\clearpage
\hypertarget{library-prerequisite-2373fe691c32d999}{}
\subsection*{\small\ttfamily\raggedright Applied\allowbreak{}Modeling\allowbreak{}Lib.\allowbreak{}Queueing.\allowbreak{}Nonpreemptive\allowbreak{}Priority\allowbreak{}Job}
\textbf{Source locator:} {\footnotesize\raggedright cited publication, lines 264-\allowbreak{}271\par}
\paragraph{Verbatim paper-source connection}
\begin{ReviewVerbatim}
A4. Service Process                                                      Vi'(Xi) = pi + vi Wi (5)
The processor is modeled as an M/M/1 queueing                          where (5) is analogous to (1) with total expected class-
system with nonpreemptive priority among job                           i user cost consisting of two components: a direct
classes. The processing requirement Hi(-) of a class-i                 payment of pi per job, and the expected class-i delay
job, ri, follows an exponential distribution with finite cost, given by vi * Wi(X).
mean ci. Processing requirements are independent           The following theorem generalizes the net value
across jobs and identically distributed within job       maximizing pricing scheme, (3), to the multiclass case.
classes.
\end{ReviewVerbatim}



\paragraph{Lean declaration type (metadata)}
\begin{ReviewVerbatim}
field identifier:
{n : ℕ} → {JobId : Type u_1} → AppliedModelingLib.Queueing.NonpreemptivePriorityJob n JobId → JobId

field priority:
{n : ℕ} → {JobId : Type u_1} → AppliedModelingLib.Queueing.NonpreemptivePriorityJob n JobId → Fin n

field arrivalTime:
{n : ℕ} → {JobId : Type u_1} → AppliedModelingLib.Queueing.NonpreemptivePriorityJob n JobId → ℝ

field serviceWork:
{n : ℕ} → {JobId : Type u_1} → AppliedModelingLib.Queueing.NonpreemptivePriorityJob n JobId → ℝ
\end{ReviewVerbatim}

\paragraph{Exact Lean library declaration}
\begin{ReviewVerbatim}
/-- A finite-trace customer, retaining its identifier, declared class,
physical arrival time, and required service work. -/
structure NonpreemptivePriorityJob (n : ℕ) (JobId : Type*) where
  identifier : JobId
  priority : Fin n
  arrivalTime : ℝ
  serviceWork : ℝ
\end{ReviewVerbatim}

\paragraph{Recorded source-to-library screening}
\textbf{Verdict:} \texttt{matches}
\\
{\raggedright
\textbf{Reason:} The expanded record retains exactly the data of one A4 queue customer needed by the nonpreemptive-\allowbreak{}priority dynamics:\allowbreak{} identity, class priority, arrival epoch, and service requirement.\allowbreak{} It introduces no scheduling behavior by itself.\allowbreak{}
\par}
\reviewmatch{Matches source input}
\noindent\textbf{Reviewer annotation}\par
\reviewerbox
\clearpage
\hypertarget{library-prerequisite-e68cb5fce009ec53}{}
\subsection*{\small\ttfamily\raggedright Applied\allowbreak{}Modeling\allowbreak{}Lib.\allowbreak{}Queueing.\allowbreak{}Nonpreemptive\allowbreak{}Priority\allowbreak{}Work\allowbreak{}State}
\textbf{Source locator:} {\footnotesize\raggedright cited publication, lines 264-\allowbreak{}271\par}
\paragraph{Verbatim paper-source connection}
\begin{ReviewVerbatim}
A4. Service Process                                                      Vi'(Xi) = pi + vi Wi (5)
The processor is modeled as an M/M/1 queueing                          where (5) is analogous to (1) with total expected class-
system with nonpreemptive priority among job                           i user cost consisting of two components: a direct
classes. The processing requirement Hi(-) of a class-i                 payment of pi per job, and the expected class-i delay
job, ri, follows an exponential distribution with finite cost, given by vi * Wi(X).
mean ci. Processing requirements are independent           The following theorem generalizes the net value
across jobs and identically distributed within job       maximizing pricing scheme, (3), to the multiclass case.
classes.
\end{ReviewVerbatim}



\paragraph{Lean declaration type (metadata)}
\begin{ReviewVerbatim}
field currentTime:
{n : ℕ} → {JobId : Type u_1} → AppliedModelingLib.Queueing.NonpreemptivePriorityWorkState n JobId → ℝ

field active:
{n : ℕ} →
  {JobId : Type u_1} →
    AppliedModelingLib.Queueing.NonpreemptivePriorityWorkState n JobId →
      Option (AppliedModelingLib.Queueing.NonpreemptivePriorityJob n JobId × ℝ)

field waiting:
{n : ℕ} →
  {JobId : Type u_1} →
    AppliedModelingLib.Queueing.NonpreemptivePriorityWorkState n JobId →
      Fin n → List (AppliedModelingLib.Queueing.NonpreemptivePriorityJob n JobId)

field completed:
{n : ℕ} →
  {JobId : Type u_1} →
    AppliedModelingLib.Queueing.NonpreemptivePriorityWorkState n JobId →
      List (AppliedModelingLib.Queueing.NonpreemptivePriorityJob n JobId × ℝ)
\end{ReviewVerbatim}

\paragraph{Exact Lean library declaration}
\begin{ReviewVerbatim}
/-- A work-level queue state.  The active pair stores the currently served
job and its residual work; `waiting` consists of one FIFO list per class; and
`completed` records each completed job with its physical completion time. -/
structure NonpreemptivePriorityWorkState (n : ℕ) (JobId : Type*) where
  currentTime : ℝ
  active : Option (NonpreemptivePriorityJob n JobId × ℝ)
  waiting : Fin n → List (NonpreemptivePriorityJob n JobId)
  completed : List (NonpreemptivePriorityJob n JobId × ℝ)
\end{ReviewVerbatim}

\paragraph{Direct retained dependencies}
\begin{itemize}\raggedright
\item \hyperlink{library-prerequisite-2373fe691c32d999}{Applied\allowbreak{}Modeling\allowbreak{}Lib.\allowbreak{}Queueing.\allowbreak{}Nonpreemptive\allowbreak{}Priority\allowbreak{}Job}
\end{itemize}
\paragraph{Recorded source-to-library screening}
\textbf{Verdict:} \texttt{matches}
\\
{\raggedright
\textbf{Reason:} The state records A4's one-\allowbreak{}server active/\allowbreak{}residual-\allowbreak{}work and class-\allowbreak{}FIFO nonpreemptive-\allowbreak{}priority queue discipline.\allowbreak{}
\par}
\reviewmatch{Matches source input}
\noindent\textbf{Reviewer annotation}\par
\reviewerbox
\clearpage
\hypertarget{library-prerequisite-837de7f0cc179814}{}
\subsection*{\small\ttfamily\raggedright Applied\allowbreak{}Modeling\allowbreak{}Lib.\allowbreak{}Queueing.\allowbreak{}admit\allowbreak{}Nonpreemptive\allowbreak{}Priority\allowbreak{}Job}
\textbf{Source locator:} {\footnotesize\raggedright cited publication, lines 264-\allowbreak{}271\par}
\paragraph{Verbatim paper-source connection}
\begin{ReviewVerbatim}
A4. Service Process                                                      Vi'(Xi) = pi + vi Wi (5)
The processor is modeled as an M/M/1 queueing                          where (5) is analogous to (1) with total expected class-
system with nonpreemptive priority among job                           i user cost consisting of two components: a direct
classes. The processing requirement Hi(-) of a class-i                 payment of pi per job, and the expected class-i delay
job, ri, follows an exponential distribution with finite cost, given by vi * Wi(X).
mean ci. Processing requirements are independent           The following theorem generalizes the net value
across jobs and identically distributed within job       maximizing pricing scheme, (3), to the multiclass case.
classes.
\end{ReviewVerbatim}



\paragraph{Lean-expanded library semantic target}
\begin{ReviewVerbatim}
type:
{n : ℕ} →
  {JobId : Type u_1} →
    AppliedModelingLib.Queueing.NonpreemptivePriorityWorkState n JobId →
      AppliedModelingLib.Queueing.NonpreemptivePriorityJob n JobId →
        AppliedModelingLib.Queueing.NonpreemptivePriorityWorkState n JobId

value:
fun {n : ℕ} {JobId : Type u_1} (state : AppliedModelingLib.Queueing.NonpreemptivePriorityWorkState n JobId)
    (job : AppliedModelingLib.Queueing.NonpreemptivePriorityJob n JobId) =>
  have prepared : AppliedModelingLib.Queueing.NonpreemptivePriorityWorkState n JobId :=
    AppliedModelingLib.Queueing.startNextNonpreemptivePriorityJob state;
  match prepared.active with
  | none =>
    { currentTime := prepared.currentTime, active := some (job, job.serviceWork), waiting := prepared.waiting,
      completed := prepared.completed }
  | some val => AppliedModelingLib.Queueing.enqueueNonpreemptivePriorityJob prepared job
\end{ReviewVerbatim}

\paragraph{Exact Lean library declaration}
\begin{ReviewVerbatim}
/-- An arrival joins service immediately only when the server is idle after
starting any previously waiting work; otherwise it joins the tail of its own
class-FIFO list. -/
noncomputable def admitNonpreemptivePriorityJob
    {n : ℕ} {JobId : Type*}
    (state : NonpreemptivePriorityWorkState n JobId)
    (job : NonpreemptivePriorityJob n JobId) :
    NonpreemptivePriorityWorkState n JobId :=
  by
    classical
    let prepared := startNextNonpreemptivePriorityJob state
    exact match prepared.active with
    | none => { prepared with active := some (job, job.serviceWork) }
    | some _ => enqueueNonpreemptivePriorityJob prepared job
\end{ReviewVerbatim}

\paragraph{Direct retained dependencies}
\begin{itemize}\raggedright
\item \hyperlink{library-prerequisite-2373fe691c32d999}{Applied\allowbreak{}Modeling\allowbreak{}Lib.\allowbreak{}Queueing.\allowbreak{}Nonpreemptive\allowbreak{}Priority\allowbreak{}Job}
\item \hyperlink{library-prerequisite-e68cb5fce009ec53}{Applied\allowbreak{}Modeling\allowbreak{}Lib.\allowbreak{}Queueing.\allowbreak{}Nonpreemptive\allowbreak{}Priority\allowbreak{}Work\allowbreak{}State}
\item \hyperlink{governing-declaration-d385a019cc9e3344}{Applied\allowbreak{}Modeling\allowbreak{}Lib.\allowbreak{}Queueing.\allowbreak{}enqueue\allowbreak{}Nonpreemptive\allowbreak{}Priority\allowbreak{}Job}
\item \hyperlink{governing-declaration-0dd8378f656d740f}{Applied\allowbreak{}Modeling\allowbreak{}Lib.\allowbreak{}Queueing.\allowbreak{}start\allowbreak{}Next\allowbreak{}Nonpreemptive\allowbreak{}Priority\allowbreak{}Job}
\end{itemize}
\paragraph{Recorded source-to-library screening}
\textbf{Verdict:} \texttt{matches}
\\
{\raggedright
\textbf{Reason:} This is A4's nonpreemptive admission and within-\allowbreak{}class FIFO enqueue operation.\allowbreak{}
\par}
\reviewmatch{Matches source input}
\noindent\textbf{Reviewer annotation}\par
\reviewerbox
\clearpage
\hypertarget{library-prerequisite-96200a424cf39dee}{}
\subsection*{\small\ttfamily\raggedright Applied\allowbreak{}Modeling\allowbreak{}Lib.\allowbreak{}Queueing.\allowbreak{}advance\allowbreak{}Nonpreemptive\allowbreak{}Priority\allowbreak{}Work\allowbreak{}State}
\textbf{Source locator:} {\footnotesize\raggedright cited publication, lines 264-\allowbreak{}271\par}
\paragraph{Verbatim paper-source connection}
\begin{ReviewVerbatim}
A4. Service Process                                                      Vi'(Xi) = pi + vi Wi (5)
The processor is modeled as an M/M/1 queueing                          where (5) is analogous to (1) with total expected class-
system with nonpreemptive priority among job                           i user cost consisting of two components: a direct
classes. The processing requirement Hi(-) of a class-i                 payment of pi per job, and the expected class-i delay
job, ri, follows an exponential distribution with finite cost, given by vi * Wi(X).
mean ci. Processing requirements are independent           The following theorem generalizes the net value
across jobs and identically distributed within job       maximizing pricing scheme, (3), to the multiclass case.
classes.
\end{ReviewVerbatim}



\paragraph{Lean-expanded library semantic target}
\begin{ReviewVerbatim}
type:
{n : ℕ} →
  {JobId : Type u_1} →
    ℕ →
      ℝ →
        AppliedModelingLib.Queueing.NonpreemptivePriorityWorkState n JobId →
          AppliedModelingLib.Queueing.NonpreemptivePriorityWorkState n JobId

value:
fun {n : ℕ} {JobId : Type u_1} (a : ℕ) (a_1 : ℝ)
    (a_2 : AppliedModelingLib.Queueing.NonpreemptivePriorityWorkState n JobId) =>
  Nat.brecOn (motive := fun (x : ℕ) =>
    ℝ →
      AppliedModelingLib.Queueing.NonpreemptivePriorityWorkState n JobId →
        AppliedModelingLib.Queueing.NonpreemptivePriorityWorkState n JobId)
    a AppliedModelingLib.Queueing.advanceNonpreemptivePriorityWorkState._f a_1 a_2
\end{ReviewVerbatim}

\paragraph{Exact Lean library declaration}
\begin{ReviewVerbatim}
/-- Advance a finite queue state to a target time, processing at most `fuel`
service completions.  With enough fuel to cover the currently resident jobs,
this is the literal nonpreemptive service evolution up to that time. -/
noncomputable def advanceNonpreemptivePriorityWorkState
    {n : ℕ} {JobId : Type*} :
    ℕ → ℝ → NonpreemptivePriorityWorkState n JobId →
      NonpreemptivePriorityWorkState n JobId
  | 0, target, state =>
      if htarget : target ≤ state.currentTime then state
      else match state.active with
      | none => { state with currentTime := target }
      | some _ => state
  | fuel + 1, target, state =>
      if htarget : target ≤ state.currentTime then state
      else match state.active with
      | none => { state with currentTime := target }
      | some (job, residual) =>
          if hcomplete : residual ≤ target - state.currentTime then
            advanceNonpreemptivePriorityWorkState fuel target
              (completeNonpreemptivePriorityWorkJob
                { state with currentTime := state.currentTime + residual })
          else
            { state with
              currentTime := target
              active := some (job, residual - (target - state.currentTime)) }
\end{ReviewVerbatim}

\paragraph{Direct retained dependencies}
\begin{itemize}\raggedright
\item \hyperlink{library-prerequisite-2373fe691c32d999}{Applied\allowbreak{}Modeling\allowbreak{}Lib.\allowbreak{}Queueing.\allowbreak{}Nonpreemptive\allowbreak{}Priority\allowbreak{}Job}
\item \hyperlink{library-prerequisite-e68cb5fce009ec53}{Applied\allowbreak{}Modeling\allowbreak{}Lib.\allowbreak{}Queueing.\allowbreak{}Nonpreemptive\allowbreak{}Priority\allowbreak{}Work\allowbreak{}State}
\item \hyperlink{governing-declaration-e7772cc6b7f51cdf}{Applied\allowbreak{}Modeling\allowbreak{}Lib.\allowbreak{}Queueing.\allowbreak{}complete\allowbreak{}Nonpreemptive\allowbreak{}Priority\allowbreak{}Work\allowbreak{}Job}
\end{itemize}
\paragraph{Recorded source-to-library screening}
\textbf{Verdict:} \texttt{matches}
\\
{\raggedright
\textbf{Reason:} This is A4's work-\allowbreak{}state evolution through a completion before selecting the next priority job.\allowbreak{}
\par}
\reviewmatch{Matches source input}
\noindent\textbf{Reviewer annotation}\par
\reviewerbox
\clearpage
\hypertarget{library-prerequisite-d49e277073e052d4}{}
\subsection*{\small\ttfamily\raggedright Applied\allowbreak{}Modeling\allowbreak{}Lib.\allowbreak{}Queueing.\allowbreak{}total\allowbreak{}Nonpreemptive\allowbreak{}Priority\allowbreak{}Work\allowbreak{}Jobs}
\textbf{Source locator:} {\footnotesize\raggedright cited publication, lines 264-\allowbreak{}271\par}
\paragraph{Verbatim paper-source connection}
\begin{ReviewVerbatim}
A4. Service Process                                                      Vi'(Xi) = pi + vi Wi (5)
The processor is modeled as an M/M/1 queueing                          where (5) is analogous to (1) with total expected class-
system with nonpreemptive priority among job                           i user cost consisting of two components: a direct
classes. The processing requirement Hi(-) of a class-i                 payment of pi per job, and the expected class-i delay
job, ri, follows an exponential distribution with finite cost, given by vi * Wi(X).
mean ci. Processing requirements are independent           The following theorem generalizes the net value
across jobs and identically distributed within job       maximizing pricing scheme, (3), to the multiclass case.
classes.
\end{ReviewVerbatim}



\paragraph{Lean-expanded library semantic target}
\begin{ReviewVerbatim}
type:
{n : ℕ} → {JobId : Type u_1} → AppliedModelingLib.Queueing.NonpreemptivePriorityWorkState n JobId → ℕ

value:
fun {n : ℕ} {JobId : Type u_1} (state : AppliedModelingLib.Queueing.NonpreemptivePriorityWorkState n JobId) =>
  (match state.active with
    | none => 0
    | some val => 1) +
    AppliedModelingLib.Queueing.totalPriorityWaitingJobs state
\end{ReviewVerbatim}

\paragraph{Exact Lean library declaration}
\begin{ReviewVerbatim}
/-- The total number of jobs currently resident in the queue, including an
active job if one exists. -/
def totalNonpreemptivePriorityWorkJobs
    {n : ℕ} {JobId : Type*}
    (state : NonpreemptivePriorityWorkState n JobId) : ℕ :=
  (match state.active with | none => 0 | some _ => 1) +
    totalPriorityWaitingJobs state
\end{ReviewVerbatim}

\paragraph{Direct retained dependencies}
\begin{itemize}\raggedright
\item \hyperlink{library-prerequisite-2373fe691c32d999}{Applied\allowbreak{}Modeling\allowbreak{}Lib.\allowbreak{}Queueing.\allowbreak{}Nonpreemptive\allowbreak{}Priority\allowbreak{}Job}
\item \hyperlink{library-prerequisite-e68cb5fce009ec53}{Applied\allowbreak{}Modeling\allowbreak{}Lib.\allowbreak{}Queueing.\allowbreak{}Nonpreemptive\allowbreak{}Priority\allowbreak{}Work\allowbreak{}State}
\item \hyperlink{governing-declaration-f2f50f38ca2b2a0b}{Applied\allowbreak{}Modeling\allowbreak{}Lib.\allowbreak{}Queueing.\allowbreak{}total\allowbreak{}Priority\allowbreak{}Waiting\allowbreak{}Jobs}
\end{itemize}
\paragraph{Recorded source-to-library screening}
\textbf{Verdict:} \texttt{matches}
\\
{\raggedright
\textbf{Reason:} This finite fuel/\allowbreak{}count helper supports A4 execution without changing the source queue discipline.\allowbreak{}
\par}
\reviewmatch{Matches source input}
\noindent\textbf{Reviewer annotation}\par
\reviewerbox
\clearpage
\hypertarget{library-prerequisite-74e6067041d9c8dc}{}
\subsection*{\small\ttfamily\raggedright Applied\allowbreak{}Modeling\allowbreak{}Lib.\allowbreak{}Queueing.\allowbreak{}run\allowbreak{}Nonpreemptive\allowbreak{}Priority\allowbreak{}Arrival\allowbreak{}Trace}
\textbf{Source locator:} {\footnotesize\raggedright cited publication, lines 264-\allowbreak{}271\par}
\paragraph{Verbatim paper-source connection}
\begin{ReviewVerbatim}
A4. Service Process                                                      Vi'(Xi) = pi + vi Wi (5)
The processor is modeled as an M/M/1 queueing                          where (5) is analogous to (1) with total expected class-
system with nonpreemptive priority among job                           i user cost consisting of two components: a direct
classes. The processing requirement Hi(-) of a class-i                 payment of pi per job, and the expected class-i delay
job, ri, follows an exponential distribution with finite cost, given by vi * Wi(X).
mean ci. Processing requirements are independent           The following theorem generalizes the net value
across jobs and identically distributed within job       maximizing pricing scheme, (3), to the multiclass case.
classes.
\end{ReviewVerbatim}



\paragraph{Lean-expanded library semantic target}
\begin{ReviewVerbatim}
type:
{n : ℕ} →
  {JobId : Type u_1} →
    AppliedModelingLib.Queueing.NonpreemptivePriorityWorkState n JobId →
      List (AppliedModelingLib.Queueing.NonpreemptivePriorityJob n JobId) →
        AppliedModelingLib.Queueing.NonpreemptivePriorityWorkState n JobId

value:
fun {n : ℕ} {JobId : Type u_1} (initial : AppliedModelingLib.Queueing.NonpreemptivePriorityWorkState n JobId)
    (jobs : List (AppliedModelingLib.Queueing.NonpreemptivePriorityJob n JobId)) =>
  List.foldl
    (fun (state : AppliedModelingLib.Queueing.NonpreemptivePriorityWorkState n JobId)
        (job : AppliedModelingLib.Queueing.NonpreemptivePriorityJob n JobId) =>
      AppliedModelingLib.Queueing.advanceThenAdmitNonpreemptivePriorityJob
        (AppliedModelingLib.Queueing.totalNonpreemptivePriorityWorkJobs state) state job)
    initial jobs
\end{ReviewVerbatim}

\paragraph{Exact Lean library declaration}
\begin{ReviewVerbatim}
/-- Execute a finite, chronologically supplied list of arrivals.  The list
order is the tie convention if two jobs are presented at the same epoch. -/
noncomputable def runNonpreemptivePriorityArrivalTrace
    {n : ℕ} {JobId : Type*}
    (initial : NonpreemptivePriorityWorkState n JobId)
    (jobs : List (NonpreemptivePriorityJob n JobId)) :
    NonpreemptivePriorityWorkState n JobId :=
  jobs.foldl (fun state job =>
    advanceThenAdmitNonpreemptivePriorityJob
      (totalNonpreemptivePriorityWorkJobs state) state job) initial
\end{ReviewVerbatim}

\paragraph{Direct retained dependencies}
\begin{itemize}\raggedright
\item \hyperlink{library-prerequisite-2373fe691c32d999}{Applied\allowbreak{}Modeling\allowbreak{}Lib.\allowbreak{}Queueing.\allowbreak{}Nonpreemptive\allowbreak{}Priority\allowbreak{}Job}
\item \hyperlink{library-prerequisite-e68cb5fce009ec53}{Applied\allowbreak{}Modeling\allowbreak{}Lib.\allowbreak{}Queueing.\allowbreak{}Nonpreemptive\allowbreak{}Priority\allowbreak{}Work\allowbreak{}State}
\item \hyperlink{governing-declaration-fe14523475acac38}{Applied\allowbreak{}Modeling\allowbreak{}Lib.\allowbreak{}Queueing.\allowbreak{}advance\allowbreak{}Then\allowbreak{}Admit\allowbreak{}Nonpreemptive\allowbreak{}Priority\allowbreak{}Job}
\item \hyperlink{library-prerequisite-d49e277073e052d4}{Applied\allowbreak{}Modeling\allowbreak{}Lib.\allowbreak{}Queueing.\allowbreak{}total\allowbreak{}Nonpreemptive\allowbreak{}Priority\allowbreak{}Work\allowbreak{}Jobs}
\end{itemize}
\paragraph{Recorded source-to-library screening}
\textbf{Verdict:} \texttt{matches}
\\
{\raggedright
\textbf{Reason:} It is the chronological finite execution of A4's nonpreemptive priority discipline.\allowbreak{}
\par}
\reviewmatch{Matches source input}
\noindent\textbf{Reviewer annotation}\par
\reviewerbox
\clearpage
\hypertarget{library-prerequisite-6e4e64d50452f911}{}
\subsection*{\small\ttfamily\raggedright Applied\allowbreak{}Modeling\allowbreak{}Lib.\allowbreak{}Queueing.\allowbreak{}canonical\allowbreak{}Stationary\allowbreak{}Priority\allowbreak{}Class\allowbreak{}Tagged\allowbreak{}Finite\allowbreak{}Window\allowbreak{}State}
\textbf{Source locator:} {\footnotesize\raggedright cited publication, lines 264-\allowbreak{}271\par}
\paragraph{Verbatim paper-source connection}
\begin{ReviewVerbatim}
A4. Service Process                                                      Vi'(Xi) = pi + vi Wi (5)
The processor is modeled as an M/M/1 queueing                          where (5) is analogous to (1) with total expected class-
system with nonpreemptive priority among job                           i user cost consisting of two components: a direct
classes. The processing requirement Hi(-) of a class-i                 payment of pi per job, and the expected class-i delay
job, ri, follows an exponential distribution with finite cost, given by vi * Wi(X).
mean ci. Processing requirements are independent           The following theorem generalizes the net value
across jobs and identically distributed within job       maximizing pricing scheme, (3), to the multiclass case.
classes.
\end{ReviewVerbatim}



\paragraph{Lean-expanded library semantic target}
\begin{ReviewVerbatim}
type:
{n : ℕ} →
  (Fin n → ℝ) →
    (i : Fin n) →
      AppliedModelingLib.Queueing.MulticlassStationaryPoissonWorkClassTaggedSample (Fin n) i →
        ℝ →
          ℝ →
            AppliedModelingLib.Queueing.NonpreemptivePriorityWorkState n
              (AppliedModelingLib.Queueing.NonpreemptivePriorityArrivalIndex n)

value:
fun {n : ℕ} (meanService : Fin n → ℝ) (i : Fin n)
    (z : AppliedModelingLib.Queueing.MulticlassStationaryPoissonWorkClassTaggedSample (Fin n) i) (a b : ℝ) =>
  have initial :
    AppliedModelingLib.Queueing.NonpreemptivePriorityWorkState n
      (AppliedModelingLib.Queueing.NonpreemptivePriorityArrivalIndex n) :=
    AppliedModelingLib.Queueing.emptyNonpreemptivePriorityWorkState a;
  have afterArrivals :
    AppliedModelingLib.Queueing.NonpreemptivePriorityWorkState n
      (AppliedModelingLib.Queueing.NonpreemptivePriorityArrivalIndex n) :=
    AppliedModelingLib.Queueing.runNonpreemptivePriorityArrivalTrace initial
      (AppliedModelingLib.Queueing.canonicalStationaryPriorityClassTaggedArrivalWindowJobs meanService i z a b);
  AppliedModelingLib.Queueing.advanceNonpreemptivePriorityWorkState
    (AppliedModelingLib.Queueing.totalNonpreemptivePriorityWorkJobs afterArrivals) b afterArrivals
\end{ReviewVerbatim}

\paragraph{Exact Lean library declaration}
\begin{ReviewVerbatim}
/-- Execute the canonically ordered finite selected/Palm arrival ledger,
starting empty at the left endpoint and serving through the right endpoint. -/
def canonicalStationaryPriorityClassTaggedFiniteWindowState
    {n : ℕ} (meanService : Fin n → ℝ) (i : Fin n)
    (z : MulticlassStationaryPoissonWorkClassTaggedSample (Fin n) i)
    (a b : ℝ) :
    NonpreemptivePriorityWorkState n (NonpreemptivePriorityArrivalIndex n) :=
  let initial := emptyNonpreemptivePriorityWorkState
    (n := n) (JobId := NonpreemptivePriorityArrivalIndex n) a
  let afterArrivals := runNonpreemptivePriorityArrivalTrace initial
    (canonicalStationaryPriorityClassTaggedArrivalWindowJobs meanService i z a b)
  advanceNonpreemptivePriorityWorkState
    (totalNonpreemptivePriorityWorkJobs afterArrivals) b afterArrivals
\end{ReviewVerbatim}

\paragraph{Direct retained dependencies}
\begin{itemize}\raggedright
\item \hyperlink{library-prerequisite-755cdb28ceaf46d2}{Applied\allowbreak{}Modeling\allowbreak{}Lib.\allowbreak{}Queueing.\allowbreak{}Multiclass\allowbreak{}Stationary\allowbreak{}Poisson\allowbreak{}Work\allowbreak{}Class\allowbreak{}Tagged\allowbreak{}Sample}
\item \hyperlink{governing-declaration-5d3724f4dc832d77}{Applied\allowbreak{}Modeling\allowbreak{}Lib.\allowbreak{}Queueing.\allowbreak{}Nonpreemptive\allowbreak{}Priority\allowbreak{}Arrival\allowbreak{}Index}
\item \hyperlink{library-prerequisite-e68cb5fce009ec53}{Applied\allowbreak{}Modeling\allowbreak{}Lib.\allowbreak{}Queueing.\allowbreak{}Nonpreemptive\allowbreak{}Priority\allowbreak{}Work\allowbreak{}State}
\item \hyperlink{library-prerequisite-96200a424cf39dee}{Applied\allowbreak{}Modeling\allowbreak{}Lib.\allowbreak{}Queueing.\allowbreak{}advance\allowbreak{}Nonpreemptive\allowbreak{}Priority\allowbreak{}Work\allowbreak{}State}
\item \hyperlink{governing-declaration-c8dff705756357c1}{Applied\allowbreak{}Modeling\allowbreak{}Lib.\allowbreak{}Queueing.\allowbreak{}canonical\allowbreak{}Stationary\allowbreak{}Priority\allowbreak{}Class\allowbreak{}Tagged\allowbreak{}Arrival\allowbreak{}Window\allowbreak{}Jobs}
\item \hyperlink{governing-declaration-0e2f4b9eb63fa294}{Applied\allowbreak{}Modeling\allowbreak{}Lib.\allowbreak{}Queueing.\allowbreak{}empty\allowbreak{}Nonpreemptive\allowbreak{}Priority\allowbreak{}Work\allowbreak{}State}
\item \hyperlink{library-prerequisite-74e6067041d9c8dc}{Applied\allowbreak{}Modeling\allowbreak{}Lib.\allowbreak{}Queueing.\allowbreak{}run\allowbreak{}Nonpreemptive\allowbreak{}Priority\allowbreak{}Arrival\allowbreak{}Trace}
\item \hyperlink{library-prerequisite-d49e277073e052d4}{Applied\allowbreak{}Modeling\allowbreak{}Lib.\allowbreak{}Queueing.\allowbreak{}total\allowbreak{}Nonpreemptive\allowbreak{}Priority\allowbreak{}Work\allowbreak{}Jobs}
\end{itemize}
\paragraph{Recorded source-to-library screening}
\textbf{Verdict:} \texttt{matches}
\\
{\raggedright
\textbf{Reason:} This finite replay applies the A4 discipline to the separately mapped A1 stationary input, without being promoted to a source formula itself.\allowbreak{}
\par}
\reviewmatch{Matches source input}
\noindent\textbf{Reviewer annotation}\par
\reviewerbox
\clearpage
\hypertarget{library-prerequisite-bb4b8c8010d5959b}{}
\subsection*{\small\ttfamily\raggedright Applied\allowbreak{}Modeling\allowbreak{}Lib.\allowbreak{}Queueing.\allowbreak{}canonical\allowbreak{}Stationary\allowbreak{}Priority\allowbreak{}Class\allowbreak{}Tagged\allowbreak{}Future\allowbreak{}Arrival\allowbreak{}Jobs}
\textbf{Source locator:} {\footnotesize\raggedright cited publication, lines 264-\allowbreak{}271\par}
\paragraph{Verbatim paper-source connection}
\begin{ReviewVerbatim}
A4. Service Process                                                      Vi'(Xi) = pi + vi Wi (5)
The processor is modeled as an M/M/1 queueing                          where (5) is analogous to (1) with total expected class-
system with nonpreemptive priority among job                           i user cost consisting of two components: a direct
classes. The processing requirement Hi(-) of a class-i                 payment of pi per job, and the expected class-i delay
job, ri, follows an exponential distribution with finite cost, given by vi * Wi(X).
mean ci. Processing requirements are independent           The following theorem generalizes the net value
across jobs and identically distributed within job       maximizing pricing scheme, (3), to the multiclass case.
classes.
\end{ReviewVerbatim}



\paragraph{Lean-expanded library semantic target}
\begin{ReviewVerbatim}
type:
{n : ℕ} →
  (Fin n → ℝ) →
    (i : Fin n) →
      AppliedModelingLib.Queueing.MulticlassStationaryPoissonWorkClassTaggedSample (Fin n) i →
        ℝ →
          List
            (AppliedModelingLib.Queueing.NonpreemptivePriorityJob n
              (AppliedModelingLib.Queueing.NonpreemptivePriorityArrivalIndex n))

value:
fun {n : ℕ} (meanService : Fin n → ℝ) (i : Fin n)
    (z : AppliedModelingLib.Queueing.MulticlassStationaryPoissonWorkClassTaggedSample (Fin n) i) (t : ℝ) =>
  List.map
    (fun (q : AppliedModelingLib.Queueing.NonpreemptivePriorityArrivalIndex n) =>
      { identifier := q, priority := q.fst,
        arrivalTime := AppliedModelingLib.Queueing.stationaryPriorityClassTaggedArrival i z q.fst q.snd,
        serviceWork :=
          AppliedModelingLib.Queueing.stationaryPriorityClassTaggedWorkRequirementAt meanService i z q.fst q.snd })
    (AppliedModelingLib.Queueing.canonicalStationaryPriorityClassTaggedFutureArrivalIndices i z t)
\end{ReviewVerbatim}

\paragraph{Exact Lean library declaration}
\begin{ReviewVerbatim}
/-- The literal finite list of post-tag priority jobs used to continue the
post-admission selected queue state. -/
noncomputable def canonicalStationaryPriorityClassTaggedFutureArrivalJobs
    {n : ℕ} (meanService : Fin n → ℝ) (i : Fin n)
    (z : MulticlassStationaryPoissonWorkClassTaggedSample (Fin n) i)
    (t : ℝ) : List (NonpreemptivePriorityJob n (NonpreemptivePriorityArrivalIndex n)) :=
  (canonicalStationaryPriorityClassTaggedFutureArrivalIndices i z t).map fun q =>
    { identifier := q
      priority := q.1
      arrivalTime := stationaryPriorityClassTaggedArrival i z q.1 q.2
      serviceWork := stationaryPriorityClassTaggedWorkRequirementAt meanService i z q.1 q.2 }
\end{ReviewVerbatim}

\paragraph{Direct retained dependencies}
\begin{itemize}\raggedright
\item \hyperlink{library-prerequisite-755cdb28ceaf46d2}{Applied\allowbreak{}Modeling\allowbreak{}Lib.\allowbreak{}Queueing.\allowbreak{}Multiclass\allowbreak{}Stationary\allowbreak{}Poisson\allowbreak{}Work\allowbreak{}Class\allowbreak{}Tagged\allowbreak{}Sample}
\item \hyperlink{governing-declaration-5d3724f4dc832d77}{Applied\allowbreak{}Modeling\allowbreak{}Lib.\allowbreak{}Queueing.\allowbreak{}Nonpreemptive\allowbreak{}Priority\allowbreak{}Arrival\allowbreak{}Index}
\item \hyperlink{library-prerequisite-2373fe691c32d999}{Applied\allowbreak{}Modeling\allowbreak{}Lib.\allowbreak{}Queueing.\allowbreak{}Nonpreemptive\allowbreak{}Priority\allowbreak{}Job}
\item \hyperlink{governing-declaration-e1de2ae4651d3437}{Applied\allowbreak{}Modeling\allowbreak{}Lib.\allowbreak{}Queueing.\allowbreak{}canonical\allowbreak{}Stationary\allowbreak{}Priority\allowbreak{}Class\allowbreak{}Tagged\allowbreak{}Future\allowbreak{}Arrival\allowbreak{}Indices}
\item \hyperlink{governing-declaration-2945df1a7a8ec970}{Applied\allowbreak{}Modeling\allowbreak{}Lib.\allowbreak{}Queueing.\allowbreak{}stationary\allowbreak{}Priority\allowbreak{}Class\allowbreak{}Tagged\allowbreak{}Arrival}
\item \hyperlink{governing-declaration-949c09de5ada557f}{Applied\allowbreak{}Modeling\allowbreak{}Lib.\allowbreak{}Queueing.\allowbreak{}stationary\allowbreak{}Priority\allowbreak{}Class\allowbreak{}Tagged\allowbreak{}Work\allowbreak{}Requirement\allowbreak{}At}
\end{itemize}
\paragraph{Recorded source-to-library screening}
\textbf{Verdict:} \texttt{matches}
\\
{\raggedright
\textbf{Reason:} It converts the mapped post-\allowbreak{}tag A1 input to jobs for the A4 priority queue, preserving class and work marks.\allowbreak{}
\par}
\reviewmatch{Matches source input}
\noindent\textbf{Reviewer annotation}\par
\reviewerbox
\clearpage
\hypertarget{library-prerequisite-cc744f6520cc7734}{}
\subsection*{\small\ttfamily\raggedright Applied\allowbreak{}Modeling\allowbreak{}Lib.\allowbreak{}Queueing.\allowbreak{}clear\allowbreak{}Nonpreemptive\allowbreak{}Priority\allowbreak{}Completion\allowbreak{}Ledger}
\textbf{Source locator:} {\footnotesize\raggedright cited publication, lines 264-\allowbreak{}271\par}
\paragraph{Verbatim paper-source connection}
\begin{ReviewVerbatim}
A4. Service Process                                                      Vi'(Xi) = pi + vi Wi (5)
The processor is modeled as an M/M/1 queueing                          where (5) is analogous to (1) with total expected class-
system with nonpreemptive priority among job                           i user cost consisting of two components: a direct
classes. The processing requirement Hi(-) of a class-i                 payment of pi per job, and the expected class-i delay
job, ri, follows an exponential distribution with finite cost, given by vi * Wi(X).
mean ci. Processing requirements are independent           The following theorem generalizes the net value
across jobs and identically distributed within job       maximizing pricing scheme, (3), to the multiclass case.
classes.
\end{ReviewVerbatim}



\paragraph{Lean-expanded library semantic target}
\begin{ReviewVerbatim}
type:
{n : ℕ} →
  {JobId : Type u_1} →
    AppliedModelingLib.Queueing.NonpreemptivePriorityWorkState n JobId →
      AppliedModelingLib.Queueing.NonpreemptivePriorityWorkState n JobId

value:
fun {n : ℕ} {JobId : Type u_1} (state : AppliedModelingLib.Queueing.NonpreemptivePriorityWorkState n JobId) =>
  { currentTime := state.currentTime, active := state.active, waiting := state.waiting, completed := [] }
\end{ReviewVerbatim}

\paragraph{Exact Lean library declaration}
\begin{ReviewVerbatim}
/-- Reset the observational completion ledger while retaining the complete
live priority queue.  This does not alter any future service or admission
decision, and is useful when a new tagged observation begins at a specified
physical epoch. -/
def clearNonpreemptivePriorityCompletionLedger
    {n : ℕ} {JobId : Type*}
    (state : NonpreemptivePriorityWorkState n JobId) :
    NonpreemptivePriorityWorkState n JobId :=
  { state with completed := [] }
\end{ReviewVerbatim}

\paragraph{Direct retained dependencies}
\begin{itemize}\raggedright
\item \hyperlink{library-prerequisite-2373fe691c32d999}{Applied\allowbreak{}Modeling\allowbreak{}Lib.\allowbreak{}Queueing.\allowbreak{}Nonpreemptive\allowbreak{}Priority\allowbreak{}Job}
\item \hyperlink{library-prerequisite-e68cb5fce009ec53}{Applied\allowbreak{}Modeling\allowbreak{}Lib.\allowbreak{}Queueing.\allowbreak{}Nonpreemptive\allowbreak{}Priority\allowbreak{}Work\allowbreak{}State}
\end{itemize}
\paragraph{Recorded source-to-library screening}
\textbf{Verdict:} \texttt{matches}
\\
{\raggedright
\textbf{Reason:} This observation-\allowbreak{}only ledger reset preserves the A4 live queue dynamics while beginning a tagged completion observation.\allowbreak{}
\par}
\reviewmatch{Matches source input}
\noindent\textbf{Reviewer annotation}\par
\reviewerbox
\clearpage
\hypertarget{library-prerequisite-795b7af1c661bc8d}{}
\subsection*{\small\ttfamily\raggedright Applied\allowbreak{}Modeling\allowbreak{}Lib.\allowbreak{}Queueing.\allowbreak{}finite\allowbreak{}Nonpreemptive\allowbreak{}Priority\allowbreak{}Queue\allowbreak{}Wait}
\textbf{Source locator:} {\footnotesize\raggedright cited publication, lines 1118-\allowbreak{}1129\par}
\paragraph{Verbatim paper-source connection}
\begin{ReviewVerbatim}
Wi(.) = The expected waiting time of a class-i job                      KLEINROCK, L. 1976. Queueing System, Vols. I and II.
         in the queue in steady state;                                       John Wiley, New York.
 Wi ( )= The expected sojourn time of a class-i job                     KNUDSEN, N. C. 1972. Individual and Social Optimiza-
         in the system in steady state.                                      tion in a Multi-Server Queue With a General Cost-
                                                                             Benefit Structure. Econometrica 40, 515-528.
For the M/M/1 queue with nonpreemptive priority                         KREPS, D. 1990. A Course in Microeconomic Theory.
discipline (Heyman and Sobel 1982, p. 435)                                 Princeton University Press, Princeton, N.J.
                                                                        MARCHAND, M. 1974. Priority Pricing. Mgmt. Sci. 20,
   _ i( L1()
 W1CX)  = W1 Wqc1+ -+
                   Ci AR   (A-)
                      cj. (-                                                  1131-1140.

[Next verbatim source excerpt]

A4. Service Process                                                      Vi'(Xi) = pi + vi Wi (5)
The processor is modeled as an M/M/1 queueing                          where (5) is analogous to (1) with total expected class-
system with nonpreemptive priority among job                           i user cost consisting of two components: a direct
classes. The processing requirement Hi(-) of a class-i                 payment of pi per job, and the expected class-i delay
job, ri, follows an exponential distribution with finite cost, given by vi * Wi(X).
mean ci. Processing requirements are independent           The following theorem generalizes the net value
across jobs and identically distributed within job       maximizing pricing scheme, (3), to the multiclass case.
classes.

[Next verbatim source excerpt]

giving rise to X = A (p) or V'(X) = 4-'(X/A) = p.                     Al. Arrival Processes

  Consider the case where the delay cost v is positive.               The arrival process of class-i jobs to the system, which
Here, the cost faced by an individual user consists of                reflects the aggregation of infinitesimal users' job flow,
two components: a direct payment of p per job, and                    follows a stationary Poisson process with a price-




                                                 This content downloaded from
                               108.30.55.84 on Mon, 24 Aug 2026 10:55:53 UTC
                                         All use subject to https://about.jstor.org/terms

[form-feed in source extraction]
                                                                                Priority Pricingfor the MIMI1 Queue / 873

dependent arrival rate Xi. The R arrival processes are                 by the system (per unit of time) to the organization,
stochastically independent.                                            where we assume for now that the system administra-
\end{ReviewVerbatim}



\paragraph{Lean-expanded library semantic target}
\begin{ReviewVerbatim}
type:
{n : ℕ} → (Fin n → ℝ) → (Fin n → ℝ) → Fin n → ℝ

value:
fun {n : ℕ} (arrivalRate meanService : Fin n → ℝ) (i : Fin n) =>
  AppliedModelingLib.Queueing.finitePriorityResidualWork arrivalRate meanService /
    (AppliedModelingLib.Queueing.finitePriorityStrictSlack meanService arrivalRate i *
      AppliedModelingLib.Queueing.finitePriorityInclusiveSlack meanService arrivalRate i)
\end{ReviewVerbatim}

\paragraph{Exact Lean library declaration}
\begin{ReviewVerbatim}
/-- Standard nonpreemptive-priority mean waiting time in queue, expressed in
terms of the residual-work numerator and the two priority-load slacks. -/
noncomputable def finiteNonpreemptivePriorityQueueWait
    {n : ℕ} (arrivalRate meanService : Fin n → ℝ) (i : Fin n) : ℝ :=
  finitePriorityResidualWork arrivalRate meanService /
    (finitePriorityStrictSlack meanService arrivalRate i *
      finitePriorityInclusiveSlack meanService arrivalRate i)
\end{ReviewVerbatim}

\paragraph{Direct retained dependencies}
\begin{itemize}\raggedright
\item \hyperlink{governing-declaration-f2b9a1af307e7469}{Applied\allowbreak{}Modeling\allowbreak{}Lib.\allowbreak{}Queueing.\allowbreak{}finite\allowbreak{}Priority\allowbreak{}Inclusive\allowbreak{}Slack}
\item \hyperlink{governing-declaration-80f4f37f124e636e}{Applied\allowbreak{}Modeling\allowbreak{}Lib.\allowbreak{}Queueing.\allowbreak{}finite\allowbreak{}Priority\allowbreak{}Residual\allowbreak{}Work}
\item \hyperlink{governing-declaration-3f5cf610a8190094}{Applied\allowbreak{}Modeling\allowbreak{}Lib.\allowbreak{}Queueing.\allowbreak{}finite\allowbreak{}Priority\allowbreak{}Strict\allowbreak{}Slack}
\end{itemize}
\paragraph{Recorded source-to-library screening}
\textbf{Verdict:} \texttt{matches}
\\
{\raggedright
\textbf{Reason:} Its residual-\allowbreak{}work numerator and strict/\allowbreak{}inclusive priority-\allowbreak{}load slacks are exactly the Appendix (A-\allowbreak{}1) mean queue-\allowbreak{}wait ingredients.\allowbreak{}
\par}
\reviewmatch{Matches source input}
\noindent\textbf{Reviewer annotation}\par
\reviewerbox
\clearpage
\hypertarget{library-prerequisite-886cbd948a5fd8b1}{}
\subsection*{\small\ttfamily\raggedright Applied\allowbreak{}Modeling\allowbreak{}Lib.\allowbreak{}Queueing.\allowbreak{}nonpreemptive\allowbreak{}Priority\allowbreak{}Recorded\allowbreak{}Completion\allowbreak{}Time}
\textbf{Source locator:} {\footnotesize\raggedright cited publication, lines 1118-\allowbreak{}1129\par}
\paragraph{Verbatim paper-source connection}
\begin{ReviewVerbatim}
Wi(.) = The expected waiting time of a class-i job                      KLEINROCK, L. 1976. Queueing System, Vols. I and II.
         in the queue in steady state;                                       John Wiley, New York.
 Wi ( )= The expected sojourn time of a class-i job                     KNUDSEN, N. C. 1972. Individual and Social Optimiza-
         in the system in steady state.                                      tion in a Multi-Server Queue With a General Cost-
                                                                             Benefit Structure. Econometrica 40, 515-528.
For the M/M/1 queue with nonpreemptive priority                         KREPS, D. 1990. A Course in Microeconomic Theory.
discipline (Heyman and Sobel 1982, p. 435)                                 Princeton University Press, Princeton, N.J.
                                                                        MARCHAND, M. 1974. Priority Pricing. Mgmt. Sci. 20,
   _ i( L1()
 W1CX)  = W1 Wqc1+ -+
                   Ci AR   (A-)
                      cj. (-                                                  1131-1140.

[Next verbatim source excerpt]

A4. Service Process                                                      Vi'(Xi) = pi + vi Wi (5)
The processor is modeled as an M/M/1 queueing                          where (5) is analogous to (1) with total expected class-
system with nonpreemptive priority among job                           i user cost consisting of two components: a direct
classes. The processing requirement Hi(-) of a class-i                 payment of pi per job, and the expected class-i delay
job, ri, follows an exponential distribution with finite cost, given by vi * Wi(X).
mean ci. Processing requirements are independent           The following theorem generalizes the net value
across jobs and identically distributed within job       maximizing pricing scheme, (3), to the multiclass case.
classes.

[Next verbatim source excerpt]

giving rise to X = A (p) or V'(X) = 4-'(X/A) = p.                     Al. Arrival Processes

  Consider the case where the delay cost v is positive.               The arrival process of class-i jobs to the system, which
Here, the cost faced by an individual user consists of                reflects the aggregation of infinitesimal users' job flow,
two components: a direct payment of p per job, and                    follows a stationary Poisson process with a price-




                                                 This content downloaded from
                               108.30.55.84 on Mon, 24 Aug 2026 10:55:53 UTC
                                         All use subject to https://about.jstor.org/terms

[form-feed in source extraction]
                                                                                Priority Pricingfor the MIMI1 Queue / 873

dependent arrival rate Xi. The R arrival processes are                 by the system (per unit of time) to the organization,
stochastically independent.                                            where we assume for now that the system administra-
\end{ReviewVerbatim}



\paragraph{Lean-expanded library semantic target}
\begin{ReviewVerbatim}
type:
{n : ℕ} →
  AppliedModelingLib.Queueing.NonpreemptivePriorityWorkState n
      (AppliedModelingLib.Queueing.NonpreemptivePriorityArrivalIndex n) →
    AppliedModelingLib.Queueing.NonpreemptivePriorityArrivalIndex n → Option ℝ

value:
fun {n : ℕ}
    (state :
      AppliedModelingLib.Queueing.NonpreemptivePriorityWorkState n
        (AppliedModelingLib.Queueing.NonpreemptivePriorityArrivalIndex n))
    (identifier : AppliedModelingLib.Queueing.NonpreemptivePriorityArrivalIndex n) =>
  Option.map Prod.snd
    (List.find?
      (fun
          (entry :
            AppliedModelingLib.Queueing.NonpreemptivePriorityJob n
                (AppliedModelingLib.Queueing.NonpreemptivePriorityArrivalIndex n) ×
              ℝ) =>
        decide (entry.1.identifier = identifier))
      state.completed.reverse)
\end{ReviewVerbatim}

\paragraph{Exact Lean library declaration}
\begin{ReviewVerbatim}
/-- The first physical completion epoch recorded for a fixed labelled job in
a finite queue state.  Completion records are prepended as service occurs, so
the ledger is read in chronological order to keep this observation unchanged
by later, unrelated continuations. -/
noncomputable def nonpreemptivePriorityRecordedCompletionTime
    {n : ℕ}
    (state : NonpreemptivePriorityWorkState n (NonpreemptivePriorityArrivalIndex n))
    (identifier : NonpreemptivePriorityArrivalIndex n) : Option ℝ :=
  (state.completed.reverse.find? fun entry => decide (entry.1.identifier = identifier)).map Prod.snd
\end{ReviewVerbatim}

\paragraph{Direct retained dependencies}
\begin{itemize}\raggedright
\item \hyperlink{governing-declaration-5d3724f4dc832d77}{Applied\allowbreak{}Modeling\allowbreak{}Lib.\allowbreak{}Queueing.\allowbreak{}Nonpreemptive\allowbreak{}Priority\allowbreak{}Arrival\allowbreak{}Index}
\item \hyperlink{library-prerequisite-2373fe691c32d999}{Applied\allowbreak{}Modeling\allowbreak{}Lib.\allowbreak{}Queueing.\allowbreak{}Nonpreemptive\allowbreak{}Priority\allowbreak{}Job}
\item \hyperlink{library-prerequisite-e68cb5fce009ec53}{Applied\allowbreak{}Modeling\allowbreak{}Lib.\allowbreak{}Queueing.\allowbreak{}Nonpreemptive\allowbreak{}Priority\allowbreak{}Work\allowbreak{}State}
\end{itemize}
\paragraph{Recorded source-to-library screening}
\textbf{Verdict:} \texttt{matches}
\\
{\raggedright
\textbf{Reason:} This finite replay completion observation supports the Appendix tagged queue-\allowbreak{}wait quantity without asserting a separate source theorem.\allowbreak{}
\par}
\reviewmatch{Matches source input}
\noindent\textbf{Reviewer annotation}\par
\reviewerbox
\clearpage
\hypertarget{library-prerequisite-3702583c174721c0}{}
\subsection*{\small\ttfamily\raggedright Applied\allowbreak{}Modeling\allowbreak{}Lib.\allowbreak{}Queueing.\allowbreak{}stationary\allowbreak{}Priority\allowbreak{}Class\allowbreak{}Tagged\allowbreak{}Job}
\textbf{Source locator:} {\footnotesize\raggedright cited publication, lines 264-\allowbreak{}271\par}
\paragraph{Verbatim paper-source connection}
\begin{ReviewVerbatim}
A4. Service Process                                                      Vi'(Xi) = pi + vi Wi (5)
The processor is modeled as an M/M/1 queueing                          where (5) is analogous to (1) with total expected class-
system with nonpreemptive priority among job                           i user cost consisting of two components: a direct
classes. The processing requirement Hi(-) of a class-i                 payment of pi per job, and the expected class-i delay
job, ri, follows an exponential distribution with finite cost, given by vi * Wi(X).
mean ci. Processing requirements are independent           The following theorem generalizes the net value
across jobs and identically distributed within job       maximizing pricing scheme, (3), to the multiclass case.
classes.
\end{ReviewVerbatim}



\paragraph{Lean-expanded library semantic target}
\begin{ReviewVerbatim}
type:
{n : ℕ} →
  (Fin n → ℝ) →
    (i : Fin n) →
      AppliedModelingLib.Queueing.MulticlassStationaryPoissonWorkClassTaggedSample (Fin n) i →
        AppliedModelingLib.Queueing.NonpreemptivePriorityJob n
          (AppliedModelingLib.Queueing.NonpreemptivePriorityArrivalIndex n)

value:
fun {n : ℕ} (meanService : Fin n → ℝ) (i : Fin n)
    (z : AppliedModelingLib.Queueing.MulticlassStationaryPoissonWorkClassTaggedSample (Fin n) i) =>
  { identifier := ⟨i, 0⟩, priority := i,
    arrivalTime := AppliedModelingLib.Queueing.stationaryPriorityClassTaggedArrival i z i 0,
    serviceWork := AppliedModelingLib.Queueing.stationaryPriorityClassTaggedWorkRequirementAt meanService i z i 0 }
\end{ReviewVerbatim}

\paragraph{Exact Lean library declaration}
\begin{ReviewVerbatim}
/-- The distinguished class-`i` customer at the Palm origin, represented as
an ordinary priority job with its original class/index identifier. -/
def stationaryPriorityClassTaggedJob
    {n : ℕ} (meanService : Fin n → ℝ) (i : Fin n)
    (z : MulticlassStationaryPoissonWorkClassTaggedSample (Fin n) i) :
    NonpreemptivePriorityJob n (NonpreemptivePriorityArrivalIndex n) :=
  { identifier := Sigma.mk i 0
    priority := i
    arrivalTime := stationaryPriorityClassTaggedArrival i z i 0
    serviceWork := stationaryPriorityClassTaggedWorkRequirementAt meanService i z i 0 }
\end{ReviewVerbatim}

\paragraph{Direct retained dependencies}
\begin{itemize}\raggedright
\item \hyperlink{library-prerequisite-755cdb28ceaf46d2}{Applied\allowbreak{}Modeling\allowbreak{}Lib.\allowbreak{}Queueing.\allowbreak{}Multiclass\allowbreak{}Stationary\allowbreak{}Poisson\allowbreak{}Work\allowbreak{}Class\allowbreak{}Tagged\allowbreak{}Sample}
\item \hyperlink{governing-declaration-5d3724f4dc832d77}{Applied\allowbreak{}Modeling\allowbreak{}Lib.\allowbreak{}Queueing.\allowbreak{}Nonpreemptive\allowbreak{}Priority\allowbreak{}Arrival\allowbreak{}Index}
\item \hyperlink{library-prerequisite-2373fe691c32d999}{Applied\allowbreak{}Modeling\allowbreak{}Lib.\allowbreak{}Queueing.\allowbreak{}Nonpreemptive\allowbreak{}Priority\allowbreak{}Job}
\item \hyperlink{governing-declaration-2945df1a7a8ec970}{Applied\allowbreak{}Modeling\allowbreak{}Lib.\allowbreak{}Queueing.\allowbreak{}stationary\allowbreak{}Priority\allowbreak{}Class\allowbreak{}Tagged\allowbreak{}Arrival}
\item \hyperlink{governing-declaration-949c09de5ada557f}{Applied\allowbreak{}Modeling\allowbreak{}Lib.\allowbreak{}Queueing.\allowbreak{}stationary\allowbreak{}Priority\allowbreak{}Class\allowbreak{}Tagged\allowbreak{}Work\allowbreak{}Requirement\allowbreak{}At}
\end{itemize}
\paragraph{Recorded source-to-library screening}
\textbf{Verdict:} \texttt{matches}
\\
{\raggedright
\textbf{Reason:} It represents the selected class customer at the Palm origin as an ordinary A4 priority job.\allowbreak{}
\par}
\reviewmatch{Matches source input}
\noindent\textbf{Reviewer annotation}\par
\reviewerbox
\clearpage
\hypertarget{library-prerequisite-098447dd1f7443fd}{}
\subsection*{\small\ttfamily\raggedright Applied\allowbreak{}Modeling\allowbreak{}Lib.\allowbreak{}Queueing.\allowbreak{}stationary\allowbreak{}Priority\allowbreak{}Class\allowbreak{}Tagged\allowbreak{}Finite\allowbreak{}Replay\allowbreak{}Post\allowbreak{}Arrival\allowbreak{}State}
\textbf{Source locator:} {\footnotesize\raggedright cited publication, lines 1118-\allowbreak{}1129\par}
\paragraph{Verbatim paper-source connection}
\begin{ReviewVerbatim}
Wi(.) = The expected waiting time of a class-i job                      KLEINROCK, L. 1976. Queueing System, Vols. I and II.
         in the queue in steady state;                                       John Wiley, New York.
 Wi ( )= The expected sojourn time of a class-i job                     KNUDSEN, N. C. 1972. Individual and Social Optimiza-
         in the system in steady state.                                      tion in a Multi-Server Queue With a General Cost-
                                                                             Benefit Structure. Econometrica 40, 515-528.
For the M/M/1 queue with nonpreemptive priority                         KREPS, D. 1990. A Course in Microeconomic Theory.
discipline (Heyman and Sobel 1982, p. 435)                                 Princeton University Press, Princeton, N.J.
                                                                        MARCHAND, M. 1974. Priority Pricing. Mgmt. Sci. 20,
   _ i( L1()
 W1CX)  = W1 Wqc1+ -+
                   Ci AR   (A-)
                      cj. (-                                                  1131-1140.

[Next verbatim source excerpt]

A4. Service Process                                                      Vi'(Xi) = pi + vi Wi (5)
The processor is modeled as an M/M/1 queueing                          where (5) is analogous to (1) with total expected class-
system with nonpreemptive priority among job                           i user cost consisting of two components: a direct
classes. The processing requirement Hi(-) of a class-i                 payment of pi per job, and the expected class-i delay
job, ri, follows an exponential distribution with finite cost, given by vi * Wi(X).
mean ci. Processing requirements are independent           The following theorem generalizes the net value
across jobs and identically distributed within job       maximizing pricing scheme, (3), to the multiclass case.
classes.

[Next verbatim source excerpt]

giving rise to X = A (p) or V'(X) = 4-'(X/A) = p.                     Al. Arrival Processes

  Consider the case where the delay cost v is positive.               The arrival process of class-i jobs to the system, which
Here, the cost faced by an individual user consists of                reflects the aggregation of infinitesimal users' job flow,
two components: a direct payment of p per job, and                    follows a stationary Poisson process with a price-




                                                 This content downloaded from
                               108.30.55.84 on Mon, 24 Aug 2026 10:55:53 UTC
                                         All use subject to https://about.jstor.org/terms

[form-feed in source extraction]
                                                                                Priority Pricingfor the MIMI1 Queue / 873

dependent arrival rate Xi. The R arrival processes are                 by the system (per unit of time) to the organization,
stochastically independent.                                            where we assume for now that the system administra-
\end{ReviewVerbatim}



\paragraph{Lean-expanded library semantic target}
\begin{ReviewVerbatim}
type:
{n : ℕ} →
  (Fin n → ℝ) →
    (i : Fin n) →
      AppliedModelingLib.Queueing.MulticlassStationaryPoissonWorkClassTaggedSample (Fin n) i →
        ℝ →
          ℝ →
            AppliedModelingLib.Queueing.NonpreemptivePriorityWorkState n
              (AppliedModelingLib.Queueing.NonpreemptivePriorityArrivalIndex n)

value:
fun {n : ℕ} (meanService : Fin n → ℝ) (i : Fin n)
    (z : AppliedModelingLib.Queueing.MulticlassStationaryPoissonWorkClassTaggedSample (Fin n) i) (older t : ℝ) =>
  have preArrival :
    AppliedModelingLib.Queueing.NonpreemptivePriorityWorkState n
      (AppliedModelingLib.Queueing.NonpreemptivePriorityArrivalIndex n) :=
    AppliedModelingLib.Queueing.canonicalStationaryPriorityClassTaggedFiniteWindowState meanService i z (-older) 0;
  have afterArrivals :
    AppliedModelingLib.Queueing.NonpreemptivePriorityWorkState n
      (AppliedModelingLib.Queueing.NonpreemptivePriorityArrivalIndex n) :=
    AppliedModelingLib.Queueing.runNonpreemptivePriorityArrivalTrace
      (AppliedModelingLib.Queueing.admitNonpreemptivePriorityJob
        (AppliedModelingLib.Queueing.clearNonpreemptivePriorityCompletionLedger preArrival)
        (AppliedModelingLib.Queueing.stationaryPriorityClassTaggedJob meanService i z))
      (AppliedModelingLib.Queueing.canonicalStationaryPriorityClassTaggedFutureArrivalJobs meanService i z t);
  AppliedModelingLib.Queueing.advanceNonpreemptivePriorityWorkState
    (AppliedModelingLib.Queueing.totalNonpreemptivePriorityWorkJobs afterArrivals) t afterArrivals
\end{ReviewVerbatim}

\paragraph{Exact Lean library declaration}
\begin{ReviewVerbatim}
/-- A completely finite two-sided replay: execute the selected/Palm ledger
from a finite past horizon to zero, clear its historical observations, admit
the distinguished customer, and continue through a finite future horizon. -/
noncomputable def stationaryPriorityClassTaggedFiniteReplayPostArrivalState
    {n : ℕ} (meanService : Fin n → ℝ) (i : Fin n)
    (z : MulticlassStationaryPoissonWorkClassTaggedSample (Fin n) i)
    (older t : ℝ) : NonpreemptivePriorityWorkState n (NonpreemptivePriorityArrivalIndex n) :=
  let preArrival := canonicalStationaryPriorityClassTaggedFiniteWindowState
    meanService i z (-older) 0
  let afterArrivals := runNonpreemptivePriorityArrivalTrace
    (admitNonpreemptivePriorityJob
      (clearNonpreemptivePriorityCompletionLedger preArrival)
      (stationaryPriorityClassTaggedJob meanService i z))
    (canonicalStationaryPriorityClassTaggedFutureArrivalJobs meanService i z t)
  advanceNonpreemptivePriorityWorkState
    (totalNonpreemptivePriorityWorkJobs afterArrivals) t afterArrivals
\end{ReviewVerbatim}

\paragraph{Direct retained dependencies}
\begin{itemize}\raggedright
\item \hyperlink{library-prerequisite-755cdb28ceaf46d2}{Applied\allowbreak{}Modeling\allowbreak{}Lib.\allowbreak{}Queueing.\allowbreak{}Multiclass\allowbreak{}Stationary\allowbreak{}Poisson\allowbreak{}Work\allowbreak{}Class\allowbreak{}Tagged\allowbreak{}Sample}
\item \hyperlink{governing-declaration-5d3724f4dc832d77}{Applied\allowbreak{}Modeling\allowbreak{}Lib.\allowbreak{}Queueing.\allowbreak{}Nonpreemptive\allowbreak{}Priority\allowbreak{}Arrival\allowbreak{}Index}
\item \hyperlink{library-prerequisite-e68cb5fce009ec53}{Applied\allowbreak{}Modeling\allowbreak{}Lib.\allowbreak{}Queueing.\allowbreak{}Nonpreemptive\allowbreak{}Priority\allowbreak{}Work\allowbreak{}State}
\item \hyperlink{library-prerequisite-837de7f0cc179814}{Applied\allowbreak{}Modeling\allowbreak{}Lib.\allowbreak{}Queueing.\allowbreak{}admit\allowbreak{}Nonpreemptive\allowbreak{}Priority\allowbreak{}Job}
\item \hyperlink{library-prerequisite-96200a424cf39dee}{Applied\allowbreak{}Modeling\allowbreak{}Lib.\allowbreak{}Queueing.\allowbreak{}advance\allowbreak{}Nonpreemptive\allowbreak{}Priority\allowbreak{}Work\allowbreak{}State}
\item \hyperlink{library-prerequisite-6e4e64d50452f911}{Applied\allowbreak{}Modeling\allowbreak{}Lib.\allowbreak{}Queueing.\allowbreak{}canonical\allowbreak{}Stationary\allowbreak{}Priority\allowbreak{}Class\allowbreak{}Tagged\allowbreak{}Finite\allowbreak{}Window\allowbreak{}State}
\item \hyperlink{library-prerequisite-bb4b8c8010d5959b}{Applied\allowbreak{}Modeling\allowbreak{}Lib.\allowbreak{}Queueing.\allowbreak{}canonical\allowbreak{}Stationary\allowbreak{}Priority\allowbreak{}Class\allowbreak{}Tagged\allowbreak{}Future\allowbreak{}Arrival\allowbreak{}Jobs}
\item \hyperlink{library-prerequisite-cc744f6520cc7734}{Applied\allowbreak{}Modeling\allowbreak{}Lib.\allowbreak{}Queueing.\allowbreak{}clear\allowbreak{}Nonpreemptive\allowbreak{}Priority\allowbreak{}Completion\allowbreak{}Ledger}
\item \hyperlink{library-prerequisite-74e6067041d9c8dc}{Applied\allowbreak{}Modeling\allowbreak{}Lib.\allowbreak{}Queueing.\allowbreak{}run\allowbreak{}Nonpreemptive\allowbreak{}Priority\allowbreak{}Arrival\allowbreak{}Trace}
\item \hyperlink{library-prerequisite-3702583c174721c0}{Applied\allowbreak{}Modeling\allowbreak{}Lib.\allowbreak{}Queueing.\allowbreak{}stationary\allowbreak{}Priority\allowbreak{}Class\allowbreak{}Tagged\allowbreak{}Job}
\item \hyperlink{library-prerequisite-d49e277073e052d4}{Applied\allowbreak{}Modeling\allowbreak{}Lib.\allowbreak{}Queueing.\allowbreak{}total\allowbreak{}Nonpreemptive\allowbreak{}Priority\allowbreak{}Work\allowbreak{}Jobs}
\end{itemize}
\paragraph{Recorded source-to-library screening}
\textbf{Verdict:} \texttt{matches}
\\
{\raggedright
\textbf{Reason:} It constructs the finite two-\allowbreak{}sided A1/\allowbreak{}A4 replay used to support the stationary tagged wait calculation.\allowbreak{}
\par}
\reviewmatch{Matches source input}
\noindent\textbf{Reviewer annotation}\par
\reviewerbox
\clearpage
\hypertarget{library-prerequisite-945b0f6ad93a6ca5}{}
\subsection*{\small\ttfamily\raggedright Applied\allowbreak{}Modeling\allowbreak{}Lib.\allowbreak{}Queueing.\allowbreak{}stationary\allowbreak{}Priority\allowbreak{}Class\allowbreak{}Tagged\allowbreak{}Finite\allowbreak{}Replay\allowbreak{}Post\allowbreak{}Arrival\allowbreak{}Response\allowbreak{}Time}
\textbf{Source locator:} {\footnotesize\raggedright cited publication, lines 1118-\allowbreak{}1129\par}
\paragraph{Verbatim paper-source connection}
\begin{ReviewVerbatim}
Wi(.) = The expected waiting time of a class-i job                      KLEINROCK, L. 1976. Queueing System, Vols. I and II.
         in the queue in steady state;                                       John Wiley, New York.
 Wi ( )= The expected sojourn time of a class-i job                     KNUDSEN, N. C. 1972. Individual and Social Optimiza-
         in the system in steady state.                                      tion in a Multi-Server Queue With a General Cost-
                                                                             Benefit Structure. Econometrica 40, 515-528.
For the M/M/1 queue with nonpreemptive priority                         KREPS, D. 1990. A Course in Microeconomic Theory.
discipline (Heyman and Sobel 1982, p. 435)                                 Princeton University Press, Princeton, N.J.
                                                                        MARCHAND, M. 1974. Priority Pricing. Mgmt. Sci. 20,
   _ i( L1()
 W1CX)  = W1 Wqc1+ -+
                   Ci AR   (A-)
                      cj. (-                                                  1131-1140.

[Next verbatim source excerpt]

A4. Service Process                                                      Vi'(Xi) = pi + vi Wi (5)
The processor is modeled as an M/M/1 queueing                          where (5) is analogous to (1) with total expected class-
system with nonpreemptive priority among job                           i user cost consisting of two components: a direct
classes. The processing requirement Hi(-) of a class-i                 payment of pi per job, and the expected class-i delay
job, ri, follows an exponential distribution with finite cost, given by vi * Wi(X).
mean ci. Processing requirements are independent           The following theorem generalizes the net value
across jobs and identically distributed within job       maximizing pricing scheme, (3), to the multiclass case.
classes.

[Next verbatim source excerpt]

giving rise to X = A (p) or V'(X) = 4-'(X/A) = p.                     Al. Arrival Processes

  Consider the case where the delay cost v is positive.               The arrival process of class-i jobs to the system, which
Here, the cost faced by an individual user consists of                reflects the aggregation of infinitesimal users' job flow,
two components: a direct payment of p per job, and                    follows a stationary Poisson process with a price-




                                                 This content downloaded from
                               108.30.55.84 on Mon, 24 Aug 2026 10:55:53 UTC
                                         All use subject to https://about.jstor.org/terms

[form-feed in source extraction]
                                                                                Priority Pricingfor the MIMI1 Queue / 873

dependent arrival rate Xi. The R arrival processes are                 by the system (per unit of time) to the organization,
stochastically independent.                                            where we assume for now that the system administra-
\end{ReviewVerbatim}



\paragraph{Lean-expanded library semantic target}
\begin{ReviewVerbatim}
type:
{n : ℕ} →
  (Fin n → ℝ) →
    (i : Fin n) →
      AppliedModelingLib.Queueing.MulticlassStationaryPoissonWorkClassTaggedSample (Fin n) i → ℝ → ℝ → Option ℝ

value:
fun {n : ℕ} (meanService : Fin n → ℝ) (i : Fin n)
    (z : AppliedModelingLib.Queueing.MulticlassStationaryPoissonWorkClassTaggedSample (Fin n) i) (older t : ℝ) =>
  AppliedModelingLib.Queueing.nonpreemptivePriorityRecordedCompletionTime
    (AppliedModelingLib.Queueing.stationaryPriorityClassTaggedFiniteReplayPostArrivalState meanService i z older t)
    ⟨i, 0⟩
\end{ReviewVerbatim}

\paragraph{Exact Lean library declaration}
\begin{ReviewVerbatim}
/-- The selected customer's completion observation in a fully finite
two-sided replay. -/
noncomputable def stationaryPriorityClassTaggedFiniteReplayPostArrivalResponseTime
    {n : ℕ} (meanService : Fin n → ℝ) (i : Fin n)
    (z : MulticlassStationaryPoissonWorkClassTaggedSample (Fin n) i)
    (older t : ℝ) : Option ℝ :=
  nonpreemptivePriorityRecordedCompletionTime
    (stationaryPriorityClassTaggedFiniteReplayPostArrivalState meanService i z older t)
    (Sigma.mk i 0)
\end{ReviewVerbatim}

\paragraph{Direct retained dependencies}
\begin{itemize}\raggedright
\item \hyperlink{library-prerequisite-755cdb28ceaf46d2}{Applied\allowbreak{}Modeling\allowbreak{}Lib.\allowbreak{}Queueing.\allowbreak{}Multiclass\allowbreak{}Stationary\allowbreak{}Poisson\allowbreak{}Work\allowbreak{}Class\allowbreak{}Tagged\allowbreak{}Sample}
\item \hyperlink{library-prerequisite-886cbd948a5fd8b1}{Applied\allowbreak{}Modeling\allowbreak{}Lib.\allowbreak{}Queueing.\allowbreak{}nonpreemptive\allowbreak{}Priority\allowbreak{}Recorded\allowbreak{}Completion\allowbreak{}Time}
\item \hyperlink{library-prerequisite-098447dd1f7443fd}{Applied\allowbreak{}Modeling\allowbreak{}Lib.\allowbreak{}Queueing.\allowbreak{}stationary\allowbreak{}Priority\allowbreak{}Class\allowbreak{}Tagged\allowbreak{}Finite\allowbreak{}Replay\allowbreak{}Post\allowbreak{}Arrival\allowbreak{}State}
\end{itemize}
\paragraph{Recorded source-to-library screening}
\textbf{Verdict:} \texttt{matches}
\\
{\raggedright
\textbf{Reason:} This is the selected customer's finite-\allowbreak{}replay completion observation supporting the Appendix stationary wait bridge.\allowbreak{}
\par}
\reviewmatch{Matches source input}
\noindent\textbf{Reviewer annotation}\par
\reviewerbox
\clearpage
\hypertarget{library-prerequisite-c3ca5cbd067c9957}{}
\subsection*{\small\ttfamily\raggedright Applied\allowbreak{}Modeling\allowbreak{}Lib.\allowbreak{}Queueing.\allowbreak{}stationary\allowbreak{}Priority\allowbreak{}Class\allowbreak{}Tagged\allowbreak{}Palm\allowbreak{}Measure}
\textbf{Source locator:} {\footnotesize\raggedright cited publication, lines 227-\allowbreak{}242\par}
\paragraph{Verbatim paper-source connection}
\begin{ReviewVerbatim}
giving rise to X = A (p) or V'(X) = 4-'(X/A) = p.                     Al. Arrival Processes

  Consider the case where the delay cost v is positive.               The arrival process of class-i jobs to the system, which
Here, the cost faced by an individual user consists of                reflects the aggregation of infinitesimal users' job flow,
two components: a direct payment of p per job, and                    follows a stationary Poisson process with a price-




                                                 This content downloaded from
                               108.30.55.84 on Mon, 24 Aug 2026 10:55:53 UTC
                                         All use subject to https://about.jstor.org/terms

[form-feed in source extraction]
                                                                                Priority Pricingfor the MIMI1 Queue / 873

dependent arrival rate Xi. The R arrival processes are                 by the system (per unit of time) to the organization,
stochastically independent.                                            where we assume for now that the system administra-

[Next verbatim source excerpt]

 Lq *)= The mean number of class-i jobs in the                          KLEINROCK, L. 1967. Optimum Bribing for Queue Posi-
         queue in steady state;                                              tion. Opns. Res. 15, 304-318.
Wi(.) = The expected waiting time of a class-i job                      KLEINROCK, L. 1976. Queueing System, Vols. I and II.
         in the queue in steady state;                                       John Wiley, New York.
 Wi ( )= The expected sojourn time of a class-i job                     KNUDSEN, N. C. 1972. Individual and Social Optimiza-
         in the system in steady state.                                      tion in a Multi-Server Queue With a General Cost-
\end{ReviewVerbatim}



\paragraph{Lean-expanded library semantic target}
\begin{ReviewVerbatim}
type:
{n : ℕ} →
  (arrivalRate : Fin n → ℝ) →
    (∀ (j : Fin n), 0 < arrivalRate j) →
      (i : Fin n) →
        MeasureTheory.Measure (AppliedModelingLib.Queueing.MulticlassStationaryPoissonWorkClassTaggedSample (Fin n) i)

value:
fun {n : ℕ} (arrivalRate : Fin n → ℝ) (harrivalRate : ∀ (j : Fin n), 0 < arrivalRate j) (i : Fin n) =>
  (AppliedModelingLib.Probability.Palm.targetPassiveTaggedArrivalAtZero
      (AppliedModelingLib.Probability.Queueing.stationaryPoissonWorkTaggedArrivalAtZero (harrivalRate i))
      (AppliedModelingLib.Queueing.multiclassStationaryPoissonWorkRestLaw arrivalRate harrivalRate i)).Ptag
\end{ReviewVerbatim}

\paragraph{Exact Lean library declaration}
\begin{ReviewVerbatim}
/-- The class-tagged Palm law obtained by inserting one class-`i` arrival at
time zero into the stationary marked-Poisson input. -/
noncomputable def stationaryPriorityClassTaggedPalmMeasure
    {n : ℕ} (arrivalRate : Fin n → ℝ)
    (harrivalRate : ∀ j, 0 < arrivalRate j) (i : Fin n) :
    Measure (MulticlassStationaryPoissonWorkClassTaggedSample (Fin n) i) :=
  (Probability.Palm.targetPassiveTaggedArrivalAtZero
    (Probability.Queueing.stationaryPoissonWorkTaggedArrivalAtZero (harrivalRate i))
    (multiclassStationaryPoissonWorkRestLaw arrivalRate harrivalRate i)).Ptag
\end{ReviewVerbatim}

\paragraph{Direct retained dependencies}
\begin{itemize}\raggedright
\item \hyperlink{governing-declaration-5a10835f452c324b}{Applied\allowbreak{}Modeling\allowbreak{}Lib.\allowbreak{}Probability.\allowbreak{}Palm.\allowbreak{}target\allowbreak{}Passive\allowbreak{}Tagged\allowbreak{}Arrival\allowbreak{}At\allowbreak{}Zero}
\item \hyperlink{governing-declaration-ab1cfe37ad8c32b0}{Applied\allowbreak{}Modeling\allowbreak{}Lib.\allowbreak{}Probability.\allowbreak{}Poisson\allowbreak{}Process.\allowbreak{}Good\allowbreak{}Suspension\allowbreak{}State}
\item \hyperlink{governing-declaration-a8385d799d0a4d16}{Applied\allowbreak{}Modeling\allowbreak{}Lib.\allowbreak{}Probability.\allowbreak{}Poisson\allowbreak{}Process.\allowbreak{}inst\allowbreak{}Measurable\allowbreak{}Space\allowbreak{}Good\allowbreak{}Suspension\allowbreak{}State}
\item \hyperlink{governing-declaration-a69bedb5b0d420c6}{Applied\allowbreak{}Modeling\allowbreak{}Lib.\allowbreak{}Probability.\allowbreak{}Queueing.\allowbreak{}Tagged\allowbreak{}Arrival\allowbreak{}At\allowbreak{}Zero}
\item \hyperlink{governing-declaration-e9116e504d73a45c}{Applied\allowbreak{}Modeling\allowbreak{}Lib.\allowbreak{}Probability.\allowbreak{}Queueing.\allowbreak{}stationary\allowbreak{}Poisson\allowbreak{}Work\allowbreak{}Tagged\allowbreak{}Arrival\allowbreak{}At\allowbreak{}Zero}
\item \hyperlink{library-prerequisite-755cdb28ceaf46d2}{Applied\allowbreak{}Modeling\allowbreak{}Lib.\allowbreak{}Queueing.\allowbreak{}Multiclass\allowbreak{}Stationary\allowbreak{}Poisson\allowbreak{}Work\allowbreak{}Class\allowbreak{}Tagged\allowbreak{}Sample}
\item \hyperlink{governing-declaration-62a268164b6ae28b}{Applied\allowbreak{}Modeling\allowbreak{}Lib.\allowbreak{}Queueing.\allowbreak{}Stationary\allowbreak{}Poisson\allowbreak{}Work\allowbreak{}Path}
\item \hyperlink{governing-declaration-9bede4eaca6a95fb}{Applied\allowbreak{}Modeling\allowbreak{}Lib.\allowbreak{}Queueing.\allowbreak{}multiclass\allowbreak{}Stationary\allowbreak{}Poisson\allowbreak{}Work\allowbreak{}Rest\allowbreak{}Law}
\end{itemize}
\paragraph{Recorded source-to-library screening}
\textbf{Verdict:} \texttt{matches}
\\
{\raggedright
\textbf{Reason:} This is A1's class-\allowbreak{}selected arrival-\allowbreak{}at-\allowbreak{}zero Palm law;\allowbreak{} it is routed as a source-\allowbreak{}model carrier, not as a separate source result.\allowbreak{}
\par}
\reviewmatch{Matches source input}
\noindent\textbf{Reviewer annotation}\par
\reviewerbox
\clearpage
\hypertarget{library-prerequisite-6b2b5577562811ed}{}
\subsection*{\small\ttfamily\raggedright Applied\allowbreak{}Modeling\allowbreak{}Lib.\allowbreak{}Queueing.\allowbreak{}stationary\allowbreak{}Priority\allowbreak{}Class\allowbreak{}Tagged\allowbreak{}Response\allowbreak{}Time}
\textbf{Source locator:} {\footnotesize\raggedright cited publication, lines 1118-\allowbreak{}1129\par}
\paragraph{Verbatim paper-source connection}
\begin{ReviewVerbatim}
Wi(.) = The expected waiting time of a class-i job                      KLEINROCK, L. 1976. Queueing System, Vols. I and II.
         in the queue in steady state;                                       John Wiley, New York.
 Wi ( )= The expected sojourn time of a class-i job                     KNUDSEN, N. C. 1972. Individual and Social Optimiza-
         in the system in steady state.                                      tion in a Multi-Server Queue With a General Cost-
                                                                             Benefit Structure. Econometrica 40, 515-528.
For the M/M/1 queue with nonpreemptive priority                         KREPS, D. 1990. A Course in Microeconomic Theory.
discipline (Heyman and Sobel 1982, p. 435)                                 Princeton University Press, Princeton, N.J.
                                                                        MARCHAND, M. 1974. Priority Pricing. Mgmt. Sci. 20,
   _ i( L1()
 W1CX)  = W1 Wqc1+ -+
                   Ci AR   (A-)
                      cj. (-                                                  1131-1140.

[Next verbatim source excerpt]

A4. Service Process                                                      Vi'(Xi) = pi + vi Wi (5)
The processor is modeled as an M/M/1 queueing                          where (5) is analogous to (1) with total expected class-
system with nonpreemptive priority among job                           i user cost consisting of two components: a direct
classes. The processing requirement Hi(-) of a class-i                 payment of pi per job, and the expected class-i delay
job, ri, follows an exponential distribution with finite cost, given by vi * Wi(X).
mean ci. Processing requirements are independent           The following theorem generalizes the net value
across jobs and identically distributed within job       maximizing pricing scheme, (3), to the multiclass case.
classes.

[Next verbatim source excerpt]

giving rise to X = A (p) or V'(X) = 4-'(X/A) = p.                     Al. Arrival Processes

  Consider the case where the delay cost v is positive.               The arrival process of class-i jobs to the system, which
Here, the cost faced by an individual user consists of                reflects the aggregation of infinitesimal users' job flow,
two components: a direct payment of p per job, and                    follows a stationary Poisson process with a price-




                                                 This content downloaded from
                               108.30.55.84 on Mon, 24 Aug 2026 10:55:53 UTC
                                         All use subject to https://about.jstor.org/terms

[form-feed in source extraction]
                                                                                Priority Pricingfor the MIMI1 Queue / 873

dependent arrival rate Xi. The R arrival processes are                 by the system (per unit of time) to the organization,
stochastically independent.                                            where we assume for now that the system administra-
\end{ReviewVerbatim}



\paragraph{Lean-expanded library semantic target}
\begin{ReviewVerbatim}
type:
{n : ℕ} →
  (Fin n → ℝ) → (i : Fin n) → AppliedModelingLib.Queueing.MulticlassStationaryPoissonWorkClassTaggedSample (Fin n) i → ℝ

value:
fun {n : ℕ} (meanService : Fin n → ℝ) (i : Fin n)
    (a : AppliedModelingLib.Queueing.MulticlassStationaryPoissonWorkClassTaggedSample (Fin n) i) =>
  AppliedModelingLib.Probability.stabilizedFiniteReplayResponse
    (fun (horizon : ℕ) (z : AppliedModelingLib.Queueing.MulticlassStationaryPoissonWorkClassTaggedSample (Fin n) i) =>
      (AppliedModelingLib.Queueing.stationaryPriorityClassTaggedFiniteReplayPostArrivalResponseTime meanService i z
            ↑horizon ↑horizon).getD
        0)
    a
\end{ReviewVerbatim}

\paragraph{Exact Lean library declaration}
\begin{ReviewVerbatim}
/-- The selected customer's stationary response is the canonical limsup of
total two-sided finite replays at matching natural past and future horizons.
On the almost-sure coalescence-and-stabilization event below, this is the
literal physical completion epoch.  Each finite replay uses zero before the
tag has completed; outside the stabilization event the limsup still provides
a total real-valued version. -/
noncomputable def stationaryPriorityClassTaggedResponseTime
    {n : ℕ} (meanService : Fin n → ℝ) (i : Fin n) :
    MulticlassStationaryPoissonWorkClassTaggedSample (Fin n) i → ℝ :=
  Probability.stabilizedFiniteReplayResponse fun horizon z =>
    (stationaryPriorityClassTaggedFiniteReplayPostArrivalResponseTime
      meanService i z (horizon : ℝ) (horizon : ℝ)).getD 0
\end{ReviewVerbatim}

\paragraph{Direct retained dependencies}
\begin{itemize}\raggedright
\item \hyperlink{library-prerequisite-4b8788f2f9337ce3}{Applied\allowbreak{}Modeling\allowbreak{}Lib.\allowbreak{}Probability.\allowbreak{}stabilized\allowbreak{}Finite\allowbreak{}Replay\allowbreak{}Response}
\item \hyperlink{library-prerequisite-755cdb28ceaf46d2}{Applied\allowbreak{}Modeling\allowbreak{}Lib.\allowbreak{}Queueing.\allowbreak{}Multiclass\allowbreak{}Stationary\allowbreak{}Poisson\allowbreak{}Work\allowbreak{}Class\allowbreak{}Tagged\allowbreak{}Sample}
\item \hyperlink{library-prerequisite-945b0f6ad93a6ca5}{Applied\allowbreak{}Modeling\allowbreak{}Lib.\allowbreak{}Queueing.\allowbreak{}stationary\allowbreak{}Priority\allowbreak{}Class\allowbreak{}Tagged\allowbreak{}Finite\allowbreak{}Replay\allowbreak{}Post\allowbreak{}Arrival\allowbreak{}Response\allowbreak{}Time}
\end{itemize}
\paragraph{Recorded source-to-library screening}
\textbf{Verdict:} \texttt{matches}
\\
{\raggedright
\textbf{Reason:} This stabilized finite-\allowbreak{}replay response observable is infrastructure for the mapped Appendix bridge and is not credited as a standalone source formula.\allowbreak{}
\par}
\reviewmatch{Matches source input}
\noindent\textbf{Reviewer annotation}\par
\reviewerbox
\clearpage
\hypertarget{library-prerequisite-8f1ec496ea0db7a0}{}
\subsection*{\small\ttfamily\raggedright Applied\allowbreak{}Modeling\allowbreak{}Lib.\allowbreak{}Queueing.\allowbreak{}stationary\allowbreak{}Priority\allowbreak{}Class\allowbreak{}Tagged\allowbreak{}Work\allowbreak{}Requirement}
\textbf{Source locator:} {\footnotesize\raggedright cited publication, lines 264-\allowbreak{}271\par}
\paragraph{Verbatim paper-source connection}
\begin{ReviewVerbatim}
A4. Service Process                                                      Vi'(Xi) = pi + vi Wi (5)
The processor is modeled as an M/M/1 queueing                          where (5) is analogous to (1) with total expected class-
system with nonpreemptive priority among job                           i user cost consisting of two components: a direct
classes. The processing requirement Hi(-) of a class-i                 payment of pi per job, and the expected class-i delay
job, ri, follows an exponential distribution with finite cost, given by vi * Wi(X).
mean ci. Processing requirements are independent           The following theorem generalizes the net value
across jobs and identically distributed within job       maximizing pricing scheme, (3), to the multiclass case.
classes.
\end{ReviewVerbatim}



\paragraph{Lean-expanded library semantic target}
\begin{ReviewVerbatim}
type:
{Class : Type u_1} →
  (Class → ℝ) →
    (i : Class) →
      ((ℤ → ℝ) × (ℤ → ℝ)) × ({ j : Class // j ≠ i } → AppliedModelingLib.Queueing.StationaryPoissonWorkPath) → ℝ

value:
fun {Class : Type u_1} (meanService : Class → ℝ) (i : Class)
    (a : ((ℤ → ℝ) × (ℤ → ℝ)) × ({ j : Class // j ≠ i } → AppliedModelingLib.Queueing.StationaryPoissonWorkPath)) =>
  meanService i * AppliedModelingLib.Queueing.multiclassStationaryPoissonWorkClassTaggedRequirement i a
\end{ReviewVerbatim}

\paragraph{Exact Lean library declaration}
\begin{ReviewVerbatim}
/-- The scaled work requirement of a class-`i` customer pinned at time zero
in the concrete multiclass Palm input. -/
def stationaryPriorityClassTaggedWorkRequirement
    (meanService : Class → ℝ) (i : Class) :
    ((ℤ → ℝ) × (ℤ → ℝ)) ×
      ({j : Class // j ≠ i} → StationaryPoissonWorkPath) → ℝ :=
  fun z => meanService i * multiclassStationaryPoissonWorkClassTaggedRequirement i z
\end{ReviewVerbatim}

\paragraph{Direct retained dependencies}
\begin{itemize}\raggedright
\item \hyperlink{governing-declaration-62a268164b6ae28b}{Applied\allowbreak{}Modeling\allowbreak{}Lib.\allowbreak{}Queueing.\allowbreak{}Stationary\allowbreak{}Poisson\allowbreak{}Work\allowbreak{}Path}
\item \hyperlink{governing-declaration-2317cefa31cefbe1}{Applied\allowbreak{}Modeling\allowbreak{}Lib.\allowbreak{}Queueing.\allowbreak{}multiclass\allowbreak{}Stationary\allowbreak{}Poisson\allowbreak{}Work\allowbreak{}Class\allowbreak{}Tagged\allowbreak{}Requirement}
\end{itemize}
\paragraph{Recorded source-to-library screening}
\textbf{Verdict:} \texttt{matches}
\\
{\raggedright
\textbf{Reason:} It is the selected class's scaled exponential work observable needed by A4's service process.\allowbreak{}
\par}
\reviewmatch{Matches source input}
\noindent\textbf{Reviewer annotation}\par
\reviewerbox
\clearpage
\hypertarget{library-prerequisite-e35fa775c8012225}{}
\subsection*{\small\ttfamily\raggedright Applied\allowbreak{}Modeling\allowbreak{}Lib.\allowbreak{}Queueing.\allowbreak{}stationary\allowbreak{}Priority\allowbreak{}Class\allowbreak{}Tagged\allowbreak{}Queue\allowbreak{}Wait}
\textbf{Source locator:} {\footnotesize\raggedright cited publication, lines 1118-\allowbreak{}1129\par}
\paragraph{Verbatim paper-source connection}
\begin{ReviewVerbatim}
Wi(.) = The expected waiting time of a class-i job                      KLEINROCK, L. 1976. Queueing System, Vols. I and II.
         in the queue in steady state;                                       John Wiley, New York.
 Wi ( )= The expected sojourn time of a class-i job                     KNUDSEN, N. C. 1972. Individual and Social Optimiza-
         in the system in steady state.                                      tion in a Multi-Server Queue With a General Cost-
                                                                             Benefit Structure. Econometrica 40, 515-528.
For the M/M/1 queue with nonpreemptive priority                         KREPS, D. 1990. A Course in Microeconomic Theory.
discipline (Heyman and Sobel 1982, p. 435)                                 Princeton University Press, Princeton, N.J.
                                                                        MARCHAND, M. 1974. Priority Pricing. Mgmt. Sci. 20,
   _ i( L1()
 W1CX)  = W1 Wqc1+ -+
                   Ci AR   (A-)
                      cj. (-                                                  1131-1140.

[Next verbatim source excerpt]

A4. Service Process                                                      Vi'(Xi) = pi + vi Wi (5)
The processor is modeled as an M/M/1 queueing                          where (5) is analogous to (1) with total expected class-
system with nonpreemptive priority among job                           i user cost consisting of two components: a direct
classes. The processing requirement Hi(-) of a class-i                 payment of pi per job, and the expected class-i delay
job, ri, follows an exponential distribution with finite cost, given by vi * Wi(X).
mean ci. Processing requirements are independent           The following theorem generalizes the net value
across jobs and identically distributed within job       maximizing pricing scheme, (3), to the multiclass case.
classes.

[Next verbatim source excerpt]

giving rise to X = A (p) or V'(X) = 4-'(X/A) = p.                     Al. Arrival Processes

  Consider the case where the delay cost v is positive.               The arrival process of class-i jobs to the system, which
Here, the cost faced by an individual user consists of                reflects the aggregation of infinitesimal users' job flow,
two components: a direct payment of p per job, and                    follows a stationary Poisson process with a price-




                                                 This content downloaded from
                               108.30.55.84 on Mon, 24 Aug 2026 10:55:53 UTC
                                         All use subject to https://about.jstor.org/terms

[form-feed in source extraction]
                                                                                Priority Pricingfor the MIMI1 Queue / 873

dependent arrival rate Xi. The R arrival processes are                 by the system (per unit of time) to the organization,
stochastically independent.                                            where we assume for now that the system administra-
\end{ReviewVerbatim}



\paragraph{Lean-expanded library semantic target}
\begin{ReviewVerbatim}
type:
{n : ℕ} →
  (Fin n → ℝ) → (i : Fin n) → AppliedModelingLib.Queueing.MulticlassStationaryPoissonWorkClassTaggedSample (Fin n) i → ℝ

value:
fun {n : ℕ} (meanService : Fin n → ℝ) (i : Fin n)
    (a : AppliedModelingLib.Queueing.MulticlassStationaryPoissonWorkClassTaggedSample (Fin n) i) =>
  AppliedModelingLib.Queueing.stationaryPriorityClassTaggedResponseTime meanService i a -
    AppliedModelingLib.Queueing.stationaryPriorityClassTaggedWorkRequirement meanService i a
\end{ReviewVerbatim}

\paragraph{Exact Lean library declaration}
\begin{ReviewVerbatim}
/-- The selected customer's queue-wait observable: its literal response time
minus its own required service.  The nonpreemptive trace is the object whose
pathwise workload decomposition will identify this difference with service
completed before the tag starts. -/
noncomputable def stationaryPriorityClassTaggedQueueWait
    {n : ℕ} (meanService : Fin n → ℝ) (i : Fin n) :
    MulticlassStationaryPoissonWorkClassTaggedSample (Fin n) i → ℝ :=
  fun z => stationaryPriorityClassTaggedResponseTime meanService i z -
    stationaryPriorityClassTaggedWorkRequirement meanService i z
\end{ReviewVerbatim}

\paragraph{Direct retained dependencies}
\begin{itemize}\raggedright
\item \hyperlink{library-prerequisite-755cdb28ceaf46d2}{Applied\allowbreak{}Modeling\allowbreak{}Lib.\allowbreak{}Queueing.\allowbreak{}Multiclass\allowbreak{}Stationary\allowbreak{}Poisson\allowbreak{}Work\allowbreak{}Class\allowbreak{}Tagged\allowbreak{}Sample}
\item \hyperlink{library-prerequisite-6b2b5577562811ed}{Applied\allowbreak{}Modeling\allowbreak{}Lib.\allowbreak{}Queueing.\allowbreak{}stationary\allowbreak{}Priority\allowbreak{}Class\allowbreak{}Tagged\allowbreak{}Response\allowbreak{}Time}
\item \hyperlink{library-prerequisite-8f1ec496ea0db7a0}{Applied\allowbreak{}Modeling\allowbreak{}Lib.\allowbreak{}Queueing.\allowbreak{}stationary\allowbreak{}Priority\allowbreak{}Class\allowbreak{}Tagged\allowbreak{}Work\allowbreak{}Requirement}
\end{itemize}
\paragraph{Recorded source-to-library screening}
\textbf{Verdict:} \texttt{matches}
\\
{\raggedright
\textbf{Reason:} Response time minus own service is the source's queue-\allowbreak{}wait quantity, matching Appendix (A-\allowbreak{}1) rather than sojourn time.\allowbreak{}
\par}
\reviewmatch{Matches source input}
\noindent\textbf{Reviewer annotation}\par
\reviewerbox
\clearpage
\hypertarget{library-prerequisite-1fb8084d5335bf95}{}
\subsection*{\small\ttfamily\raggedright Applied\allowbreak{}Modeling\allowbreak{}Lib.\allowbreak{}Queueing.\allowbreak{}stationary\allowbreak{}Priority\allowbreak{}Work\allowbreak{}Requirement}
\textbf{Source locator:} {\footnotesize\raggedright cited publication, lines 264-\allowbreak{}271\par}
\paragraph{Verbatim paper-source connection}
\begin{ReviewVerbatim}
A4. Service Process                                                      Vi'(Xi) = pi + vi Wi (5)
The processor is modeled as an M/M/1 queueing                          where (5) is analogous to (1) with total expected class-
system with nonpreemptive priority among job                           i user cost consisting of two components: a direct
classes. The processing requirement Hi(-) of a class-i                 payment of pi per job, and the expected class-i delay
job, ri, follows an exponential distribution with finite cost, given by vi * Wi(X).
mean ci. Processing requirements are independent           The following theorem generalizes the net value
across jobs and identically distributed within job       maximizing pricing scheme, (3), to the multiclass case.
classes.
\end{ReviewVerbatim}



\paragraph{Lean-expanded library semantic target}
\begin{ReviewVerbatim}
type:
{Class : Type u_1} →
  (Class → ℝ) → Class → (Class → AppliedModelingLib.Probability.PoissonProcess.GoodSuspensionState × (ℤ → ℝ)) → ℤ → ℝ

value:
fun {Class : Type u_1} (meanService : Class → ℝ) (i : Class)
    (a : Class → AppliedModelingLib.Probability.PoissonProcess.GoodSuspensionState × (ℤ → ℝ)) (a_1 : ℤ) =>
  meanService i * AppliedModelingLib.Probability.PoissonProcess.multiclassStationaryPoissonWorkRequirement i a a_1
\end{ReviewVerbatim}

\paragraph{Exact Lean library declaration}
\begin{ReviewVerbatim}
/-- The work requirement of a class-labelled arrival when its class has the
given positive mean service requirement.  It scales the concrete unit-rate
exponential mark on the stationary multiclass input carrier. -/
def stationaryPriorityWorkRequirement
    (meanService : Class → ℝ) (i : Class) :
    (Class → (Probability.PoissonProcess.GoodSuspensionState × (ℤ → ℝ))) → ℤ → ℝ :=
  fun omega n => meanService i *
    Probability.PoissonProcess.multiclassStationaryPoissonWorkRequirement i omega n
\end{ReviewVerbatim}

\paragraph{Direct retained dependencies}
\begin{itemize}\raggedright
\item \hyperlink{governing-declaration-ab1cfe37ad8c32b0}{Applied\allowbreak{}Modeling\allowbreak{}Lib.\allowbreak{}Probability.\allowbreak{}Poisson\allowbreak{}Process.\allowbreak{}Good\allowbreak{}Suspension\allowbreak{}State}
\item \hyperlink{governing-declaration-13879bfc74e6d5d0}{Applied\allowbreak{}Modeling\allowbreak{}Lib.\allowbreak{}Probability.\allowbreak{}Poisson\allowbreak{}Process.\allowbreak{}multiclass\allowbreak{}Stationary\allowbreak{}Poisson\allowbreak{}Work\allowbreak{}Requirement}
\end{itemize}
\paragraph{Recorded source-to-library screening}
\textbf{Verdict:} \texttt{matches}
\\
{\raggedright
\textbf{Reason:} It supplies A4's class-\allowbreak{}scaled unit-\allowbreak{}exponential service-\allowbreak{}work mark.\allowbreak{}
\par}
\reviewmatch{Matches source input}
\noindent\textbf{Reviewer annotation}\par
\reviewerbox
\clearpage
\hypertarget{library-prerequisite-ea60ce4afba27cd7}{}
\subsection*{\small\ttfamily\raggedright Applied\allowbreak{}Modeling\allowbreak{}Lib.\allowbreak{}Queueing.\allowbreak{}stationary\allowbreak{}Priority\allowbreak{}Work\allowbreak{}Requirement\_\allowbreak{}has\allowbreak{}Law}
\textbf{Source locator:} {\footnotesize\raggedright cited publication, lines 264-\allowbreak{}271\par}
\paragraph{Verbatim paper-source connection}
\begin{ReviewVerbatim}
A4. Service Process                                                      Vi'(Xi) = pi + vi Wi (5)
The processor is modeled as an M/M/1 queueing                          where (5) is analogous to (1) with total expected class-
system with nonpreemptive priority among job                           i user cost consisting of two components: a direct
classes. The processing requirement Hi(-) of a class-i                 payment of pi per job, and the expected class-i delay
job, ri, follows an exponential distribution with finite cost, given by vi * Wi(X).
mean ci. Processing requirements are independent           The following theorem generalizes the net value
across jobs and identically distributed within job       maximizing pricing scheme, (3), to the multiclass case.
classes.
\end{ReviewVerbatim}



\paragraph{Lean declaration type (metadata)}
\begin{ReviewVerbatim}
∀ {Class : Type u_1} [inst : Fintype Class] (arrivalRate meanService : Class → ℝ),
  (∀ (i : Class), 0 < arrivalRate i) →
    (∀ (i : Class), 0 < meanService i) →
      ∀ (i : Class) (n : ℤ),
        ProbabilityTheory.HasLaw
          (fun (omega : Class → AppliedModelingLib.Probability.PoissonProcess.GoodSuspensionState × (ℤ → ℝ)) =>
            AppliedModelingLib.Queueing.stationaryPriorityWorkRequirement meanService i omega n)
          (ProbabilityTheory.expMeasure (meanService i)⁻¹)
          (AppliedModelingLib.Probability.PoissonProcess.multiclassStationaryPoissonWorkMeasure arrivalRate)
\end{ReviewVerbatim}

\paragraph{Exact Lean library declaration}
\begin{ReviewVerbatim}
/-- A class-specific scaled work mark has the exponential law whose rate is
the reciprocal of its mean service requirement. -/
theorem stationaryPriorityWorkRequirement_hasLaw
    (arrivalRate meanService : Class → ℝ)
    (harrivalRate : ∀ i, 0 < arrivalRate i)
    (hmeanService : ∀ i, 0 < meanService i)
    (i : Class) (n : ℤ) :
    HasLaw (fun omega : Class →
      (Probability.PoissonProcess.GoodSuspensionState × (ℤ → ℝ)) =>
        stationaryPriorityWorkRequirement meanService i omega n)
      (expMeasure (meanService i)⁻¹)
      (Probability.PoissonProcess.multiclassStationaryPoissonWorkMeasure arrivalRate) := by
  have hscale :
      Measure.map (fun x : ℝ => meanService i * x) (expMeasure (1 : ℝ)) =
        expMeasure (meanService i)⁻¹ := by
    calc
      Measure.map (fun x : ℝ => meanService i * x) (expMeasure (1 : ℝ)) =
          Measure.map (fun x : ℝ => x / (meanService i)⁻¹) (expMeasure (1 : ℝ)) := by
            congr 1
            funext x
            field_simp [ne_of_gt (hmeanService i)]
      _ = expMeasure (meanService i)⁻¹ := by
            exact Probability.map_div_unitExpMeasure_eq_expMeasure
              (inv_pos.mpr (hmeanService i))
  exact HasLaw.comp
    ⟨by fun_prop, hscale⟩
    (Probability.PoissonProcess.multiclassStationaryPoissonWorkRequirement_hasLaw
      arrivalRate harrivalRate i n)
\end{ReviewVerbatim}

\paragraph{Direct retained dependencies}
\begin{itemize}\raggedright
\item \hyperlink{governing-declaration-ab1cfe37ad8c32b0}{Applied\allowbreak{}Modeling\allowbreak{}Lib.\allowbreak{}Probability.\allowbreak{}Poisson\allowbreak{}Process.\allowbreak{}Good\allowbreak{}Suspension\allowbreak{}State}
\item \hyperlink{governing-declaration-a8385d799d0a4d16}{Applied\allowbreak{}Modeling\allowbreak{}Lib.\allowbreak{}Probability.\allowbreak{}Poisson\allowbreak{}Process.\allowbreak{}inst\allowbreak{}Measurable\allowbreak{}Space\allowbreak{}Good\allowbreak{}Suspension\allowbreak{}State}
\item \hyperlink{library-prerequisite-dba97dbb9d8c7c91}{Applied\allowbreak{}Modeling\allowbreak{}Lib.\allowbreak{}Probability.\allowbreak{}Poisson\allowbreak{}Process.\allowbreak{}multiclass\allowbreak{}Stationary\allowbreak{}Poisson\allowbreak{}Work\allowbreak{}Measure}
\item \hyperlink{library-prerequisite-1fb8084d5335bf95}{Applied\allowbreak{}Modeling\allowbreak{}Lib.\allowbreak{}Queueing.\allowbreak{}stationary\allowbreak{}Priority\allowbreak{}Work\allowbreak{}Requirement}
\end{itemize}
\paragraph{Recorded source-to-library screening}
\textbf{Verdict:} \texttt{matches}
\\
{\raggedright
\textbf{Reason:} It establishes the A4 exponential work law at reciprocal mean-\allowbreak{}service rate.\allowbreak{}
\par}
\reviewmatch{Matches source input}
\noindent\textbf{Reviewer annotation}\par
\reviewerbox
\clearpage
\hypertarget{paper-prerequisite-section-5ef5ef0364b6939c}{}
\section*{Paper-specific semantic prerequisites}
\hypertarget{paper-prerequisite-60ffddd9622d7d65}{}
\subsection*{\small\ttfamily\raggedright Mendelson\allowbreak{}Whang1990\allowbreak{}Priority\allowbreak{}Pricing.\allowbreak{}priority\allowbreak{}Time\allowbreak{}Dependent\allowbreak{}Price}
\noindent\textbf{Paper-local declaration:}\par
{\footnotesize\ttfamily\raggedright Mendelson\allowbreak{}Whang1990\allowbreak{}Priority\allowbreak{}Pricing.\allowbreak{}priority\allowbreak{}Time\allowbreak{}Dependent\allowbreak{}Price\par}
\textbf{Source locator:} {\footnotesize\raggedright cited publication, lines 760-\allowbreak{}805\par}
\paragraph{Verbatim paper-source connection}
\begin{ReviewVerbatim}
and incentive-compatibility in the general setting                   E[pt+(ri)] = E[Ai Ti + 112B *Ti r]
with heterogeneous service requirements. What price                                 = Ai * ci + B * c?,
structure is optimal and incentive compatible? The
                                                                    since E[r,] = 2ci . It follows that
specific PTD pricing scheme we propose is given by
p+(t) = (p+(t), p2+(t), . . ., pZ(t)) with                          p* = E[pt((ri)] for alli= 1, 2, ...,R- 1.
                                                                    Likewise
Pti(t) = Ait + V12Bt2 (19)

where                                                                A
                                                                     PR= -=-a'R
                                                                              -2CR V XkR
                                                                                 CR+(E    2= E[p+(TR)]
                                                                                       - -)R E[R(R)
                                                                           SR- SR k=1 Sk-l Sk

                                                                    which completes the proof of optimality.
 B        =     E     .k    *     Xk(20)
       k= 1 Sk-l Sk                                                    b. Next we prove incentive-compatibility, that is,
                                                                    class-i users will find it in their best interest to join
                                                                    priority-class i. As in Theorem 2, we prove local
 Ai 2+ 2 + - __ 2) (21)                                             optimality and then transitivity.
         z-1 9 k=i+l 'kISk k Sk-1SIk
                                                                       i. Local optimality: If a class-i job (i = 1,
with Ai = EL cIX*, ai = ViXA AR, Si = REk=1     ckXCk
                                        - 1) joins priority-class i, its expected total cost
So = 0, S) = 1 - Si, and where X * are evaluated at the
                                                                     is E[pt((ri)] + vi * (Wi7 + ci), where Wi' denotes
optimum of Theorem 1, and summations from R to
                                                                    the expected queueing time (excluding service) for
(R - 1) are interpreted as zero (see the Appendix for
                                                                    priority-class i. On the other hand, if the user chooses
a summary of the notation). Expression (21) may also
                                                                    priority-class (i + 1), the expected total cost of the job
be written in the equivalent recursive form
                                                                     is given by E[pl+(Tri)] + vi * (Wl1+1 + ci). Thus, for
                                                                     all i= 1, 2, . . ., R - I wehave
Ai=Ai-,- a- -2 + -2ai_) i=2,3 ... R
                                                                     ELpo(Tri)] + vi Wi7- EE[p +I(ri)] + viWi'+ I

with                                                                            r a, ai+ + lAR AR
                                                                           =Ni I 37 +1 13i23+1 CI -S
           aR
AR = - -2                                                                   _AR
        SR-1 * SR
                                                                                          - 1 i S [civS3i+1 + civi+1 IX1 1i-I
\end{ReviewVerbatim}



\paragraph{Lean-expanded paper semantic target}
\begin{ReviewVerbatim}
type:
ℝ → ℝ → ℝ → ℝ

value:
fun (linearCoefficient quadraticCoefficient processingTime : ℝ) =>
  linearCoefficient * processingTime + 1 / 2 * quadraticCoefficient * processingTime ^ 2
\end{ReviewVerbatim}


\paragraph{Recorded source-to-declaration screening}
\textbf{Verdict:} \texttt{matches}
\\
{\raggedright
\textbf{Reason:} The expanded Lean value is A\_\allowbreak{}i t + one-\allowbreak{}half B t squared, exactly Equation (19)'s processing-\allowbreak{}time-\allowbreak{}dependent charge for one declared priority class.\allowbreak{}
\par}
\reviewmatch{Matches source input}
\noindent\textbf{Reviewer annotation}\par
\reviewerbox
\clearpage
\hypertarget{paper-prerequisite-ae05cc1b41c33470}{}
\subsection*{\small\ttfamily\raggedright Mendelson\allowbreak{}Whang1990\allowbreak{}Priority\allowbreak{}Pricing.\allowbreak{}priority\allowbreak{}Time\allowbreak{}Linear\allowbreak{}Coefficient}
\noindent\textbf{Paper-local declaration:}\par
{\footnotesize\ttfamily\raggedright Mendelson\allowbreak{}Whang1990\allowbreak{}Priority\allowbreak{}Pricing.\allowbreak{}priority\allowbreak{}Time\allowbreak{}Linear\allowbreak{}Coefficient\par}
\textbf{Source locator:} {\footnotesize\raggedright cited publication, lines 760-\allowbreak{}805\par}
\paragraph{Verbatim paper-source connection}
\begin{ReviewVerbatim}
and incentive-compatibility in the general setting                   E[pt+(ri)] = E[Ai Ti + 112B *Ti r]
with heterogeneous service requirements. What price                                 = Ai * ci + B * c?,
structure is optimal and incentive compatible? The
                                                                    since E[r,] = 2ci . It follows that
specific PTD pricing scheme we propose is given by
p+(t) = (p+(t), p2+(t), . . ., pZ(t)) with                          p* = E[pt((ri)] for alli= 1, 2, ...,R- 1.
                                                                    Likewise
Pti(t) = Ait + V12Bt2 (19)

where                                                                A
                                                                     PR= -=-a'R
                                                                              -2CR V XkR
                                                                                 CR+(E    2= E[p+(TR)]
                                                                                       - -)R E[R(R)
                                                                           SR- SR k=1 Sk-l Sk

                                                                    which completes the proof of optimality.
 B        =     E     .k    *     Xk(20)
       k= 1 Sk-l Sk                                                    b. Next we prove incentive-compatibility, that is,
                                                                    class-i users will find it in their best interest to join
                                                                    priority-class i. As in Theorem 2, we prove local
 Ai 2+ 2 + - __ 2) (21)                                             optimality and then transitivity.
         z-1 9 k=i+l 'kISk k Sk-1SIk
                                                                       i. Local optimality: If a class-i job (i = 1,
with Ai = EL cIX*, ai = ViXA AR, Si = REk=1     ckXCk
                                        - 1) joins priority-class i, its expected total cost
So = 0, S) = 1 - Si, and where X * are evaluated at the
                                                                     is E[pt((ri)] + vi * (Wi7 + ci), where Wi' denotes
optimum of Theorem 1, and summations from R to
                                                                    the expected queueing time (excluding service) for
(R - 1) are interpreted as zero (see the Appendix for
                                                                    priority-class i. On the other hand, if the user chooses
a summary of the notation). Expression (21) may also
                                                                    priority-class (i + 1), the expected total cost of the job
be written in the equivalent recursive form
                                                                     is given by E[pl+(Tri)] + vi * (Wl1+1 + ci). Thus, for
                                                                     all i= 1, 2, . . ., R - I wehave
Ai=Ai-,- a- -2 + -2ai_) i=2,3 ... R
                                                                     ELpo(Tri)] + vi Wi7- EE[p +I(ri)] + viWi'+ I

with                                                                            r a, ai+ + lAR AR
                                                                           =Ni I 37 +1 13i23+1 CI -S
           aR
AR = - -2                                                                   _AR
        SR-1 * SR
                                                                                          - 1 i S [civS3i+1 + civi+1 IX1 1i-I
\end{ReviewVerbatim}



\paragraph{Lean-expanded paper semantic target}
\begin{ReviewVerbatim}
type:
{n : ℕ} → (Fin n → ℝ) → (Fin n → ℝ) → (Fin n → ℝ) → Fin n → ℝ

value:
fun {n : ℕ} (arrivalRate delayCost meanService : Fin n → ℝ) (i : Fin n) =>
  MendelsonWhang1990PriorityPricing.priorityLinearNumerator arrivalRate delayCost meanService i /
      (AppliedModelingLib.Queueing.finitePriorityStrictSlack meanService arrivalRate i *
        AppliedModelingLib.Queueing.finitePriorityInclusiveSlack meanService arrivalRate i ^ 2) +
    ∑ k : Fin n,
      if i < k then
        (1 /
              (AppliedModelingLib.Queueing.finitePriorityStrictSlack meanService arrivalRate k ^ 2 *
                AppliedModelingLib.Queueing.finitePriorityInclusiveSlack meanService arrivalRate k) +
            1 /
              (AppliedModelingLib.Queueing.finitePriorityStrictSlack meanService arrivalRate k *
                AppliedModelingLib.Queueing.finitePriorityInclusiveSlack meanService arrivalRate k ^ 2)) *
          MendelsonWhang1990PriorityPricing.priorityLinearNumerator arrivalRate delayCost meanService k
      else 0
\end{ReviewVerbatim}

\paragraph{Direct retained dependencies}
\begin{itemize}\raggedright
\item \hyperlink{governing-declaration-f2b9a1af307e7469}{Applied\allowbreak{}Modeling\allowbreak{}Lib.\allowbreak{}Queueing.\allowbreak{}finite\allowbreak{}Priority\allowbreak{}Inclusive\allowbreak{}Slack}
\item \hyperlink{governing-declaration-3f5cf610a8190094}{Applied\allowbreak{}Modeling\allowbreak{}Lib.\allowbreak{}Queueing.\allowbreak{}finite\allowbreak{}Priority\allowbreak{}Strict\allowbreak{}Slack}
\item \hyperlink{governing-declaration-96f4c7bfcd6904cc}{Mendelson\allowbreak{}Whang1990\allowbreak{}Priority\allowbreak{}Pricing.\allowbreak{}priority\allowbreak{}Linear\allowbreak{}Numerator}
\end{itemize}
\paragraph{Recorded source-to-declaration screening}
\textbf{Verdict:} \texttt{matches}
\\
{\raggedright
\textbf{Reason:} Its strict/\allowbreak{}inclusive slack terms and lower-\allowbreak{}priority suffix are Equation (21)'s A\_\allowbreak{}i coefficient.\allowbreak{}
\par}
\reviewmatch{Matches source input}
\noindent\textbf{Reviewer annotation}\par
\reviewerbox
\clearpage
\hypertarget{paper-prerequisite-1840b9f7b498b425}{}
\subsection*{\small\ttfamily\raggedright Mendelson\allowbreak{}Whang1990\allowbreak{}Priority\allowbreak{}Pricing.\allowbreak{}priority\allowbreak{}Time\allowbreak{}Quadratic\allowbreak{}Coefficient}
\noindent\textbf{Paper-local declaration:}\par
{\footnotesize\ttfamily\raggedright Mendelson\allowbreak{}Whang1990\allowbreak{}Priority\allowbreak{}Pricing.\allowbreak{}priority\allowbreak{}Time\allowbreak{}Quadratic\allowbreak{}Coefficient\par}
\textbf{Source locator:} {\footnotesize\raggedright cited publication, lines 760-\allowbreak{}805\par}
\paragraph{Verbatim paper-source connection}
\begin{ReviewVerbatim}
and incentive-compatibility in the general setting                   E[pt+(ri)] = E[Ai Ti + 112B *Ti r]
with heterogeneous service requirements. What price                                 = Ai * ci + B * c?,
structure is optimal and incentive compatible? The
                                                                    since E[r,] = 2ci . It follows that
specific PTD pricing scheme we propose is given by
p+(t) = (p+(t), p2+(t), . . ., pZ(t)) with                          p* = E[pt((ri)] for alli= 1, 2, ...,R- 1.
                                                                    Likewise
Pti(t) = Ait + V12Bt2 (19)

where                                                                A
                                                                     PR= -=-a'R
                                                                              -2CR V XkR
                                                                                 CR+(E    2= E[p+(TR)]
                                                                                       - -)R E[R(R)
                                                                           SR- SR k=1 Sk-l Sk

                                                                    which completes the proof of optimality.
 B        =     E     .k    *     Xk(20)
       k= 1 Sk-l Sk                                                    b. Next we prove incentive-compatibility, that is,
                                                                    class-i users will find it in their best interest to join
                                                                    priority-class i. As in Theorem 2, we prove local
 Ai 2+ 2 + - __ 2) (21)                                             optimality and then transitivity.
         z-1 9 k=i+l 'kISk k Sk-1SIk
                                                                       i. Local optimality: If a class-i job (i = 1,
with Ai = EL cIX*, ai = ViXA AR, Si = REk=1     ckXCk
                                        - 1) joins priority-class i, its expected total cost
So = 0, S) = 1 - Si, and where X * are evaluated at the
                                                                     is E[pt((ri)] + vi * (Wi7 + ci), where Wi' denotes
optimum of Theorem 1, and summations from R to
                                                                    the expected queueing time (excluding service) for
(R - 1) are interpreted as zero (see the Appendix for
                                                                    priority-class i. On the other hand, if the user chooses
a summary of the notation). Expression (21) may also
                                                                    priority-class (i + 1), the expected total cost of the job
be written in the equivalent recursive form
                                                                     is given by E[pl+(Tri)] + vi * (Wl1+1 + ci). Thus, for
                                                                     all i= 1, 2, . . ., R - I wehave
Ai=Ai-,- a- -2 + -2ai_) i=2,3 ... R
                                                                     ELpo(Tri)] + vi Wi7- EE[p +I(ri)] + viWi'+ I

with                                                                            r a, ai+ + lAR AR
                                                                           =Ni I 37 +1 13i23+1 CI -S
           aR
AR = - -2                                                                   _AR
        SR-1 * SR
                                                                                          - 1 i S [civS3i+1 + civi+1 IX1 1i-I
\end{ReviewVerbatim}



\paragraph{Lean-expanded paper semantic target}
\begin{ReviewVerbatim}
type:
{n : ℕ} → (Fin n → ℝ) → (Fin n → ℝ) → (Fin n → ℝ) → ℝ

value:
fun {n : ℕ} (arrivalRate delayCost meanService : Fin n → ℝ) =>
  ∑ k : Fin n,
    delayCost k * arrivalRate k /
      (AppliedModelingLib.Queueing.finitePriorityStrictSlack meanService arrivalRate k *
        AppliedModelingLib.Queueing.finitePriorityInclusiveSlack meanService arrivalRate k)
\end{ReviewVerbatim}

\paragraph{Direct retained dependencies}
\begin{itemize}\raggedright
\item \hyperlink{governing-declaration-f2b9a1af307e7469}{Applied\allowbreak{}Modeling\allowbreak{}Lib.\allowbreak{}Queueing.\allowbreak{}finite\allowbreak{}Priority\allowbreak{}Inclusive\allowbreak{}Slack}
\item \hyperlink{governing-declaration-3f5cf610a8190094}{Applied\allowbreak{}Modeling\allowbreak{}Lib.\allowbreak{}Queueing.\allowbreak{}finite\allowbreak{}Priority\allowbreak{}Strict\allowbreak{}Slack}
\end{itemize}
\paragraph{Recorded source-to-declaration screening}
\textbf{Verdict:} \texttt{matches}
\\
{\raggedright
\textbf{Reason:} Its common finite sum is Equation (20)'s B coefficient.\allowbreak{}
\par}
\reviewmatch{Matches source input}
\noindent\textbf{Reviewer annotation}\par
\reviewerbox
\clearpage
\hypertarget{paper-prerequisite-d3f248a0395e1c76}{}
\subsection*{\small\ttfamily\raggedright Mendelson\allowbreak{}Whang1990\allowbreak{}Priority\allowbreak{}Pricing.\allowbreak{}heterogeneous\allowbreak{}Priority\allowbreak{}Time\allowbreak{}Dependent\allowbreak{}Price}
\noindent\textbf{Paper-local declaration:}\par
{\footnotesize\ttfamily\raggedright Mendelson\allowbreak{}Whang1990\allowbreak{}Priority\allowbreak{}Pricing.\allowbreak{}heterogeneous\allowbreak{}Priority\allowbreak{}Time\allowbreak{}Dependent\allowbreak{}Price\par}
\textbf{Source locator:} {\footnotesize\raggedright cited publication, lines 760-\allowbreak{}805\par}
\paragraph{Verbatim paper-source connection}
\begin{ReviewVerbatim}
and incentive-compatibility in the general setting                   E[pt+(ri)] = E[Ai Ti + 112B *Ti r]
with heterogeneous service requirements. What price                                 = Ai * ci + B * c?,
structure is optimal and incentive compatible? The
                                                                    since E[r,] = 2ci . It follows that
specific PTD pricing scheme we propose is given by
p+(t) = (p+(t), p2+(t), . . ., pZ(t)) with                          p* = E[pt((ri)] for alli= 1, 2, ...,R- 1.
                                                                    Likewise
Pti(t) = Ait + V12Bt2 (19)

where                                                                A
                                                                     PR= -=-a'R
                                                                              -2CR V XkR
                                                                                 CR+(E    2= E[p+(TR)]
                                                                                       - -)R E[R(R)
                                                                           SR- SR k=1 Sk-l Sk

                                                                    which completes the proof of optimality.
 B        =     E     .k    *     Xk(20)
       k= 1 Sk-l Sk                                                    b. Next we prove incentive-compatibility, that is,
                                                                    class-i users will find it in their best interest to join
                                                                    priority-class i. As in Theorem 2, we prove local
 Ai 2+ 2 + - __ 2) (21)                                             optimality and then transitivity.
         z-1 9 k=i+l 'kISk k Sk-1SIk
                                                                       i. Local optimality: If a class-i job (i = 1,
with Ai = EL cIX*, ai = ViXA AR, Si = REk=1     ckXCk
                                        - 1) joins priority-class i, its expected total cost
So = 0, S) = 1 - Si, and where X * are evaluated at the
                                                                     is E[pt((ri)] + vi * (Wi7 + ci), where Wi' denotes
optimum of Theorem 1, and summations from R to
                                                                    the expected queueing time (excluding service) for
(R - 1) are interpreted as zero (see the Appendix for
                                                                    priority-class i. On the other hand, if the user chooses
a summary of the notation). Expression (21) may also
                                                                    priority-class (i + 1), the expected total cost of the job
be written in the equivalent recursive form
                                                                     is given by E[pl+(Tri)] + vi * (Wl1+1 + ci). Thus, for
                                                                     all i= 1, 2, . . ., R - I wehave
Ai=Ai-,- a- -2 + -2ai_) i=2,3 ... R
                                                                     ELpo(Tri)] + vi Wi7- EE[p +I(ri)] + viWi'+ I

with                                                                            r a, ai+ + lAR AR
                                                                           =Ni I 37 +1 13i23+1 CI -S
           aR
AR = - -2                                                                   _AR
        SR-1 * SR
                                                                                          - 1 i S [civS3i+1 + civi+1 IX1 1i-I
\end{ReviewVerbatim}



\paragraph{Lean-expanded paper semantic target}
\begin{ReviewVerbatim}
type:
{n : ℕ} → (Fin n → ℝ) → (Fin n → ℝ) → (Fin n → ℝ) → Fin n → ℝ → ℝ

value:
fun {n : ℕ} (arrivalRate delayCost meanService : Fin n → ℝ) (priorityClass : Fin n) (processingTime : ℝ) =>
  MendelsonWhang1990PriorityPricing.priorityTimeDependentPrice
    (MendelsonWhang1990PriorityPricing.priorityTimeLinearCoefficient arrivalRate delayCost meanService priorityClass)
    (MendelsonWhang1990PriorityPricing.priorityTimeQuadraticCoefficient arrivalRate delayCost meanService)
    processingTime
\end{ReviewVerbatim}

\paragraph{Direct retained dependencies}
\begin{itemize}\raggedright
\item \hyperlink{paper-prerequisite-60ffddd9622d7d65}{Mendelson\allowbreak{}Whang1990\allowbreak{}Priority\allowbreak{}Pricing.\allowbreak{}priority\allowbreak{}Time\allowbreak{}Dependent\allowbreak{}Price}
\item \hyperlink{paper-prerequisite-ae05cc1b41c33470}{Mendelson\allowbreak{}Whang1990\allowbreak{}Priority\allowbreak{}Pricing.\allowbreak{}priority\allowbreak{}Time\allowbreak{}Linear\allowbreak{}Coefficient}
\item \hyperlink{paper-prerequisite-1840b9f7b498b425}{Mendelson\allowbreak{}Whang1990\allowbreak{}Priority\allowbreak{}Pricing.\allowbreak{}priority\allowbreak{}Time\allowbreak{}Quadratic\allowbreak{}Coefficient}
\end{itemize}
\paragraph{Recorded source-to-declaration screening}
\textbf{Verdict:} \texttt{matches}
\\
{\raggedright
\textbf{Reason:} Equations (19)-\allowbreak{}-\allowbreak{}(21) give the declared-\allowbreak{}class charge as the mapped linear coefficient times service time plus one half the common quadratic coefficient times its square.\allowbreak{}
\par}
\reviewmatch{Matches source input}
\noindent\textbf{Reviewer annotation}\par
\reviewerbox
\clearpage
\hypertarget{paper-prerequisite-49957d909aa24bc5}{}
\subsection*{\small\ttfamily\raggedright Mendelson\allowbreak{}Whang1990\allowbreak{}Priority\allowbreak{}Pricing.\allowbreak{}priority\allowbreak{}Queueing\allowbreak{}Time}
\noindent\textbf{Paper-local declaration:}\par
{\footnotesize\ttfamily\raggedright Mendelson\allowbreak{}Whang1990\allowbreak{}Priority\allowbreak{}Pricing.\allowbreak{}priority\allowbreak{}Queueing\allowbreak{}Time\par}
\textbf{Source locator:} {\footnotesize\raggedright cited publication, lines 1118-\allowbreak{}1129\par}
\paragraph{Verbatim paper-source connection}
\begin{ReviewVerbatim}
Wi(.) = The expected waiting time of a class-i job                      KLEINROCK, L. 1976. Queueing System, Vols. I and II.
         in the queue in steady state;                                       John Wiley, New York.
 Wi ( )= The expected sojourn time of a class-i job                     KNUDSEN, N. C. 1972. Individual and Social Optimiza-
         in the system in steady state.                                      tion in a Multi-Server Queue With a General Cost-
                                                                             Benefit Structure. Econometrica 40, 515-528.
For the M/M/1 queue with nonpreemptive priority                         KREPS, D. 1990. A Course in Microeconomic Theory.
discipline (Heyman and Sobel 1982, p. 435)                                 Princeton University Press, Princeton, N.J.
                                                                        MARCHAND, M. 1974. Priority Pricing. Mgmt. Sci. 20,
   _ i( L1()
 W1CX)  = W1 Wqc1+ -+
                   Ci AR   (A-)
                      cj. (-                                                  1131-1140.

[Next verbatim source excerpt]

A4. Service Process                                                      Vi'(Xi) = pi + vi Wi (5)
The processor is modeled as an M/M/1 queueing                          where (5) is analogous to (1) with total expected class-
system with nonpreemptive priority among job                           i user cost consisting of two components: a direct
classes. The processing requirement Hi(-) of a class-i                 payment of pi per job, and the expected class-i delay
job, ri, follows an exponential distribution with finite cost, given by vi * Wi(X).
mean ci. Processing requirements are independent           The following theorem generalizes the net value
across jobs and identically distributed within job       maximizing pricing scheme, (3), to the multiclass case.
classes.

[Next verbatim source excerpt]

giving rise to X = A (p) or V'(X) = 4-'(X/A) = p.                     Al. Arrival Processes

  Consider the case where the delay cost v is positive.               The arrival process of class-i jobs to the system, which
Here, the cost faced by an individual user consists of                reflects the aggregation of infinitesimal users' job flow,
two components: a direct payment of p per job, and                    follows a stationary Poisson process with a price-




                                                 This content downloaded from
                               108.30.55.84 on Mon, 24 Aug 2026 10:55:53 UTC
                                         All use subject to https://about.jstor.org/terms

[form-feed in source extraction]
                                                                                Priority Pricingfor the MIMI1 Queue / 873

dependent arrival rate Xi. The R arrival processes are                 by the system (per unit of time) to the organization,
stochastically independent.                                            where we assume for now that the system administra-
\end{ReviewVerbatim}



\paragraph{Lean-expanded paper semantic target}
\begin{ReviewVerbatim}
type:
{n : ℕ} → (Fin n → ℝ) → (Fin n → ℝ) → Fin n → ℝ

value:
fun {n : ℕ} (arrivalRate meanService : Fin n → ℝ) (i : Fin n) =>
  AppliedModelingLib.Queueing.finiteNonpreemptivePriorityQueueWait arrivalRate meanService i
\end{ReviewVerbatim}

\paragraph{Direct retained dependencies}
\begin{itemize}\raggedright
\item \hyperlink{library-prerequisite-795b7af1c661bc8d}{Applied\allowbreak{}Modeling\allowbreak{}Lib.\allowbreak{}Queueing.\allowbreak{}finite\allowbreak{}Nonpreemptive\allowbreak{}Priority\allowbreak{}Queue\allowbreak{}Wait}
\end{itemize}
\paragraph{Recorded source-to-declaration screening}
\textbf{Verdict:} \texttt{matches}
\\
{\raggedright
\textbf{Reason:} It names the Appendix (A-\allowbreak{}1) stationary nonpreemptive-\allowbreak{}priority M/\allowbreak{}M/\allowbreak{}1 queue-\allowbreak{}wait expression with the source priority ordering.\allowbreak{}
\par}
\reviewmatch{Matches source input}
\noindent\textbf{Reviewer annotation}\par
\reviewerbox
\clearpage
\hypertarget{paper-prerequisite-69133735a09302ef}{}
\subsection*{\small\ttfamily\raggedright Mendelson\allowbreak{}Whang1990\allowbreak{}Priority\allowbreak{}Pricing.\allowbreak{}homogeneous\allowbreak{}Priority\allowbreak{}Queueing\allowbreak{}Time}
\noindent\textbf{Paper-local declaration:}\par
{\footnotesize\ttfamily\raggedright Mendelson\allowbreak{}Whang1990\allowbreak{}Priority\allowbreak{}Pricing.\allowbreak{}homogeneous\allowbreak{}Priority\allowbreak{}Queueing\allowbreak{}Time\par}
\textbf{Source locator:} {\footnotesize\raggedright cited publication, lines 387-\allowbreak{}438\par}
\paragraph{Verbatim paper-source connection}
\begin{ReviewVerbatim}
Consider a system with R job classes that satisfy
and assigning the highest priority (= priority 1) to                   assumptions A1-A5 of Section 2, and assume in ad-
class 1, the second-highest priority to class 2, and so                dition that their service requirements are homogene-
on, down to class R. In the special case of equal service              ous across job classes, i.e., c, = C2 = ... = CR = 1,
requirements, c, = C2 = ... = CR, and priorities are                   where the last equality can be assumed without loss
determined by the ranking of delay costs. This policy
                                                                       of generality. In this case, jobs are characterized by
is known to be optimal in the case of fixed arrival                    their delay cost parameters, and the optimal non-
rates with the objective of minimizing the expected
                                                                       preemptive priority scheme orders jobs by decreasing
overall system delay cost. We claim that the same
                                                                       delay costs (i.e., vu > V2 > ... > VR), assigning a higher
scheduling policy is optimal in the current model,                     priority to job classes i with a higher delay cost vi.
where the arrival-rate and priority decisions are jointly              Under this scheme, the expected number of class-i
made to maximize the expected net value. To see this,
                                                                       jobs and the expected waiting time in the system for
note that any priority assignment can be expressed as
                                                                       a class-i job at X are, respectively, given by equation
a permutation r of {1, 2, . . ., R }. Let 9Q denote the                (A- 1) in the Appendix, namely
set of permutations of 1, 2, . . ., R}, and ro be the
permutation of (1, 2, ..., R) which corresponds to                                   XiSR

the c-A rule. The net value maximization problem                                     S/-i=
                                                                                     Ni-l *Si S S+ Xi (8)




                                                  This content downloaded from
                                108.30.55.84 on Mon, 24 Aug 2026 10:55:53 UTC
                                          All use subject to https://about.jstor.org/terms

[form-feed in source extraction]
                                                                                  PriorityPricingfortheMIM/1 Queue / 875

and                                                                   attempt to maximize the objective function of the
                                                                      organization as a whole. Specifically, a user whose job
 Wi(S ) = - _+ Wi (X ) + 1 (9)                                        belongs to class i may find it in his self-interest to
           Si- I .S
                                                                       classify a job into a different priority class, j, either in
                         order
 where SO = 0, Si = EL Xk,  Si to gain
                                  = 1  priority
                                           - Si and reduce
                                                     for   the all
                                                               expected
                                                                     i delay
                                                                        = 1,
2, 3,... , R and Wiq denotes cost
                              the   when  j < i, or to reduce
                                       expected               the service charge
                                                            waiting          timewhen
                             j > i. This would result in a divergence between users'
of a class-i job in the queue. Also, using (9)
\end{ReviewVerbatim}



\paragraph{Lean-expanded paper semantic target}
\begin{ReviewVerbatim}
type:
{n : ℕ} → (Fin n → ℝ) → Fin n → ℝ

value:
fun {n : ℕ} (arrivalRate : Fin n → ℝ) (i : Fin n) =>
  MendelsonWhang1990PriorityPricing.priorityQueueingTime arrivalRate (fun (x : Fin n) => 1) i
\end{ReviewVerbatim}

\paragraph{Direct retained dependencies}
\begin{itemize}\raggedright
\item \hyperlink{paper-prerequisite-49957d909aa24bc5}{Mendelson\allowbreak{}Whang1990\allowbreak{}Priority\allowbreak{}Pricing.\allowbreak{}priority\allowbreak{}Queueing\allowbreak{}Time}
\end{itemize}
\paragraph{Recorded source-to-declaration screening}
\textbf{Verdict:} \texttt{matches}
\\
{\raggedright
\textbf{Reason:} It is the Appendix (A-\allowbreak{}1) priority queue-\allowbreak{}wait formula specialized to homogeneous unit mean service.\allowbreak{}
\par}
\reviewmatch{Matches source input}
\noindent\textbf{Reviewer annotation}\par
\reviewerbox
\clearpage
\hypertarget{paper-prerequisite-3db4e2685735dcc0}{}
\subsection*{\small\ttfamily\raggedright Mendelson\allowbreak{}Whang1990\allowbreak{}Priority\allowbreak{}Pricing.\allowbreak{}homogeneous\allowbreak{}Priority\allowbreak{}Price}
\noindent\textbf{Paper-local declaration:}\par
{\footnotesize\ttfamily\raggedright Mendelson\allowbreak{}Whang1990\allowbreak{}Priority\allowbreak{}Pricing.\allowbreak{}homogeneous\allowbreak{}Priority\allowbreak{}Price\par}
\textbf{Source locator:} {\footnotesize\raggedright cited publication, lines 448-\allowbreak{}460\par}
\paragraph{Verbatim paper-source connection}
\begin{ReviewVerbatim}
Using Theorem 1, the optimal prices are given by
                                                                         In this section, we consider the family of priority-
       R Xk*Vk                                                        dependent (PD) pricing schemes, where each job is
      k=l I k * k-I                                                   charged only on the basis of the priority class selected.
                                                                      A PD pricing scheme is specified by a price vector
          RXk Vk Wk + Xk+1 Vk+l Wkq+ (                                P = (PI, P2, .. ., PR); for a given p, the arrival-rate
                      + E - - O ~~~~~~(0)                             vector and priority assignment (?X (p), r(p)) constitute
          k=i         Ok

                                                                      an (symmetric Nash) equilibrium if each class-i user's
where we define XA+I = VR+I = WR = 0.
                                                                       action maximizes the expected net value, taking all
  Equation (10) clearly demonstrates the delay cost-
\end{ReviewVerbatim}



\paragraph{Lean-expanded paper semantic target}
\begin{ReviewVerbatim}
type:
{n : ℕ} → (Fin n → ℝ) → (Fin n → ℝ) → Fin n → ℝ

value:
fun {n : ℕ} (arrivalRate delayCost : Fin n → ℝ) (i : Fin n) =>
  ∑ k : Fin n,
      arrivalRate k * delayCost k /
        (AppliedModelingLib.Queueing.finitePriorityInclusiveSlack (fun (x : Fin n) => 1) arrivalRate k *
          AppliedModelingLib.Queueing.finitePriorityStrictSlack (fun (x : Fin n) => 1) arrivalRate k) +
    ∑ k : Fin n,
      if i ≤ k then
        (arrivalRate k * delayCost k * MendelsonWhang1990PriorityPricing.homogeneousPriorityQueueingTime arrivalRate k +
            MendelsonWhang1990PriorityPricing.prioritySuccessorOrZero
              (fun (j : Fin n) =>
                arrivalRate j * delayCost j *
                  MendelsonWhang1990PriorityPricing.homogeneousPriorityQueueingTime arrivalRate j)
              k) /
          AppliedModelingLib.Queueing.finitePriorityInclusiveSlack (fun (x : Fin n) => 1) arrivalRate k
      else 0
\end{ReviewVerbatim}

\paragraph{Direct retained dependencies}
\begin{itemize}\raggedright
\item \hyperlink{governing-declaration-f2b9a1af307e7469}{Applied\allowbreak{}Modeling\allowbreak{}Lib.\allowbreak{}Queueing.\allowbreak{}finite\allowbreak{}Priority\allowbreak{}Inclusive\allowbreak{}Slack}
\item \hyperlink{governing-declaration-3f5cf610a8190094}{Applied\allowbreak{}Modeling\allowbreak{}Lib.\allowbreak{}Queueing.\allowbreak{}finite\allowbreak{}Priority\allowbreak{}Strict\allowbreak{}Slack}
\item \hyperlink{paper-prerequisite-69133735a09302ef}{Mendelson\allowbreak{}Whang1990\allowbreak{}Priority\allowbreak{}Pricing.\allowbreak{}homogeneous\allowbreak{}Priority\allowbreak{}Queueing\allowbreak{}Time}
\item \hyperlink{governing-declaration-f2dbc4d52705c26d}{Mendelson\allowbreak{}Whang1990\allowbreak{}Priority\allowbreak{}Pricing.\allowbreak{}priority\allowbreak{}Successor\allowbreak{}Or\allowbreak{}Zero}
\end{itemize}
\paragraph{Recorded source-to-declaration screening}
\textbf{Verdict:} \texttt{matches}
\\
{\raggedright
\textbf{Reason:} Its expanded two-\allowbreak{}sum expression implements Equation (10), including the source's zero-\allowbreak{}padded successor convention.\allowbreak{}
\par}
\reviewmatch{Matches source input}
\noindent\textbf{Reviewer annotation}\par
\reviewerbox
\clearpage
\hypertarget{paper-prerequisite-9c41d2efaf9290fc}{}
\subsection*{\small\ttfamily\raggedright Mendelson\allowbreak{}Whang1990\allowbreak{}Priority\allowbreak{}Pricing.\allowbreak{}homogeneous\allowbreak{}Priority\allowbreak{}Sojourn\allowbreak{}Time}
\noindent\textbf{Paper-local declaration:}\par
{\footnotesize\ttfamily\raggedright Mendelson\allowbreak{}Whang1990\allowbreak{}Priority\allowbreak{}Pricing.\allowbreak{}homogeneous\allowbreak{}Priority\allowbreak{}Sojourn\allowbreak{}Time\par}
\textbf{Source locator:} {\footnotesize\raggedright cited publication, lines 387-\allowbreak{}438\par}
\paragraph{Verbatim paper-source connection}
\begin{ReviewVerbatim}
Consider a system with R job classes that satisfy
and assigning the highest priority (= priority 1) to                   assumptions A1-A5 of Section 2, and assume in ad-
class 1, the second-highest priority to class 2, and so                dition that their service requirements are homogene-
on, down to class R. In the special case of equal service              ous across job classes, i.e., c, = C2 = ... = CR = 1,
requirements, c, = C2 = ... = CR, and priorities are                   where the last equality can be assumed without loss
determined by the ranking of delay costs. This policy
                                                                       of generality. In this case, jobs are characterized by
is known to be optimal in the case of fixed arrival                    their delay cost parameters, and the optimal non-
rates with the objective of minimizing the expected
                                                                       preemptive priority scheme orders jobs by decreasing
overall system delay cost. We claim that the same
                                                                       delay costs (i.e., vu > V2 > ... > VR), assigning a higher
scheduling policy is optimal in the current model,                     priority to job classes i with a higher delay cost vi.
where the arrival-rate and priority decisions are jointly              Under this scheme, the expected number of class-i
made to maximize the expected net value. To see this,
                                                                       jobs and the expected waiting time in the system for
note that any priority assignment can be expressed as
                                                                       a class-i job at X are, respectively, given by equation
a permutation r of {1, 2, . . ., R }. Let 9Q denote the                (A- 1) in the Appendix, namely
set of permutations of 1, 2, . . ., R}, and ro be the
permutation of (1, 2, ..., R) which corresponds to                                   XiSR

the c-A rule. The net value maximization problem                                     S/-i=
                                                                                     Ni-l *Si S S+ Xi (8)




                                                  This content downloaded from
                                108.30.55.84 on Mon, 24 Aug 2026 10:55:53 UTC
                                          All use subject to https://about.jstor.org/terms

[form-feed in source extraction]
                                                                                  PriorityPricingfortheMIM/1 Queue / 875

and                                                                   attempt to maximize the objective function of the
                                                                      organization as a whole. Specifically, a user whose job
 Wi(S ) = - _+ Wi (X ) + 1 (9)                                        belongs to class i may find it in his self-interest to
           Si- I .S
                                                                       classify a job into a different priority class, j, either in
                         order
 where SO = 0, Si = EL Xk,  Si to gain
                                  = 1  priority
                                           - Si and reduce
                                                     for   the all
                                                               expected
                                                                     i delay
                                                                        = 1,
2, 3,... , R and Wiq denotes cost
                              the   when  j < i, or to reduce
                                       expected               the service charge
                                                            waiting          timewhen
                             j > i. This would result in a divergence between users'
of a class-i job in the queue. Also, using (9)
\end{ReviewVerbatim}



\paragraph{Lean-expanded paper semantic target}
\begin{ReviewVerbatim}
type:
{n : ℕ} → (Fin n → ℝ) → Fin n → ℝ

value:
fun {n : ℕ} (arrivalRate : Fin n → ℝ) (i : Fin n) =>
  MendelsonWhang1990PriorityPricing.prioritySojournTime arrivalRate (fun (x : Fin n) => 1) i
\end{ReviewVerbatim}

\paragraph{Direct retained dependencies}
\begin{itemize}\raggedright
\item \hyperlink{governing-declaration-703b87082556903f}{Mendelson\allowbreak{}Whang1990\allowbreak{}Priority\allowbreak{}Pricing.\allowbreak{}priority\allowbreak{}Sojourn\allowbreak{}Time}
\end{itemize}
\paragraph{Recorded source-to-declaration screening}
\textbf{Verdict:} \texttt{matches}
\\
{\raggedright
\textbf{Reason:} Equation (9) is the homogeneous queue wait plus the unit service requirement.\allowbreak{}
\par}
\reviewmatch{Matches source input}
\noindent\textbf{Reviewer annotation}\par
\reviewerbox
\clearpage
\hypertarget{paper-prerequisite-53f58e4a79f3cd3c}{}
\subsection*{\small\ttfamily\raggedright Mendelson\allowbreak{}Whang1990\allowbreak{}Priority\allowbreak{}Pricing.\allowbreak{}homogeneous\allowbreak{}Priority\allowbreak{}Expected\allowbreak{}Cost}
\noindent\textbf{Paper-local declaration:}\par
{\footnotesize\ttfamily\raggedright Mendelson\allowbreak{}Whang1990\allowbreak{}Priority\allowbreak{}Pricing.\allowbreak{}homogeneous\allowbreak{}Priority\allowbreak{}Expected\allowbreak{}Cost\par}
\textbf{Source locator:} {\footnotesize\raggedright cited publication, lines 448-\allowbreak{}460\par}
\paragraph{Verbatim paper-source connection}
\begin{ReviewVerbatim}
Using Theorem 1, the optimal prices are given by
                                                                         In this section, we consider the family of priority-
       R Xk*Vk                                                        dependent (PD) pricing schemes, where each job is
      k=l I k * k-I                                                   charged only on the basis of the priority class selected.
                                                                      A PD pricing scheme is specified by a price vector
          RXk Vk Wk + Xk+1 Vk+l Wkq+ (                                P = (PI, P2, .. ., PR); for a given p, the arrival-rate
                      + E - - O ~~~~~~(0)                             vector and priority assignment (?X (p), r(p)) constitute
          k=i         Ok

                                                                      an (symmetric Nash) equilibrium if each class-i user's
where we define XA+I = VR+I = WR = 0.
                                                                       action maximizes the expected net value, taking all
  Equation (10) clearly demonstrates the delay cost-
\end{ReviewVerbatim}



\paragraph{Lean-expanded paper semantic target}
\begin{ReviewVerbatim}
type:
{n : ℕ} → (Fin n → ℝ) → (Fin n → ℝ) → Fin n → Fin n → ℝ

value:
fun {n : ℕ} (arrivalRate delayCost : Fin n → ℝ) (i selectedPriority : Fin n) =>
  MendelsonWhang1990PriorityPricing.homogeneousPriorityPrice arrivalRate delayCost selectedPriority +
    delayCost i * MendelsonWhang1990PriorityPricing.homogeneousPrioritySojournTime arrivalRate selectedPriority
\end{ReviewVerbatim}

\paragraph{Direct retained dependencies}
\begin{itemize}\raggedright
\item \hyperlink{paper-prerequisite-3db4e2685735dcc0}{Mendelson\allowbreak{}Whang1990\allowbreak{}Priority\allowbreak{}Pricing.\allowbreak{}homogeneous\allowbreak{}Priority\allowbreak{}Price}
\item \hyperlink{paper-prerequisite-9c41d2efaf9290fc}{Mendelson\allowbreak{}Whang1990\allowbreak{}Priority\allowbreak{}Pricing.\allowbreak{}homogeneous\allowbreak{}Priority\allowbreak{}Sojourn\allowbreak{}Time}
\end{itemize}
\paragraph{Recorded source-to-declaration screening}
\textbf{Verdict:} \texttt{matches}
\\
{\raggedright
\textbf{Reason:} The Section 2 private priority cost is Equation (10)'s declared-\allowbreak{}priority price plus true-\allowbreak{}class delay cost times the declared-\allowbreak{}priority sojourn time.\allowbreak{}
\par}
\reviewmatch{Matches source input}
\noindent\textbf{Reviewer annotation}\par
\reviewerbox
\clearpage
\hypertarget{paper-prerequisite-0efbdff1257f8766}{}
\subsection*{\small\ttfamily\raggedright Mendelson\allowbreak{}Whang1990\allowbreak{}Priority\allowbreak{}Pricing.\allowbreak{}homogeneous\allowbreak{}Priority\allowbreak{}Population}
\noindent\textbf{Paper-local declaration:}\par
{\footnotesize\ttfamily\raggedright Mendelson\allowbreak{}Whang1990\allowbreak{}Priority\allowbreak{}Pricing.\allowbreak{}homogeneous\allowbreak{}Priority\allowbreak{}Population\par}
\textbf{Source locator:} {\footnotesize\raggedright cited publication, lines 387-\allowbreak{}438\par}
\paragraph{Verbatim paper-source connection}
\begin{ReviewVerbatim}
Consider a system with R job classes that satisfy
and assigning the highest priority (= priority 1) to                   assumptions A1-A5 of Section 2, and assume in ad-
class 1, the second-highest priority to class 2, and so                dition that their service requirements are homogene-
on, down to class R. In the special case of equal service              ous across job classes, i.e., c, = C2 = ... = CR = 1,
requirements, c, = C2 = ... = CR, and priorities are                   where the last equality can be assumed without loss
determined by the ranking of delay costs. This policy
                                                                       of generality. In this case, jobs are characterized by
is known to be optimal in the case of fixed arrival                    their delay cost parameters, and the optimal non-
rates with the objective of minimizing the expected
                                                                       preemptive priority scheme orders jobs by decreasing
overall system delay cost. We claim that the same
                                                                       delay costs (i.e., vu > V2 > ... > VR), assigning a higher
scheduling policy is optimal in the current model,                     priority to job classes i with a higher delay cost vi.
where the arrival-rate and priority decisions are jointly              Under this scheme, the expected number of class-i
made to maximize the expected net value. To see this,
                                                                       jobs and the expected waiting time in the system for
note that any priority assignment can be expressed as
                                                                       a class-i job at X are, respectively, given by equation
a permutation r of {1, 2, . . ., R }. Let 9Q denote the                (A- 1) in the Appendix, namely
set of permutations of 1, 2, . . ., R}, and ro be the
permutation of (1, 2, ..., R) which corresponds to                                   XiSR

the c-A rule. The net value maximization problem                                     S/-i=
                                                                                     Ni-l *Si S S+ Xi (8)




                                                  This content downloaded from
                                108.30.55.84 on Mon, 24 Aug 2026 10:55:53 UTC
                                          All use subject to https://about.jstor.org/terms

[form-feed in source extraction]
                                                                                  PriorityPricingfortheMIM/1 Queue / 875

and                                                                   attempt to maximize the objective function of the
                                                                      organization as a whole. Specifically, a user whose job
 Wi(S ) = - _+ Wi (X ) + 1 (9)                                        belongs to class i may find it in his self-interest to
           Si- I .S
                                                                       classify a job into a different priority class, j, either in
                         order
 where SO = 0, Si = EL Xk,  Si to gain
                                  = 1  priority
                                           - Si and reduce
                                                     for   the all
                                                               expected
                                                                     i delay
                                                                        = 1,
2, 3,... , R and Wiq denotes cost
                              the   when  j < i, or to reduce
                                       expected               the service charge
                                                            waiting          timewhen
                             j > i. This would result in a divergence between users'
of a class-i job in the queue. Also, using (9)
\end{ReviewVerbatim}



\paragraph{Lean-expanded paper semantic target}
\begin{ReviewVerbatim}
type:
{n : ℕ} → (Fin n → ℝ) → Fin n → ℝ

value:
fun {n : ℕ} (arrivalRate : Fin n → ℝ) (i : Fin n) =>
  arrivalRate i * MendelsonWhang1990PriorityPricing.homogeneousPrioritySojournTime arrivalRate i
\end{ReviewVerbatim}

\paragraph{Direct retained dependencies}
\begin{itemize}\raggedright
\item \hyperlink{paper-prerequisite-9c41d2efaf9290fc}{Mendelson\allowbreak{}Whang1990\allowbreak{}Priority\allowbreak{}Pricing.\allowbreak{}homogeneous\allowbreak{}Priority\allowbreak{}Sojourn\allowbreak{}Time}
\end{itemize}
\paragraph{Recorded source-to-declaration screening}
\textbf{Verdict:} \texttt{matches}
\\
{\raggedright
\textbf{Reason:} Equation (8) is the homogeneous Little's-\allowbreak{}law population:\allowbreak{} class arrival rate multiplied by its sojourn time.\allowbreak{}
\par}
\reviewmatch{Matches source input}
\noindent\textbf{Reviewer annotation}\par
\reviewerbox
\clearpage
\hypertarget{paper-prerequisite-6c93de2469b53cb2}{}
\subsection*{\small\ttfamily\raggedright Mendelson\allowbreak{}Whang1990\allowbreak{}Priority\allowbreak{}Pricing.\allowbreak{}individual\allowbreak{}Demand\allowbreak{}Condition}
\noindent\textbf{Paper-local declaration:}\par
{\footnotesize\ttfamily\raggedright Mendelson\allowbreak{}Whang1990\allowbreak{}Priority\allowbreak{}Pricing.\allowbreak{}individual\allowbreak{}Demand\allowbreak{}Condition\par}
\textbf{Source locator:} {\footnotesize\raggedright cited publication, lines 258-\allowbreak{}268\par}
\paragraph{Verbatim paper-source connection}
\begin{ReviewVerbatim}
relationship (1). Let the system administrator set the
A3. Delay Cost Structure
                                                                       R-dimensional price vector p = (PI, P2, P3, .., PR),
Each class-i job is characterized by a delay cost of vi                where pi is the price charged to a class-i job. Then, the
per unit time.                                                         class-i demand relationship is given by

A4. Service Process                                                      Vi'(Xi) = pi + vi Wi (5)
The processor is modeled as an M/M/1 queueing                          where (5) is analogous to (1) with total expected class-
system with nonpreemptive priority among job                           i user cost consisting of two components: a direct
classes. The processing requirement Hi(-) of a class-i                 payment of pi per job, and the expected class-i delay
job, ri, follows an exponential distribution with finite cost, given by vi * Wi(X).

[Next verbatim source excerpt]

A5. Decision Structure
                                                                      price per class-i job is given by
We assume the individual decision structure: users do
not collude, and each maximizes his expected net
                                                                                 R aow.
                                                                        Pi =, vjXi, for i = 1, 2, ..., R (6)
benefit.
\end{ReviewVerbatim}



\paragraph{Lean-expanded paper semantic target}
\begin{ReviewVerbatim}
type:
{Class : Type u_1} →
  [Fintype Class] →
    (Class → ℝ → ℝ) → (Class → ℝ) → (Class → ℝ) → (Class → ℝ) → (Class → (Class → ℝ) → ℝ) → Class → ℝ → Prop

value:
fun {Class : Type u_1} [Fintype Class] (value : Class → ℝ → ℝ) (delayCost flow price : Class → ℝ)
    (waitingTime : Class → (Class → ℝ) → ℝ) (i : Class) (marginalValue : ℝ) =>
  HasDerivAt (value i) marginalValue (flow i) ∧ marginalValue = price i + delayCost i * waitingTime i flow
\end{ReviewVerbatim}


\paragraph{Recorded source-to-declaration screening}
\textbf{Verdict:} \texttt{matches}
\\
{\raggedright
\textbf{Reason:} The Lean predicate records both parts of the source condition:\allowbreak{} the displayed number is the derivative of class i's value at its flow, and Equation (5) equates it to direct price plus per-\allowbreak{}time delay cost times expected delay.\allowbreak{} This is the atomistic A5 decision equation, not a collusive first-\allowbreak{}order condition.\allowbreak{}
\par}
\reviewmatch{Matches source input}
\noindent\textbf{Reviewer annotation}\par
\reviewerbox
\clearpage
\hypertarget{paper-prerequisite-7dbaca4145f3a4a5}{}
\subsection*{\small\ttfamily\raggedright Mendelson\allowbreak{}Whang1990\allowbreak{}Priority\allowbreak{}Pricing.\allowbreak{}net\allowbreak{}Value}
\noindent\textbf{Paper-local declaration:}\par
{\footnotesize\ttfamily\raggedright Mendelson\allowbreak{}Whang1990\allowbreak{}Priority\allowbreak{}Pricing.\allowbreak{}net\allowbreak{}Value\par}
\textbf{Source locator:} {\footnotesize\raggedright cited publication, lines 195-\allowbreak{}200\par}
\paragraph{Verbatim paper-source connection}
\begin{ReviewVerbatim}
Two alternative decision structures may be envi-                   p, represented by the demand relationship (1), we turn
sioned in this setting. Either users set the aggregate                to the problem of maximizing the expected net value
arrival rate collusively, acting in effect as a single                of the system as a whole. The net value of the system
decision maker; or users make individual decisions on                 at arrival rate X is given by V(X) - v * X * W(X),
whether or not to join the system, and the overall                    where the first term represents its expected gross value,
arrival rate is determined by aggregating the results ofand the second represents the expected delay cost.

[Next verbatim source excerpt]

                                                                       belongs. Under a given prionity scheme, this problem
The value function of the class-i job stream can be                    can be formulated as
expressed as Vi (Xi), and the total value is given by the                        F       R
sum of Vi(Xi) over all classes, that is, the expected                   max V(X) - E viLi(X)J (4)
(gross) value function is given by                                          X        i=l


           R                                                           where Li is the mean number of class-i jobs in the
V(X)= E V1(Xi)                                                         system in steady state, viLi is the expected delay cost
           i=l1
                                                                       incurred by class-i jobs per unit of time, and V(2X) =
where Vi(Xi) are differentiable, nondecreasing and                     E 1 Vi(Xi) is the total gross value per unit of time.
\end{ReviewVerbatim}



\paragraph{Lean-expanded paper semantic target}
\begin{ReviewVerbatim}
type:
{Class : Type u_1} → [Fintype Class] → (Class → ℝ → ℝ) → (Class → ℝ) → (Class → (Class → ℝ) → ℝ) → (Class → ℝ) → ℝ

value:
fun {Class : Type u_1} [Fintype Class] (value : Class → ℝ → ℝ) (delayCost : Class → ℝ)
    (waitingTime : Class → (Class → ℝ) → ℝ) (flow : Class → ℝ) =>
  AppliedModelingLib.finiteDelayNetValue value delayCost waitingTime flow
\end{ReviewVerbatim}

\paragraph{Direct retained dependencies}
\begin{itemize}\raggedright
\item \hyperlink{governing-declaration-971183e1bc53251f}{Applied\allowbreak{}Modeling\allowbreak{}Lib.\allowbreak{}finite\allowbreak{}Delay\allowbreak{}Net\allowbreak{}Value}
\end{itemize}
\paragraph{Recorded source-to-declaration screening}
\textbf{Verdict:} \texttt{matches}
\\
{\raggedright
\textbf{Reason:} The expanded finite sum is gross value V\_\allowbreak{}i(lambda\_\allowbreak{}i) less v\_\allowbreak{}i lambda\_\allowbreak{}i W\_\allowbreak{}i(lambda) for every class.\allowbreak{} Since lambda\_\allowbreak{}i W\_\allowbreak{}i is the steady-\allowbreak{}state class population by Little's law, this is exactly Equation (4)'s total value minus total delay cost.\allowbreak{}
\par}
\reviewmatch{Matches source input}
\noindent\textbf{Reviewer annotation}\par
\reviewerbox
\clearpage
\hypertarget{paper-prerequisite-54c5682d350b5d3d}{}
\subsection*{\small\ttfamily\raggedright Mendelson\allowbreak{}Whang1990\allowbreak{}Priority\allowbreak{}Pricing.\allowbreak{}priority\allowbreak{}Cheating\allowbreak{}Penalty}
\noindent\textbf{Paper-local declaration:}\par
{\footnotesize\ttfamily\raggedright Mendelson\allowbreak{}Whang1990\allowbreak{}Priority\allowbreak{}Pricing.\allowbreak{}priority\allowbreak{}Cheating\allowbreak{}Penalty\par}
\textbf{Source locator:} {\footnotesize\raggedright cited publication, lines 921-\allowbreak{}937\par}
\paragraph{Verbatim paper-source connection}
\begin{ReviewVerbatim}
arrive during the service time of the tagged job totals                 where ll'(k) = vi(Wk - W') + ci(Ak- Ai). The term
v1 Xt2/2. The impact of this delayed startup of the                     Ili(k) can be regarded as the cheating penalty incurred
class-I busy period propagates to all class-I jobs that                 when a class-i user chooses priority-class k. Theorem
arrive by the end of the busy period, and the total                     3 demonstrates that II'(k), as a discrete function of k,
externality cost inflicted on class- I jobs is given by                 has a unique minimum value of zero at k = i. Hence,
                                                                        any class-i user choosing a different priority class will
    = 2 (Xlt + CIX2t + C2X3t + . v.X-t2                                 incur a strictly positive penalty. Another desirable
              2         2                                               property of the cheating penalty fl'(k), given by the
                                                                        following theorem, is that for all user classes i,
The externality Ek on class (k ? 1) by the tagged job
                                                                        the cheating penalty II'(k) increases as the selected
can be computed in a similar way, but each service
                                                                        priority-class, k, gets farther away from the correct
period must now be extended by a factor of /lk-l
                                               user class, i, in either direction. That is, the cheating
because the service process is subject to interruptions
                                                                        penalty increases with the degree of the cheating.
\end{ReviewVerbatim}



\paragraph{Lean-expanded paper semantic target}
\begin{ReviewVerbatim}
type:
{n : ℕ} → (Fin n → ℝ) → (Fin n → ℝ) → (Fin n → ℝ) → (Fin n → ℝ) → Fin n → Fin n → ℝ

value:
fun {n : ℕ} (meanService delayCost queueingTime linearCoefficient : Fin n → ℝ) (customer declaredPriority : Fin n) =>
  meanService customer * (linearCoefficient declaredPriority - linearCoefficient customer) +
    delayCost customer * (queueingTime declaredPriority - queueingTime customer)
\end{ReviewVerbatim}


\paragraph{Recorded source-to-declaration screening}
\textbf{Verdict:} \texttt{matches}
\\
{\raggedright
\textbf{Reason:} The Lean value is c\_\allowbreak{}i(A\_\allowbreak{}k-\allowbreak{}A\_\allowbreak{}i)+v\_\allowbreak{}i(W\_\allowbreak{}k-\allowbreak{}W\_\allowbreak{}i), exactly the source's displayed cheating penalty for a true class i declaring priority k.\allowbreak{} The common expected quadratic component cancels, leaving the linear-\allowbreak{}price and queue-\allowbreak{}wait differences shown in the source.\allowbreak{}
\par}
\reviewmatch{Matches source input}
\noindent\textbf{Reviewer annotation}\par
\reviewerbox
\clearpage
\hypertarget{paper-prerequisite-476808203e3008c1}{}
\subsection*{\small\ttfamily\raggedright Mendelson\allowbreak{}Whang1990\allowbreak{}Priority\allowbreak{}Pricing.\allowbreak{}priority\allowbreak{}Time\allowbreak{}Dependent\allowbreak{}Expected\allowbreak{}Cost}
\noindent\textbf{Paper-local declaration:}\par
{\footnotesize\ttfamily\raggedright Mendelson\allowbreak{}Whang1990\allowbreak{}Priority\allowbreak{}Pricing.\allowbreak{}priority\allowbreak{}Time\allowbreak{}Dependent\allowbreak{}Expected\allowbreak{}Cost\par}
\textbf{Source locator:} {\footnotesize\raggedright cited publication, lines 326-\allowbreak{}362\par}
\paragraph{Verbatim paper-source connection}
\begin{ReviewVerbatim}
Hence                                                                  becomes

                                                                                   r     R   R
            a
                                                                        maxf j Vi(Xi) - viLi(, r)t.
    P~E
     d oi 3ax,
                                                                       If X * and r* are the optimal arrival-rate vector and
      + vi *[Wi(X*)+x* Wi(X*)]-viWi(X*)                                scheduling permutation, respectively, then by the opti-
                                                                       mality of the c-, rule for fixed X = X * we have
           R a~~~i x
      R                                                                             R            R

   - JX,*- W (X*)                                                       max{ Vi(Xi) - viLi(X,(r)
      j=1 J axj

since for] i d                                                                   R           R


                                                                          = E Vi(Xi) -E viLi(2*, r*)
                                                                               i=l       i=l


1a-
 Ax.Li (2X) a - xa
          axi    iW xi
                    = -i a 'J/(X).                                               R           R


                                                                           X Vi(X*) - E viLi(X, rO)
  Note that when R = 1, (6) reduces to (3). Also, as
in (3), the right-hand side of (6) represents the expected                                   R   R

cost inflicted on all other users as a consequence of                      < max Vi(xi) - viLi(X, rO)}
an infinitesimal unit increase in the job flow of a class-
i user. Thus, similar to (3), we may call the right-hand               which demonstrates that the same scheduling rate is
side of (6) the externality cost. The result of Theorem                also optimal in the current setting.

[Next verbatim source excerpt]

                                                              Consider a system with R job classes that satisfy
class-i jobs. As before, all aggregate class characteristics
                                                            assumptions A1-A5 of Section 1, where class-i jobs
are common knowledge, whereas individual job char-
                                                            have exponential processing requirements with mean
acteristics are known only to the user. Thus, both the
                                                           ci. We assume that

[Next verbatim source excerpt]

  Under this general framework, an equilibrium can                     geneous service requirements. The answer is negative,
be defined in a similar way to (12) and (13). Let WJi                  as the following example demonstrates. Consider the
denote the expected total waiting time in the system                   case of two job classes (R = 2) and let p(t) = (pI (t),
for a class-i job picking priority-class j. Then, ( X (p),            P2(t)) be given by (16). Ifp(t) is optimal and incentive-
f(p)) is an equilibrium with respect to the PTD pricing compatible, optimality implies that for a class-i job,
scheme p(t) if for each i = 1, 2, . . ., R              we just have E[pi(ri)] = p*, where p* are given by
                                                                       Theorem 1. Also from (16)
E[pr,(p)(ri)] + viWri,(p)(2(p), r(p))

   = min E[pj(ri)] + viWj(X(p), r(p))}                                 E[p2(T )] =- E[pl(Tri)] =- P
                                                                                             a   a


and                                                                          and                          (17)

Vi'(Xi(p)) = E[prj(p)(Ti)] + Vi Wr,(p)(X (p), r(p))
                                                                       E[pl(T2)] = a * E[p2(T2)] = a *2
where Ti is the actual processing requirement of a
                                                                      where a = K1/K2.
class-i job.
                                                                         In order for the pricing schedule p(t) to be incentive-
  A PTD pricing scheme p0(t) will be called optimal
                                                                       compatible, it should satisfy
and incentive-compatible if it satisfies the following
conditions:
                                                                       E[pI(TI)] + v* W < S E[p2(rl)] + V1 WY
   1. Incentive-compatibility: For all i = 1, 2, . .. , R
                                                                       and
   ri(p?)                   =        i      (14)
                                                                       E[pi(T2)] + v2 * Wl q E[p2(T2)] + V2 WY-
      2.       Optimality:                                       The                 resulting                         arrival
                                                                       Hence using (17), incentive-compatibility implies
  the expected                                          net                value                 of        the            system
is, for all i= 1, 2, .. ., R
                                                                       pi +V* 1 S-*p* + VI * W
  X(pO)                =     _X*                (15)                                         a



where X * achieves the overall optimum.                                      and                          (18)

[Next verbatim source excerpt]

                                                              Consider a system with R job classes that satisfy
class-i jobs. As before, all aggregate class characteristics
                                                            assumptions A1-A5 of Section 1, where class-i jobs
are common knowledge, whereas individual job char-
                                                            have exponential processing requirements with mean
acteristics are known only to the user. Thus, both the
                                                           ci. We assume that
\end{ReviewVerbatim}



\paragraph{Lean-expanded paper semantic target}
\begin{ReviewVerbatim}
type:
{n : ℕ} → (Fin n → ℝ) → (Fin n → ℝ) → (Fin n → ℝ) → (Fin n → ℝ → ℝ) → Fin n → Fin n → ℝ

value:
fun {n : ℕ} (arrivalRate delayCost meanService : Fin n → ℝ) (pricing : Fin n → ℝ → ℝ)
    (customer declaredPriority : Fin n) =>
  ∫ (processingTime : ℝ),
      pricing declaredPriority processingTime ∂ProbabilityTheory.expMeasure (meanService customer)⁻¹ +
    delayCost customer *
      (MendelsonWhang1990PriorityPricing.priorityQueueingTime arrivalRate meanService declaredPriority +
        meanService customer)
\end{ReviewVerbatim}

\paragraph{Direct retained dependencies}
\begin{itemize}\raggedright
\item \hyperlink{paper-prerequisite-49957d909aa24bc5}{Mendelson\allowbreak{}Whang1990\allowbreak{}Priority\allowbreak{}Pricing.\allowbreak{}priority\allowbreak{}Queueing\allowbreak{}Time}
\end{itemize}
\paragraph{Recorded source-to-declaration screening}
\textbf{Verdict:} \texttt{matches}
\\
{\raggedright
\textbf{Reason:} The selected source bundle defines the PTD equilibrium by comparing E[p\_\allowbreak{}j(T\_\allowbreak{}i)] + v\_\allowbreak{}i W\_\allowbreak{}j\textasciicircum{}i(X(p), r(p)) across priority classes j, identifies T\_\allowbreak{}i as class i's actual processing requirement, calls W the expected total time in system, and states that class-\allowbreak{}i processing requirements are exponential with mean c\_\allowbreak{}i.\allowbreak{} The displayed Lean target integrates the declared-\allowbreak{}priority charge against expMeasure (meanService customer)⁻¹ and adds delayCost customer times the declared-\allowbreak{}priority queueing time plus that customer's mean service time.\allowbreak{} Its displayed queueing and sojourn prerequisites make the latter the same total-\allowbreak{}time component, so the formula preserves the source's direct-\allowbreak{}charge and delay-\allowbreak{}cost terms on the displayed positive-\allowbreak{}mean stationary uses.\allowbreak{}
\par}
\reviewmatch{Matches source input}
\noindent\textbf{Reviewer annotation}\par
\reviewerbox
\clearpage
\hypertarget{paper-prerequisite-84dca75fd93aa845}{}
\subsection*{\small\ttfamily\raggedright Mendelson\allowbreak{}Whang1990\allowbreak{}Priority\allowbreak{}Pricing.\allowbreak{}stable\allowbreak{}Priority\allowbreak{}Flow\allowbreak{}Domain}
\noindent\textbf{Paper-local declaration:}\par
{\footnotesize\ttfamily\raggedright Mendelson\allowbreak{}Whang1990\allowbreak{}Priority\allowbreak{}Pricing.\allowbreak{}stable\allowbreak{}Priority\allowbreak{}Flow\allowbreak{}Domain\par}
\textbf{Source locator:} {\footnotesize\raggedright cited publication, lines 1118-\allowbreak{}1129\par}
\paragraph{Verbatim paper-source connection}
\begin{ReviewVerbatim}
Wi(.) = The expected waiting time of a class-i job                      KLEINROCK, L. 1976. Queueing System, Vols. I and II.
         in the queue in steady state;                                       John Wiley, New York.
 Wi ( )= The expected sojourn time of a class-i job                     KNUDSEN, N. C. 1972. Individual and Social Optimiza-
         in the system in steady state.                                      tion in a Multi-Server Queue With a General Cost-
                                                                             Benefit Structure. Econometrica 40, 515-528.
For the M/M/1 queue with nonpreemptive priority                         KREPS, D. 1990. A Course in Microeconomic Theory.
discipline (Heyman and Sobel 1982, p. 435)                                 Princeton University Press, Princeton, N.J.
                                                                        MARCHAND, M. 1974. Priority Pricing. Mgmt. Sci. 20,
   _ i( L1()
 W1CX)  = W1 Wqc1+ -+
                   Ci AR   (A-)
                      cj. (-                                                  1131-1140.

[Next verbatim source excerpt]

A4. Service Process                                                      Vi'(Xi) = pi + vi Wi (5)
The processor is modeled as an M/M/1 queueing                          where (5) is analogous to (1) with total expected class-
system with nonpreemptive priority among job                           i user cost consisting of two components: a direct
classes. The processing requirement Hi(-) of a class-i                 payment of pi per job, and the expected class-i delay
job, ri, follows an exponential distribution with finite cost, given by vi * Wi(X).
mean ci. Processing requirements are independent           The following theorem generalizes the net value
across jobs and identically distributed within job       maximizing pricing scheme, (3), to the multiclass case.
classes.

[Next verbatim source excerpt]

giving rise to X = A (p) or V'(X) = 4-'(X/A) = p.                     Al. Arrival Processes

  Consider the case where the delay cost v is positive.               The arrival process of class-i jobs to the system, which
Here, the cost faced by an individual user consists of                reflects the aggregation of infinitesimal users' job flow,
two components: a direct payment of p per job, and                    follows a stationary Poisson process with a price-




                                                 This content downloaded from
                               108.30.55.84 on Mon, 24 Aug 2026 10:55:53 UTC
                                         All use subject to https://about.jstor.org/terms

[form-feed in source extraction]
                                                                                Priority Pricingfor the MIMI1 Queue / 873

dependent arrival rate Xi. The R arrival processes are                 by the system (per unit of time) to the organization,
stochastically independent.                                            where we assume for now that the system administra-
\end{ReviewVerbatim}



\paragraph{Lean-expanded paper semantic target}
\begin{ReviewVerbatim}
type:
{n : ℕ} → (Fin n → ℝ) → (Fin n → ℝ) → Prop

value:
fun {n : ℕ} (meanService a : Fin n → ℝ) =>
  {flow : Fin n → ℝ |
      (∀ (j : Fin n), 0 < meanService j) ∧ (∀ (j : Fin n), 0 ≤ flow j) ∧ ∑ j : Fin n, flow j * meanService j < 1}
    a
\end{ReviewVerbatim}


\paragraph{Recorded source-to-declaration screening}
\textbf{Verdict:} \texttt{matches}
\\
{\raggedright
\textbf{Reason:} cited publication:\allowbreak{}1118-\allowbreak{}1129 gives the stationary nonpreemptive-\allowbreak{}priority M/\allowbreak{}M/\allowbreak{}1 time formula with S\_\allowbreak{}i as cumulative offered work and its slack as 1-\allowbreak{}S\_\allowbreak{}i;\allowbreak{} the target is precisely the nonnegative-\allowbreak{}rate, positive-\allowbreak{}mean, total-\allowbreak{}load-\allowbreak{}below-\allowbreak{}one stationary domain needed for that formula.\allowbreak{}
\par}
\reviewmatch{Matches source input}
\noindent\textbf{Reviewer annotation}\par
\reviewerbox
\clearpage
\hypertarget{paper-prerequisite-3ae6b5293fffafb7}{}
\subsection*{\small\ttfamily\raggedright Mendelson\allowbreak{}Whang1990\allowbreak{}Priority\allowbreak{}Pricing.\allowbreak{}priority\allowbreak{}Time\allowbreak{}Dependent\allowbreak{}Pricing\allowbreak{}Optimal\allowbreak{}And\allowbreak{}Incentive\allowbreak{}Compatible\allowbreak{}At}
\noindent\textbf{Paper-local declaration:}\par
{\footnotesize\ttfamily\raggedright Mendelson\allowbreak{}Whang1990\allowbreak{}Priority\allowbreak{}Pricing.\allowbreak{}priority\allowbreak{}Time\allowbreak{}Dependent\allowbreak{}Pricing\allowbreak{}Optimal\allowbreak{}And\allowbreak{}Incentive\allowbreak{}Compatible\allowbreak{}At\par}
\textbf{Source locator:} {\footnotesize\raggedright cited publication, lines 326-\allowbreak{}362\par}
\paragraph{Verbatim paper-source connection}
\begin{ReviewVerbatim}
Hence                                                                  becomes

                                                                                   r     R   R
            a
                                                                        maxf j Vi(Xi) - viLi(, r)t.
    P~E
     d oi 3ax,
                                                                       If X * and r* are the optimal arrival-rate vector and
      + vi *[Wi(X*)+x* Wi(X*)]-viWi(X*)                                scheduling permutation, respectively, then by the opti-
                                                                       mality of the c-, rule for fixed X = X * we have
           R a~~~i x
      R                                                                             R            R

   - JX,*- W (X*)                                                       max{ Vi(Xi) - viLi(X,(r)
      j=1 J axj

since for] i d                                                                   R           R


                                                                          = E Vi(Xi) -E viLi(2*, r*)
                                                                               i=l       i=l


1a-
 Ax.Li (2X) a - xa
          axi    iW xi
                    = -i a 'J/(X).                                               R           R


                                                                           X Vi(X*) - E viLi(X, rO)
  Note that when R = 1, (6) reduces to (3). Also, as
in (3), the right-hand side of (6) represents the expected                                   R   R

cost inflicted on all other users as a consequence of                      < max Vi(xi) - viLi(X, rO)}
an infinitesimal unit increase in the job flow of a class-
i user. Thus, similar to (3), we may call the right-hand               which demonstrates that the same scheduling rate is
side of (6) the externality cost. The result of Theorem                also optimal in the current setting.

[Next verbatim source excerpt]

                                                              Consider a system with R job classes that satisfy
class-i jobs. As before, all aggregate class characteristics
                                                            assumptions A1-A5 of Section 1, where class-i jobs
are common knowledge, whereas individual job char-
                                                            have exponential processing requirements with mean
acteristics are known only to the user. Thus, both the
                                                           ci. We assume that

[Next verbatim source excerpt]

  Under this general framework, an equilibrium can                     geneous service requirements. The answer is negative,
be defined in a similar way to (12) and (13). Let WJi                  as the following example demonstrates. Consider the
denote the expected total waiting time in the system                   case of two job classes (R = 2) and let p(t) = (pI (t),
for a class-i job picking priority-class j. Then, ( X (p),            P2(t)) be given by (16). Ifp(t) is optimal and incentive-
f(p)) is an equilibrium with respect to the PTD pricing compatible, optimality implies that for a class-i job,
scheme p(t) if for each i = 1, 2, . . ., R              we just have E[pi(ri)] = p*, where p* are given by
                                                                       Theorem 1. Also from (16)
E[pr,(p)(ri)] + viWri,(p)(2(p), r(p))

   = min E[pj(ri)] + viWj(X(p), r(p))}                                 E[p2(T )] =- E[pl(Tri)] =- P
                                                                                             a   a


and                                                                          and                          (17)

Vi'(Xi(p)) = E[prj(p)(Ti)] + Vi Wr,(p)(X (p), r(p))
                                                                       E[pl(T2)] = a * E[p2(T2)] = a *2
where Ti is the actual processing requirement of a
                                                                      where a = K1/K2.
class-i job.
                                                                         In order for the pricing schedule p(t) to be incentive-
  A PTD pricing scheme p0(t) will be called optimal
                                                                       compatible, it should satisfy
and incentive-compatible if it satisfies the following
conditions:
                                                                       E[pI(TI)] + v* W < S E[p2(rl)] + V1 WY
   1. Incentive-compatibility: For all i = 1, 2, . .. , R
                                                                       and
   ri(p?)                   =        i      (14)
                                                                       E[pi(T2)] + v2 * Wl q E[p2(T2)] + V2 WY-
      2.       Optimality:                                       The                 resulting                         arrival
                                                                       Hence using (17), incentive-compatibility implies
  the expected                                          net                value                 of        the            system
is, for all i= 1, 2, .. ., R
                                                                       pi +V* 1 S-*p* + VI * W
  X(pO)                =     _X*                (15)                                         a



where X * achieves the overall optimum.                                      and                          (18)

[Next verbatim source excerpt]

                                                              Consider a system with R job classes that satisfy
class-i jobs. As before, all aggregate class characteristics
                                                            assumptions A1-A5 of Section 1, where class-i jobs
are common knowledge, whereas individual job char-
                                                            have exponential processing requirements with mean
acteristics are known only to the user. Thus, both the
                                                           ci. We assume that
\end{ReviewVerbatim}



\paragraph{Lean-expanded paper semantic target}
\begin{ReviewVerbatim}
type:
{n : ℕ} →
  (Fin n → ℝ → ℝ) →
    (Fin n → ℝ) →
      (Fin n → ℝ) →
        (meanService : Fin n → ℝ) →
          (Fin n → ℝ) → (∀ (j : Fin n), 0 < meanService j) → (Fin n → ℝ → ℝ) → (Fin n → Fin n) → Prop

value:
fun {n : ℕ} (value : Fin n → ℝ → ℝ) (arrivalRate delayCost meanService marginalValue : Fin n → ℝ)
    (hmeanService : ∀ (j : Fin n), 0 < meanService j) (pricing : Fin n → ℝ → ℝ) (assignedPriority : Fin n → Fin n) =>
  IsMaxOn
      (fun (pair : (Fin n → ℝ) × Equiv.Perm (Fin n)) =>
        MendelsonWhang1990PriorityPricing.priorityScheduleNetValue value pair.1 delayCost meanService pair.2)
      ((MendelsonWhang1990PriorityPricing.stablePriorityFlowDomain meanService).prod Set.univ)
      (arrivalRate, Equiv.refl (Fin n)) ∧
    (∀ (i : Fin n),
        HasDerivAt (value i) (marginalValue i) (arrivalRate i) ∧
          marginalValue i =
            MendelsonWhang1990PriorityPricing.priorityTimeDependentExpectedCost arrivalRate delayCost meanService
              pricing i (assignedPriority i)) ∧
      (∀ (i selectedPriority : Fin n),
          MendelsonWhang1990PriorityPricing.priorityTimeDependentExpectedCost arrivalRate delayCost meanService pricing
              i (assignedPriority i) ≤
            MendelsonWhang1990PriorityPricing.priorityTimeDependentExpectedCost arrivalRate delayCost meanService
              pricing i selectedPriority) ∧
        ∀ (i : Fin n), assignedPriority i = i
\end{ReviewVerbatim}

\paragraph{Direct retained dependencies}
\begin{itemize}\raggedright
\item \hyperlink{governing-declaration-2d648b9cf4c8cf80}{Mendelson\allowbreak{}Whang1990\allowbreak{}Priority\allowbreak{}Pricing.\allowbreak{}priority\allowbreak{}Schedule\allowbreak{}Net\allowbreak{}Value}
\item \hyperlink{paper-prerequisite-476808203e3008c1}{Mendelson\allowbreak{}Whang1990\allowbreak{}Priority\allowbreak{}Pricing.\allowbreak{}priority\allowbreak{}Time\allowbreak{}Dependent\allowbreak{}Expected\allowbreak{}Cost}
\item \hyperlink{paper-prerequisite-84dca75fd93aa845}{Mendelson\allowbreak{}Whang1990\allowbreak{}Priority\allowbreak{}Pricing.\allowbreak{}stable\allowbreak{}Priority\allowbreak{}Flow\allowbreak{}Domain}
\end{itemize}
\paragraph{Recorded source-to-declaration screening}
\textbf{Verdict:} \texttt{matches}
\\
{\raggedright
\textbf{Reason:} The selected source bundle states the equilibrium cost-\allowbreak{}minimization condition over every priority j and the marginal-\allowbreak{}user equality V\_\allowbreak{}i'(X\_\allowbreak{}i(p)) = E[p\_\allowbreak{}\{r\_\allowbreak{}i(p)\}(T\_\allowbreak{}i)] + v\_\allowbreak{}i W\_\allowbreak{}\{r\_\allowbreak{}i(p)\}\textasciicircum{}i(X(p), r(p));\allowbreak{} it then requires incentive compatibility r\_\allowbreak{}i(p⁰) = i and optimality X(p⁰) = X*, where X* achieves the overall optimum.\allowbreak{} The displayed Lean target jointly maximizes the displayed stationary net-\allowbreak{}value objective over stable flow vectors and priority permutations at (arrivalRate, identity), requires the derivative/\allowbreak{}equal-\allowbreak{}expected-\allowbreak{}cost equality for each class, minimizes that expected cost over every selected priority, and requires assignedPriority i = i.\allowbreak{} The displayed stable-\allowbreak{}flow and priority-\allowbreak{}sojourn prerequisites restrict this comparison to the source's exponential stationary priority-\allowbreak{}queue domain, so all source conditions in the supplied bundle are represented without a changed formula, quantifier, or conclusion.\allowbreak{}
\par}
\reviewmatch{Matches source input}
\noindent\textbf{Reviewer annotation}\par
\reviewerbox
\clearpage
\hypertarget{paper-prerequisite-33c9dc55ddc7d5b6}{}
\subsection*{\small\ttfamily\raggedright Mendelson\allowbreak{}Whang1990\allowbreak{}Priority\allowbreak{}Pricing.\allowbreak{}source\allowbreak{}Delay\allowbreak{}Cost\allowbreak{}Conditions}
\noindent\textbf{Paper-local declaration:}\par
{\footnotesize\ttfamily\raggedright Mendelson\allowbreak{}Whang1990\allowbreak{}Priority\allowbreak{}Pricing.\allowbreak{}source\allowbreak{}Delay\allowbreak{}Cost\allowbreak{}Conditions\par}
\textbf{Source locator:} {\footnotesize\raggedright cited publication, lines 259-\allowbreak{}262\par}
\paragraph{Verbatim paper-source connection}
\begin{ReviewVerbatim}
A3. Delay Cost Structure
                                                                       R-dimensional price vector p = (PI, P2, P3, .., PR),
Each class-i job is characterized by a delay cost of vi                where pi is the price charged to a class-i job. Then, the
per unit time.                                                         class-i demand relationship is given by
\end{ReviewVerbatim}



\paragraph{Lean-expanded paper semantic target}
\begin{ReviewVerbatim}
type:
{Class : Type u_1} → [Fintype Class] → (Class → ℝ) → Prop

value:
fun {Class : Type u_1} [Fintype Class] (delayCost : Class → ℝ) => ∀ (i : Class), 0 ≤ delayCost i
\end{ReviewVerbatim}


\paragraph{Recorded source-to-declaration screening}
\textbf{Verdict:} \texttt{matches}
\\
{\raggedright
\textbf{Reason:} The source calls v\_\allowbreak{}i a delay cost per unit time and explicitly introduces the one-\allowbreak{}class delay cost as positive.\allowbreak{} Lean records the corresponding nonnegative cost-\allowbreak{}domain convention class by class, including the zero-\allowbreak{}cost boundary;\allowbreak{} this weak closure does not omit or reverse any source restriction and makes the formal theorem at least as general as the stated cost model.\allowbreak{}
\par}
\reviewmatch{Matches source input}
\noindent\textbf{Reviewer annotation}\par
\reviewerbox
\clearpage
\hypertarget{paper-prerequisite-9c435671d506207c}{}
\subsection*{\small\ttfamily\raggedright Mendelson\allowbreak{}Whang1990\allowbreak{}Priority\allowbreak{}Pricing.\allowbreak{}source\allowbreak{}Value\allowbreak{}Function\allowbreak{}Conditions}
\noindent\textbf{Paper-local declaration:}\par
{\footnotesize\ttfamily\raggedright Mendelson\allowbreak{}Whang1990\allowbreak{}Priority\allowbreak{}Pricing.\allowbreak{}source\allowbreak{}Value\allowbreak{}Function\allowbreak{}Conditions\par}
\textbf{Source locator:} {\footnotesize\raggedright cited publication, lines 244-\allowbreak{}257\par}
\paragraph{Verbatim paper-source connection}
\begin{ReviewVerbatim}
A2. Value Function
                                                                       belongs. Under a given prionity scheme, this problem
The value function of the class-i job stream can be                    can be formulated as
expressed as Vi (Xi), and the total value is given by the                        F       R
sum of Vi(Xi) over all classes, that is, the expected                   max V(X) - E viLi(X)J (4)
(gross) value function is given by                                          X        i=l


           R                                                           where Li is the mean number of class-i jobs in the
V(X)= E V1(Xi)                                                         system in steady state, viLi is the expected delay cost
           i=l1
                                                                       incurred by class-i jobs per unit of time, and V(2X) =
where Vi(Xi) are differentiable, nondecreasing and                     E 1 Vi(Xi) is the total gross value per unit of time.
concave functions of their arguments Xi for all i.                       Consider the multiple class analog of the demand
\end{ReviewVerbatim}



\paragraph{Lean-expanded paper semantic target}
\begin{ReviewVerbatim}
type:
{Class : Type u_1} → [Fintype Class] → (Class → ℝ → ℝ) → Prop

value:
fun {Class : Type u_1} [Fintype Class] (value : Class → ℝ → ℝ) =>
  (∀ (i : Class), DifferentiableOn ℝ (value i) (Set.Ici 0)) ∧
    (∀ (i : Class), MonotoneOn (value i) (Set.Ici 0)) ∧ ∀ (i : Class), ConcaveOn ℝ (Set.Ici 0) (value i)
\end{ReviewVerbatim}


\paragraph{Recorded source-to-declaration screening}
\textbf{Verdict:} \texttt{matches}
\\
{\raggedright
\textbf{Reason:} The expanded Lean predicate now includes all three properties stated in Assumption A2—differentiability, nondecreasingness, and concavity—for every class value function on the nonnegative flow domain.\allowbreak{} No A2 property is left only in prose or in a theorem-\allowbreak{}local premise.\allowbreak{}
\par}
\reviewmatch{Matches source input}
\noindent\textbf{Reviewer annotation}\par
\reviewerbox
\clearpage
\hypertarget{paper-prerequisite-816477dc3af70c36}{}
\subsection*{\small\ttfamily\raggedright Mendelson\allowbreak{}Whang1990\allowbreak{}Priority\allowbreak{}Pricing.\allowbreak{}stationary\allowbreak{}Class\allowbreak{}Tagged\allowbreak{}Expected\allowbreak{}Priority\allowbreak{}Time\allowbreak{}Dependent\allowbreak{}Price}
\noindent\textbf{Paper-local declaration:}\par
{\footnotesize\ttfamily\raggedright Mendelson\allowbreak{}Whang1990\allowbreak{}Priority\allowbreak{}Pricing.\allowbreak{}stationary\allowbreak{}Class\allowbreak{}Tagged\allowbreak{}Expected\allowbreak{}Priority\allowbreak{}Time\allowbreak{}Dependent\allowbreak{}Price\par}
\textbf{Source locator:} {\footnotesize\raggedright cited publication, lines 760-\allowbreak{}805\par}
\paragraph{Verbatim paper-source connection}
\begin{ReviewVerbatim}
and incentive-compatibility in the general setting                   E[pt+(ri)] = E[Ai Ti + 112B *Ti r]
with heterogeneous service requirements. What price                                 = Ai * ci + B * c?,
structure is optimal and incentive compatible? The
                                                                    since E[r,] = 2ci . It follows that
specific PTD pricing scheme we propose is given by
p+(t) = (p+(t), p2+(t), . . ., pZ(t)) with                          p* = E[pt((ri)] for alli= 1, 2, ...,R- 1.
                                                                    Likewise
Pti(t) = Ait + V12Bt2 (19)

where                                                                A
                                                                     PR= -=-a'R
                                                                              -2CR V XkR
                                                                                 CR+(E    2= E[p+(TR)]
                                                                                       - -)R E[R(R)
                                                                           SR- SR k=1 Sk-l Sk

                                                                    which completes the proof of optimality.
 B        =     E     .k    *     Xk(20)
       k= 1 Sk-l Sk                                                    b. Next we prove incentive-compatibility, that is,
                                                                    class-i users will find it in their best interest to join
                                                                    priority-class i. As in Theorem 2, we prove local
 Ai 2+ 2 + - __ 2) (21)                                             optimality and then transitivity.
         z-1 9 k=i+l 'kISk k Sk-1SIk
                                                                       i. Local optimality: If a class-i job (i = 1,
with Ai = EL cIX*, ai = ViXA AR, Si = REk=1     ckXCk
                                        - 1) joins priority-class i, its expected total cost
So = 0, S) = 1 - Si, and where X * are evaluated at the
                                                                     is E[pt((ri)] + vi * (Wi7 + ci), where Wi' denotes
optimum of Theorem 1, and summations from R to
                                                                    the expected queueing time (excluding service) for
(R - 1) are interpreted as zero (see the Appendix for
                                                                    priority-class i. On the other hand, if the user chooses
a summary of the notation). Expression (21) may also
                                                                    priority-class (i + 1), the expected total cost of the job
be written in the equivalent recursive form
                                                                     is given by E[pl+(Tri)] + vi * (Wl1+1 + ci). Thus, for
                                                                     all i= 1, 2, . . ., R - I wehave
Ai=Ai-,- a- -2 + -2ai_) i=2,3 ... R
                                                                     ELpo(Tri)] + vi Wi7- EE[p +I(ri)] + viWi'+ I

with                                                                            r a, ai+ + lAR AR
                                                                           =Ni I 37 +1 13i23+1 CI -S
           aR
AR = - -2                                                                   _AR
        SR-1 * SR
                                                                                          - 1 i S [civS3i+1 + civi+1 IX1 1i-I
\end{ReviewVerbatim}



\paragraph{Lean-expanded paper semantic target}
\begin{ReviewVerbatim}
type:
{n : ℕ} → (arrivalRate : Fin n → ℝ) → (Fin n → ℝ) → (Fin n → ℝ) → (∀ (j : Fin n), 0 < arrivalRate j) → Fin n → Fin n → ℝ

value:
fun {n : ℕ} (arrivalRate delayCost meanService : Fin n → ℝ) (harrivalRate : ∀ (j : Fin n), 0 < arrivalRate j)
    (customer declaredPriority : Fin n) =>
  ∫ (omega :
    ((ℤ → ℝ) × (ℤ → ℝ)) × ({ j : Fin n // j ≠ customer } → AppliedModelingLib.Queueing.StationaryPoissonWorkPath)),
    MendelsonWhang1990PriorityPricing.heterogeneousPriorityTimeDependentPrice arrivalRate delayCost meanService
      declaredPriority
      (AppliedModelingLib.Queueing.stationaryPriorityClassTaggedWorkRequirement meanService customer
        omega) ∂(AppliedModelingLib.Probability.Palm.targetPassiveTaggedArrivalAtZero
        (AppliedModelingLib.Probability.Queueing.stationaryPoissonWorkTaggedArrivalAtZero (harrivalRate customer))
        (AppliedModelingLib.Queueing.multiclassStationaryPoissonWorkRestLaw arrivalRate harrivalRate customer)).Ptag
\end{ReviewVerbatim}

\paragraph{Direct retained dependencies}
\begin{itemize}\raggedright
\item \hyperlink{governing-declaration-5a10835f452c324b}{Applied\allowbreak{}Modeling\allowbreak{}Lib.\allowbreak{}Probability.\allowbreak{}Palm.\allowbreak{}target\allowbreak{}Passive\allowbreak{}Tagged\allowbreak{}Arrival\allowbreak{}At\allowbreak{}Zero}
\item \hyperlink{governing-declaration-ab1cfe37ad8c32b0}{Applied\allowbreak{}Modeling\allowbreak{}Lib.\allowbreak{}Probability.\allowbreak{}Poisson\allowbreak{}Process.\allowbreak{}Good\allowbreak{}Suspension\allowbreak{}State}
\item \hyperlink{governing-declaration-a8385d799d0a4d16}{Applied\allowbreak{}Modeling\allowbreak{}Lib.\allowbreak{}Probability.\allowbreak{}Poisson\allowbreak{}Process.\allowbreak{}inst\allowbreak{}Measurable\allowbreak{}Space\allowbreak{}Good\allowbreak{}Suspension\allowbreak{}State}
\item \hyperlink{governing-declaration-a69bedb5b0d420c6}{Applied\allowbreak{}Modeling\allowbreak{}Lib.\allowbreak{}Probability.\allowbreak{}Queueing.\allowbreak{}Tagged\allowbreak{}Arrival\allowbreak{}At\allowbreak{}Zero}
\item \hyperlink{governing-declaration-e9116e504d73a45c}{Applied\allowbreak{}Modeling\allowbreak{}Lib.\allowbreak{}Probability.\allowbreak{}Queueing.\allowbreak{}stationary\allowbreak{}Poisson\allowbreak{}Work\allowbreak{}Tagged\allowbreak{}Arrival\allowbreak{}At\allowbreak{}Zero}
\item \hyperlink{governing-declaration-62a268164b6ae28b}{Applied\allowbreak{}Modeling\allowbreak{}Lib.\allowbreak{}Queueing.\allowbreak{}Stationary\allowbreak{}Poisson\allowbreak{}Work\allowbreak{}Path}
\item \hyperlink{governing-declaration-9bede4eaca6a95fb}{Applied\allowbreak{}Modeling\allowbreak{}Lib.\allowbreak{}Queueing.\allowbreak{}multiclass\allowbreak{}Stationary\allowbreak{}Poisson\allowbreak{}Work\allowbreak{}Rest\allowbreak{}Law}
\item \hyperlink{library-prerequisite-8f1ec496ea0db7a0}{Applied\allowbreak{}Modeling\allowbreak{}Lib.\allowbreak{}Queueing.\allowbreak{}stationary\allowbreak{}Priority\allowbreak{}Class\allowbreak{}Tagged\allowbreak{}Work\allowbreak{}Requirement}
\item \hyperlink{paper-prerequisite-d3f248a0395e1c76}{Mendelson\allowbreak{}Whang1990\allowbreak{}Priority\allowbreak{}Pricing.\allowbreak{}heterogeneous\allowbreak{}Priority\allowbreak{}Time\allowbreak{}Dependent\allowbreak{}Price}
\end{itemize}
\paragraph{Recorded source-to-declaration screening}
\textbf{Verdict:} \texttt{matches}
\\
{\raggedright
\textbf{Reason:} cited publication:\allowbreak{}760-\allowbreak{}805 specifies p\_\allowbreak{}i\textasciicircum{}+(t)=A\_\allowbreak{}i t+(1/\allowbreak{}2)B t\textasciicircum{}2 and evaluates E[p\_\allowbreak{}i\textasciicircum{}+(T\_\allowbreak{}i)] for the class-\allowbreak{}specific exponential requirement;\allowbreak{} the target is that same PTD charge averaged over the stationary class-\allowbreak{}tagged service requirement, with the declared priority selecting i.\allowbreak{}
\par}
\reviewmatch{Matches source input}
\noindent\textbf{Reviewer annotation}\par
\reviewerbox

\clearpage
\hypertarget{source-claim-0d52797e81514cb8}{}
\hypertarget{source-section-f10b5f68e4acf3dc}{}
\section*{Source claims}
\subsection*{Review row 1}
\textbf{Source locator:} {\footnotesize\raggedright cited publication, lines 272-\allowbreak{}305\par}
\paragraph{Verbatim source input}
\begin{ReviewVerbatim}
Theorem 1. Under the above assumptions, the optimal
A5. Decision Structure
                                                                      price per class-i job is given by
We assume the individual decision structure: users do
not collude, and each maximizes his expected net
                                                                                 R aow.
                                                                        Pi =, vjXi, for i = 1, 2, ..., R (6)
benefit.

  The specific assumptions of exponential service                      where Wj is the expected delay of a class-j job, and
time distributions and nonpreemptive priority among                    the derivatives a Wj/0 Xi are evaluated at X = X *, the
job classes will be invoked only from the next section;                arrival-rate vector that maximizes (4).

the results derived in this section are quite general,
                                                                       Proof. The ith first-order condition for maximizing
and apply to any service-time distribution and head-
                                                                       (4) is
of-the-line priority discipline. We assume throughout
                                                                                     R
that the joining and pliority decisions are made by
                                                                          Vi'            (Xi)=>           viyj.         (7)
atomistic individual users, each maximizing the                                          j=I A
expected net benefit. This individual decision struc-
                                                                       Furthermore, the class-i demand relationship is again
ture is in sharp contrast to the class (or collusive)
                                                                       given by (5)
decision structure, where each class is represented by
a single decision maker who makes the arrival-rate                      Vi/(Xi) = Pi + Vi * wi
and priority decisions for the class so as to maximize
                                                                       Combining the above equations we obtain, at
the expected net value of the class as a whole. Clearly,
                                                                        X= A*
pricing decisions for our individual decision structure
                                                                                 R a
\end{ReviewVerbatim}


\paragraph{Formalized review target}
The archival source above is not asserted equivalent to this formalized target.
\begin{ReviewVerbatim}
Theorem 1's externality-price identity holds for each selected class whose welfare-maximizing arrival rate is strictly positive, under the paper's stated value, delay-cost, demand, and differentiability conditions.
\end{ReviewVerbatim}

\textbf{Archival source anchor:} {\footnotesize\raggedright cited publication, lines 272-\allowbreak{}305\par}
\paragraph{Expanded PaperInterface specification}
\begin{ReviewVerbatim}
∀ {Class : Type u_1} [inst : Fintype Class] (value : Class → ℝ → ℝ) (delayCost flow price : Class → ℝ)
  (waitingTime : Class → (Class → ℝ) → ℝ) (i : Class) (marginalValue : ℝ) (partialWaiting : Class → ℝ),
  MendelsonWhang1990PriorityPricing.sourceValueFunctionConditions value →
    MendelsonWhang1990PriorityPricing.sourceDelayCostConditions delayCost →
      (∀ (j : Class), 0 ≤ flow j) →
        0 < flow i →
          IsMaxOn (MendelsonWhang1990PriorityPricing.netValue value delayCost waitingTime)
              (Set.Ici fun (x : Class) => 0) flow →
            MendelsonWhang1990PriorityPricing.individualDemandCondition value delayCost flow price waitingTime i
                marginalValue →
              (∀ (j : Class),
                  HasDerivAt (fun (t : ℝ) => waitingTime j (Function.update flow i t)) (partialWaiting j) (flow i)) →
                price i = MendelsonWhang1990PriorityPricing.externalityPrice delayCost flow partialWaiting
\end{ReviewVerbatim}

\paragraph{Retained governing declarations}
\begin{itemize}\raggedright
\item \hyperlink{governing-declaration-806f5a55d6021334}{Mendelson\allowbreak{}Whang1990\allowbreak{}Priority\allowbreak{}Pricing.\allowbreak{}externality\allowbreak{}Price}
\item \hyperlink{paper-prerequisite-6c93de2469b53cb2}{Mendelson\allowbreak{}Whang1990\allowbreak{}Priority\allowbreak{}Pricing.\allowbreak{}individual\allowbreak{}Demand\allowbreak{}Condition}
\item \hyperlink{paper-prerequisite-7dbaca4145f3a4a5}{Mendelson\allowbreak{}Whang1990\allowbreak{}Priority\allowbreak{}Pricing.\allowbreak{}net\allowbreak{}Value}
\item \hyperlink{governing-declaration-3df5d95507228c45}{Mendelson\allowbreak{}Whang1990\allowbreak{}Priority\allowbreak{}Pricing.\allowbreak{}paper\allowbreak{}Interface\allowbreak{}Classical\allowbreak{}Decidable\allowbreak{}Eq}
\item \hyperlink{paper-prerequisite-33c9dc55ddc7d5b6}{Mendelson\allowbreak{}Whang1990\allowbreak{}Priority\allowbreak{}Pricing.\allowbreak{}source\allowbreak{}Delay\allowbreak{}Cost\allowbreak{}Conditions}
\item \hyperlink{paper-prerequisite-9c435671d506207c}{Mendelson\allowbreak{}Whang1990\allowbreak{}Priority\allowbreak{}Pricing.\allowbreak{}source\allowbreak{}Value\allowbreak{}Function\allowbreak{}Conditions}
\end{itemize}
\paragraph{Lean-elaborated claim roles}\begin{ReviewVerbatim}
1. parameter: Type u_1
2. parameter: Fintype Class
3. parameter: Class → ℝ → ℝ
4. parameter: Class → ℝ
5. parameter: Class → ℝ
6. parameter: Class → ℝ
7. parameter: Class → (Class → ℝ) → ℝ
8. parameter: Class
9. parameter: ℝ
10. parameter: Class → ℝ
11. assumption: MendelsonWhang1990PriorityPricing.sourceValueFunctionConditions value
12. assumption: MendelsonWhang1990PriorityPricing.sourceDelayCostConditions delayCost
13. assumption: ∀ (j : Class), 0 ≤ flow j
14. assumption: 0 < flow i
15. assumption: IsMaxOn (MendelsonWhang1990PriorityPricing.netValue value delayCost waitingTime) (Set.Ici fun (x : Class) => 0) flow
16. assumption: MendelsonWhang1990PriorityPricing.individualDemandCondition value delayCost flow price waitingTime i marginalValue
17. assumption: ∀ (j : Class), HasDerivAt (fun (t : ℝ) => waitingTime j (Function.update flow i t)) (partialWaiting j) (flow i)
18. conclusion: price i = MendelsonWhang1990PriorityPricing.externalityPrice delayCost flow partialWaiting
\end{ReviewVerbatim}

\paragraph{Lean proof endpoint (built separately)}\begin{ReviewVerbatim}
MendelsonWhang1990PriorityPricing.theoremOneExternalityPricing
\end{ReviewVerbatim}

\paragraph{Recorded source-to-Spec screening}
\textbf{Verdict:} \texttt{matches\_approved\_corrected\_target}
\\
{\raggedright
\textbf{Reason:} Theorem 1's displayed source statement and proof quantify classes i,j = 1,.\allowbreak{}.\allowbreak{}.\allowbreak{},R and use the finite sum over j, while the target quantifies an arbitrary finite Class through [Fintype Class] and concludes the same class-\allowbreak{}i price as the sum of delayCost j * flow j * partialWaiting j.\allowbreak{} The source proof combines the class-\allowbreak{}i first-\allowbreak{}order condition with the individual demand equality;\allowbreak{} the target makes those two displayed ingredients explicit, and the approved corrected-\allowbreak{}target record expressly restricts the selected optimal coordinate to 0 < flow i, the condition under which that ordinary first-\allowbreak{}order equality is used.\allowbreak{} No DecidableEq Class premise is needed to represent the source's finite class domain:\allowbreak{} the displayed local paperInterfaceClassicalDecidableEq supplies Classical.\allowbreak{}decEq solely for Function.\allowbreak{}update in the coordinate-\allowbreak{}derivative expression, so its omission from the public target neither changes the finite-\allowbreak{}class model domain nor adds or removes a source-\allowbreak{}side economic assumption.\allowbreak{}
\par}
\reviewmatch{Matches formalized target}
\noindent\textbf{Reviewer annotation}\par
\reviewerbox
\clearpage
\hypertarget{source-claim-86fb5cbbc8012c1d}{}
\section*{Review row 2}
\textbf{Source locator:} {\footnotesize\raggedright cited publication, lines 1123-\allowbreak{}1130\par}
\paragraph{Verbatim source input}
\begin{ReviewVerbatim}
For the M/M/1 queue with nonpreemptive priority                         KREPS, D. 1990. A Course in Microeconomic Theory.
discipline (Heyman and Sobel 1982, p. 435)                                 Princeton University Press, Princeton, N.J.
                                                                        MARCHAND, M. 1974. Priority Pricing. Mgmt. Sci. 20,
   _ i( L1()
 W1CX)  = W1 Wqc1+ -+
                   Ci AR   (A-)
                      cj. (-                                                  1131-1140.
                Xi        i-,           i
\end{ReviewVerbatim}



\paragraph{Expanded PaperInterface specification}
\begin{ReviewVerbatim}
∀ {n : ℕ} (arrivalRate meanService : Fin n → ℝ) (harrivalRate : ∀ (j : Fin n), 0 < arrivalRate j),
  (∀ (j : Fin n), 0 < meanService j) →
    ∑ j : Fin n, arrivalRate j * meanService j < 1 →
      ∀ (i : Fin n),
        ∫ (z : AppliedModelingLib.Queueing.MulticlassStationaryPoissonWorkClassTaggedSample (Fin n) i),
              AppliedModelingLib.Queueing.stationaryPriorityClassTaggedResponseTime meanService i
                z ∂AppliedModelingLib.Queueing.stationaryPriorityClassTaggedPalmMeasure arrivalRate harrivalRate i =
            MendelsonWhang1990PriorityPricing.prioritySojournTime arrivalRate meanService i ∧
          MendelsonWhang1990PriorityPricing.priorityPopulation arrivalRate meanService i =
              arrivalRate i * MendelsonWhang1990PriorityPricing.prioritySojournTime arrivalRate meanService i ∧
            MendelsonWhang1990PriorityPricing.priorityPopulation arrivalRate meanService i / arrivalRate i =
              MendelsonWhang1990PriorityPricing.prioritySojournTime arrivalRate meanService i
\end{ReviewVerbatim}

\paragraph{Retained governing declarations}
\begin{itemize}\raggedright
\item \hyperlink{governing-declaration-ab1cfe37ad8c32b0}{Applied\allowbreak{}Modeling\allowbreak{}Lib.\allowbreak{}Probability.\allowbreak{}Poisson\allowbreak{}Process.\allowbreak{}Good\allowbreak{}Suspension\allowbreak{}State}
\item \hyperlink{governing-declaration-a8385d799d0a4d16}{Applied\allowbreak{}Modeling\allowbreak{}Lib.\allowbreak{}Probability.\allowbreak{}Poisson\allowbreak{}Process.\allowbreak{}inst\allowbreak{}Measurable\allowbreak{}Space\allowbreak{}Good\allowbreak{}Suspension\allowbreak{}State}
\item \hyperlink{library-prerequisite-755cdb28ceaf46d2}{Applied\allowbreak{}Modeling\allowbreak{}Lib.\allowbreak{}Queueing.\allowbreak{}Multiclass\allowbreak{}Stationary\allowbreak{}Poisson\allowbreak{}Work\allowbreak{}Class\allowbreak{}Tagged\allowbreak{}Sample}
\item \hyperlink{governing-declaration-62a268164b6ae28b}{Applied\allowbreak{}Modeling\allowbreak{}Lib.\allowbreak{}Queueing.\allowbreak{}Stationary\allowbreak{}Poisson\allowbreak{}Work\allowbreak{}Path}
\item \hyperlink{library-prerequisite-c3ca5cbd067c9957}{Applied\allowbreak{}Modeling\allowbreak{}Lib.\allowbreak{}Queueing.\allowbreak{}stationary\allowbreak{}Priority\allowbreak{}Class\allowbreak{}Tagged\allowbreak{}Palm\allowbreak{}Measure}
\item \hyperlink{library-prerequisite-6b2b5577562811ed}{Applied\allowbreak{}Modeling\allowbreak{}Lib.\allowbreak{}Queueing.\allowbreak{}stationary\allowbreak{}Priority\allowbreak{}Class\allowbreak{}Tagged\allowbreak{}Response\allowbreak{}Time}
\item \hyperlink{governing-declaration-65f779b7d3134146}{Mendelson\allowbreak{}Whang1990\allowbreak{}Priority\allowbreak{}Pricing.\allowbreak{}priority\allowbreak{}Population}
\item \hyperlink{governing-declaration-703b87082556903f}{Mendelson\allowbreak{}Whang1990\allowbreak{}Priority\allowbreak{}Pricing.\allowbreak{}priority\allowbreak{}Sojourn\allowbreak{}Time}
\end{itemize}
\paragraph{Lean-elaborated claim roles}\begin{ReviewVerbatim}
1. parameter: ℕ
2. parameter: Fin n → ℝ
3. parameter: Fin n → ℝ
4. assumption: ∀ (j : Fin n), 0 < arrivalRate j
5. assumption: ∀ (j : Fin n), 0 < meanService j
6. assumption: ∑ j : Fin n, arrivalRate j * meanService j < 1
7. parameter: Fin n
8. conclusion: ∫ (z : AppliedModelingLib.Queueing.MulticlassStationaryPoissonWorkClassTaggedSample (Fin n) i),
      AppliedModelingLib.Queueing.stationaryPriorityClassTaggedResponseTime meanService i
        z ∂AppliedModelingLib.Queueing.stationaryPriorityClassTaggedPalmMeasure arrivalRate harrivalRate i =
    MendelsonWhang1990PriorityPricing.prioritySojournTime arrivalRate meanService i ∧
  MendelsonWhang1990PriorityPricing.priorityPopulation arrivalRate meanService i =
      arrivalRate i * MendelsonWhang1990PriorityPricing.prioritySojournTime arrivalRate meanService i ∧
    MendelsonWhang1990PriorityPricing.priorityPopulation arrivalRate meanService i / arrivalRate i =
      MendelsonWhang1990PriorityPricing.prioritySojournTime arrivalRate meanService i
\end{ReviewVerbatim}

\paragraph{Lean proof endpoint (built separately)}\begin{ReviewVerbatim}
MendelsonWhang1990PriorityPricing.appendixPriorityQueueingTime
\end{ReviewVerbatim}

\paragraph{Recorded source-to-Spec screening}
\textbf{Verdict:} \texttt{matches}
\\
{\raggedright
\textbf{Reason:} The Appendix A-\allowbreak{}1 source item is the stationary M/\allowbreak{}M/\allowbreak{}1 nonpreemptive-\allowbreak{}priority mean-\allowbreak{}time formula.\allowbreak{}  The target instantiates that formula with positive class-\allowbreak{}specific exponential service means and stable offered load;\allowbreak{} its population and population-\allowbreak{}to-\allowbreak{}rate conjuncts only make the corresponding Little-\allowbreak{}law notation explicit, with positivity making the division valid.\allowbreak{}
\par}
\reviewmatch{Matches source input}
\noindent\textbf{Reviewer annotation}\par
\reviewerbox
\clearpage
\hypertarget{source-claim-35dff0bacecfebeb}{}
\section*{Review row 3}
\textbf{Source locator:} {\footnotesize\raggedright cited publication, lines 550-\allowbreak{}559\par}
\paragraph{Verbatim source input}
\begin{ReviewVerbatim}
Theorem 2. Let p* be the optimal PD price vector                        an equality.
given by (10). Then, p* is incentive-compatible as well. This proves the local optimality of choosing priority
That is                                               class i, that is

                                                                        TCQ = pi* + vi Wi-;-< p, I + vi Wi-1 = TC-'
p*I + vi W, CPk* + Vi Wk
                                                                        and
                    for any k E{ 1, 2, . . ., R -{i

                                                                        TC' = p* + vi Wi < p+I + vi W1+ I = TQ+l.
\end{ReviewVerbatim}


\paragraph{Formalized review target}
The archival source above is not asserted equivalent to this formalized target.
\begin{ReviewVerbatim}
Theorem 2's strict priority incentive comparison holds on the homogeneous stable full-positive-flow interior domain, so each class strictly prefers its assigned priority to every distinct priority.
\end{ReviewVerbatim}

\textbf{Archival source anchor:} {\footnotesize\raggedright cited publication, lines 550-\allowbreak{}559\par}
\paragraph{Expanded PaperInterface specification}
\begin{ReviewVerbatim}
∀ {n : ℕ} (arrivalRate delayCost : Fin n → ℝ),
  (∀ (j : Fin n), 0 < arrivalRate j) →
    ∑ j : Fin n, arrivalRate j < 1 →
      (∀ {i k : Fin n}, i < k → delayCost k < delayCost i) →
        ∀ (i selectedPriority : Fin n),
          selectedPriority ≠ i →
            MendelsonWhang1990PriorityPricing.homogeneousPriorityExpectedCost arrivalRate delayCost i i <
              MendelsonWhang1990PriorityPricing.homogeneousPriorityExpectedCost arrivalRate delayCost i selectedPriority
\end{ReviewVerbatim}

\paragraph{Retained governing declarations}
\begin{itemize}\raggedright
\item \hyperlink{paper-prerequisite-53f58e4a79f3cd3c}{Mendelson\allowbreak{}Whang1990\allowbreak{}Priority\allowbreak{}Pricing.\allowbreak{}homogeneous\allowbreak{}Priority\allowbreak{}Expected\allowbreak{}Cost}
\end{itemize}
\paragraph{Lean-elaborated claim roles}\begin{ReviewVerbatim}
1. parameter: ℕ
2. parameter: Fin n → ℝ
3. parameter: Fin n → ℝ
4. assumption: ∀ (j : Fin n), 0 < arrivalRate j
5. assumption: ∑ j : Fin n, arrivalRate j < 1
6. assumption: ∀ {i k : Fin n}, i < k → delayCost k < delayCost i
7. parameter: Fin n
8. parameter: Fin n
9. assumption: selectedPriority ≠ i
10. conclusion: MendelsonWhang1990PriorityPricing.homogeneousPriorityExpectedCost arrivalRate delayCost i i <
  MendelsonWhang1990PriorityPricing.homogeneousPriorityExpectedCost arrivalRate delayCost i selectedPriority
\end{ReviewVerbatim}

\paragraph{Lean proof endpoint (built separately)}\begin{ReviewVerbatim}
MendelsonWhang1990PriorityPricing.theoremTwoPriorityIncentiveCompatibility
\end{ReviewVerbatim}

\paragraph{Recorded source-to-Spec screening}
\textbf{Verdict:} \texttt{matches\_approved\_corrected\_target}
\\
{\raggedright
\textbf{Reason:} The source's strict assigned-\allowbreak{}priority comparison for every distinct priority is the Spec's strict homogeneous expected-\allowbreak{}cost inequality on the recorded stable full-\allowbreak{}positive homogeneous domain.\allowbreak{}
\par}
\reviewmatch{Matches formalized target}
\noindent\textbf{Reviewer annotation}\par
\reviewerbox
\clearpage
\hypertarget{source-claim-a044a878df73f034}{}
\section*{Review row 4}
\textbf{Source locator:} {\footnotesize\raggedright cited publication, lines 806-\allowbreak{}846\par}
\paragraph{Verbatim source input}
\begin{ReviewVerbatim}
Theorem 3. The PTD pricing scheme (19) is optimal
                                                                                                          + viSi(Si+1 - Si-1)]
and incentive-compatible.
                                                                               AR

Proof. a. We first prove the optimality of (19). It is                         s1~~~-5vl+*i+ X* + civi+i X*'
sufficient to prove that for all i = 1, 2, ...,
E[pt+(ri)] = pi*, where p* are given by Theorem 1.                             *i-)              i+         ki
For all i = 1, 2,.. .,R -1, we have from Theorem 1

                                                                            _AR
P' E VkXk-
      k=1 xi


                                                                            * [i+1(ciVi+1 - CE+lv1) CkXk)j <
       =[332 + 1E (-!2 + -2 _ ak) ]
          i (i k=i+l S)k-lSk Ek k-I
                                                                     since by construction ci vi+1 < ci+1 vi. Thus

                                                                     E[pt(ri)] + iviW < E[p ++I(ri)] + viWlq+
            k= I k-1 i2
        -A1 -. c,
                                                                    i.e., any class-i job will prefer priority-class i over
        =Ai * ci + B * c?I                                          priority-class (i + 1).




                                               This content downloaded from
                             108.30.55.84 on Mon, 24 Aug 2026 10:55:53 UTC
                                       All use subject to https://about.jstor.org/terms

[form-feed in source extraction]
880 / Mendelson and Whang

  In a similar way, for i = 2, 3,..., R                                which completes the proof of transitivity for i < j <
                                                                       k. Transitivity for i > j > k follows in a similar way.
{E[pt(ri)] + vj WI - {E[pti (ri)] + vi Wiq Ij
                                                                         Incentive-compatibility follows by repeated appli-
                                                                       cation of the transitivity principle, showing that
       = ai- + a + l | [ AR hAR]                                       TC' < TC' for all i = 1, 2, ...,5 R, j E {1, 25 ....5
                                                                       R}-{i}.

[Next verbatim source excerpt]

  Under this general framework, an equilibrium can                     geneous service requirements. The answer is negative,
be defined in a similar way to (12) and (13). Let WJi                  as the following example demonstrates. Consider the
denote the expected total waiting time in the system                   case of two job classes (R = 2) and let p(t) = (pI (t),
for a class-i job picking priority-class j. Then, ( X (p),            P2(t)) be given by (16). Ifp(t) is optimal and incentive-
f(p)) is an equilibrium with respect to the PTD pricing compatible, optimality implies that for a class-i job,
scheme p(t) if for each i = 1, 2, . . ., R              we just have E[pi(ri)] = p*, where p* are given by
                                                                       Theorem 1. Also from (16)
E[pr,(p)(ri)] + viWri,(p)(2(p), r(p))

   = min E[pj(ri)] + viWj(X(p), r(p))}                                 E[p2(T )] =- E[pl(Tri)] =- P
                                                                                             a   a


and                                                                          and                          (17)

Vi'(Xi(p)) = E[prj(p)(Ti)] + Vi Wr,(p)(X (p), r(p))
                                                                       E[pl(T2)] = a * E[p2(T2)] = a *2
where Ti is the actual processing requirement of a
                                                                      where a = K1/K2.
class-i job.
                                                                         In order for the pricing schedule p(t) to be incentive-
  A PTD pricing scheme p0(t) will be called optimal
                                                                       compatible, it should satisfy
and incentive-compatible if it satisfies the following
conditions:
                                                                       E[pI(TI)] + v* W < S E[p2(rl)] + V1 WY
   1. Incentive-compatibility: For all i = 1, 2, . .. , R
                                                                       and
   ri(p?)                   =        i      (14)
                                                                       E[pi(T2)] + v2 * Wl q E[p2(T2)] + V2 WY-
      2.       Optimality:                                       The                 resulting                         arrival
                                                                       Hence using (17), incentive-compatibility implies
  the expected                                          net                value                 of        the            system
is, for all i= 1, 2, .. ., R
                                                                       pi +V* 1 S-*p* + VI * W
  X(pO)                =     _X*                (15)                                         a



where X * achieves the overall optimum.                                      and                          (18)

[Next verbatim source excerpt]

Hence                                                                  becomes

                                                                                   r     R   R
            a
                                                                        maxf j Vi(Xi) - viLi(, r)t.
    P~E
     d oi 3ax,
                                                                       If X * and r* are the optimal arrival-rate vector and
      + vi *[Wi(X*)+x* Wi(X*)]-viWi(X*)                                scheduling permutation, respectively, then by the opti-
                                                                       mality of the c-, rule for fixed X = X * we have
           R a~~~i x
      R                                                                             R            R

   - JX,*- W (X*)                                                       max{ Vi(Xi) - viLi(X,(r)
      j=1 J axj

since for] i d                                                                   R           R


                                                                          = E Vi(Xi) -E viLi(2*, r*)
                                                                               i=l       i=l


1a-
 Ax.Li (2X) a - xa
          axi    iW xi
                    = -i a 'J/(X).                                               R           R


                                                                           X Vi(X*) - E viLi(X, rO)
  Note that when R = 1, (6) reduces to (3). Also, as
in (3), the right-hand side of (6) represents the expected                                   R   R

cost inflicted on all other users as a consequence of                      < max Vi(xi) - viLi(X, rO)}
an infinitesimal unit increase in the job flow of a class-
i user. Thus, similar to (3), we may call the right-hand               which demonstrates that the same scheduling rate is
side of (6) the externality cost. The result of Theorem                also optimal in the current setting.

[Next verbatim source excerpt]

                                                              Consider a system with R job classes that satisfy
class-i jobs. As before, all aggregate class characteristics
                                                            assumptions A1-A5 of Section 1, where class-i jobs
are common knowledge, whereas individual job char-
                                                            have exponential processing requirements with mean
acteristics are known only to the user. Thus, both the
                                                           ci. We assume that
\end{ReviewVerbatim}


\paragraph{Formalized review target}
The archival source above is not asserted equivalent to this formalized target.
\begin{ReviewVerbatim}
Theorem 3's priority-and-time-dependent price implements the welfare-maximizing arrival vector and gives every class a strict assigned-priority preference on the heterogeneous stationary full-positive-flow interior domain.
\end{ReviewVerbatim}

\textbf{Archival source anchor:} {\footnotesize\raggedright cited publication, lines 806-\allowbreak{}846\par}
\paragraph{Expanded PaperInterface specification}
\begin{ReviewVerbatim}
∀ {n : ℕ} (value : Fin n → ℝ → ℝ) (arrivalRate delayCost meanService marginalValue : Fin n → ℝ),
  MendelsonWhang1990PriorityPricing.sourceValueFunctionConditions value →
    MendelsonWhang1990PriorityPricing.sourceDelayCostConditions delayCost →
      ∀ (harrivalRate : ∀ (j : Fin n), 0 < arrivalRate j) (hmeanService : ∀ (j : Fin n), 0 < meanService j),
        ∑ j : Fin n, arrivalRate j * meanService j < 1 →
          (∀ {i k : Fin n}, i < k → meanService i * delayCost k < meanService k * delayCost i) →
            (∀ (i : Fin n), HasDerivAt (value i) (marginalValue i) (arrivalRate i)) →
              IsMaxOn
                  (fun (flow : Fin n → ℝ) =>
                    MendelsonWhang1990PriorityPricing.netValue value delayCost
                      (fun (k : Fin n) (rates : Fin n → ℝ) =>
                        MendelsonWhang1990PriorityPricing.prioritySojournTime rates meanService k)
                      flow)
                  (MendelsonWhang1990PriorityPricing.stablePriorityFlowDomain meanService) arrivalRate →
                MendelsonWhang1990PriorityPricing.priorityTimeDependentPricingOptimalAndIncentiveCompatibleAt value
                    arrivalRate delayCost meanService marginalValue hmeanService
                    (MendelsonWhang1990PriorityPricing.heterogeneousPriorityTimeDependentPrice arrivalRate delayCost
                      meanService)
                    id ∧
                  (∀ (i : Fin n),
                      MendelsonWhang1990PriorityPricing.stationaryClassTaggedExpectedPriorityTimeDependentPrice
                          arrivalRate delayCost meanService harrivalRate i i =
                        MendelsonWhang1990PriorityPricing.externalityPrice delayCost arrivalRate
                          (MendelsonWhang1990PriorityPricing.prioritySojournDerivative arrivalRate meanService i)) ∧
                    ∀ (i selectedPriority : Fin n),
                      selectedPriority ≠ i →
                        0 <
                          MendelsonWhang1990PriorityPricing.priorityCheatingPenalty meanService delayCost
                            (MendelsonWhang1990PriorityPricing.priorityQueueingTime arrivalRate meanService)
                            (MendelsonWhang1990PriorityPricing.priorityTimeLinearCoefficient arrivalRate delayCost
                              meanService)
                            i selectedPriority
\end{ReviewVerbatim}

\paragraph{Retained governing declarations}
\begin{itemize}\raggedright
\item \hyperlink{governing-declaration-806f5a55d6021334}{Mendelson\allowbreak{}Whang1990\allowbreak{}Priority\allowbreak{}Pricing.\allowbreak{}externality\allowbreak{}Price}
\item \hyperlink{paper-prerequisite-d3f248a0395e1c76}{Mendelson\allowbreak{}Whang1990\allowbreak{}Priority\allowbreak{}Pricing.\allowbreak{}heterogeneous\allowbreak{}Priority\allowbreak{}Time\allowbreak{}Dependent\allowbreak{}Price}
\item \hyperlink{paper-prerequisite-7dbaca4145f3a4a5}{Mendelson\allowbreak{}Whang1990\allowbreak{}Priority\allowbreak{}Pricing.\allowbreak{}net\allowbreak{}Value}
\item \hyperlink{paper-prerequisite-54c5682d350b5d3d}{Mendelson\allowbreak{}Whang1990\allowbreak{}Priority\allowbreak{}Pricing.\allowbreak{}priority\allowbreak{}Cheating\allowbreak{}Penalty}
\item \hyperlink{paper-prerequisite-49957d909aa24bc5}{Mendelson\allowbreak{}Whang1990\allowbreak{}Priority\allowbreak{}Pricing.\allowbreak{}priority\allowbreak{}Queueing\allowbreak{}Time}
\item \hyperlink{governing-declaration-41afc2346e61f660}{Mendelson\allowbreak{}Whang1990\allowbreak{}Priority\allowbreak{}Pricing.\allowbreak{}priority\allowbreak{}Sojourn\allowbreak{}Derivative}
\item \hyperlink{governing-declaration-703b87082556903f}{Mendelson\allowbreak{}Whang1990\allowbreak{}Priority\allowbreak{}Pricing.\allowbreak{}priority\allowbreak{}Sojourn\allowbreak{}Time}
\item \hyperlink{paper-prerequisite-3ae6b5293fffafb7}{Mendelson\allowbreak{}Whang1990\allowbreak{}Priority\allowbreak{}Pricing.\allowbreak{}priority\allowbreak{}Time\allowbreak{}Dependent\allowbreak{}Pricing\allowbreak{}Optimal\allowbreak{}And\allowbreak{}Incentive\allowbreak{}Compatible\allowbreak{}At}
\item \hyperlink{paper-prerequisite-ae05cc1b41c33470}{Mendelson\allowbreak{}Whang1990\allowbreak{}Priority\allowbreak{}Pricing.\allowbreak{}priority\allowbreak{}Time\allowbreak{}Linear\allowbreak{}Coefficient}
\item \hyperlink{paper-prerequisite-33c9dc55ddc7d5b6}{Mendelson\allowbreak{}Whang1990\allowbreak{}Priority\allowbreak{}Pricing.\allowbreak{}source\allowbreak{}Delay\allowbreak{}Cost\allowbreak{}Conditions}
\item \hyperlink{paper-prerequisite-9c435671d506207c}{Mendelson\allowbreak{}Whang1990\allowbreak{}Priority\allowbreak{}Pricing.\allowbreak{}source\allowbreak{}Value\allowbreak{}Function\allowbreak{}Conditions}
\item \hyperlink{paper-prerequisite-84dca75fd93aa845}{Mendelson\allowbreak{}Whang1990\allowbreak{}Priority\allowbreak{}Pricing.\allowbreak{}stable\allowbreak{}Priority\allowbreak{}Flow\allowbreak{}Domain}
\item \hyperlink{paper-prerequisite-816477dc3af70c36}{Mendelson\allowbreak{}Whang1990\allowbreak{}Priority\allowbreak{}Pricing.\allowbreak{}stationary\allowbreak{}Class\allowbreak{}Tagged\allowbreak{}Expected\allowbreak{}Priority\allowbreak{}Time\allowbreak{}Dependent\allowbreak{}Price}
\end{itemize}
\paragraph{Lean-elaborated claim roles}\begin{ReviewVerbatim}
1. parameter: ℕ
2. parameter: Fin n → ℝ → ℝ
3. parameter: Fin n → ℝ
4. parameter: Fin n → ℝ
5. parameter: Fin n → ℝ
6. parameter: Fin n → ℝ
7. assumption: MendelsonWhang1990PriorityPricing.sourceValueFunctionConditions value
8. assumption: MendelsonWhang1990PriorityPricing.sourceDelayCostConditions delayCost
9. assumption: ∀ (j : Fin n), 0 < arrivalRate j
10. assumption: ∀ (j : Fin n), 0 < meanService j
11. assumption: ∑ j : Fin n, arrivalRate j * meanService j < 1
12. assumption: ∀ {i k : Fin n}, i < k → meanService i * delayCost k < meanService k * delayCost i
13. assumption: ∀ (i : Fin n), HasDerivAt (value i) (marginalValue i) (arrivalRate i)
14. assumption: IsMaxOn
  (fun (flow : Fin n → ℝ) =>
    MendelsonWhang1990PriorityPricing.netValue value delayCost
      (fun (k : Fin n) (rates : Fin n → ℝ) => MendelsonWhang1990PriorityPricing.prioritySojournTime rates meanService k)
      flow)
  (MendelsonWhang1990PriorityPricing.stablePriorityFlowDomain meanService) arrivalRate
15. conclusion: MendelsonWhang1990PriorityPricing.priorityTimeDependentPricingOptimalAndIncentiveCompatibleAt value arrivalRate
    delayCost meanService marginalValue hmeanService
    (MendelsonWhang1990PriorityPricing.heterogeneousPriorityTimeDependentPrice arrivalRate delayCost meanService) id ∧
  (∀ (i : Fin n),
      MendelsonWhang1990PriorityPricing.stationaryClassTaggedExpectedPriorityTimeDependentPrice arrivalRate delayCost
          meanService harrivalRate i i =
        MendelsonWhang1990PriorityPricing.externalityPrice delayCost arrivalRate
          (MendelsonWhang1990PriorityPricing.prioritySojournDerivative arrivalRate meanService i)) ∧
    ∀ (i selectedPriority : Fin n),
      selectedPriority ≠ i →
        0 <
          MendelsonWhang1990PriorityPricing.priorityCheatingPenalty meanService delayCost
            (MendelsonWhang1990PriorityPricing.priorityQueueingTime arrivalRate meanService)
            (MendelsonWhang1990PriorityPricing.priorityTimeLinearCoefficient arrivalRate delayCost meanService) i
            selectedPriority
\end{ReviewVerbatim}

\paragraph{Lean proof endpoint (built separately)}\begin{ReviewVerbatim}
MendelsonWhang1990PriorityPricing.theoremThreePriorityTimeDependentPricing
\end{ReviewVerbatim}

\paragraph{Recorded source-to-Spec screening}
\textbf{Verdict:} \texttt{matches\_approved\_corrected\_target}
\\
{\raggedright
\textbf{Reason:} Theorem 3 states that the Equation (19) PTD schedule is optimal and incentive-\allowbreak{}compatible, and its supplied proof excerpt identifies each truthful expected charge with the Theorem 1 externality price and derives strict preference over every distinct priority.\allowbreak{}  The target expresses those conclusions through the source's equilibrium conditions and the recorded stable full-\allowbreak{}positive interior scope, without attributing strict zero-\allowbreak{}flow comparisons to the archival statement.\allowbreak{}
\par}
\reviewmatch{Matches formalized target}
\noindent\textbf{Reviewer annotation}\par
\reviewerbox
\clearpage
\hypertarget{source-claim-324a284201e724f2}{}
\section*{Review row 5}
\textbf{Source locator:} {\footnotesize\raggedright cited publication, lines 940-\allowbreak{}985\par}
\paragraph{Verbatim source input}
\begin{ReviewVerbatim}
Theorem 4. For all user classes i = 1, 2, ..., R,
job is, therefore, given by
                                                                        the cheating penalty ll'(k) satisfies 0 = H'(i) <
  R       R   t2                                        HI .(i + 1) < H 1.(i + 2) < ... < H' .(R) and H` (l) >
  E       Ek        =     E      -   Vk           Xkl(Sk-I
                                                        Hli(2) >3k)
                                                                  ... > H'(i) = 0. Equivalently, ll'(k) is a
k=l k=- 2
                                                                        unimodalfunction of k.
which is the quadratic term of (19). The linear term
of the proposed price schedule can be interpreted in a                  Proof. Suppose that j > i. Then, by the definition of
similar way as the externality cost inflicted on other                  ll'(k), we have
jobs arriving before the tagged job starts its service:
the first term of Ai t, ait/(Si_1 ?) is the externality cost            IH.(j + 1) - HIi(j)
inflicted by the tagged job on the jobs of its own
classes (i.e., class-i), while                                              = v1(J471+ - Wj) -c1(Aj-A1+1) (24)

                                                                        But we know from Theorem 3 that
     k=i+l k-I k Sk- k
                                                                        IIji(j + 1) = vj(JW,7 - W,) - cj(Aj - Aj+l) > 0
represents the externality cost inflicted on the lower
class jobs. To further decompose the second term, the                  or
first part

      R


      t , E ak
    k=i+l Sk_1 Sk
                                                                         WUq>- W,q >-(A --Aj+1). (25)
                                         It follows from equations (24) and (25) that
is incurred by the lower class jobs which have arrived
before the tagged job arrives, and the second part
          R   1
                                                                        HIi(j + 1) - Hi(j) > c _ C)(Aj - Aj+1) > 0
t - E 3 2
    k=i+l Ak-lk
                                                                       where the last inequality follows from vi/ci > vj/cj and
is incurred by the lower class jobs coming after the                   Aj > Aj+1. This shows that 0 = H'(i) < IIi(i + 1) <
                                                                        ... < li(R).
tagged job, but before the tagged job starts its service.
                                                                            A similar argument for the case j < i implies
Altogether, the price represents the conditional expec-
tation of the tagged job's externality on the system.
                                                                        fIi(l) > WI(2) > . . . > iii(i) = O.
  From a class-i user's viewpoint, the total expected
\end{ReviewVerbatim}


\paragraph{Formalized review target}
The archival source above is not asserted equivalent to this formalized target.
\begin{ReviewVerbatim}
Theorem 4's cheating penalty is zero at the assigned priority and strictly increases away from it on the stationary full-positive-flow interior domain.
\end{ReviewVerbatim}

\textbf{Archival source anchor:} {\footnotesize\raggedright cited publication, lines 940-\allowbreak{}985\par}
\paragraph{Expanded PaperInterface specification}
\begin{ReviewVerbatim}
∀ {n : ℕ} (arrivalRate delayCost meanService : Fin n → ℝ),
  (∀ (j : Fin n), 0 < arrivalRate j) →
    (∀ (j : Fin n), 0 < meanService j) →
      ∑ j : Fin n, arrivalRate j * meanService j < 1 →
        (∀ {i k : Fin n}, i < k → meanService i * delayCost k < meanService k * delayCost i) →
          (∀ (i : Fin n),
              MendelsonWhang1990PriorityPricing.priorityCheatingPenalty meanService delayCost
                  (MendelsonWhang1990PriorityPricing.priorityQueueingTime arrivalRate meanService)
                  (MendelsonWhang1990PriorityPricing.priorityTimeLinearCoefficient arrivalRate delayCost meanService) i
                  i =
                0) ∧
            (∀ {i j k : Fin n},
                i < j →
                  j < k →
                    MendelsonWhang1990PriorityPricing.priorityCheatingPenalty meanService delayCost
                        (MendelsonWhang1990PriorityPricing.priorityQueueingTime arrivalRate meanService)
                        (MendelsonWhang1990PriorityPricing.priorityTimeLinearCoefficient arrivalRate delayCost
                          meanService)
                        i j <
                      MendelsonWhang1990PriorityPricing.priorityCheatingPenalty meanService delayCost
                        (MendelsonWhang1990PriorityPricing.priorityQueueingTime arrivalRate meanService)
                        (MendelsonWhang1990PriorityPricing.priorityTimeLinearCoefficient arrivalRate delayCost
                          meanService)
                        i k) ∧
              ∀ {i j k : Fin n},
                k < i →
                  j < k →
                    MendelsonWhang1990PriorityPricing.priorityCheatingPenalty meanService delayCost
                        (MendelsonWhang1990PriorityPricing.priorityQueueingTime arrivalRate meanService)
                        (MendelsonWhang1990PriorityPricing.priorityTimeLinearCoefficient arrivalRate delayCost
                          meanService)
                        i k <
                      MendelsonWhang1990PriorityPricing.priorityCheatingPenalty meanService delayCost
                        (MendelsonWhang1990PriorityPricing.priorityQueueingTime arrivalRate meanService)
                        (MendelsonWhang1990PriorityPricing.priorityTimeLinearCoefficient arrivalRate delayCost
                          meanService)
                        i j
\end{ReviewVerbatim}

\paragraph{Retained governing declarations}
\begin{itemize}\raggedright
\item \hyperlink{paper-prerequisite-54c5682d350b5d3d}{Mendelson\allowbreak{}Whang1990\allowbreak{}Priority\allowbreak{}Pricing.\allowbreak{}priority\allowbreak{}Cheating\allowbreak{}Penalty}
\item \hyperlink{paper-prerequisite-49957d909aa24bc5}{Mendelson\allowbreak{}Whang1990\allowbreak{}Priority\allowbreak{}Pricing.\allowbreak{}priority\allowbreak{}Queueing\allowbreak{}Time}
\item \hyperlink{paper-prerequisite-ae05cc1b41c33470}{Mendelson\allowbreak{}Whang1990\allowbreak{}Priority\allowbreak{}Pricing.\allowbreak{}priority\allowbreak{}Time\allowbreak{}Linear\allowbreak{}Coefficient}
\end{itemize}
\paragraph{Lean-elaborated claim roles}\begin{ReviewVerbatim}
1. parameter: ℕ
2. parameter: Fin n → ℝ
3. parameter: Fin n → ℝ
4. parameter: Fin n → ℝ
5. assumption: ∀ (j : Fin n), 0 < arrivalRate j
6. assumption: ∀ (j : Fin n), 0 < meanService j
7. assumption: ∑ j : Fin n, arrivalRate j * meanService j < 1
8. assumption: ∀ {i k : Fin n}, i < k → meanService i * delayCost k < meanService k * delayCost i
9. conclusion: (∀ (i : Fin n),
    MendelsonWhang1990PriorityPricing.priorityCheatingPenalty meanService delayCost
        (MendelsonWhang1990PriorityPricing.priorityQueueingTime arrivalRate meanService)
        (MendelsonWhang1990PriorityPricing.priorityTimeLinearCoefficient arrivalRate delayCost meanService) i i =
      0) ∧
  (∀ {i j k : Fin n},
      i < j →
        j < k →
          MendelsonWhang1990PriorityPricing.priorityCheatingPenalty meanService delayCost
              (MendelsonWhang1990PriorityPricing.priorityQueueingTime arrivalRate meanService)
              (MendelsonWhang1990PriorityPricing.priorityTimeLinearCoefficient arrivalRate delayCost meanService) i j <
            MendelsonWhang1990PriorityPricing.priorityCheatingPenalty meanService delayCost
              (MendelsonWhang1990PriorityPricing.priorityQueueingTime arrivalRate meanService)
              (MendelsonWhang1990PriorityPricing.priorityTimeLinearCoefficient arrivalRate delayCost meanService) i k) ∧
    ∀ {i j k : Fin n},
      k < i →
        j < k →
          MendelsonWhang1990PriorityPricing.priorityCheatingPenalty meanService delayCost
              (MendelsonWhang1990PriorityPricing.priorityQueueingTime arrivalRate meanService)
              (MendelsonWhang1990PriorityPricing.priorityTimeLinearCoefficient arrivalRate delayCost meanService) i k <
            MendelsonWhang1990PriorityPricing.priorityCheatingPenalty meanService delayCost
              (MendelsonWhang1990PriorityPricing.priorityQueueingTime arrivalRate meanService)
              (MendelsonWhang1990PriorityPricing.priorityTimeLinearCoefficient arrivalRate delayCost meanService) i j
\end{ReviewVerbatim}

\paragraph{Lean proof endpoint (built separately)}\begin{ReviewVerbatim}
MendelsonWhang1990PriorityPricing.theoremFourCheatingPenaltyMonotonicity
\end{ReviewVerbatim}

\paragraph{Recorded source-to-Spec screening}
\textbf{Verdict:} \texttt{matches\_approved\_corrected\_target}
\\
{\raggedright
\textbf{Reason:} The source's zero truthful penalty and two strict monotonicity chains are preserved by the Spec.\allowbreak{} Its full-\allowbreak{}positive stationary interior domain is the map-\allowbreak{}pinned correction, with no archival-\allowbreak{}equivalence claim.\allowbreak{}
\par}
\reviewmatch{Matches formalized target}
\noindent\textbf{Reviewer annotation}\par
\reviewerbox

\end{document}
